\documentclass[9pt,twoside]{extarticle}
\usepackage{Breviary-Appendix}
%%
\begin{document}
\section*{\center \textsc{Appendix Verna}\\ \large \textit{\textred{Ex Breviario Carmelitarum 1866 - Pars Verna}}}
\begin{multicols}{2} % Start the two-column layout
%This first office is the same as in the Pars Hiemalis, but much less complete. This office in the Verna ends with Lectio viij.
\renewcommand{\thepage}{\arabic{page}*}
\subsubsection*{\center \textsc{NOTANDUM}}
\begin{tightlist}{}
\item \textred{Per totam Quadrages. quando fit de Festo, legitur semper Lectio nona de Homil. Feriæ occurr. et fit commem. ejusdem in Laud. et Vesperis, de quocumque fiat Offic. fitque hæc commemoratio ante commemorationem cujuscumque simplicissimi.}
\item \textred{Officia Passionis Domini Nostri Jesu Christi, non ultra Quadragesimam transferenda. (Ex decreto S.R.C. 31 Jan. 1879.)}
\end{tightlist}
\begin{center}\textcolor{red}{\rule{\linewidth}{0.5pt}} \end{center}
\fancyhead[C]{\textsc{Comm. Passionis D.N.J.C.}} % Book title in header center for even pages
\subsection*{\center \textsc{FERIA TERTIA}\\
\textred{Post Domin. Sexages.}}
\subsection*{\center \textsc{Commem. Passionis D.N.J.C.}}
\subsubsection*{\center \textred{Duplex minus. 1 classis.}}
%\begin{center}\textcolor{red}{\rule{\linewidth}{0.5pt}} \end{center}

\subsubsection*{\center \textsc{AD VESPERAS.}}
\begin{tightlist}{}
\item \textred{Aña.} Cálicem salutáris accípiam et nomen Dómini invocábo.
\item \textred{Psalmi de Offic. præc.}
\end{tightlist}
\subsubsection*{\textred{Capitulum. \hfill Philip. 2.}}
\lettrineThree{F}ratres, Hoc enim sentíte in vobis, quod et in Christo Jesu, qui, cum in forma Dei esset, non rapínam arbitrátus est, esse se æquálem Deo~: sed semetípsum exinanívit, formam servi accípiens, in similitúdinem hóminum factus, et hábitu invéntus ut homo. Humiliávit semetípsum factus obédiens usque ad mortem, mortem autem crucis.

\begin{tightlist}{}
\item \textred{℟.}~Crucem tuam adorámus, Dómine, sanctámque recólimus passiónem~: \textred{*} Ecce enim propter lignum venit gáudium in univérso mundo. \textred{℣.}~Deus misereátur nostri, et benedícat nobis, illúminet vultum suum super nos, et misereátur nostri. \textred{E}cce enim. Glória Patri. \textred{E}cce enim.
\end{tightlist}

\subsubsection*{\center \textred{Hymnus.}}
\lettrine{M}œréntes óculi spárgite lácrymas,\\
Et luctu résonent íntima córdium,\\
Illátas réfero Númini ab ímpiis\\
Pœnas, et fera vúlnera.\\
\\
\textred{A}ccíncta heu! gládiis turba satéllitum\\
Arréptum Dóminum fústibus ímpetit;\\
Nunc cædit cólaphis, nunc quatit hórridis\\
Divínum caput íctibus.\\
\\
\textred{H}aud finis scéleri~: tráditur ímprobo\\
Christus carnífici; nec mora, bárbarus\\
In Regem súperum non timet ímpio\\
Ausu, vértere déxteram.\\
\\
\textred{A}udíte, o pópuli~: Numen amábile,\\
Manánte ex húmeris úndique sánguine,\\
Lictóris rábidi sústinet ímpetum,\\
Et vocem premit ínnocens.\\
\\
\textred{Q}uis non illácrymet? Jam nova cóndidit\\
Torménta indómitæ gentis iníquitas;\\
Infígit cérebro proh dolor! éfferis\\
Sertum vépribus ásperum.\\
\\
\textred{H}eu crimen! tráhitur fúnibus ímprobis\\
Funésti Dóminus supplícii ad locum;\\
Illic oppétiit fúnere, spíritum\\
Patri restítuens suo.\\
\\
\textred{P}asso pro míseris tam fera vúlnera,\\
In terris résonet débita glória;\\
Sacrúmque assídue Nomen in æthera\\
Humánum genus éfferat.\\
Amen.

\begin{tightlist}{}
\item \textred{℣.}~Oblátus est quia ipse vóluit. \textred{℟.}~Et livóre ejus sanáti sumus.
\item \textred{Ad Magnificat, Aña.} O vos omnes, qui transítis per viam, atténdite, et vidéte, si est dolor sicut dolor meus.
\end{tightlist}

\subsubsection*{\center \textred{Oratio.}}
\lettrine{O}mnípotens sempitérne Deus, qui humáno géneri, ad imitándum humilitátis exémplum, Salvatórem nostrum carnem súmere, et crucem subíre fecísti~: concéde propítius; ut, sicut solémnem Commemoratiónem Passiónis ejus celebrámus, ita et patiéntiæ ipsíus habére documénta, et resurrectiónis consórtia mereámur. Per eúmdem Dóminum.

\begin{tightlist}{}
\item \textred{Hæc Oratio dicitur etiam ad Laud. Horas, et ij. Vesp.}
\item \textred{Reliqua omnia ad Matut. Laud. et Horas ut in Offic. vot., [90]*, exceptis hic sign.}
\end{tightlist}

\subsubsection*{\center \textsc{IN SECUNDO NOCTURNO.}}

\subsubsection*{\center Sermo sancti Augustíni Epíscopi.}
\subsubsection*{\center \textred{Sermo 41. de Passione Domini.\\ Lectio iv.}}
\lettrineThree{A}dmonet nos, fratres caríssimi, ad solemnitátem Domínicæ Passiónis dies ipse, in quo eam nec muta eleménta tacuérunt. Célebret eam lux fídei, linguis hóminum, quam conclamavérunt étiam siléntia tenebrárum. Hódie Dóminus noster in statéra crucis prétium nostræ salútis appéndit, et una morte univérsum mundum, sicut ómnium Cónditor, ita ómnium Reparátor absólvit. Indubitánter enim credámus, quod totum mundum redémit, qui plus dedit, quam totus mundus valéret. Méritum enim redémptæ mercédis dígnitas insígnis prétii supergréssa est. Inter redémptum et rediméntem dispensátio fuit, compensátio non fuit. Qui ergo non habébat peccáta própria, digne delévit aliéna. Solus hic pia víctima pro ómnibus cécidit, ut omnes leváret. Et quia débitum solus non hábuit, recte fœnus misericórdiæ pro debitóribus erogávit.

\begin{tightlist}{}
\item \textred{℟.}~Ego propter te Chananæórum reges percússi, et dedi tibi sceptrum regále; \textred{*} Et tu dedísti cápiti meo spíneam corónam, et percussísti arúndine caput meum. \textred{℣.}~Pópule meus, quid feci tibi, aut in quo contristávi te? respónde mihi. \textred{E}t tu dedísti cápiti meo.
\end{tightlist}

\subsubsection*{\center\textred{Lectio v.}}
\lettrine{P}erpénde inter hæc, qui talem pro nobis dedit pecúniam, qualem a nobis sit exactúrus usúram. In hac ítaque die fides prophéticæ annuntiatiónis impléta est, ita dicéntis~: Corpus meum dedi percutiéntibus~: et genas meas velléntibus; fáciem meam non avérti a fœditáte sputórum. Suscépit mala nostra, ut tribúeret bona sua. Hinc intelligámus, quantum hóminem dilígere dignátus sit ante culpam, quem sic díligit post ruínam. Agnósce, homo, quantum váleas et quantum débeas~: et dum tantam redemptiónis tuæ pérspicis dignitátem, ipse tibi indícito peccándi pudórem. Ecce, pro ímpio píetas flagellátur, pro stulto sapiéntia illúditur, pro mendáce véritas necátur, damnátur justítia pro iníquo, misericórdia affícitur pro crudéli, pro mísero replétur sincéritas acéto, inebriátur felle dulcédo, addícitur innocéntia pro reo, móritur vita pro mórtuo.

\begin{tightlist}{}
\item \textred{℟.}~Ego duxi te per desértum quadragínta annis, et manna cibávi te; \textred{*} Et tu me cecidísti álapis, et flagéllis. \textred{℣.}~Pópule meus, quid feci tibi, aut in quo contristávi te? respónde mihi. \textred{E}t tu me.
\end{tightlist}

\subsubsection*{\center\textred{Lectio vj.}}
\lettrine{E}xpávit scelus hóminum natúra rerum, et quem creatúra rebéllis non agnóscit, eum mundi Dóminum tremens terra testátur, et cœli Regem sol fúgiens confitétur. Clámyde coccínea indúitur, quia sánguine Mártyrum suórum Ecclésiæ corpus ornátur. Coróna spínea cápiti ejus impónitur, quia púnctio peccatórum nostrórum, quorum remissióne Redemptóris glória strúitur, áridis tríbulis comparátur. Studeámus nunc e divérso, ut membrórum vita cápitis sit coróna. Quod vero sitíre se in cruce pósitus dicit, fidem incrédulæ gentis concupíscit. Sed e contra acétum malítiæ pórrigunt, quia vinum sapiéntiæ, quod a Deo accéperant, peccándo corrúperant. Velum templi scínditur, quia Synagóga honóre nudátur, observátio antíqua dissólvitur, Ecclésiæ únitas præmonstrátur. Monuménta aperiúntur, quia mortis jura jure superántur.

\begin{tightlist}{}
\item \textred{℟.}~Ego exaltávi te magna virtúte~: \textred{*} Et tu me suspendísti in patíbulo crucis, in qua expándi manus meas ad pópulum non credéntem et contradicéntem. \textred{℣.}~Pópule meus, quid feci tibi, aut in quo contristávi te? respónde mihi. \textred{E}t.Glória Patri. \textred{E}t.
\end{tightlist}
\subsubsection*{\center \textsc{IN TERTIO NOCTURNO.}}
%\begin{tightlist}{}
%\item \textred{Aña.} Fílii hóminum dentes eórum arma, et sagíttæ, et linga eórum gládius acútus. 
%\item \textred{Ps.} Miserére mei Deus, miserére mei. 63.
%\item \textred{Aña.} Intendérunt arcum rem amáram, ut sagíttent in occúltis immaculátum.
%\item \textred{Ps.} Exáudi Deus oratiónem. 67.
%\item \textred{Aña.} Factus sum sicut homo sine adjutório inter mórtuos liber.
%\item \textred{Ps.} Dómine Deus salútis meæ. 97.
%\item \textred{℣.}~Factus est obédiens usque ad mortem. \textred{℟.}~Mortem autem crucis.
%\end{tightlist}

\subsubsection*{\center Léctio sancti Evangélii secúndum Joánnem.}
\subsubsection*{\center \textred{Lectio vij. \hfill Cap. 19.}}
\lettrine{I}n illo témpore~: Sciens Jesus, quia ómnia consummáta sunt, ut consummarétur scriptúra, dixit~: Sítio. Et réliqua.

\begin{center}Homilía sancti Joánnis Chrysóstomi.\\
\textred{Hom. 84. in c. 19. Joann.}\end{center}
\lettrineThree{S}ciens Jesus ómnia consummáta, ait, Sítio~: étiam hoc verbo prophetíam implens. Tu autem adstántium pravitátem intéllige; nam licet innúmeros habeámus inimícos, qui nos gravíssime persecúti sint, cum tamen intérimi eos vidémus, fléctimur misericórdia. Hi vero nil moti, magis efferántur, et irrídent, acetúmque in spóngia pórrigunt, ut bibat tamquam reus~: nam propter hoc hyssópum ádditur, quo gustáto, mórtuus est. Vides quam non perturbáte, et quam poténter ómnia transigéntem? Quod et prósequens osténdit; cum enim ómnia consummáta essent, inclináto cápite, quod affíxum non erat, trádidit spíritum, hoc est exspirávit.

\begin{tightlist}{}
\item \textred{℟.}~Ego te potávi aqua salútis de petra~: \textred{*} Et tu me potásti felle, et acéto. \textred{℣.}~Pópule meus, quid feci tibi, aut in quo contristávi te? respónde mihi. \textred{E}t tu.
\end{tightlist}

\subsubsection*{\center\textred{Lectio viij.}}
\lettrine{A}tqui non post inclinátum caput exspirátur; sed post exspiratiónem caput inclinátur. Hic autem contra áccidit, per quod osténdit eum Evangelísta esse ómnium Dóminum. Judæi autem camélum devorántes, et cúlicem excolántes, tantum nefas ausi, de die consúlant~: Quóniam Parascéve erat, ut non remanérent in cruce córpora sábbato, rogavérunt Pilátum, ut frangeréntur eórum crura. Vides quanti sit véritas? Per ea, quæ illi fáciunt, prophetía implétur; ália enim ex hoc manifestáta est. Venérunt ergo mílites, et aliórum fregérunt crura, non Christi; sed ad Judæórum grátiam conciliándam latus ejus láncea aperuérunt, et mórtuo adhuc insúltant. O péssimam voluntátem, ac scelestíssimam!

\begin{tightlist}{}
\item \textred{℟.}~Ego ante te apérui mare~: \textred{*} Et tu aperuísti láncea latus Salvatóri tuo, qui redémit mundum in sánguine suo. \textred{℣.}~Pópule meus, quid feci tibi, aut in quo contristávi te? respónde mihi. \textred{E}t tu.
\end{tightlist}

\begin{tightlist}{}
\item \textred{Ad Vesperas, Añæ Laudum. Psalmi, Hymnus ℣. et Aña ad Magnificat, ut in Officio votivo, infra.} [97*]
\end{tightlist}
%The rubric at the end says to take the rest for Vespers and such from the votive office in the back, and does not have anything for the minor hours this day.
\begin{center}\textcolor{red}{\rule{\linewidth}{0.5pt}} \end{center}   

\subsection*{\center \textsc{FERIA SEXTA}\\
\textred{Post Domin. 1. Quadrag.}}
\subsection*{\center \textsc{Lanceæ et Clavorum D.N.J.C.}}
\subsubsection*{\center \textred{Duplex minus. 1 classis.}}
%\begin{center}\textcolor{red}{\rule{\linewidth}{0.5pt}} \end{center}
\subsubsection*{\center \textsc{AD VESPERAS.}}
\fancyhead[C]{\textsc{Lanceæ et Clavorum D.N.J.C.}}
    \begin{tightlist}{}
\item \textred{Aña.} Unus mílitum láncea latus ejus apéruit, et contínuo exívit sanguis et aqua.
\item \textred{Psalmi de Offic. præced.}
\end{tightlist}

\subsubsection*{\textred{Capitulum. \hfill 1. Joan. 5.}}
\lettrineThree{C}aríssimi~: Quis est, qui vincit mundum, nisi qui credit, quóniam Jesus est fílius Dei? Hic est, qui venit per aquam, et sánguinem, Jesus Christus, non in aqua solum, sed in aqua, et sánguine.

\begin{tightlist}{}
\item \textred{℟.}~Adorámus te, Christe, et benedícimus tibi, quia per sanctam Crucem tuam redemísti mundum; \textred{*} Miserére nostri, qui passus es pro nobis. \textred{℣.}~In ligno pendens, nostræ salútis viam Verbum Patris invénit, in qua Redémptor noster, mortis víncula rumpens cállidam áspidem superávit. \textred{M}iserére. Glória Patri. \textred{M}iserére.
\end{tightlist}

\subsubsection*{\center \textred{Hymnus.}}
\lettrine{Q}uænam lingua tibi, o Láncea, débitas\\
Grates pro mérito est apta repéndere?\\
Christi vivíficum namque áperis latus,\\
Unde Ecclésia náscitur.\\
\\
\textred{H}æc est Heva viri de látere éxiens,\\
Olli membra gravis dum sopor óccupat;\\
Hanc quippe alter Adam, corde scaténtibus\\
Unda et sánguine, prócreat.\\
\\
\textred{O} Clavi! æqua manet vos quoque grátia,\\
Christi quando sacris ártubus ínsiti,\\
Delétum Dómini sánguine fígitis\\
Mortis chirógraphum cruci.\\
\\
\textred{T}e, Jesu, Súperi láudibus éfferant,\\
Qui Clavórum áditus, sígnaque Lánceæ\\
In cœlo rétines, vivus ubi ímperas\\
Cum Patre atque Paráclito.\\
Amen.

\begin{tightlist}{}
\item \textred{℣.}~Fodérunt manus meas et pedes meos. \textred{℟.}~Dinumeravérunt ómnia ossa mea.
\item \textred{Ad Magnif. Aña.} Delens quod advérsus nos erat, chirógraphum decréti, quod erat contrárium nobis, et ipsum tulit de médio, affígens illud cruci.
\end{tightlist}

\subsubsection*{\center \textred{Oratio.}}
\lettrine{D}eus qui in assúmptæ carnis infirmitáte Clavis affígi, et Láncea vulnerári pro mundi salúte voluísti~: concéde propítius, ut qui eorúmdem Clavórum, et Lánceæ solémnia venerámur in terris, de glorióso victóriæ tuæ triúmpho gratulémur in cœlis. Qui vivis.

\begin{tightlist}{}
\item \textred{Ad Completorium, et per Horas in fine Hymnorum dicitur.}
\item \textred{C}lavis foráto et Láncea,\\
Jesu, tibi siti glória,\\
Cum Patre et almo Spíritu\\
Nunc et per omne sæculum.\\
Amen.
\item \textred{Ad Nunc dimittis, Aña.} Hoc signum Crucis * erit in cælo, cum Dóminus ad judicándum vénerit.
\end{tightlist}

\subsubsection*{\center \textsc{AD MATUTINUM.}}
\begin{tightlist}{}
\item \textred{Invit.} Christum Láncea et Clavis vulnerátum, \textred{*} Veníte adorémus.
\item \textred{Ps.} Veníte, exsultémus. 2.
\end{tightlist}

\subsubsection*{\center \textred{Hymnus.}}
\lettrine{S}alvéte Clavi et Láncea,\\
Squalóre nuper óbsita,\\
Quæ mersa Christi córpori,\\
Almo rubétis sánguine.\\
\\
\textred{V}os ad scelus Judáica\\
Elégerat perfídia~:\\
Sed in minístra grátiæ\\
Vos vertit e cœlo Deus.\\
\\
\textred{N}am quot sacrátis ártubus,\\
Sculpsístis olim vúlnera,\\
E tot reclúsis fóntibus\\
Dona éffluunt cœléstia.\\
\\
\textred{C}lavis foráto et Láncea,\\
Jesu, tibi sit glória,\\
Cum Patre et almo Spíritu\\
Nunc et per omne sæculum.\\
Amen.

\begin{tightlist}{}
\item \textred{Si omissus sit Hymnus ad Vesperas, dicatur ad Matutinum, et Hymnus Matutini conjungatur cum Hymno ad Laudes.}
\end{tightlist}

\subsubsection*{\center  \textsc{IN PRIMO NOCTURNO.}}
\begin{tightlist}{}
\item \textred{Aña.} Ad Jesum autem cum veníssent, non fregérunt ejus crura, sed unus mílitum Láncea latus ejus apéruit.
\item \textred{Ps. 1.} Beátus vir, [76] \textred{cum reliq. j. et ij. Noct. de Communi Confess. Pont.}
\item \textred{Aña.} Exívit sanguis et aqua; et qui vidit, testimónium perhíbuit, et verum est testimónium ejus.
%\item \textred{Ps. 2.} Quare fremuérunt.
\item \textred{Aña.} Alia Scriptúra dicit~: Vidébunt in quem transfixérunt.
%\item \textred{Ps. 3.} Dómine quid.
\item \textred{℣.}~Fodérunt manus meas et pedes meos. \textred{℟.}~Dinumeravérunt ómnia ossa mea.
\end{tightlist}

\subsubsection*{\center De Zacharía Prophéta.}
\subsubsection*{\center \textred{Lectio j. \hfill Cap. 12.}}
\lettrineThree{H}æc dicit Dóminus~: Effúndam super domum David, et super habitatóres Jerúsalem spíritum grátiæ, et precum~: et adspícient ad me, quem confixérunt. Et plangent eum planctu quasi super unigénitum, et dolébunt super eum, ut doléri solet in morte primogéniti. In die illa magnus erit planctus in Jerúsalem, sicut planctus Adadrémmon in campo Magéddon. Et planget terra, famíliæ et famíliæ seórsum~: famíliæ domus David seórsum, et mulíeres eórum seórsum~: famíliæ domus Nathan seórsum, et mulíeres eórum seórsum~: famíliæ domus Levi seórsum, et mulíeres eórum seórsum~: famíliæ Sémei seórsum, et mulíeres eórum seórsum. Omnes famíliæ réliquæ, famíliæ et famílæ seórsum, et mulíeres eórum seórsum. Hæc dicit Dóminus Deus~: convertímini ad me, et salvi éritis.

\begin{tightlist}{}
\item \textred{℟.}~Ad Jesum autem cum veníssent, ut vidérent eum jam mórtuum, non fregérunt ejus crura, \textred{*} Sed unus mílitum Láncea latus ejus apéruit. \textred{℣.}~Et qui vidit, testimónium perhíbuit~: et verum est testimónium ejus. \textred{S}ed.
\end{tightlist}

\subsubsection*{\center \textred{Lectio ij. \hfill Cap. 13.}}
\lettrine{I}n die illa erit fons patens dómui David, et habitántibus Jerúsalem in ablutiónem peccatóris, et menstruátæ. Et erit in die illa, dicit Dóminus exercítuum~: Dispérdam nómina idolórum de terra, et non memorabúntur ultra~: et pseudoprophétas, et spíritum immúndum áuferam de terra. Et erit, cum prophetáverit quíspiam ultra, dicent ei pater ejus, et mater ejus, qui genuérunt eum~: Non vives, quia mendácium locútus es in nómine Dómini. Et confígent eum pater ejus, et mater ejus, genitóres ejus, cum prophetáverit. Et erit~: In die illa confundéntur prophétæ, unusquísque ex visióne sua cum prophetáverit~: nec operiéntur pállio sáccino, ut mentiántur~: sed dicet~: Non sum prophéta, homo agrícola ego sum~: quóniam Adam exémplum meum ab adolescéntia mea. Hæc dicit Dóminus.

\begin{tightlist}{}
\item \textred{℟.}~Unus mílitum Láncea latus ejus apéruit. \textred{*} Et contínuo exívit sanguis, et aqua. \textred{℣.}~In die illa erit fons patens dómui David, et habitántibus Jerúsalem. \textred{E}t.
\end{tightlist}

\subsubsection*{\center \textred{Lectio iij.}}
\lettrine{E}t dicétur ei~: Quid sunt plagæ istæ in médio mánuum tuárum? Et dicet~: His plagátus sum in domo eórum qui dilegébant me. Frámea suscitáre super pastórem meum, et super virum cohæréntem mihi, dicit Dóminus exercítuum~: pércute pastórem, et dispergéntur oves, et convértam manum meam ad párvulos. Et erunt in omni terra, dicit Dóminus~: partes duæ in ea dispergéntur et defícient; et tértia pars relinquétur in ea. Et ducam tértiam partem per ignem, et uram eos sicut úritur argéntum, et probábo eos sicut probátur aurum. Ipse vocábit nomen meum, et ego exáudiam eum. Dicam~: Pópulus meus es; et ipse dicet~: Dóminus Deus meus. Hæc dicit Dóminus.

\begin{tightlist}{}
\item \textred{℟.}~In die illa erit fons patens dómui David, et habitántibus Jerúsalem \textred{*} In ablutiónem peccatóris, et menstruátæ. \textred{℣.}~Hic est, qui venit per aquam, et sánguinem Jesus Christus. \textred{I}n. Glória Patri. \textred{I}n.
\end{tightlist}

\subsubsection*{\center \textsc{IN SECUNDO NOCTURNO.}}
\begin{tightlist}{}
\item \textred{Aña.} Adspícient ad me~: quem confixérunt, et plangent eum planctu quasi super Unigénitum.
%\item \textred{Ps. 4.} Cum invocárem. [79]
\item \textred{Aña.} Quid sunt plagæ istæ in médio mánuum tuárum?
%\item \textred{Ps. 5.} Verba mea.
\item \textred{Aña.} His plagátus sum in domo eórum, qui diligébant me.
%\item \textred{Ps. 8.} Dómine Dóminus.
\item \textred{℣.}~Quem tu percussísti persecúti sunt. Et super dolórem vúlnerum meórum addidérunt.
\end{tightlist}

\subsubsection*{\center Sermo Innocéntii Papæ Sexti.}
\subsubsection*{\textred{Decret. de Fest. Lanc. et Clav. Dom. \hfill Lectio~iv.}}
\lettrineThree{I}n Redemptóris nostri Dómini Jesu Christi sacratíssima passióne sic nos gloriári opórtet, ut ipsíus passiónis cuncta mystéria dinumerántes et mérita, síngulis étiam ejus salutáribus instruméntis gloriémur. Inter quæ illud celébriter memorándum est, quod ipse Salvátor, emísso in cruce jam spíritu, sustínuit perforári Láncea laus suum, ut, inde Sánguinis et aquæ profluéntibus undis, formarétur única, et immaculáta, ac virgo sancta mater Ecclésia sponsa sua. O beatíssima ipsíus sacri láteris apertúra, unde nobis tot, et tanta divínæ pietátis dona fluxérunt! O felix Láncea, quæ tot bona nobis effícere, et ad tanti triúmphi glóriam méruit superáddi.

\begin{tightlist}{}
\item \textred{℟.}~Hic est, qui venit per aquam, et sánguinem Jesus Christus~: \textred{*} Non in aqua solum, sed in aqua et sánguine. \textred{℣.}~Tres sunt, qui testimónium dant in terra, spíritus, aqua, et sanguis, et hi tres unum sunt. \textred{N}on.
\end{tightlist}

\subsubsection*{\center \textred{Lectio v.}}
\lettrine{H}æc latus ipsum aperiéndo sacratíssimas jánuas nobis regni cœléstis apéruit. Hæc vúlnera nostra sanávit, vitámque nobis réddidit et salútem. Hæc innóxium transfigéndo, illíus sánguine culpas nostras abstérsit~: et demum ejúsdem sacratíssimis undis irrigáta, cæcitátis nostræ ténebras sústulit, et nos ipsíus divínæ pietátis flúviis mundávit. Illi étiam dulces Clavi, quibus ipse Salvátor eídem cruci fuit affíxus, quique non solum immaculáto respérgi sánguine, et molem ferre tanti pónderis meruérunt, sed et nos étiam per ipsórum salutíferas plagas dulcédinem tantam ipsíus divínæ caritátis accépimus, ut manus nostræ a peccáti solútæ néxibus, pedésque nostri a mortis láqueis fúerint liberáti, sunt devotíssime recoléndi.

\begin{tightlist}{}
\item \textred{℟.}~Quóniam circumdedérunt me canes multi, concílium malignántium obsédit me~: \textred{*} Fodérunt manus meas, et pedes meos. \textred{℣.}~Dicétur ei~: Quid sunt plagæ istæ in médio mánuum tuárum? \textred{F}odérunt.
\end{tightlist}

\subsubsection*{\center \textred{Lectio vj.}}
\lettrine{Q}uid enim vúlnere, et plagis hujúsmodi sánctius? Quid eis salúbrius, ex quibus salus nostra procéssit, et in quibus assídue se possunt curáre ánimæ devotórum? Licet ígitur Láncea, et Clavi prædícti, alíaque ipsíus passiónis salutária instruménta sint a cunctis fidélibus Christi ubílibet veneránda, et de ipsa étiam passióne in eádem Ecclésia solémnia annis síngulis Offícia celebréntur, et fiant; dignum tamen, et convéniens reputámus, si de ipsíus passiónis speciálibus instruméntis, et præsértim in pártibus, in quibus instruménta ipsa dicúntur habéri, solémne, ac speciále festum celebrétur, et fiat~: Nosque illos fidéles Christi, qui áliqua ex instruméntis ipsíus habére se gaudent in eórum devotióne divínis offíciis, atque munéribus speciáliter foveámus.

\begin{tightlist}{}
\item \textred{℟.}~Dicétur ei~: Quid sunt plagæ istæ in médio mánuum tuárum? Et dicet~: \textred{*} His plagátus sum in domo eórum, qui diligébant me. \textred{℣.}~Nisi vídero in mánibus ejus fixúram clavórum, non credam. \textred{H}is. Glória Patri. \textred{H}is.
\end{tightlist}

\subsubsection*{\center \textsc{IN TERTIO NOCTURNO.}}
\begin{tightlist}{}
\item \textred{Aña.} Fodérunt manus meas et pedes meos, dinumeravérunt ómnia ossa mea.
\item \textred{Ps. 95} Cantáte. [179]
\item \textred{Aña.} Nisi vídero in mánibus ejus fixúram clavórum, et mittam dígitum meum in locum clavórum et mittam manum meam in latus ejus, non credam.
\item \textred{Ps. 96.} Dóminus regnávit. [180]
\item \textred{Aña.} Infer dígitum tuum huc, et vide manus meas, et affer manum tuam, et mitte in latus meum.
\item \textred{Ps. 97.} Cantáte. [180]
\item \textred{℣.}~Ipse vulnerátus est propter iniquitátes nostras. \textred{℟.}~Attrítus est propter scélera nostra.
\end{tightlist}

\subsubsection*{\center Léctio sancti Evangélii secúndum Joánnem.}
\subsubsection*{\center \textred{Lectio vij. \hfill Cap. 19.}}
\lettrine{I}n illo témpore~: Sciens Jesus, quia ómnia consummáta sunt, ut consummarétur Scriptúra, dixit~: Sítio. Et réliqua.

\begin{center}Homilía sancti Augustíni Epíscopi.\\ \textred{Tract.~120.~in~Joann.}\end{center}
\lettrineThree{U}nus mílitum Láncea latus ejus apéruit, et contínuo exívit sanguis, et aqua. Vigilánti verbo Evangelísta usus est, ut non díceret, latus ejus percússit, aut vulnerávit, aut quid áliud, sed apéruit~: ut illic quodámmodo vitæ óstium panderétur, unde Sacraménta Ecclésiæ manavérunt, sine quibus ad vitam, quæ vera vita est, non intrátur. Ille sanguis, qui fusus est, in remissiónem fusus est peccatórum. Aqua illa salutáre témperat póculum~: hæc et lavácrum præstat, et potum. Hoc prænuntiábat, quod Nœ in látere arcæ óstium fácere jussus est, qua intrárent animália, quæ non erant dilúvio peritúra, quibus præfigurabátur Ecclésia.

\begin{tightlist}{}
\item \textred{℟.}~Facta sunt hæc, ut Scriptúra implerétur~: Os non comminuétis ex eo. Et íterum ália Scriptúra dicit~: \textred{*} Vidébunt in quem transfixérunt. \textred{℣.}~Effúndam super domum David, et super habitatóres Jerúsalem spíritum grátiæ, et precum. \textred{V}idébunt.
\end{tightlist}

\subsubsection*{\center \textred{Lectio viij.}}
\lettrine{P}ropter hoc prima múlier facta est de látere viri dormiéntis, et appelláta est vita, matérque vivórum. Magnum quippe significávit bonum, ante magnum prævaricatiónis malum. Hic secúndus Adam, inclináto cápite, in cruce dormívit; ut inde formarétur ei conjux, quæ de látere dormiéntis efflúxit. O mors, unde mórtui revivíscunt! Quid isto sánguine múndius? Quid vúlnere isto salúbrius? Facta sunt enim hæc, inquit, ut Scriptúra implerétur~: Os non comminuétis ex eo. Et íterum ália Scriptúra dicit~: Vidébunt in quem transfixérunt.

\begin{tightlist}{}
\item \textred{℟.}~Effúndam super domum David, et super habitatóres Jerúsalem spíritum grátiæ, et precum~: \textred{*} Et adspícient ad me, quem confixérunt. \textred{℣.}~Et plangent eum planctu quasi super unigénitum. \textred{E}t adspícient.
\item \textred{Lectio ix. legitur de Homilia Feriæ occurr.}
\item \textred{℟. ix.} Adorámus te, Christe, et benedícimus tibi, quia per sanctam Crucem tuam redemísti mundum; \textred{*} Miserére nostri, qui passus es pro nobis. \textred{℣.}~In ligno pendens, nostræ salútis viam Verbum Patris invénit, in qua Redémptor noster, mortis víncula rumpens, cállidam áspidem superávit. \textred{M}iserére. Glória Patri. \textred{M}iserére.
\item Te Deum laudámus. 6.
\item \textred{Sacerd. ℣.} Adorámus te Christe, et benedícimus tibi. \textred{℟.}~Qui vulnerátus es propter iniquitátes nostras.
\end{tightlist}

\subsubsection*{\center \textsc{AD LAUDES.}}
\begin{center}\textred{et per Horas, Aña.}\end{center}
\lettrine{U}nus mílitum Láncea latus ejus apéruit, et contínuo exívit sanguis et aqua.

\begin{tightlist}{}
\item \textred{Ps. de Dominica.} 27.
\item \textred{Aña.} Fodérunt manus meas, et pedes meos~: dinumeravérunt ómnia ossa mea.
\item \textred{Aña.} Tres sunt, qui testimónium dant in terra~: spíritus, aqua, et sanguis; et hi tres unum sunt.
\item \textred{Aña.} Quid turbáti estis, et cogitatiónes ascéndunt in corda vestra? Vidéte manus meas, et pedes meos, quia ego ipse sum.
\item \textred{Aña.} Infer dígitum tuum huc et vide manus meas, et affer manum tuam, et mitte in latus meum.
\end{tightlist}

\subsubsection*{\textred{Capitulum. \hfill 1. Joan. 5.}}
\lettrineThree{C}aríssimi~: Quis est qui vincit mundum, nisi qui credit quóniam Jesus est Fílius Dei? Hic est qui venit per aquam et sánguinem, Jesus Christus, non in aqua solum, sed in aqua et sánguine.

\subsubsection*{\center \textred{Hymnus.}}
\lettrine{V}inctam ergo Christi sánguine\\
Convértite in me cúspidem~:\\
Feríte cor, pedes, manus,\\
Pœnam a nocénte súmite.\\
\\
\textred{A}t quæso culpis débitas,\\
Quas jure plagas fígitis,\\
Cruóre divíno íllitæ\\
Fiant medéla spíritus.\\
\\
\textred{S}it gressus ad malum ímpotens,\\
Manus nocére désinant,\\
Omnísque corde e sáucio\\
Profánus ardor éxeat.\\
\\
\textred{C}lavis foráto et Láncea,\\
Jesu, tibi sit glória,\\
Cum Patre et almo Spíritu\\
Nunc et per omne sæculum.\\
Amen.

\begin{tightlist}{}
\item \textred{℣.}~Disciplína pacis nostræ super eum. \textred{℟.}~Cujus livóre sanáti sumus.
\item \textred{Ad Benedictus, Aña.} Inglórius erit inter viros adspéctus ejus, et forma ejus inter fílios hóminum. Iste aspérget gentes multas.
\end{tightlist}

\subsubsection*{\center \textred{Oratio.}}
\lettrine{D}eus qui in assúmptæ carnis infirmitáte Clavis affígi, et Láncea vulnerári pro mundi salúte voluísti~: concéde propítius, ut qui eorúmdem Clavórum, et Lánceæ solémnia venerámur in terris, de glorióso victóriæ tuæ triúmpho gratulémur in cœlis. Qui vivis.

\subsubsection*{\center \textsc{AD PRIMAM.}}
\noindent \textred{℣.}~Qui nos salvásti in Cruce.

\subsubsection*{\center \textsc{AD TERTIAM.}}
\subsubsection*{\textred{Capitulum. \hfill 1. Joan. 5.}}
\lettrine{C}aríssimi~: Quis est qui vincit mundum, nisi qui credit quóniam Jesus est Fílius Dei? Hic est qui venit per aquam et sánguinem, Jesus Christus, non in aqua solum, sed in aqua et sánguine.

\begin{tightlist}{}
%\item \textred{Capitulum.} Caríssimi, \textred{ut supra.}
\item \textred{℟.}~Fodérunt manus meas, \textred{*} Et pedes meos. \textred{F}odérunt. \textred{℣.}~Dinumeravérunt ómnia ossa mea. \textred{E}t pedes. Glória Patri. \textred{F}odérunt.
\item \textred{℣.}~Quem tu percussísti, persecúti sunt. \textred{℟.}~Et super dolórem vúlnerum meórum addidérunt.
\end{tightlist}

\subsubsection*{\center \textsc{AD SEXTAM.}}
\subsubsection*{\textred{Capitulum. \hfill 1. Petr. 2.}}
\lettrine{F}ratres, Christus passus est pro nobis, vobis relínquens exémplum, ut sequámini vestígia ejus. Qui peccátum non fecit, nec invéntus est dolus in ore ejus.

\begin{tightlist}{}
\item \textred{℟.}~Quem tu percussísti, \textred{*} Persecúti sunt. \textred{Q}uem. \textred{℣.}~Et super dolórem vúlnerum meórum addidérunt. \textred{P}ersecúti. Glória Patri. \textred{Q}uem.
\item \textred{℣.}~Ipse vulnerátus est propter iniquitátes nostras. \textred{℟.}~Attrítus est propter scélera nostra.
\end{tightlist}

\subsubsection*{\center \textsc{AD NONAM.}}
\subsubsection*{\textred{Capitulum. \hfill 1. Petr. 2.}}
\lettrine{Q}ui peccáta nostra ipse pértulit in córpore suo super lignum, ut peccátis mórtui justítiæ vivámus; cujus livóre sanáti estis.

\begin{tightlist}{}
\item \textred{℟.}~Ipse vulnerátus est \textred{*} Propter iniquitátes nostras. \textred{I}pse. \textred{℣.}~Attrítus est propter scélera nostra. \textred{P}ropter. Glória Patri. \textred{I}pse.
\item \textred{℣.}~Disciplína pacis nostræ super eum. \textred{℟.}~Cujus livóre sanáti sumus.
\end{tightlist}

\subsubsection*{\center \textsc{AD VESPERAS.}}
\begin{tightlist}{}
\item \textred{Añæ Laudum, Psalmi de Dom. sed loco ultimi, Ps. 115.} Crédidi. 91.
\end{tightlist}

\subsubsection*{\textred{Capitulum. \hfill 1. Joan. 5.}}
\lettrineThree{C}aríssimi~: Quis est, qui vincit mundum, nisi qui credit, quóniam Jesus est fílius Dei? Hic est, qui venit per aquam, et sánguinem, Jesus Christus, non in aqua solum, sed in aqua, et sánguine.

\subsubsection*{\center \textred{Hymnus.}}
\lettrine{Q}uænam lingua tibi, o Láncea, débitas\\
Grates pro mérito est apta repéndere?\\
Christi vivíficum namque áperis latus,\\
Unde Ecclésia náscitur.\\
\\
\textred{H}æc est Heva viri de látere éxiens,\\
Olli membra gravis dum sopor óccupat;\\
Hanc quippe alter Adam, corde scaténtibus\\
Unda et sánguine, prócreat.\\
\\
\textred{O} Clavi! æqua manet vos quoque grátia,\\
Christi quando sacris ártubus ínsiti,\\
Delétum Dómini sánguine fígitis\\
Mortis chirógraphum cruci.\\
\\
\textred{T}e, Jesu, Súperi láudibus éfferant,\\
Qui Clavórum áditus, sígnaque Lánceæ\\
In cœlo rétines, vivus ubi ímperas\\
Cum Patre atque Paráclito.\\
Amen.

\begin{tightlist}{}
\item \textred{℣.}~Fodérunt manus meas et pedes meos. \textred{℟.}~Dinumeravérunt ómnia ossa mea.
%\item \textred{Capit. Hymn. et ℣.ut in j. Vesp.}
\item \textred{Ad Magnif. Aña.} Vere languóres nostros ipse tulit, et dolóres nostros ipse portávit, et nos putávimus eum quasi leprósum, et percússum a Deo, et humiliátum.
\end{tightlist}

\begin{center}\textcolor{red}{\rule{\linewidth}{0.5pt}} \end{center}

\subsection*{\center \textsc{FERIA SEXTA}\\
\textred{post Domin. II. Quadrag.}}
\subsection*{\center \textsc{SS. Sindonis D.N.J.C.}}
\subsubsection*{\center \textred{Duplex minus. 1 classis.}}
%\begin{center}\textcolor{red}{\rule{\linewidth}{0.5pt}} \end{center}
\subsubsection*{\center \textsc{AD VESPERAS.}}
\fancyhead[C]{\textsc{SS. Sindonis D.N.J.C.}}
    \begin{tightlist}{}
\item \textred{Aña.} Joseph nóbilis decúrio, vir bonus et justus, et ipse dives, erat exspéctans regnum Dei.
\item \textred{Ps. de Officio præc.}
\end{tightlist}

\subsubsection*{\textred{Capitulum. \hfill Isaiæ 63.}}
\lettrineThree{Q}uis est iste, qui venit de Edom tinctis véstibus de Bosra? Iste formósus in stola sua grádiens in multitúdine fortitúdinis suæ? Ego, qui loquor justítiam, et propugnátor sum ad salvándum.

\begin{tightlist}{}
\item \textred{℟.}~Te, sacra Sindon, devóte venerámur, qua Salvátor noster de cruce depósitus invólvi dignátus est, \textred{*} Quæ Viri dolórum refers imáginem, ipsíus cruóre delineátam. \textred{℣.}~Salve Sindon pretiosíssima, ex membris Dómini tamquam margarítis ornáta. \textred{Q}uæ. Glória Patri. \textred{Q}uæ.
\end{tightlist}

\subsubsection*{\center \textred{Hymnus.}}
\lettrine{G}lóriam sacræ celebrémus omnes\\
Síndonis~: lætis recolámus hymnis,\\
Et piis votis monuménta nostræ\\
Certa salútis.\\
\\
\textred{Q}uæ refert semper veneránda Sindon\\
Sánguine impréssis decoráta signis,\\
Dum cruce ex alta tulit involútum\\
Corpus Iésu.\\
\\
\textred{R}eddit hæc sævos ánimo dolóres,\\
Quos tulit, casum miserátus Adæ,\\
Christus humáni géneris Redémptor\\
Morte perémpta.\\
\\
\textred{S}áucium ferro latus, atque palmas\\
Et pedes clavis, laceráta flagris\\
Membra, et infíxam cápiti corónam\\
Monstrat imágo.\\
\\
\textred{Q}uis pius siccis óculis, et absque\\
Intimi cordis gémitu, notáta\\
Vívaque indígnæ simulácra mortis\\
Cérnere possit?\\
\\
\textred{N}ostra cum solum tibi, Christe, culpa\\
Causa tantórum fúerit malórum\\
Nostra debétur tibi vita~: vitam\\
Dédimus ipsam.\\
\\
\textred{S}it tibi, Fili, decus atque virtus,\\
Qui tuo mundum rédimis cruóre;\\
Quique cum summo Genitóre, et almo\\
Flámine regnas.\\
Amen.

\begin{tightlist}{}
\item \textred{℣.}~Tuam Síndonem venerámur, Dómine. \textred{℟.}~Tuam recólimus passiónem.
\item \textred{Ad Magnificat, Aña.} Joseph vir bonus, et justus accéssit ad Pilátum, et pétiit corpus Jesu, quo accépto, invólvit illud in Síndone munda.
\end{tightlist}

\subsubsection*{\center \textred{Oratio.}}
\lettrine{D}eus, qui nobis in sancta Síndone, qua corpus tuum sacratíssimum e cruce depósitum a Joseph involútum fuit, passiónis tuæ vestígia reliquísti~: concéde propítius, ut per mortem et sepultúram tuam, ad resurrectiónis glóriam perducámur. Qui vivis et regnas.

\begin{tightlist}{}
\item \textred{In fine Hymnorum Completorii et Horarum dicitur ℣.} Glória tibi Dómine, qui nos salvásti in Cruce. 
\item \textred{Ad Nunc dimittis, Aña.} Hoc signum Crucis * erit in cælo, cum Dóminus ad judicándum vénerit.
\end{tightlist}

\subsubsection*{\center \textsc{AD MATUTINUM.}}
\begin{tightlist}{}
\item \textred{Invitat.} Christum Dóminum, qui passiónis suæ memóriam in sacra Síndone rénovat, \textred{*} Veníte adorémus.
\item \textred{Ps.} Veníte exsultémus. 2.
\end{tightlist}

\subsubsection*{\center \textred{Hymnus.}}
\lettrine{M}ystérium mirábile\\
Hac luce nobis pánditur,\\
Verus Dei quum Fílius\\
Mortem cruéntam sústinet.\\
\\
\textred{C}ausam tuéndo sérvuli\\
Reíque formam véstiens,\\
Pro servo herus suspénditur\\
Pro sonte justus pléctitur.\\
\\
\textred{N}ecis manent insígnia\\
Impréssa sacra in Síndone,\\
Quæ post triúmphum nóbilem\\
Corpus cruéntum invólverat.\\
\\
\textred{S}unt mortis hæc, et tártari,\\
Mundíque victi insígnia,\\
Trophæa sunt hæc ínclyta\\
Ductóris invictíssimi.\\
\\
\textred{D}ebémus ergo hanc grátiam\\
Nostræ salútis Víndici,\\
Ut dæmonis contra dolos\\
Hac militémus téssera.\\
\\
\textred{V}itæ vetústæ mórtui\\
Surgámus in vitam novam;\\
Christum secúti per crucem\\
Christi fruémur glória.\\
\\
\textred{G}lória tibi, Dómine,\\
Qui nos salvásti in cruce\\
Cum Patre, et sancto Spíritu\\
In sempitérna sæcula\\
Amen.

\subsubsection*{\center \textsc{IN PRIMO NOCTURNO.}}
\begin{tightlist}{}
\item \textred{Aña.} Rubrum est induméntum tuum, et vestiménta tua sicut calcántium in torculári.
\item \textred{Ps. 4.} Cum invocárem. 56.
\item \textred{Aña.} Aspérsus est sanguis super vestiménta mea; ómnia induménta mea inquinávi.
\item \textred{Ps. 15.} Consérva me. 132.
\item \textred{Aña.} Divisérunt sibi vestiménta mea, et super vestem meam misérunt sortem.
\item \textred{Ps. 14.} Dómine quis habitábit. 65.
\item \textred{℣.}~Tuam Síndonem venerámur, Dómine. \textred{℟.}~Tuam recólimus passiónem.
\end{tightlist}

\subsubsection*{\center De Isaía Prophéta.}
\begin{center}\textred{Lectio j. \hfill Cap. 53.}\end{center}
\lettrineThree{Q}uis crédidit audítui nostro? et bráchium Dómini cui revelátum est?
Et ascéndet sicut virgúltum coram eo, et sicut radix de terra sitiénti~: non est spécies ei, neque decor~: et vídimus eum, et non erat aspéctus, et desiderávimus eum:
despéctum, et novíssimum virórum, virum dolórum, et sciéntem infirmitátem~: et quasi abscónditus vultus ejus et despéctus, unde nec reputávimus eum.
Vere languóres nostros ipse tulit, et dolóres nostros ipse portávit~: et nos putávimus eum quasi leprósum, et percússum a Deo et humiliátum.
Ipse autem vulnerátus est propter iniquitátes nostras~: attrítus est propter scélera nostra~: disciplína pacis nostræ super eum, et livóre ejus sanáti sumus. Hæc dicit Dóminus.

\begin{tightlist}{}
\item \textred{℟.}~Ecce vídimus eum non habéntem spéciem, neque decórem~: adspéctus ejus in eo non est; Hic peccáta nostra portávit, et pro nobis dolet, ipse autem vulnerátus est propter iniquitátes nostras~: \textred{*} Cujus livóre sanáti sumus. \textred{℣.}~Vere languóres nostros ipse tulit, et dolóres nostros ipse portávit. \textred{C}ujus.
\end{tightlist}

\subsubsection*{\center \textred{Lectio ij.}}
\lettrine{O}mnes nos quasi oves errávimus, unusquísque in viam suam declinávit~: et pósuit Dóminus in eo iniquitátem ómnium nostrum.
Oblátus est quia ipse vóluit, et non apéruit os suum~: sicut ovis ad occisiónem ducétur, et quasi agnus coram tondénte se obmutéscet, et non apériet os suum.
De angústia, et de judício sublátus est~: generatiónem ejus quis enarrábit? quia abscíssus est de terra vivéntium~: propter scelus pópuli mei percússi eum.
Et dabit ímpios pro sepultúra, et dívitem pro morte sua, eo quod iniquitátem non fécerit, neque dolus fúerit in ore ejus. Hæc dicit Dóminus.

\begin{tightlist}{}
\item \textred{℟.}~Quid sunt plagæ istæ in médio mánuum tuárum? \textred{*} Et dicet; his plagátus sum in domo eórum qui diligébant me. \textred{℣.}~Frámea suscitáre super pastórem meum, et super virum cohæréntem mihi. \textred{E}t dicet.
\end{tightlist}

\subsubsection*{\center \textred{Lectio iij.}}
\lettrine{E}t Dóminus vóluit contérere eum in infirmitáte. Si posúerit pro peccáto ánimam suam, vidébit semen longǽvum, et volúntas Dómini in manu ejus dirigétur.
Pro eo quod laborávit ánima ejus, vidébit et saturábitur. In sciéntia sua justificábit ipse justus servus meus multos, et iniquitátes eórum ipse portábit.
Ideo dispértiam ei plúrimos, et fórtium dívidet spólia, pro eo quod trádidit in mortem ánimam suam, et cum scelerátis reputátus est, et ipse peccáta multórum tulit, et pro transgressóribus rogávit. Hæc dicit Dóminus.

\begin{tightlist}{}
\item \textred{℟.}~Tuam Síndonem venerámur Dómine, tuam gloriósam recólimus passiónem~: \textred{*} Miserére nobis qui passus es pro nobis. \textred{℣.}~Salva præséntem catérvam in tuis hódie láudibus congregátam. \textred{M}iserére. Glória Patri. \textred{M}iserére.
\end{tightlist}

\subsubsection*{\center \textsc{IN SECUNDO NOCTURNO.}}
\begin{tightlist}{}
\item \textred{Aña.} Quasi abscónditus vultus ejus, et despéctus~: unde nec reputávimus eum.
\item \textred{Ps. 23.} Dómini est terra. 79.
\item \textred{Aña.} Omnes vidéntes me derisérunt me, locúti sunt lábiis, et movérunt caput.
\item \textred{Ps. 26.} Dóminus illuminátio. 84.
\item \textred{Aña.} Trádidit in mortem ánimam suam, et ipse peccáta multórum tulit.
\item \textred{Ps. 29.} Exaltábo te Dómine. 70.
\item \textred{℣.}~Adorámus te, Christe, et benedícimus tibi. \textred{℟.}~Quia per mortem tuam redemísti mundum.
\end{tightlist}

\subsubsection*{\center Sermo sancti Ambrósii Epíscopi.}
\subsubsection*{\center \textred{Lectio iv. \hfill In Luc. c. 23.}}
\lettrineThree{Q}uid sibi vult, quod non Apóstoli, sed Joseph et Nicodémus, ut Joánnes dicit, Christum sepéliunt? Unus justus et constans, alter, qui erat magíster in Israel. Talis enim Christi est sepultúra quæ justítiam, magisteriúmque hábeat. Obstrúitur ígitur tergiversándi locus, et doméstico Judæi testimónio revincúntur. Nam si Apóstoli sepelíssent, dícerent útique non sepúltum, quem sepúltum raptum esse dixérunt. Justus autem Christi corpus óperit Síndone~: ínnocens ungit unguénto. Namque hæc non otióse distíncta reperímus, quia justítia vestit Ecclésiam, innocéntia grátiam subminístrat.

\begin{tightlist}{}
\item \textred{℟.}~O admirábilis Sindon! \textred{*} In qua involútus est thesáurus noster, redémptio captivórum. \textred{℣.}~Gaudet totus mundus, qui redémptus est Sánguine Dómini sui. \textred{I}n qua.
\end{tightlist}

\subsubsection*{\center \textred{Lectio v.}}
\lettrine{V}esti ergo et tu Dómini corpus glória sua, ut et ipse sis justus. Etsi mórtuum credis, óperi tamen divinitátis própriæ plenitúdine. Unge illud myrrha et áloe, ut bonus odor Christi sis. Bonum línteum misit Joseph ille vir justus, et fortásse illud, quod Petrus vidit e cœlo ad se esse demíssum, in quo erant génera quadrúpedum, et ferárum, et vólucrum ad similitúdinem gentílium figuráta. Mystico ígitur unguénto illo pístico consepelítur Ecclésia, quæ diversitátem populórum fídei suæ communióne sociávit.

\begin{tightlist}{}
\item \textred{℟.}~Tulérunt túnicam Joseph fratres ejus, et in sánguine hœdi, quem occíderant, tinxérunt, mitténtes qui ferrent ad patrem, et dícerent~: \textred{*} Hanc invénimus, vide num túnica fílii tui sit an non. \textred{℣.}~Quam cum agnovísset Pater, ait~: Túnica fílii mei est~: fera péssima devorávit eum. \textred{H}anc.
\end{tightlist}

\subsubsection*{\center \textred{Lectio vj.}}
\lettrine{H}unc Joseph justum Lucas dixit, Matthæus dívitem, et mérito dives hoc loco dícitur, ubi corpus súscipit Christi~: suscipiéndo enim dívitem, nescívit fídei paupertátem. Dives est ergo, qui justus est~: justus ígitur Síndone invólvit; Israelíta vero et divérsos miscet virtútum odóres, et áloes mittit quasi libras centum, hoc est perféctæ fídei quantitátem. Et ligavérunt corpus Jesu juxta consuetúdinem speciálem Judæórum, non útique notis perfídiæ, sed fídei ligatúris~: et posuérunt in horto, cui frequénter Ecclésia comparátur, quæ diversórum habet poma meritórum, florésque virtútum.

\begin{tightlist}{}
\item \textred{℟.}~Christus passus est pro nobis, vobis relínquens exémplum, ut sequámini vestígia ejus, \textred{*} Qui peccátum non fecit, nec invéntus est dolus in ore ejus. \textred{℣.}~Qui cum maledicerétur non maledicébat, cum paterétur non comminabátur. \textred{Q}ui peccátum. Glória Patri. \textred{Q}ui.
\end{tightlist}

\subsubsection*{\center \textsc{IN TERTIO NOCTURNO.}}
\begin{tightlist}{}
\item \textred{Aña.} Caro mea requiéscet in spe, quia non dabis sanctum tuum vidére corruptiónem.
\item \textred{Ps. 53.} Deus in nómine. 160.
\item \textred{Aña.} Convertísti planctum meum in gáudium mihi, conscidísti saccum meum, et circumdedísti me lætítia.
\item \textred{Ps. 75.} Notus in Judæa. 202.
\item \textred{Aña.} Factus sum sicut homo sine adjutório inter mórtuos liber.
\item \textred{Ps. 87.} Dómine Deus. 287.
\item \textred{℣.}~Omnis terra adóret te, et psallat tibi. \textred{℟.}~Psalmum dicat nómini tuo, Dómine.
\end{tightlist}

\subsubsection*{\center Léctio sancti Evangélii secúndum Marcum.}
\subsubsection*{\center \textred{Lectio vij. \hfill Cap. 15.}}
\lettrine{I}n illo témpore~: Cum jam sero esset factum, quia erat parascéve, quod est ante Sábbatum, venit Joseph ab Arimathæa nóbilis decúrio, qui et ipse erat exspéctans regnum Dei. Et réliqua.

\begin{center}Homilía Venerábilis Bedæ Presbyteri.\\
\textred{In Evang. Marci.}\end{center}
\lettrineThree{P}arascéve græce, latíne præparátio dícitur, quo nómine Judæi, qui inter Græcos morabántur, sextam Sábbati appellábant, eo quod in illa, ea, quæ requiéi Sábbati necessária essent, præparáre solérent~: juxta hoc quod de manna quondam præcéptum est~: Sexta autem die colligétis duplum. Quia ergo sexta die factus est homo, et tota est mundi creatúra perfécta; séptima autem Cónditor ab ópere suo requiévit, unde et hanc, Sábbatum, hoc est réquiem vóluit appellári~: recte Salvátor eádem sexta die crucifíxus humánæ restauratiónis implévit arcánum~: ideóque cum accepísset acétum, dixit~: Consummátum est~: hoc est sextæ diéi quod pro mundi refectióne suscépi, jam totum est opus explétum Sábbato autem in sepúlcro requiéscens, resurrectiónis, quæ octáva die ventúra erat, exspectábat evéntum.

\begin{tightlist}{}
\item \textred{℟.}~Mercátus est Joseph Síndonem mundam, ut invólveret in ea corpus Dómini~: \textred{*} Venit ergo, et tulit corpus ejus. \textred{℣.}~Rogávit Pilátum Joseph ab Arimathæa, ut tóllerat corpus Jesu. \textred{V}enit.
\end{tightlist}

\subsubsection*{\center \textred{Lectio viij.}}
\lettrine{V}enit, inquit, Joseph ab Arimathæa nóbilis decúrio, qui et ipse erat exspéctans regnum Dei, et audácter introívit ad Pilátum, et pétiit corpus Jesu. Magnæ quidem Joseph iste dignitátis ad sæculum, sed majóris apud Deum mériti fuísse laudátur. Talem namque exsístere decébat eum, qui corpus Dómini sepelíret; qui et per justítiam meritórum, tali ministério dignus esset, et per nobilitátem poténtiæ sæculáris facultátem posset obtinére ministrándi~: non enim quílibet ignótus, aut medíocris ad præsidem accédere, et crucifíxi corpus póterat impetráre.

\begin{tightlist}{}
\item \textred{℟.}~Hæc est Sindon digníssima, \textred{*} In qua salútis Auctor de cruce depósitus invólvi dignátus est. \textred{℣.}~Ut nos consepúlti cum eo, véteris hóminis exúviis depósitis, innocéntiæ Síndone tegerémur. \textred{I}n qua.
\end{tightlist}

\begin{tightlist}{}
\item \textred{Lectio ix. de Homilia Feriæ occurrentis.}
\item \textred{℟.}~Te sacra Sindon, devóte venerámur, qua Salvátor noster de cruce depósitus invólvi dignátus est, \textred{*} Quæ Viri dolórum refers imáginem, ipsíus cruóre delineátam. \textred{℣.}~Salve Sindon pretiosíssima, ex membris Dómini tamquam margarítis ornáta. \textred{Q}uæ Viri. Glória Patri. \textred{Q}uæ Viri.
\item Te Deum laudámus. 6.
\item \textred{Sacerd ℣.} Tuam Síndonem venerámur, Dómine. \textred{℟.}~Miserére nobis, qui passus es pro nobis.
\end{tightlist}

\subsubsection*{\center \textsc{AD LAUDES.}}
\begin{center}\textred{er per Horas, Aña.}\end{center}
\lettrine{J}oseph nóbilis decúrio vir bonus, et justus, et ipse dives, erat exspéctans regnum Dei.

\begin{tightlist}{}
\item \textred{Ps. de Dominica.} 27.
\item \textred{Aña.} Hic audácter introívit ad Pilátum, et pétiit corpus Jesu.
\item \textred{Aña.} Cum cognovísset Pilátus a centurióne, quod jam mórtuus esset, donávit corpus Jesu.
\item \textred{Aña.} Joseph autem mercátus Síndonem, et depónens eum, invólvit Síndone.
\item \textred{Aña.} Pósuit eum in monuménto, in quo nondum quisquam pósitus fúerat.
\end{tightlist}

\subsubsection*{\textred{Capitulum. \hfill Isaiæ 63.}}
\lettrineThree{Q}uis est iste qui venit de Edom, tinctis véstibus de Bosra? Iste formósus in stola sua grádiens in multitúdine fortitúdinis suæ? Ego qui loquor justítiam, et propugnátor sum ad salvándum.

\subsubsection*{\center \textred{Hymnus.}}
\lettrine{J}esu dulcis amor meus,\\
Ac si præsens sis, accédo~:\\
Te compléctor cum afféctu\\
Tuórum memor vúlnerum.\\
\\
\textred{O} quam nudum hic te cerno,\\
Vulnerátum et disténtum,\\
Inquinátum, involútum\\
In hoc sacráto tégmine.\\
\\
\textred{S}alve caput cruentátum\\
Spinis, cujus dulcis vultus\\
Immutávit suum florem,\\
Quem cœli tremit cúria.\\
\\
\textred{S}alve latus Salvatóris,\\
Salve mitis apertúra\\
Super rosam rubicúnda,\\
Medéla salutífera.\\
\\
\textred{M}anus sanctæ vos avéte,\\
Diris clavis perforátæ;\\
Ne repéllas me, Salvátor\\
De tuis sanctis pédibus.\\
Amen.

\begin{tightlist}{}
\item \textred{℣.}~Dóminus regnávit, decórem índuit. \textred{℟.}~Induit Dóminus fortitúdinem, et præcínxit se virtúte.
\item \textred{Ad Benedictus, Aña.} Joseph nóbilis decúrio, exspéctans et ipse regnum Dei, mercátus est Síndonem, et depónens corpus Jesu, invólvit illud in Síndone.
\end{tightlist}

\subsubsection*{\center \textred{Oratio.}}
\lettrine{D}eus, qui nobis in sancta Síndone, qua corpus tuum sacratíssimum e cruce depósitum a Joseph involútum fuit, passiónis tuæ vestígia reliquísti~: concéde propítius, ut per mortem et sepultúram tuam, ad resurrectiónis glóriam perducámur. Qui vivis et regnas.

\subsubsection*{\center \textsc{AD PRIMAM.}}
\noindent \textred{℣.}~Qui nos salvásti in Cruce.

\subsubsection*{\center \textsc{AD TERTIAM.}}
\subsubsection*{\textred{Capitulum. \hfill Isaiæ 63.}}
\lettrine{Q}uis est iste qui venit de Edom, tinctis véstibus de Bosra? Iste formósus in stola sua grádiens in multitúdine fortitúdinis suæ? Ego qui loquor justítiam, et propugnátor sum ad salvándum.

\begin{tightlist}{}
%\item \textred{Capit.} Quis est iste.
\item \textred{℟.}~Tuam Síndonem \textred{*} Venerámur, Dómine. \textred{T}uam. \textred{℣.}~Tuam recólimus passiónem. \textred{V}enerámur. Glória Patri. \textred{T}uam.
\item \textred{℣.}~Adorámus te Christe, et benedícimus tibi. \textred{℟.}~Quia per mortem tuam redemísti mundum.
\end{tightlist}

\subsubsection*{\center \textsc{AD SEXTAM.}}
\subsubsection*{\textred{Capitulum. \hfill Isai. 63.}}
\lettrine{Q}uare ergo rubrum est induméntum tuum, et vestiménta tua sicut calcántium in torculári? Torcúlar calcávi solus, et de Géntibus non est vir mecum.

\begin{tightlist}{}
\item \textred{℟.}~Adorámus te Christe, \textred{*} Et benedícimus tibi. \textred{A}dorámus. \textred{℣.}~Quia per mortem tuam redemísti mundum. \textred{E}t benedícimus. Glória Patri. \textred{A}dorámus.
\item \textred{℣.}~Omnis terra adóret te, et psallat tibi. \textred{℟.}~Psalmum dicat nómini tuo, Dómine.
\end{tightlist}

\subsubsection*{\center \textsc{AD NONAM.}}
\subsubsection*{\textred{Capitulum. \hfill Isai. 63.}}
\lettrine{C}ircumspéxi, et non erat auxiliátor, quæsívi, et non fuit qui adjuváret~: et salvávit mihi bráchium meum, et indignátio mea ipsa auxiliáta est mihi.

\begin{tightlist}{}
\item \textred{℟.}~Omnis terra adóret te, \textred{*} Et psallat tibi. \textred{O}mnis. \textred{℣.}~Psalmum dicat nómini tuo, Dómine. \textred{E}t. Glória Patri. \textred{O}mnis.
\item \textred{℣.}~Dóminus regnávit, decorem índuit. \textred{℟.}~Induit Dóminus fortitúdinem, et præcínxit se virtúte.
\end{tightlist}

\subsubsection*{\center \textsc{AD II. VESPERAS.}}
\begin{tightlist}{}
\item \textred{Añæ Laudum, Ps. de Dóm.} 52. \textred{sed loco ultimi, Ps. 141.} Voce mea. 245.
\end{tightlist}

\subsubsection*{\textred{Capitulum. \hfill Isaiæ 63.}}
\lettrineThree{Q}uis est iste, qui venit de Edom tinctis véstibus de Bosra? Iste formósus in stola sua grádiens in multitúdine fortitúdinis suæ? Ego, qui loquor justítiam, et propugnátor sum ad salvándum.

\subsubsection*{\center \textred{Hymnus.}}
\lettrine{G}lóriam sacræ celebrémus omnes\\
Síndonis~: lætis recolámus hymnis,\\
Et piis votis monuménta nostræ\\
Certa salútis.\\
\\
\textred{Q}uæ refert semper veneránda Sindon\\
Sánguine impréssis decoráta signis,\\
Dum cruce ex alta tulit involútum\\
Corpus Iésu.\\
\\
\textred{R}eddit hæc sævos ánimo dolóres,\\
Quos tulit, casum miserátus Adæ,\\
Christus humáni géneris Redémptor\\
Morte perémpta.\\
\\
\textred{S}áucium ferro latus, atque palmas\\
Et pedes clavis, laceráta flagris\\
Membra, et infíxam cápiti corónam\\
Monstrat imágo.\\
\\
\textred{Q}uis pius siccis óculis, et absque\\
Intimi cordis gémitu, notáta\\
Vívaque indígnæ simulácra mortis\\
Cérnere possit?\\
\\
\textred{N}ostra cum solum tibi, Christe, culpa\\
Causa tantórum fúerit malórum\\
Nostra debétur tibi vita~: vitam\\
Dédimus ipsam.\\
\\
\textred{S}it tibi, Fili, decus atque virtus,\\
Qui tuo mundum rédimis cruóre;\\
Quique cum summo Genitóre, et almo\\
Flámine regnas.\\
Amen.

\begin{tightlist}{}
\item \textred{℣.}~Tuam Síndonem venerámur, Dómine. \textred{℟.}~Tuam recólimus passiónem.
%\item \textred{Capit. Hymnus et ℣. ut in primis Vesp.}
\item \textred{Ad Magnif. Aña.} Homo quidam dives ab Arimathæa, nómine Joseph, accépto córpore Jesu invólvit illud in Síndone munda.
\end{tightlist}

\begin{center}\textcolor{red}{\rule{\linewidth}{0.5pt}} \end{center}

\subsection*{\center \textsc{FERIA SEXTA}\\
\textred{post Domin. III. Quadrag.}}
\subsection*{\center \textsc{SS. Quinque Vulnerum D.N.J.C.}}
\subsubsection*{\center \textred{Duplex minus. 1. classis.}}
%\begin{center}\textcolor{red}{\rule{\linewidth}{0.5pt}} \end{center}
\subsubsection*{\center \textsc{AD VESPERAS.}}
\fancyhead[C]{\textsc{SS. Quinque Vulnerum D.N.J.C.}}
\begin{tightlist}{}
\item \textred{Aña.} Vere languóres nostros ipse tulit, et dolóres nostros ipse portávit.
\item \textred{Ps. de Officio præc.}
\end{tightlist}

\subsubsection*{\textred{Capitulum. \hfill Isaiæ 53.}}
\lettrineThree{C}hristus vulnerátus est propter iniquitátes nostras, attrítus est propter scélera nostra~: cujus livóre sanáti sumus.

\begin{tightlist}{}
\item \textred{℟.}~Omnes amíci mei dereliquérunt me, et prævaluérunt insidiántes mihi; \textred{*} Et terribílibus óculis plaga crudéli percussérunt me. \textred{℣.}~Dedérunt in escam meam fel, et potavérunt me acéto. \textred{E}t terribílibus. Glória Patri. \textred{E}t.
\end{tightlist}

\subsubsection*{\center \textred{Hymnus.}}
\lettrine{P}ange, lingua, gloriósi\\
Prǽlium certáminis,\\
Et super Crucis trophǽum\\
Dic triúmphum nóbilem~:\\
Quáliter Redémptor orbis\\
Immolátus vícerit.\\
\\
\textred{D}e paréntis protoplásti\\
Fraude Factor cóndolens,\\
Quando pomi noxiális\\
Morsu in mortem córruit~:\\
Ipse lignum tunc notávit,\\
Damna ligni ut sólveret.\\
\\
\textred{H}oc opus nostræ salútis\\
Ordo depopóscerat;\\
Multifórmis proditóris\\
Ars ut artem fálleret,\\
Et medélam ferret inde,\\
Hostis unde lǽserat.\\
\\
\textred{Q}uando venit ergo sacri\\
Plenitúdo témporis,\\
Missus est ab arce Patris\\
Natus, orbis Cónditor~:\\
Atque ventre virgináli\\
Caro factus pródiit.\\
\\
\textred{V}agit infans inter arcta\\
Cónditus præsépia~:\\
Membra pannis involúta\\
Virgo mater álligat,\\
Et manus, pedésque et crura\\
Stricta cingit fáscia.\\
\\
\textred{G}lória et honor Deo\\
Usquequáque altíssimo,\\
Una Patri, Filióque,\\
Inclyto Paráclito~:\\
Cui laus est, et potéstas\\
Per ætérna sǽcula.\\
Amen.
%\item \textred{Hymn.} Pange lingua gloriósi, \textred{ut in Dom. Passionis.} 387.
\begin{tightlist}{}
\item \textred{℣.}~Fodérunt manus meas et pedes meos. \textred{℟.}~Dinumeravérunt ómnia ossa mea.
\item \textred{Ad Magnificat, Aña.} Dum in Cruce pendéret Unigénitus Dei Fílius, et ab ómnibus subsannarétur, Mater ejus Virgo María ipsum verum Deum et hóminem cóndolens venerabátur.
\end{tightlist}

\subsubsection*{\center \textred{Oratio.}}
\lettrine{D}eus, qui Unigéniti Fílii tui passióne, et per quinque Vúlnera ejus Sánguinis effusióne humánam natúram peccáto pérditam reparásti~: tríbue nobis, quæsumus; ut qui ab eo suscépta Vúlnera venerámur in terris, ejúsdem pretiosíssimi Sánguinis fructum cónsequi mereámur in cœlis. Per eúmdem Dóminum.

\begin{tightlist}{}
\item \textred{In fine Hymnor. Completorii et Horarum dic. ℣.} Glória tibi Dómine, Qui nos salvásti in Cruce.
\item \textred{Ad Nunc dimittis, Aña.} Hoc signum Crucis * erit in cælo, cum Dóminus ad judicándum vénerit.
\end{tightlist}

\subsubsection*{\center \textsc{AD MATUTINUM.}}
\begin{tightlist}{}
\item \textred{Invitatorium.} Christum in Cruce confíxum, quinque plagis vulnerátum, \textred{*} Veníte adorémus.
\item \textred{Ps.} Veníte exsultémus. 2.
\item \textred{Hymn.} Pange lingua. \textred{ut supra ad j. Vesp.}
\end{tightlist}

\subsubsection*{\center \textsc{IN PRIMO NOCTURNO.}}
\begin{tightlist}{}
\item \textred{Aña.} Quis est qui venit de Edom, tinctis véstibus de Bosra?
\item \textred{Ps. 10.} In Dómino confído. 26.
\item \textred{Aña.} Quare rubrum est induméntum tuum, et vestiménta tua sicut calcántium in torculári?
\item \textred{Ps. 14.} Dómine quis. 65.
\item \textred{Aña.} Torcúlar calcávi solus, et de géntibus non fuit vir mecum.
\item \textred{Ps. 23.} Dómini est terra. 79.
\item \textred{℣.}~Fodérunt manus meas et pedes meos. \textred{℟.}~Dinumeravérunt ómnia ossa mea.
\end{tightlist}

\subsubsection*{\center De Isaía Prophéta.}
\begin{center}\textred{Lectio j. \hfill Cap. 53.}\end{center}
\lettrineThree{Q}uis crédidit audítui nostro? et bráchium Dómini cui revelátum est?
Et ascéndet sicut virgúltum coram eo, et sicut radix de terra sitiénti~: non est spécies ei, neque decor~: et vídimus eum, et non erat aspéctus, et desiderávimus eum:
despéctum, et novíssimum virórum, virum dolórum, et sciéntem infirmitátem~: et quasi abscónditus vultus ejus et despéctus, unde nec reputávimus eum.
Vere languóres nostros ipse tulit, et dolóres nostros ipse portávit~: et nos putávimus eum quasi leprósum, et percússum a Deo et humiliátum.
Ipse autem vulnerátus est propter iniquitátes nostras~: attrítus est propter scélera nostra~: disciplína pacis nostræ super eum, et livóre ejus sanáti sumus. Hæc dicit Dóminus.

\begin{tightlist}{}
\item \textred{℟.}~Gratificávit nos Deus in dilécto Fílio suo, \textred{*} In quo habémus redemptiónem per sánguinem ejus, in remissiónem peccatórum. \textred{℣.}~Ecce venit plenitúdo témporis, in quo misit Deus Fílium suum in terras. \textred{I}n quo.
\end{tightlist}

\subsubsection*{\center \textred{Lectio ij.}}
\lettrine{O}mnes nos quasi oves errávimus, unusquísque in viam suam declinávit~: et pósuit Dóminus in eo iniquitátem ómnium nostrum.
Oblátus est quia ipse vóluit, et non apéruit os suum~: sicut ovis ad occisiónem ducétur, et quasi agnus coram tondénte se obmutéscet, et non apériet os suum.
De angústia, et de judício sublátus est~: generatiónem ejus quis enarrábit? quia abscíssus est de terra vivéntium~: propter scelus pópuli mei percússi eum.
Et dabit ímpios pro sepultúra, et dívitem pro morte sua, eo quod iniquitátem non fécerit, neque dolus fúerit in ore ejus.
Et Dóminus vóluit contérere eum in infirmitáte. Hæc dicit Dóminus.

\begin{tightlist}{}
\item \textred{℟.}~Nos, qui aliquándo erámus longe, facti sumus prope per sánguinem Jesu Christi~: \textred{*} Ipse est pax nostra, qui fecit útraque unum. \textred{℣.}~A Dómino factum est istud, et est mirábile in óculis nostris. \textred{I}pse.
\end{tightlist}

\subsubsection*{\center \textred{Lectio iij. \hfill Cap. 63.}}
\lettrine{Q}uis est iste, qui venit de Edom, tinctis véstibus de Bosra? iste formósus in stola sua, grádiens in multitúdine fortitúdinis suæ. Ego, qui loquor justítiam, et propugnátor sum ad salvándum.
Quare ergo rubrum est induméntum tuum, et vestiménta tua sicut calcántium in torculári?
Tórcular calcávi solus, et de géntibus non est vir mecum~: calcávi eos in furóre meo, et conculcávi eos in ira mea~: et aspérsus est sanguis eórum super vestiménta mea, et ómnia induménta mea inquinávi.
Dies enim ultiónis in corde meo, annus redemptiónis meæ venit.
Circumspéxi, et non erat auxiliátor~: quæsívi, et non fuit qui adjuváret. Hæc dicit Dóminus.

\begin{tightlist}{}
\item \textred{℟.}In primogénito ex mórtuis complácuit omnem plenitúdinem divinitátis inhabitáre, et per eum reconciliáre ómnia in ipsum~: \textred{*} Pacíficans per sánguinem crucis ejus, sive quæ in cœlis, sive quæ in terris sunt. \textred{℣.}~Ipse est caput córporis Ecclésiæ, in ómnibus primátum tenens. \textred{P}acíficans. Glória Patri. \textred{P}acíficans.
\end{tightlist}

\subsubsection*{\center \textsc{IN SECUNDO NOCTURNO.}}
\begin{tightlist}{}
\item \textred{Aña.} Fodérunt manus meas, et pedes meos; dinumeravérunt ómnia ossa mea.
\item \textred{Ps. 21.} Deus Deus meus. 230.
\item \textred{Aña.} A planta pedis usque ad vérticem non est in eo sánitas.
\item \textred{Ps. 65.} Jubiláte Deo. 177.
\item \textred{Aña.} Vulnus, et livor et plaga tumens non est circumligáta, nec curáta medicámine, neque fota óleo.
\item \textred{Ps. 75.} Notus in Judæa. 202.
\item \textred{℣.}~Vere languóres nostros ipse tulit. \textred{℟.}~Et dolóres nostros ipse portávit.
\end{tightlist}

\subsubsection*{\center De Sermóne sancti Bernárdi Abbátis.}
\subsubsection*{\center \textred{Lib. de Pass. Dom. c. 41.\\
Lectio iv.}}
\lettrineThree{C}aptus ergo Jesus, post plúrimas illusiónes tam Judæórum, quam géntium, post sánguinis plures effusiónes, clavis immítibus mánibus simul et pédibus perforátur, et confígitur ligno crucis Salvátor noster, mitíssimus Jesus. Intuére et réspice rosam passiónis sanguíneæ, quómodo rubet in indícium ardentíssimæ caritátis. Conténdunt pássio et cáritas~: ista ut plus árdeat, illa ut plus rúbeat.

\begin{tightlist}{}
\item \textred{℟.}~Per mortem destrúxit Jesus eum, qui mortis habébat impérium~: \textred{*} Ut liberáret eos, qui timóre mortis obnóxii erant servitúti.  \textred{℣.}~Unde débuit per ómnia frátribus assimilári, ut miséricors fíeret. \textred{U}t.
\end{tightlist}

\subsubsection*{\center \textred{Lectio v.}}
\lettrine{V}ide quómodo hoc flore rosæ florúerit óptima vitis nostra, rubicúndus Jesus. Vide totum corpus, sícubi rosæ sanguíneæ florem non invénias. Inspice manum unam et álteram, si florem rosæ non invénias in utráque. Inspice pedem unum et álterum; numquid non rósei? Inspice láteris apertúram, quia nec illa caret rosa, quamvis ipsa subrúbea sit, propter mixtúram aquæ; quia, sicut narrat Evangelísta, cum unus mílitum láncea latus ejus perforásset, exívit sanguis et aqua.

\begin{tightlist}{}
\item \textred{℟.}~Per hóminem mors, et per hóminem resurréctio mortuórum~: \textred{*} Et sicut in Adam omnes peccavérunt, ita in Christo omnes vivificabúntur. \textred{℣.}~Novíssime autem inimíca destruétur mors. \textred{E}t sicut.
\end{tightlist}

\subsubsection*{\center Ex eódem divo Bernárdo.}
\subsubsection*{\center \textred{Ser. de Pass. Dom.\\
Lectio vj.}}
\lettrine{R}espice, Dómine, sancte Pater, de sanctuário tuo, et de excélso cœlórum habitáculo, et intuére hanc sacrosánctam hóstiam quam tibi offert magnus Póntifex noster, sanctus puer tuus, Dóminus Jesus, pro peccátis fratrum suórum, et esto placábilis super multitúdine malítiæ nostræ. Cognósce, Pater, túnicam fílii tui Joseph. Heu! fera péssima devorávit eum, et conculcávit in furóre suo vestiméntum ejus~: ecce quinque scissúras lamentábiles in eo relíquit.

\begin{tightlist}{}
\item \textred{℟.}~Vidémus Jesum propter passíonem mortis, glória et honóre coronátum~: \textred{*} Ut pro ómnibus gustáret mortem. \textred{℣.}~Decébat eum, qui multos fílios in glóriam addúxerat, auctórem salútis eórum per passiónem consummári. \textred{U}t pro. Glória Patri. \textred{U}t pro.
\end{tightlist}

\subsubsection*{\center \textsc{IN TERTIO NOCTURNO.}}
\begin{tightlist}{}
\item \textred{Aña.} Quid sunt plagæ istæ in médio mánuum tuárum?
\item \textred{Ps. 95.} Cantáte Dómino. [179]
\item \textred{Aña.} His plagátus sum in domo eórum, qui me diligébant.
\item \textred{Ps. 96.} Dóminus regnávit. [180]
\item \textred{Aña.} In amaritúdine sermo meus est, et manus plagæ meæ aggraváta est super gémitum meum.
\item \textred{Ps. 97.} Cantáte Dómino. [180]
\item \textred{℣.}~Ipse vulnerátus est propter iniquitátes nostras. \textred{℟.}~Attrítus est propter scélera nostra.
\end{tightlist}

\subsubsection*{\center Léctio sancti Evangélii secúndum Joánnem.}
\subsubsection*{\center \textred{Lectio vij. \hfill Cap. 19.}}
\lettrine{I}n illo témpore~: Sciens Jesus quia ómnia consummáta sunt ut consummarétur Scriptúra, dixit~: Sítio. Et réliqua.

\begin{center}Homilía sancti Cyrílli Epíscopi Alexandríni.\\
\textred{Lib. 12. Comm. in Joann.}\end{center}
\lettrineThree{C}um ab immánibus Judæis omnis in Christum expléta esset impíetas, nec quidquam jam ad summam crudelitátem réliquum esset, ad extrémum caro Christi sanctíssima próprium quiddam ac naturále rursus pátitur. Multis enim ac váriis dolóribus exsiccáta, siti torquétur. Graves síquidem dolóres magnam ad sitim commovéndam vim habent, naturáli quodam et inexplicábili calóre húmidum consuméntes, et ígneis ardóribus præcórdia uréntes.

\begin{tightlist}{}
\item \textred{℟.}~Manifestáta est grátia Dei per illuminatiónem Salvatóris nostri Jesu Christi~: \textred{*} Qui destrúxit mortem, illuminávit autem vitam in incorruptiónem. \textred{℣.}~Liberábit et vocávit nos vocatióne sua sancta per grátiam, quæ data est nobis in Christo. \textred{Q}ui.
\end{tightlist}

\subsubsection*{\center \textred{Lectio viij.}}
\lettrine{C}onsummátum esse ait impietátis Judáicæ modum, et impoténtem sæviéndi ánimum. Quid enim Judæi intentátum reliquérunt, aut quid non extrémæ crudelitátis adhibuérunt? Quod supplícii aut contuméliæ genus prætermisérunt? Quámobrem jure, Consummátum est, ait; sed hora jam eum ad prædicándum spirítibus, qui apud ínferos erant invitánte, (advénit enim ut mortuórum ac vivórum dominarétur) ipsam pro nobis mortem suscépit, commúnem hanc natúræ nostræ passiónem sustínuit in carne sua. Cum autem Jesum caput inclinásse comperíssent, et jam eum exspirásse putárent, frustra ejus crura confríngi exístimant, sed cum adhuc mórtuum esse non nihil diffíderent, láncea latus ejus perfódiunt, unde cruor aqua mixtus scatúriit, quod Eulogíæ mysticæ et sancti baptísmatis imágo quædam erat atque primítiæ. Christi enim est vere et a Christo sanctum baptísma, et mysticæ Eulogíæ vis ex sancta ejus nobis carne pródiit.

\begin{tightlist}{}
\item \textred{℟.}~Dignus es, Dómine, accípere librum, et aperíre signácula ejus, quóniam occísus es, et redemísti nos Deo. \textred{*} In sánguine tuo. \textred{℣.}~Fecísti enim nos Deo nostro regnum et Sacerdótes. \textred{I}n sánguine.
\item \textred{Lectio ix. de Homilia Feriæ occurrentis.}
\item \textred{℟. ix.} Omnes amíci mei dereliquérunt me, et prævaluérunt insidiántes mihi~: \textred{*} Et terribílibus óculis plaga crudéli percussérunt me. \textred{℣.}~Dedérunt in escam meam fel, et potavérunt me acéto. \textred{E}t. Glória Patri. \textred{E}t.
\item Te Deum laudámus. 6.
\item \textred{Sacerd. ℣.} Oblátus est quia ipse vóluit. \textred{℟.}~Et non apéruit os suum.
\end{tightlist}

\subsubsection*{\center \textsc{AD LAUDES.}}
\begin{center}\textred{et per Horas, Aña.}\end{center}
\lettrine{V}ere languóres nostros ipse tulit, et dolóres nostros ipse portávit.

\begin{tightlist}{}
\item \textred{Ps. de Dominica.} 27.
\item \textred{Aña.} Ipse autem vulnerátus est propter iniquitátes nostras, attrítus est propter scélera nostra.
\item \textred{Aña.} Alligávit Dóminus plagam pópuli sui, et percussúram ejus sanávit.
\item \textred{Aña.} Omnis qui transíbit, stupébit super omnes plagas ejus.
\item \textred{Aña.} O vos omnes, qui transítis per viam, atténdite, et vidéte si est dolor sicut dolor meus.
\end{tightlist}

\subsubsection*{\textred{Capitulum. \hfill Is. 53.}}
\lettrineThree{C}hristus vulnerátus est propter iniquitátes nostras, attrítus est propter scélera nostra~: cujus livóre sanáti sumus.

\subsubsection*{\center \textred{Hymnus.}}
\lettrine{L}ustris se qui jam peráctis,\\
Tempus implens córporis,\\
Se volénte natus ad hoc\\
Passióni déditus,\\
Agnus in Cruce levátur\\
Immolátus stípite.\\
\\
\textred{H}ic acétum, fel, arúndo,\\
Sputa, clavi láncea;\\
Mite corpus perforátur,\\
Sanguis, unda prófluit~:\\
Terra, pontus, astra, mundus,\\
Quo lavántur flúmine!\\
\\
\textred{C}rux fidélis inter omnes\\
Arbor una nóbilis~:\\
Nulla silva talem profert\\
Fronde, flore, gérmine~:\\
Dulce lignum, dulces clavos,\\
Dulce pondus sústinet.\\
\\
\textred{F}lecte ramos arbor alta,\\
Tensa laxa víscera;\\
Et rigor lentéscat ille,\\
Quem dedit natívitas~:\\
Ut supérni membra Regis\\
Miti tendas stípite.\\
\\
\textred{S}ola digna tu fuísti\\
Ferre sæcli prétium~:\\
Atque portum præparáre\\
Nauta mundo náufrago~:\\
Quem sacer cruor perúnxit,\\
Fusus Agni córpore.\\
\\
\textred{G}lória et honor Deo\\
Usquequáque altíssimo;\\
Una Patri, Filióque,\\
Inclyto Paráclito~:\\
Cui laus est, et potéstas\\
Per ætérna sǽcula.\\
Amen.

\begin{tightlist}{}
%\item \textred{Hymnus.} Lustris sex qui jam peráctis, \textred{ut in Dom. Passionis.} 393.
\item \textred{℣.}~Vidébunt in quem transfixérunt. \textred{℟.}~Et dolébunt super eum, ut in morte primogéniti.
\item \textred{Ad Bened. Aña.} Unus mílitum láncea latus ejus apéruit, et contínuo exívit sanguis et aqua.
\end{tightlist}

\subsubsection*{\center \textred{Oratio.}}
\lettrine{D}eus, qui Unigéniti Fílii tui passióne, et per quinque Vúlnera ejus Sánguinis effusióne humánam natúram peccáto pérditam reparásti~: tríbue nobis, quæsumus; ut qui ab eo suscépta Vúlnera venerámur in terris, ejúsdem pretiosíssimi Sánguinis fructum cónsequi mereámur in cœlis. Per eúmdem Dóminum.

\subsubsection*{\center \textsc{AD PRIMAM. 
\noindent \textred{℣.}~Qui nos salvásti in Cruce.

\subsubsection*{\center \textsc{AD TERTIAM.}}
\subsubsection*{\textred{Capitulum. \hfill Is. 53.}}
\lettrine{C}hristus vulnerátus est propter iniquitátes nostras, attrítus est propter scélera nostra~: cujus livóre sanáti sumus.

\begin{tightlist}{}
%\item \textred{Capit.} Christus vulnerátus, \textred{ut supra.}
\item \textred{℟.}~Fodérunt manus meas, \textred{*} Et pedes meos. \textred{F}odérunt. \textred{℣.}~Dinumeravérunt ómnia ossa mea. \textred{E}t. Glória Patri. \textred{F}odérunt.
\item \textred{℣.}~Vere languóres nostros ipse tulit. \textred{℟.}~Et dolóres nostros ipse portávit.
\end{tightlist}

\subsubsection*{\center \textsc{AD SEXTAM.}}
\subsubsection*{\textred{Capitulum. \hfill Is. 53.}}
\lettrine{F}aciem meam non avérti ab increpántibus, et conspuéntibus in me~: Dóminus Deus auxiliátor meus, et ídeo non sum confúsus.

\begin{tightlist}{}
\item \textred{℟.}~Vere languóres nostros \textred{*} Ipse tulit. \textred{V}ere \textred{℣.}~Et dolóres nostros ipse portávit. \textred{I}pse. Glória Patri. \textred{V}ere.
\item \textred{℣.}~Ipse vulnerátus est propter iniquitátes nostras. \textred{℟.}~Attrítus est propter scélera nostra.
\end{tightlist}

\subsubsection*{\center \textsc{AD NONAM.}}
\subsubsection*{\textred{Capitulum. \hfill Is. 53.}}
\lettrine{T}rádidit Christus in mortem ánimam suam, et cum scelerátis reputátus est; et ipse peccáta multórum tulit, et pro transgressóribus rogávit.

\begin{tightlist}{}
\item \textred{℟.}~Ipse vulnerátus est \textred{*} Propter iniquitátes nostras. \textred{I}pse. \textred{℣.}~Attrítus est propter scélera nostra. \textred{P}ropter. Glória Patri. \textred{I}pse.
\item \textred{℣.}~Vidébunt in quem transfixérunt. \textred{℟.}~Et dolébunt super eum, ut in morte primogéniti.
\end{tightlist}

\subsubsection*{\center \textsc{AD II. VESPERAS.}}
\begin{tightlist}{}
\item \textred{Añæ Laudum,}
\item \textred{Ps. 115.} Crédidi. 91. 
\item \textred{Ps. 119.} Ad Dóminum. 92. 
\item \textred{Ps. 139.} Eripe me. 244. 
\item \textred{Ps. 140.} Dómine clamávi. 244. 
\item \textred{Ps. 141.} Voce mea. 245.
\end{tightlist}

\subsubsection*{\textred{Capitulum. \hfill Isaiæ 53.}}
\lettrineThree{C}hristus vulnerátus est propter iniquitátes nostras, attrítus est propter scélera nostra~: cujus livóre sanáti sumus.

\subsubsection*{\center \textred{Hymnus.}}
\lettrine{P}ange, lingua, gloriósi\\
Prǽlium certáminis,\\
Et super Crucis trophǽum\\
Dic triúmphum nóbilem~:\\
Quáliter Redémptor orbis\\
Immolátus vícerit.\\
\\
\textred{D}e paréntis protoplásti\\
Fraude Factor cóndolens,\\
Quando pomi noxiális\\
Morsu in mortem córruit~:\\
Ipse lignum tunc notávit,\\
Damna ligni ut sólveret.\\
\\
\textred{H}oc opus nostræ salútis\\
Ordo depopóscerat;\\
Multifórmis proditóris\\
Ars ut artem fálleret,\\
Et medélam ferret inde,\\
Hostis unde lǽserat.\\
\\
\textred{Q}uando venit ergo sacri\\
Plenitúdo témporis,\\
Missus est ab arce Patris\\
Natus, orbis Cónditor~:\\
Atque ventre virgináli\\
Caro factus pródiit.\\
\\
\textred{V}agit infans inter arcta\\
Cónditus præsépia~:\\
Membra pannis involúta\\
Virgo mater álligat,\\
Et manus, pedésque et crura\\
Stricta cingit fáscia.\\
\\
\textred{G}lória et honor Deo\\
Usquequáque altíssimo,\\
Una Patri, Filióque,\\
Inclyto Paráclito~:\\
Cui laus est, et potéstas\\
Per ætérna sǽcula.\\
Amen.
%\item \textred{Hymn.} Pange lingua gloriósi, \textred{ut in Dom. Passionis.} 387.
\begin{tightlist}{}
\item \textred{℣.}~Fodérunt manus meas et pedes meos. \textred{℟.}~Dinumeravérunt ómnia ossa mea.
%\item \textred{Capit. Hymn. et ℣. ut in j. Vesperis.}
\item \textred{Ad Magnificat, Aña.} Ego sum vestra redémptio~: manus meæ, quæ vos fecérunt, clavis confíxæ sunt~: propter vos flagéllis cæsus sum, spinis coronátus sum~: aquam pétii pendens, et acétum porrexérunt~: in escam meam fel dedérunt, et in latus lánceam~: mórtuus et sepúltus resurréxi~: vobíscum sum, et vivo in ætérnum.
\end{tightlist}

\begin{center}\textcolor{red}{\rule{\linewidth}{0.5pt}} \end{center}

%\subsection*{\center \textsc{\textred{F}esta \textred{A}prilis.}}
%
%\begin{center}\textcolor{red}{\rule{\linewidth}{0.5pt}} \end{center}
\subsubsection*{\center \textred{In Festo}}
\subsection*{\center \textsc{Coronæ Domini}}
\subsubsection*{\center \textred{Duplex minus 1 classis.}}
%Breviary pg. 607 and following
\begin{tightlist}{}
\item \textred{Prima Feria sexta post Octavam Paschæ non impedita Festo novem Lectionum occurr. fit Offic. de Corona Spinea Domini nostri Jesu Christi sub ritu Duplici minori 1. classis.}
\item \textred{Quod si usque ad Ascensionem Domini omnes Feriæ sextæ simili Festo impeditæ sint, fiat Officium de Corona Domini intra istud tempus Feria sexta licet impedita Festo inferioris ritus.}
\end{tightlist}

\subsubsection*{\center \textsc{AD VESPERAS.}}
\fancyhead[C]{\textsc{Coronæ Spineæ D.N.J.C.}}
\begin{tightlist}{}
\item \textred{Aña.} Mílites præsidis suscipiéntes Jesum in prætórium, congregavérunt ad eum univérsam cohórtem, allelúja.
\item \textred{Psalmi de Officio præc.}
\end{tightlist}

\subsubsection*{\textred{Capitulum. \hfill Cant. 3.}}
\lettrineThree{E}gredímini, et vidéte fíliæ Sion regem Salomónem in diadémate, quo coronávit illum mater sua in die desponsatiónis illíus, et in die lætítiæ cordis ejus.

\begin{tightlist}{}
\item \textred{℟.}~O Corónæ trophæum ínclytum, sub quo vita vicit intéritum~: \textred{*} Tu tuórum solámen mílitum, et perénnis corónæ méritum, allelúja. \textred{℣.}~O præcláræ signum victóriæ, tuos serva servos in ácie. \textred{T}u. Glória Patri. \textred{T}u.
\end{tightlist}

\subsubsection*{\center \textred{Hymnus.}}
\lettrine{V}exílla Regis pródeunt,\\
Fulget Crucis mystérium,\\
Quo carne carnis Cónditor\\
Suspénsus est patíbulo.\\
\\
\textred{Q}uo vulnerátus ínsuper\\
Mucróne diro lánceæ,\\
Ut nos laváret crímine,\\
Manávit unda et sánguine.\\
\\
\textred{I}mpléta sunt quæ cóncinit\\
David fidéli cármine,\\
Dicens~: In natiónibus\\
Regnávit a ligno Deus.\\
\\
\textred{A}rbor decóra et fúlgida,\\
Ornáta Regis púrpura,\\
Elécta digno stípite\\
Tam sancta membra tángere.\\
\\
\textred{B}eáta, cujus bráchiis\\
Sæcli pepéndit prétium,\\
Statéra facta córporis,\\
Prædámque tulit tártari.\\
\\
\textred{O} Crux ave spes única,\\
In hac triúmphi glória\\
Auge piis justítiam,\\
Reísque dona véniam.\\
\\
\textred{T}e summa Deus Trínitas,\\
Colláudet omnis spíritus~:\\
Quos per Crucis mystérium\\
Salvas, rege per sæcula.\\
Amen.

\begin{tightlist}{}
\item \textred{℣.}~Tuam corónam adorámus Dómine. \textred{℟.}~Tuum gloriósum recólimus triúmphum, allelúja.
\item \textred{Ad Magnif. Aña.} Apprehéndit Pilátus Jesum et flagellávit~: et mílites plecténtes corónam de spinis, imposuérunt cápiti ejus, et veste purpúrea circumdedérunt eum, allelúja, allelúja.
\end{tightlist}

\subsubsection*{\center \textred{Oratio.}}
\lettrine{P}ræsta quæsumus omnípotens Deus, ut qui in memóriam Passiónis Dómini nostri Jesu Christi Corónam ejus spíneam venerámur in terris, ab ipso glória et honóre coronári mereámur in cœlis. Qui tecum vivit.

\begin{tightlist}{}
\item \textred{In Hymnis Complet. et Horar. dicuntur versus} \textred{℣.}~Quæsumus Auctor ómnium, \textred{et} \textred{℣.}~Glória tibi Dómine, qui nos salvásti in Cruce.
\item \textred{Ad Nunc dimittis, Aña.} Allelúja~: Hoc signum Crucis.
\end{tightlist}

\subsubsection*{\center \textsc{AD MATUTINUM.}}
\begin{tightlist}{}
\item \textred{Invit.} Christum Regem pro nobis spinis coronátum, \textred{*} Veníte, adorémus, allelúja.
\item \textred{Psalmus.} Veníte, exsultémus. 2.
\item \textred{Hymnus.} Pange lingua, \textred{ut in Dom. Passionis.} 387.
\end{tightlist}

\subsubsection*{\center \textsc{IN PRIMO NOCTURNO.}}
\begin{tightlist}{}
\item \textred{Aña.} Maledícta terra in ópere tuo, spinas, et tríbulos germinábit tibi, allelúja.
\item \textred{Ps. 10.} In Dómino. 26. 
\item \textred{Ps. 14.} Dómine, quis habitábit. 65. 
\item \textred{Ps. 23.} Dómini est terra. 79. 
\item \textred{℣.}~Tuam corónam adorámus Dómine. \textred{℟.}~Tuum gloriósum recólimus triúmphum, allelúja.
\end{tightlist}

\subsubsection*{\center De Isaía Prophéta.}
\subsubsection*{\center \textred{Lectio j. \hfill Cap. 53.}}
\lettrineThree{Q}uis crédidit audítui nostro? et bráchium Dómini cui revelátum est?
Et ascéndet sicut virgúltum coram eo, et sicut radix de terra sitiénti~: non est spécies ei, neque decor~: et vídimus eum, et non erat aspéctus, et desiderávimus eum:
despéctum, et novíssimum virórum, virum dolórum, et sciéntem infirmitátem~: et quasi abscónditus vultus ejus et despéctus, unde nec reputávimus eum. Hæc dicit Dóminus.

\begin{tightlist}{}
\item \textred{℟.}~Ecce vídimus eum non habéntem spéciem, neque decórem~: adspéctus ejus in eo non est; hic peccáta nostra portávit, et pro nobis dolet, ipse autem vulnerátus est propter iniquitátes nostras, \textred{*} Cujus livóre sanáti sumus, allelúja. \textred{℣.}~Vere languóres nostros ipse tulit, et dolóres nostros ipse portávit. \textred{C}ujus.
\end{tightlist}

\subsubsection*{\center \textred{Lectio ij.}}
\lettrine{V}ere languóres nostros ipse tulit, et dolóres nostros ipse portávit~: et nos putávimus eum quasi leprósum, et percússum a Deo et humiliátum.
Ipse autem vulnerátus est propter iniquitátes nostras~: attrítus est propter scélera nostra~: disciplína pacis nostræ super eum, et livóre ejus sanáti sumus.
Omnes nos quasi oves errávimus, unusquísque in viam suam declinávit~: et pósuit Dóminus in eo iniquitátem ómnium nostrum. Hæc dicit Dóminus.

\begin{tightlist}{}
\item \textred{℟.}~Quid sunt plagæ istæ in médio mánuum tuárum? \textred{*} Et dicet~: His plagátus sum in domo eórum, qui diligébant me, allelúja. \textred{℣.}~Frámea suscitáre super pastórem meum, et super virum cohæréntem mihi. \textred{E}t.
\end{tightlist}

\subsubsection*{\center De Zacharía Prophéta.}
\subsubsection*{\center \textred{Lectio iij. \hfill Cap. 12. et 13.}}
\lettrine{E}t erit in die illa~: quæram contérere omnes gentes, quæ véniunt contra Jerúsalem.
Et effúndam super domum David, et super habitatóres Jerúsalem spíritum grátiæ, et precum~: et adspícient ad me, quem confixérunt~: et plangent eum planctu quasi super unigénitum, et dolébunt super eum, ut doléri solet in morte primogéniti.
In die illa magnus erit planctus in Jerúsalem. Et dicétur ei~: Quid sunt plagæ istæ in médio mánuum tuárum? Et dicet~: His plagátus sum in domo eórum, qui diligébant me.
Frámea suscitáre super pastórem meum, et super virum cohæréntem mihi, dicit Dóminus exercítuum~: pércute pastórem, et dispergéntur oves~: et convértam manum meam ad párvulos. Hæc dicit Dóminus.

\begin{tightlist}{}
\item \textred{℟.}~Tuam corónam spíneam venerámur Dómine, tuam gloriósam recólimus passiónem~: \textred{*} Miserére nobis, qui passus es pro nobis, allelúja. \textred{℣.}~Salva præséntem catérvam in tuis hódie láudibus congregátam. \textred{M}iserére. Glória Patri. \textred{M}iserére.
\end{tightlist}

\subsubsection*{\center \textsc{IN SECUNDO NOCTURNO.}}
\begin{tightlist}{}
\item \textred{Aña.} Pro fruménto ortus est mihi tríbulus, et pro hórdeo spina, allelúja.
\item \textred{Ps. 46.} Omnes Gentes. 72. 
\item \textred{Ps. 65.} Jubiláte. 177. 
\item \textred{Ps. 75.} Notus in Judæa. 202.
\item \textred{℣.}~Glória et honóre coronásti eum Dómine. \textred{℟.}~Et constituísti eum super ópera mánuum tuárum, allelúja.
\end{tightlist}

\subsubsection*{\center Ex Tractátu sancti Athanásii de Passióne et Cruce Dómini.}
\subsubsection*{\center \textred{Lectio iv.}}
\lettrineThree{C}um exuerétur Dóminus, eo ipso invisibíliter erigebátur de diábolo trophæum~: et quo illud maniféstum fieret, narrátur a Joánne indútus fuísse púrpura, a Matthæo vero et Marco, chlámyde coccínea, coronámque illi ex spinis super caput esse pósitam, et cálamum in manu datum~: et ita cum eum produxíssent, génibus flexis adorásse. Miráculum novum et incredíbile, et magnæ sine dúbio victóriæ insígne. Quem enim ut hóminem condemnáverant, eum mórtuum deínde ut Deum adoravérunt~: et quem aspernabántur ut vilem et contemptíbilem, eum confitebántur Regem~: cui túnicam, et vestem detráxerant, eum púrpura circumdedérunt, et quem ignorántes contuméliis afficiébant, eum Prophétam étiam noléntes appellavérunt~: et quem per ludíbrium, subsannationémque pulsábant, ei triumphália addidérunt ornaménta, coccíneam chlámydem, et corónam spíneam, et cálamum. Nam tamétsi ea quasi per contuméliam agébant, efficiebátur tamen inde, ut ab eo laténter spoliaréntur.

\begin{tightlist}{}
\item \textred{℟.}~Coróna áurea super caput ejus, \textred{*} Expréssa signo sanctitátis, glória honóris, et ope fortitúdinis, allelúja. \textred{℣.}~Quóniam prævenísti eum in benedictiónibus dulcédinis, posuísti in cápite ejus corónam de lápide pretióso. \textred{E}xpréssa.
\end{tightlist}

\subsubsection*{\center Sermo sancti Ambrósii Epíscopi}
\subsubsection*{\center \textred{Serm. 73. de S. Cypriano.\\
Lectio v.}}
\lettrine{V}ideámus diligénter, ne ex pálmite óperis nostri non nascátur dulcédinis uva, sed fellis; ne dicátur nobis, quod Judæis dictum est per Prophétam~: Sustínui ut fáceret uvam, fecit autem spinas; et non justítiam, sed clamórem. Quam dura condítio, ut posséssio, quæ Dómino suo suavitátis fructus afférre débuit, spinárum illum asperitáte compúnxerit~: hoc est, cum Judæi Salvatórem debúerint tota fídei devotióne suscípere, spinárum eum coronavérunt passióne. Quæ coróna quidem, quantum in Judæis est, erat injuriárum contumélia~: quantum in Dómino, erat coróna virtútum. Compúngunt ergo Judæi Dóminum dum corónant, compúngunt dum crucifígunt. Unde ait Scriptúra sancta~: Vidébunt in quem compunxérunt.

\begin{tightlist}{}
\item \textred{℟.}~Omnes amíci mei dereliquérunt me, et prævaluérunt insidiántes mihi, trádidit me, quem diligébam, \textred{*} Et terribílibus óculis plaga crudéli percutiéntes, subsannábant mihi, allelúja. \textred{℣.}~Inter iníquos projecérunt me, et non pepercérunt ánimæ meæ. \textred{E}t.
\end{tightlist}

\subsubsection*{\center \textred{Lectio vj.}}
\lettrine{I}gitur fratres vidéte, ne et nobis dicátur~: Sustínui ut fáceret uvam, fecit autem spinas~: hoc est, ne spinas, quas Dómino Judæi cruéntis imposuérunt mánibus, nos iníquis sénsibus apponámus~: et compungámus eum, non asperitáte séntium, sed óperum pravitáte. Sunt enim spinæ cordis, quæ verbum Dei étiam vulneráre consuevérunt~: de quibus méminit in Evangélio Salvátor, dicens~: Seminántis semen cecidísse inter spinas, et crevísse eas, ac suffocásse, quod satum est. Ergo fratres vidéte, ne vínea vestra non uvas, sed spinas fáciat~: ne vindémia vestra non vinum áfferat, sed acétum. Quisquis enim vindémiam legit, et paupéribus non largítur~: ille acétum cólligit, non vinum. Quisquis messes recóndit, et indigéntibus non minístrat, iste non alimóniæ fructus repósuit, sed tríbulos avarítiæ congregávit. Nam proptérea ait Scriptúra de terra nostra~: Spinas et tríbulos germinábit tibi.

\begin{tightlist}{}
\item \textred{℟.}~Vínea mea elécta, ego te plantávi~: \textred{*} Quómodo convérsa es in amaritúdinem, ut me spinis coronáres? allelúja. \textred{℣.}~Sepívi te, et lápides elégi ex te, et ædificávi turrim. \textred{Q}uómodo. Glória Patri. \textred{Q}uómodo.
\end{tightlist}

\subsubsection*{\center \textsc{IN TERTIO NOCTURNO.}}
\begin{tightlist}{}
\item \textred{Aña.} Graváta est super me manus tua, convérsus sum in ærúmna mea, dum confígitur spina, allelúja.
\item \textred{Ps. 95.} Cantáte Dómino. [179] 
\item \textred{Ps. 96.} Dóminus regnávit. [180] 
\item \textred{Ps. 97.} Cantáte. [180]
\item \textred{℣.}~Eris coróna glóriæ in manu Dómini. \textred{℟.}~Et diadéma regni in manu Dei tui, allelúja.
\end{tightlist}

\subsubsection*{\center Lectio sancti Evangélii secúndum Joánnem.}
\subsubsection*{\center \textred{Lectio vij. \hfill Cap. 19.}}
\lettrine{I}n illo témpore~: Apprehéndit Pilátus Jesum, et flagellávit, et mílites plecténtes corónam de spinis, imposuérunt cápiti ejus. Et réliqua.

\begin{center}Homilía sancti Augustíni Epíscopi.\\
\textred{Tract. 116. in Joann.}\end{center}
\lettrineThree{H}oc Pilátus non ob áliud fecísse credéndus est, nisi ut ejus injúriis Judæi satiáti suffícere sibi existimárent, et usque ad ejus mortem sævíre desísterent. Ad hoc pértinet, quod idem præses cohórtem suam étiam permísit fácere, quæ sequúntur, aut fortássis et jussit, quamvis hoc Evangelísta tacúerit~: dixit enim quid deínde fécerint mílites, Pilátum tamen id jussísse non dicit.

\begin{tightlist}{}
\item \textred{℟.}~Christus passus est pro nobis, vobis relínquens exémplum, ut sequámini vestígia ejus~: \textred{*} Qui peccátum non fecit; nec invéntus est dolus in ore ejus, allelúja. \textred{℣.}~Qui cum maledicerétur, non maledicébat; cum paterétur, non comminabátur. \textred{Q}ui. 
\item \textred{Si nona Lectio debet legi de aliquo Sancto, duæ sequentes Lectiones simul conjunctæ legantur pro octava.}
\end{tightlist}

\subsubsection*{\center \textred{Lectio viij.}}
\lettrine{E}t mílites, inquit, plecténtes corónam de spinis, imposuérunt cápiti ejus, et veste purpúrea circumdedérunt eum, et veniébant ad eum, et dicébant~: Ave Rex Judæórum, et dabant ei álapas. Sic implebántur, quæ de se prædíxerat Christus~: sic Mártyres informabántur ad ómnia, quæ persecutóri libuísset fácere, perferénda~: sic paulísper occultáta treménda poténtia, commendabátur prius imitánda patiéntia~: sic regnum, quod de hoc mundo non erat, supérbum mundum non atrocitáte pugnándi, sed patiéndi humilitáte vincébat~: sic illud granum multiplicándum seminabátur horríbili contumélia, ut mirábili pulluláret in glória.

\begin{tightlist}{}
\item \textred{℟.}~Mílites præsidis suscipiéntes Jesum in prætórium congregavérunt ad eum univérsam cohórtem; et exuéntes eum, chlámydem coccíneam circumdedérunt ei~: \textred{*} Et plecténtes corónam de spinis, posuérunt super caput ejus, allelúja, allelúja, allelúja. \textred{℣.}~Genu flexo ante eum, illudébant ei dicéntes~: Ave Rex Judæórum. \textred{E}t.
\end{tightlist}

\subsubsection*{\center \textred{Lectio ix.}} 
\lettrine{E}xiit ergo Jesus portans corónam spíneam, et purpúreum vestiméntum~: et dicit illis Pilátus~: Ecce homo. Hic appáret, non ignoránte Piláto hæc a milítibus facta, sive jússerit ea, sive permíserit, ea scílicet causa, quam supra díximus, ut hæc ejus ludíbria inimíci libentíssime bíberent, et ultérius sánguinem non sitírent. Egréditur ad eos Jesus portans spíneam corónam, et purpúreum vestiméntum, non clarus império, sed plenus oppróbrio, et dícitur eis~: Ecce homo, si Regi invidétis, jam párcite, quia dejéctum vidétis. Flagellátus est, spinis coronátus est, ludibriósa veste amíctus est, amáris convítiis illúsus est, álapis cæsus est; fervet ignomínia, frigéscat invídia. Sed non frigéscit~: inardéscit pótius, et incréscit.

\begin{tightlist}{}
\item \textred{℟.}~O Corónæ trophæum ínclytum, sub quo vita vicit intéritum~: \textred{*} Tu tuórum solámen mílitum, et perénnis corónæ méritum, allelúja. \textred{℣.}~O præcláræ signum victóriæ, tuos serva servos in ácie. \textred{T}u. Glória Patri. \textred{T}u.
\item Te Deum laudámus. 6.
\item \textred{Sacerd. ℣.} Posuérunt advérsum me mala pro bonis. \textred{℟.}~Et ódium pro dilectióne mea, allelúja.
\end{tightlist}

\subsubsection*{\center \textsc{AD LAUDES.}}
\begin{center}\textred{et per Horas, Aña.}\end{center}
\lettrine{M}ilites præsidis suscipiéntes Jesum in prætórium, congregavérunt ad eum univérsam cohórtem, allelúja.

\begin{tightlist}{}
\item \textred{Ps. de Dominica.} 27.
\item \textred{Aña.} Et exuéntes eum, chlámydem coccíneam circumdedérunt ei, allelúja.
\item \textred{Aña.} Et plecténtes corónam de spinis, posuérunt super caput ejus, et arúndinem in déxtera ejus, allelúja.
\item \textred{Aña.} Et genu flexo ante eum, illudébant ei, dicéntes~: Ave Rex Judæórum, allelúja.
\item \textred{Aña.} Et exspuéntes in eum, accepérunt arúndinem, et percutiébant caput ejus, allelúja.
\end{tightlist}

\subsubsection*{\textred{Capitulum. \hfill Cant. 3.}}
\lettrineThree{E}gredímini, et vidéte fíliæ Sion regem Salomónem in diadémate, quo coronávit illum mater sua in die desponsatiónis illíus, et in die lætítiæ cordis ejus.

\begin{tightlist}{}
\item \textred{Hymnus.} Lustris sex, \textred{ut in Dom. Passionis.} 393.
\item \textred{℣.}~Dóminus regnávit. \textred{℟.}~Decórem índuit, allelúja.
\item \textred{Ad Bened. Aña.} Mílites índuunt Jesum púrpura, et impónunt ei plecténtes spíneam corónam. Et cœpérunt salutáre eum~: Ave Rex Judæórum, allelúja.
\end{tightlist}

\subsubsection*{\center \textred{Oratio.}}
\lettrine{P}ræsta quæsumus omnípotens Deus, ut qui in memóriam Passiónis Dómini nostri Jesu Christi Corónam ejus spíneam venerámur in terris, ab ipso glória et honóre coronári mereámur in cœlis. Qui tecum vivit.

\subsubsection*{\center \textsc{AD PRIMAM.}}
\noindent \textred{℣.}~Qui nos salvásti in Cruce.

\subsubsection*{\center \textsc{AD TERTIAM.}}
\subsubsection*{\textred{Capitulum. \hfill Cant. 3.}}
\lettrine{E}gredímini, et vidéte fíliæ Sion regem Salomónem in diadémate, quo coronávit illum mater sua in die desponsatiónis illíus, et in die lætítiæ cordis ejus.

\begin{tightlist}{}
%\item \textred{Capit. Egredímini, ut supra.}
\item \textred{℟.}~Tuam corónam adorámus Dómine, \textred{*} Allelúja, allelúja. \textred{T}uam. \textred{℣.}~Tuum gloriósum recólimus triúmphum. \textred{A}llelúja, allelúja. Glória Patri. \textred{T}uam.
\item \textred{℣.}~Glória et honóre coronásti eum Dómine. \textred{℟.}~Et constituísti eum super ópera mánuum tuárum, allelúja.
\end{tightlist}

\subsubsection*{\center \textsc{AD SEXTAM.}}
\subsubsection*{\textred{Capitulum. \hfill Apoc. 6.}}
\lettrine{V}idi, et ecce equus albus, et qui sedébat super illum, habébat arcum~: et data est ei coróna, et exívit vincens, ut vínceret.

\begin{tightlist}{}
\item \textred{℟.}~Glória et honóre coronásti eum Dómine, \textred{*} Allelúja, allelúja. \textred{G}lória et honóre. \textred{℣.}~Et constituísti eum super ópera mánuum tuárum. \textred{A}llelúja, allelúja. Glória Patri. \textred{G}lória et.
\item \textred{℣.}~Eris coróna glóriæ in manu Dómini. \textred{℟.}~Et diadéma regni in manu Dei tui, allelúja.
\end{tightlist}

\subsubsection*{\center \textsc{AD NONAM.}}
\subsubsection*{\textred{Capitulum. \hfill Isa. 28.}}
\lettrine{I}n illa die erit Dóminus exercítuum coróna glóriæ, et sertum exsultatiónis resíduo pópuli sui.

\begin{tightlist}{}
\item \textred{℟.}~Eris coróna glóriæ in manu Dómini, \textred{*} Allelúja, allelúja. \textred{E}ris. \textred{℣.}~Et diadéma regni in manu Dei tui. \textred{A}llelúja, allelúja. Glória Patri. \textred{E}ris.
\item \textred{℣.}~Dóminus regnávit. \textred{℟.}~Decórem índuit, allelúja.
\end{tightlist}

\subsubsection*{\center \textsc{AD VESPERAS.}}
\begin{tightlist}{}
\item \textred{Añæ Laudum, Psalmi de j. Vesperis Apostolorum, [2].}
\end{tightlist}

\subsubsection*{\textred{Capitulum. \hfill Cant. 3.}}
\lettrineThree{E}gredímini, et vidéte fíliæ Sion regem Salomónem in diadémate, quo coronávit illum mater sua in die desponsatiónis illíus, et in die lætítiæ cordis ejus.

\subsubsection*{\center \textred{Hymnus.}}
\lettrine{V}exílla Regis pródeunt,\\
Fulget Crucis mystérium,\\
Quo carne carnis Cónditor\\
Suspénsus est patíbulo.\\
\\
\textred{Q}uo vulnerátus ínsuper\\
Mucróne diro lánceæ,\\
Ut nos laváret crímine,\\
Manávit unda et sánguine.\\
\\
\textred{I}mpléta sunt quæ cóncinit\\
David fidéli cármine,\\
Dicens~: In natiónibus\\
Regnávit a ligno Deus.\\
\\
\textred{A}rbor decóra et fúlgida,\\
Ornáta Regis púrpura,\\
Elécta digno stípite\\
Tam sancta membra tángere.\\
\\
\textred{B}eáta, cujus bráchiis\\
Sæcli pepéndit prétium,\\
Statéra facta córporis,\\
Prædámque tulit tártari.\\
\\
\textred{O} Crux ave spes única,\\
In hac triúmphi glória\\
Auge piis justítiam,\\
Reísque dona véniam.\\
\\
\textred{T}e summa Deus Trínitas,\\
Colláudet omnis spíritus~:\\
Quos per Crucis mystérium\\
Salvas, rege per sæcula.\\
Amen.

\begin{tightlist}{}
\item \textred{℣.}~Tuam corónam adorámus Dómine. \textred{℟.}~Tuum gloriósum recólimus triúmphum, allelúja.
%\item \textred{Capit. Hymn. et ℣. ut in j. Vesp.}
\item \textred{Ad Magnificat, Aña.} O corónæ miránda dígnitas, mortis pœnas repéllis débitas. Tu spes lapsis, infírmis fírmitas, tu restáuras corónas pérditas, allelúja, allelúja, allelúja.
\end{tightlist}

\begin{center}\textcolor{red}{\rule{\linewidth}{0.5pt}} \end{center}

\subsubsection*{\center \textred{Die V. Maii.}}
\subsection*{\center \textsc{S. Angeli, Martyris Ordinis Nostri.}}
\subsubsection*{\center \textred{Duplex Majus 2. Classis.}} 
\begin{center} \textred{Vide festum supra, 892.}\end{center}

\begin{center}\textcolor{red}{\rule{0.5\linewidth}{0.5pt}} \end{center}

\subsubsection*{\center \textsc{INFRA OCTAVAM\\ Sancti Angeli.}}
\fancyhead[C]{\textsc{Infra Octavam S. Angeli.}}

\begin{tightlist}{}
\item \textred{Fit de ea comm., usque ad Vig. Pent. exclusive, in Laud. per Añam, quæ habetur ad Benedictus. 901.; et in Vesp. per Añam., quæ habetur ad Magnificat de ij. Vesp. Festi, 903., exceptis Vesp. Sabbati, in quibus dicitur Aña. quæ habetur ad Magnificat in j. Vesp. 893., et dicitur Oratio Festi.}
\end{tightlist}

\begin{center}\textcolor{red}{\rule{\linewidth}{0.5pt}} \end{center}

\subsubsection*{\center \textred{Die XII. Maii.}}
\subsection*{\center \textsc{Octava S. Angeli.}}
\subsubsection*{\center \textred{Duplex minus.}}
\begin{tightlist}{}
\item \textred{Omnia ut in Festo, 892. præter hic signata.}
\item \textred{In j. Vesp. si non præcedat Festum Duplex, Aña prima Laudum, Psalmi de Officio præc.}
\item \textred{In j. Noct. Lect. de Script. occurr. Sed si venerit in Feria ij. vel iv. Rogationum, legantur Lectiones ut in die.}
\end{tightlist}

\subsubsection*{\center \textsc{IN SECUNDO NOCTURNO.}}
\fancyhead[C]{\textsc{12 Maii - Dies Octava S. Angeli.}}
\subsubsection*{\center Sermo sancti Augustíni Epíscopi.}
\subsubsection*{\center \textred{Serm. de S. Cypriano 113.\\ Lectio iv.}}
\lettrineThree{S}píritus sanctus dóceat nos in hac hora, quæ opórteat dícere. Dictúri enim sumus áliquid de laude Angeli gloriosíssimi Mártyris, cujus natálem hódie, sicut nostis, celebrámus. Quod nomen sic frequéntat Ecclésia, id est, natáles, ut natáles vocet pretiósas Mártyrum mortes. Quid est hoc fratres? quando natus est, ignorámus~: et quia hódie passus est, natálem ejus hódie celebrámus. Sed illum diem non celebrarémus, etsi nossémus. Illo enim die traxit originále peccátum~: isto autem die vicit omne peccátum.

\begin{tightlist}{}
\item \textred{℟.}~Christus ex desérto Angelum mittit ad prædicándum, dicens~: Argue, íncrepa, et clama in abominatiónem impiórum~: adínstar Joánnis Baptístæ \textred{*} Martyrio coronáberis, allelúja. \textred{℣.}~Erit dulce mihi pro tua veritáte mori. \textred{M}artyrio.
\end{tightlist}

\subsubsection*{\center \textred{Lectio v.}}
\lettrine{I}stum nobis festum diem pássio beatíssimi Angeli mártyris fecit~: cujus nos victóriæ celébritas in istum locum devotíssimos congregávit. Sed celebrátio solemnitátis mártyrum, imitátio debet esse virtútum. Fácile est honórem mártyris celebráre, magnum est fidem atque patiéntiam mártyris imitári. Hoc sic agámus, ut illud optémus~: hoc sic celebrémus, ut illud pótius diligámus. Quid laudémus in fide mártyris? quia usque ad mortem certávit pro veritáte, et ídeo vicit. Blandiéntem mundum contémpsit, sæviénti non cessit~: ídeo victor ad Deum accéssit. Abúndant in isto sæculo erróres, et terróres. Beatíssimus martyr erróres sapiéntia, terróres patiéntia superávit.

\begin{tightlist}{}
\item \textred{℟.}~Oránte sancto Angelo trecénti infidéles óculis capti sunt, \textred{*} Et póstea visum recepérunt, allelúja. \textred{℣.}~Quicúmque vestrum, dixit Angelus, credet in Jesum Christum, spíritus et córporis lumen recípiet. \textred{E}t.
\end{tightlist}

\subsubsection*{\center \textred{Lectio vj.}}
\lettrine{C}ontémnite ergo sæculum, Christiáni, contémnite sæculum. Contémnite; contempsérunt Mártyres, contempsérunt Apóstoli, contémpsit beátus Angelus, cujus hódie memóriam celebrámus. Dívites esse vultis, honoráti esse vultis, sani esse vultis~: totum ille contémpsit, ad cujus memóriam convenístis. Quid, óbsecro, tantum amátis, quod contémpsit, quem sic honorátis? Quem si ista non contempsísset, non útique sic honorarétis? Quare te invénio eárum rerum amatórem, quarum veneráris contemptórem? Certe illum, si hæc amáret, non veneraréris. Et tu noli amáre~: non enim intrávit et óstium contra te clausit. Contémne et tu, et intra post illum.

\begin{tightlist}{}
\item \textred{℟.}~Instar Patris sui Elíæ, Angelus fílium víduæ revocávit ad vitam, \textred{*} In virtúte passiónis Dómini nostri Jesu Christi, allelúja. \textred{℣.}~Scitóte, dixit adoléscens, me sancti Angeli précibus ab utráque morte liberátum. \textred{I}n. Glória Patri. \textred{I}n.
\end{tightlist}

\subsubsection*{\center \textsc{IN TERTIO NOCTURNO.}}
\subsubsection*{\center Lectio sancti Evangélii secúndum Matthæum.}
\subsubsection*{\center \textred{Lectio vij. \hfill Cap. 16.}}
\lettrine{I}n illo témpore~: Dixit Jesus discípulis suis~: Si quis vult post me veníre, ábneget semetípsum, et tollat crucem suam, et sequátur me. Et réliqua.

\begin{center}Homilía sancti Augustíni Epíscopi.\\
\textred{Serm. 47. de divers.}\end{center}
\lettrineThree{D}urum vidétur, et grave, quod Dóminus imperávit~: ut, si quis eum vult sequi, ábneget semetípsum~: sed non est durum, nec grave, quod ille ímperat, qui ádjuvat, ut fiat quod ímperat. Nam et illud verum est, quod ei dícitur in Psalmo~: Propter verba labiórum tuórum ego custodívi vias duras. Quidquid enim durum est in præcéptis, ut sit leve cáritas facit. Nóvimus, quanta ipse amor facit. Quid autem est, Neget se? Non præsúmat de se, séntiat se hóminem, et respíciat dictum Prophéticum~: Maledíctus omnis, qui spem suam ponit in hómine. Subdúcat se sibi sed non deórsum versus. Subdúcat se sibi, ut hæreat Deo.

\begin{tightlist}{}
\item \textred{℟.}~Sanctus Joánnes Baptísta Angelo dixit~: Lætáre quia sánguinem tuum pro veritáte effúndes, et Angeli te exspéctant, \textred{*} In cœlésti pátria, allelúja. \textred{℣.}~Intrépidus Christi miles parta victória perducéris ad glóriam. \textred{I}n.
\end{tightlist}

\subsubsection*{\center \textred{Lectio viij.}}
\lettrine{T}ollat crucem suam, et sequátur me. Quo sequéndus est Dóminus? Quo iit, nóvimus~: resurréxit enim, et ascéndit in cœlum~: illo sequéndus est. Plane desperándum non est, quia ipse promíssit, non quia homo áliquid potest. Jam quare desperémus, si membra illíus cápitis sumus? Bonum est illo eum sequi, sed vidéndum est qua. Etenim verba ista Dóminus Jesus non tunc dicébat, quando a mórtuis jam resurréxerat. Nondum erat passus, ventúrus erat ad crucem, ventúrus ad exhonoratiónem, ad contumélias, ad flagélla, ad spinas, ad vúlnera, ad insultatiónes, oppróbria, mortem. Quasi exasperáta est via~: pigrum te facis, non vis sequi~: Séquere. Nam quis non velit ire ad exaltatiónem? Omnes deléctat celsitúdo; sed humílitas gradus est.

\begin{tightlist}{}
\item \textred{℟.}~Moriénte Angelo audíta est vox de cœlo dicens~: Veni Angele ad regnum tibi parátum a constitutióne mundi, ut inter Sanctos collocéris, \textred{*} Et cum Martyribus in ætérnum vivas, allelúja. \textred{℣.}~Tu in supérna Jerúsalem corónam glóriæ meruísti. \textred{E}t.
\end{tightlist}

\subsubsection*{\center \textred{Lectio ix.}}
\lettrine{Q}uid est~: Tollat crucem suam? Ferat quidquid moléstum est~: sic me sequátur. Cum enim cœperit me móribus et præcéptis meis sequi, multos habébit contradictóres, multos habébit prohibitóres, multos habébit dissuasóres, et hos de ipsis quasi comítibus Christi. Cum Christo ambulábant, qui cæcos clamáre prohibébant. Sive ergo minas, sive blandiménta, sive quáslibet prohibitiónes, si sequi vis, in crucem verte; tólera, porta, noli succúmbere. Vidéntur his verbis Dómini exhortáta martyria; si persecútio est, nonne pro Christo debent cuncta contémni? Amátur mundus, sed præponátur, a quo factus est mundus. 

\begin{tightlist}{}
\item \textred{℟.}~In desérto orábat Angelus~: Clementíssime Jesu, civitátem istam Jerúsalem imménsæ bonitátis tuæ óculo réspice, \textred{*} In qua cruci affíxus sánguinem tuum effúndere voluísti, allelúja. \textred{℣.}~Miserére, miserére, et propítius esto propter nomen tuum, ne semper in ea gloriéntur inimíci tui. \textred{I}n. Glória Patri. \textred{I}n.
\item Te Deum laudámus. 6.
\item \textred{Sacerd. ℣.} Stola jucunditátis índuit eum Dóminus. \textred{℣.}~Et corónam pulchritúdinis pósuit super caput ejus.
\item \textred{In Vesp. com. seq. (Die XIII Maii. S. Roberti Bellarmino, 937.)}
\end{tightlist}

\begin{center}\textcolor{red}{\rule{\linewidth}{0.5pt}} \end{center}

\subsubsection*{\center \textred{Die XVI. Maii.}}
\subsection*{\center \textsc{S. Simonis Anglici, Confess. Ordinis Nostri.}}
\subsubsection*{\center \textred{Duplex Majus S. C.}}

\begin{center}\textcolor{red}{\rule{0.5\linewidth}{0.5pt}} \end{center}

\subsubsection*{\center \textsc{INFRA OCTAVAM\\ Sancti Simonis.}}
\fancyhead[C]{\textsc{Infra Octavam S. Simonis.}}
\begin{tightlist}{}
\item \textred{Fit de ea comm. usque ad Vig. Pent. exclusive, in Laud. per Añam. ad Benedictus de Laud. 951. et in Vesp. per Añam ad Magnif. de ij. Vesp. Festi, 954. exceptis Vesperis Sabbati, in quibus fit comm. per Añam ad Magnif. de j. Vesp. Festi. 944.}
\end{tightlist}

\begin{center}\textcolor{red}{\rule{\linewidth}{0.5pt}} \end{center}

\subsubsection*{\center \textred{Die XXI. Maii.}}
\subsection*{\center \textsc{Translationis\\ S. Joannis a Cruce}}
\subsubsection*{\center \textred{Confessoris O.N.\\ Duplex minus.}}

\begin{tightlist}{}
\item \textred{Omnia de Comm. unius Conf. non Pont. [97] præter hic posita.}
\end{tightlist}

\subsubsection*{\center \textsc{IN I. VESPERIS.}}
\fancyhead[C]{\textsc{21 Maii - Translationis S. Joannis a Cruce.}}
\begin{tightlist}{}
\item \textred{Aña.} Addúcam eum * in montem sanctum meum, et lætificábo eum in domo oratiónis meæ.
\item \textred{Psalmi de Officio præc.}
\end{tightlist}

\subsubsection*{\center \textred{Hymnus.}}
\lettrine{R}egis ætérni generóse miles,\\
Ordinis nostri cólumen Joánnes,\\
Quos tuis læti méritis dicámus,\\
Súscipe plausus.\\
\\
\textred{O} satis felix! Speculátor alti\\
Núminis, votis animóque Martyr,\\
Pœnitens Virgo, memoránde Vates,\\
Mystice Doctor.\\
\\
\textred{S}æpe cum Christo et Genitríce pura\\
Verba commísces, superísque vivens~:\\
Hinc fluit candor, tua quo refúlgent\\
Scripta, perénnis.\\
\\
\textred{M}ente perfúsa rádiis ab alto,\\
Montis ascénsum, tenebrásque noctis,\\
Et facem vivam récolens amóris,\\
Alta revélas.\\
\\
\textred{T}e sacri verbi retegénte sensus,\\
Visa calígo procul exsuláre,\\
Dum facis lumen rútilans oríri\\
Noctis ab umbris.\\
\\
\textred{P}érsonent laudes Tríadi Joánnis,\\
Cui Pater lumen, tribuítque dona\\
Spíritus, vitæ réserat volúmen\\
Cœlicus Agnus.\\
Amen.

\begin{tightlist}{}
\item \textred{℣.}~Ora pro nobis sancte pater Joánnes. \textred{℟.}~Ut digni efficiámur promissiónibus Christi, allelúja.
\item \textred{Ad Magnificat, Aña.} Collaudábunt multi sapiéntiam ejus, et usque in sæculum non delébitur, allelúja.
\end{tightlist}

\subsubsection*{\center \textred{Oratio.}}
\lettrine{D}eus, qui sanctum Joánnem Confessórem tuum perféctæ sui abnegatiónis, et crucis amatórem exímium effecísti~: concéde; ut cujus hódie Translatiónem cólimus, ejúsdem imitatióni júgiter inhæréntes, glóriam assequámur ætérnam. Per Dóminum.

\subsubsection*{\center \textsc{AD MATUTINUM.}}
\subsubsection*{\center \textred{Hymnus.}}
\lettrine{D}iem Joánnes ádvehit,\\
Crucis micántem glória,\\
Cui fixus hæsit spíritu,\\
Fixámque gessit córpori.\\
\\
\textred{N}on probra, nec ludíbria\\
Fames, flagélla, víncula,\\
Sedáre jam valent sitim,\\
Qua dura fortis áppetit.\\
\\
\textred{H}æc poscit ille gáudia,\\
Hæc vult labóris præmia,\\
Hæc vota, confíxi Ducis\\
Reférre signum pérpetim.\\
\\
\textred{N}octem volútans mysticam\\
In nocte lucem pércipit,\\
De luce flammam súscipit,\\
Carméli ut agmen ínstruat.\\
\\
\textred{T}e, Christe, laudent Cœlites,\\
Tibíque corda dédita\\
Curant, Joánne prævio,\\
Ad vera vitæ gáudia.\\
Amen.

\begin{tightlist}{}
\item \textred{In j. Noct. Lect. de Scriptura occurrente, quæ si desint, legantur sequentes:}
\end{tightlist}

\subsubsection*{\center De Epístola beáti Pauli Apóstoli ad Philippénses.}
\subsubsection*{\center \textred{Lectio j. \hfill Cap. 1. et 2.}}
\lettrineThree{F}ratres~: Vobis donátum est pro Christo, non solum ut in eum credátis, sed ut étiam pro illo patiámini:
idem certámen habéntes, quale et vidístis in me, et nunc audístis de me. Si qua ergo consolátio in Christo, si quod solátium caritátis, si qua socíetas spíritus, si qua víscera miseratiónis:
impléte gáudium meum ut idem sapiátis, eámdem caritátem habéntes, unánimes, idípsum sentiéntes,
nihil per contentiónem, neque per inánem glóriam~: sed in humilitáte superióres sibi ínvicem arbitrántes,
non quæ sua sunt sínguli considerántes, sed ea, quæ aliórum.

\begin{tightlist}{}
\item \textred{℟.}~Euge serve. [99]
\end{tightlist}

\subsubsection*{\center \textred{Lectio ij.}}
\lettrine{V}erúmtamen ad quod pervénimus ut idem sapiámus, et in eádem permaneámus régula.
Imitatóres mei estóte fratres, et observáte eos qui ita ámbulant, sicut habétis formam nostram.
Multi enim ámbulant, quos sæpe dicébam vobis (nunc autem et flens dico) inimícos crucis Christi:
quorum finis intéritus~: quorum Deus venter est~: et glória in confusióne ipsórum, qui terréna sápiunt.
Nostra autem conversátio in cælis est~: unde étiam Salvatórem exspectámus Dóminum nostrum Jesum Christum, qui reformábit corpus humilitátis nostræ, configurátum córpori claritátis suæ, secúndum operatiónem, qua étiam possit subjícere sibi ómnia.

\begin{tightlist}{}
\item \textred{℟.}~Justus germinábit. [100]
\end{tightlist}

\subsubsection*{\center \textred{Lectio iij.}}
\lettrine{N}ihil sollíciti sitis~: sed in omni oratióne, et obsecratióne, cum gratiárum actióne petitiónes vestræ innotéscant apud Deum.
Et pax Dei, quæ exsúperat omnem sensum, custódiat corda vestra, et intelligéntias vestras in Christo Jesu.
De cétero fratres, quæcúmque sunt vera, quæcúmque pudíca, quæcúmque justa, quæcúmque sancta, quæcúmque amabília, quæcúmque bonæ famæ, si qua virtus, si qua laus disciplínæ, hæc cogitáte.
Quæ et didicístis, et accepístis, et audístis, et vidístis in me, hæc ágite~: et Deus pacis erit vobíscum.

\begin{tightlist}{}
\item \textred{℟.}~Iste cognóvit. [100]
\end{tightlist}

\subsubsection*{\center \textsc{IN SECUNDO NOCTURNO.}}
\subsubsection*{\center \textred{Lectio iv.}}
\lettrineThree{J}oánnes a Cruce Ordinis Fratrum beátæ Maríæ Vírginis de Monte Carmélo admiránda innocéntia, assídua rerum divinárum contemplatióne, aspérrimo vitæ rigóre, summísque virtútibus ínclytus, strictióris disciplínæ promovéndæ ardóre veheménter accénsus, plane cœlitus datus est. Vir quidem Deo percárus, virtútibus, miráculis, nec non prophetíæ, ac mysticæ theologíæ arcána scripto explicándi dono, æque ac sancta Terésia, insígnis fuit. Tandem diro morbo, et in crure quinque plagis sánie manántibus, patientíssime tolerátis, inter sodáles pie collacrymántes in Ubédæ cœnóbio puríssimam efflávit ánimam in ósculo Dómini.

\begin{tightlist}{}
\item \textred{℟.}~Desidérium. [101]
\end{tightlist}

\subsubsection*{\center \textred{Lectio v.}}
\lettrine{A}d servi Dei exúvias, quasi odóre perfúsas et deínde corruptiónis expértes, eúmdem véluti plácido somno sopítum referéntes, pópuli eas osculántis multitúdo innúmera catervátim conflúxit, véstium, et ipsíus étiam córporis partículas veneratiónis instínctu auférre conáta. Venerábile corpus religióse tandem ibídem cónditum, et póstea ex Ubéda in Segoviénsis cœnóbii ecclésiam asportátum, sub humo in sacéllo beátæ Maríæ Vírginis de Monte Carmélo diu asservátum fuit. Tanti viri virtútum et prodigiórum fama percrebrescénte, iisdémque canónice probátis, ipsum Clemens Décimus Beatórum catálogo adscrípsit.

\begin{tightlist}{}
\item \textred{℟.}~Amávit eum. [101]
\end{tightlist}

\subsubsection*{\center \textred{Lectio vj.}}
\lettrine{I}dem corpus ex húmili, et antíquo túmulo sublátum est, ac incorrúptum, palpábile, mirificúmque emíttens odórem repértum, duodécimo Kaléndas Júnias subsequéntes lóculo pretióse ornáto inclúsum, in ara cappéllæ magnificentíssime exstrúctæ repósitum est. Quo autem rei gestæ pia et festíva pompa deínde respondéret, præscrípta die univérsus clerus sæculáris, et gubernátor cum magistrátu, generáli institúta supplicatióne, ad præfátam cœnóbii ecclésiam se contulérunt, ubi ad novam aram incruénto sacrifício solémni ritu obláto, sacri pígnoris exaltátio, máximis étiam lætítiæ signis éditis, compleméntum hábuit. Postquam vero ipse Dei fámulus, multis a Deo ejus intercessióne, patrátis miráculis, a Benedícto Décimo tértio in Sanctórum album relátus est, solémnis hujus Translatiónis festum Clemens Duodécimus Póntifex Máximus univérso Carmelitárum Ordini offício próprio celebrándum concéssit.

\begin{tightlist}{}
\item \textred{℟.}~Iste sanctus. [101]
\end{tightlist}

\subsubsection*{\center \textsc{IN TERTIO NOCTURNO.}}
\subsubsection*{\center Léctio sancti Evangélii secúndum Lucam.}
\subsubsection*{\center \textred{Lectio vij. \hfill Cap. 11.}}
\lettrine{I}n illo témpore~: Dixit Jesus turbis~: Nemo lucérnam accéndit, et in abscóndito ponit, neque sub módio, sed supra candelábrum; ut, qui ingrediúntur, lumen vídeant. Et réliqua.

\begin{center}Homilía sancti Joánnis Chrysóstomi.\\
\textred{Hom. 10. in cap. 5. Matth.}\end{center}
\lettrineThree{P}er hanc comparatiónem vult osténdere, quare Sanctos suos ipse Christus maniféstat, et non sinit eos esse abscónsos, id est, quia nec qui accéndunt lucérnam, ídeo eam accéndunt, ut ponant sub módio, sed super candelábrum, ut lúceat ómnibus, qui in domo sunt. Qui sunt accensóres lucérnæ? Pater, et Fílius. Quæ est illa lucérna? Verbum divínum, de quo dictum est~: Lucérna pédibus meis verbum tuum. Quod est candelábrum? Ecclésia, quæ bájulat verbum vitæ; consequénter dicit et Paulus~: inter quos lucétis, sicut luminária in sæculo, verbum vitæ continéntes. Consequénter et omnis ecclesiásticus vir habens verbum Dei, dícitur candelábrum.

\begin{tightlist}{}
\item \textred{℟.}~Magnificávit. [102]
\end{tightlist}

\subsubsection*{\center \textred{Lectio viij.}}
\lettrine{M}odii autem sunt hómines mundiáles, vácui a Deo, et ab ómnibus, quæ sunt Dei. Dicúntur autem módii, quóniam sicut módii désuper quidem vácui sunt, desúbtus autem pleni~: sic omnes mundi amatóres, et carnis, in spirituálibus quidem et divínis, quæ superióra recte esse dicúntur, vácui sunt, et insensáti~: in inferióribus autem, id est, in mundiálibus et terrénis aliquátenus pleni vidéntur, et sapiéntes. Proptérea sicut in módio si volúeris áliquid repónere, necésse est, ut cadat deórsum ad inferióra ejus, et tenebrósa~: sic hómini carnáli, et mundiáli, et vácuo, si díxeris áliquid, quod est Dei, non illud tenet.

\begin{tightlist}{}
\item \textred{℟.}~Vir Israelíta. [102]
\end{tightlist}

\subsubsection*{\center \textred{Lectio ix.}}
\lettrine{N}unc autem sciens, quóniam ille est laudábilis, qui facit quod docet, addit et dicit~: Sic lúceat lux vestra coram homínibus, ut vídeant ópera vestra bona, et magníficent Patrem vestrum, qui in cœlis est~: id est, sic illumináte, non ut verba vestra tantúmmodo áudiant hómines, sed ut ópera vestra bona vídeant, ut quos illuminavéritis per verbum, condiátis per exémpla óperum vestrórum~: quóniam qui docet, et facit quod docet, vere ille docet~: qui autem non facit quod docet, non álium docet, sed seípsum condémnat. Et mélius est fácere, et non docére, quam docére, et non fácere; quóniam qui facit, et si tacúerit, áliquos córrigit suo exémplo.

\begin{tightlist}{}
\item \textred{℟.}~Sancte Joánnes. [103]
\item Te Deum laudámus. 6.
\item \textred{Sacerd. ℣.} Os justi meditábitur sapiéntiam. \textred{℟.}~Et lingua ejus loquétur judícium, allelúja.
\end{tightlist}

\subsubsection*{\center \textsc{AD LAUDES.}}
\begin{center}\textred{et per Horas, Aña.}\end{center}
\lettrine{A}ddúcam eum in montem sanctum meum, et lætificábo eum in domo oratiónis meæ, allelúja.

\begin{tightlist}{}
\item \textred{Ps. de Dominica.} 27.
\item \textred{Aña.} Apéruit os suum in oratióne, et Dóminus spíritu intelligéntiæ replévit illum, allelúja.
\item \textred{Aña.} Cibávit illum Dóminus pane vitæ et intelléctus, et aqua sapiéntiæ salutáris potávit illum, allelúja.
\item \textred{Aña.} Vir sápiens plebem suam érudit, et fructus sensus illíus fidéles sunt, allelúja.
\item \textred{Aña.} Ipse tamquam imbres mittet elóquia sapiéntiæ suæ, allelúja.
\end{tightlist}

\subsubsection*{\center \textred{Hymnus.}}
\lettrine{D}um Crucem gestat Dóminus, Joánni\\
Alloquens offert méritas corónas,\\
Is crucis tantum cálices amáros\\
Súmere póscit.\\
\\
\textred{S}unt pati ac sperni sua vota, merces;\\
Flamma sed nunquam saturátur ingens;\\
Emori nulli sub honóre notus,\\
Instat, habétque.\\
\\
\textred{P}érgraves tandem Crucis hos labóres\\
Morte conclúdit, meritúmque pandit\\
Finis extrémi globus igne flagrans\\
Lapsus Olympo.\\
\\
\textred{F}it micans æther rádiis, cubíle,\\
Lámpadum lumen tenebréscit omne,\\
Testis est fragrans odor ipse vitæ\\
Córpore sparsus.\\
\\
\textred{S}umma laus Patri, genitóque Verbo,\\
Et tibi compar, utriúsque Flamen,\\
Quo Beatórum pia turba semper\\
Páscitur igne.\\
Amen.

\begin{tightlist}{}
\item \textred{℣.}~Justus germinábit sicut lílium. \textred{℟.}~Et florébit in ætérnum ante Dóminum, allelúja.
\item \textred{Ad Benedictus, Aña.} Dedit illi sciéntiam Sanctórum, honestávit illum in labóribus, et complévit labóres illíus, allelúja.
\end{tightlist}

\subsubsection*{\center \textred{Oratio.}}
\lettrine{D}eus, qui sanctum Joánnem Confessórem tuum perféctæ sui abnegatiónis, et crucis amatórem exímium effecísti~: concéde; ut cujus hódie Translatiónem cólimus, ejúsdem imitatióni júgiter inhæréntes, glóriam assequámur ætérnam. Per Dóminum.

\begin{tightlist}{}
\item \textred{Comm. Oct. S. Simonis., Aña.} Beátus Simon * quasi stella matutína in diébus suis luxit, et quasi sol sic effúlsit in templo Dei, sui memóriam relínquens pósteris in benedictióne, allelúja.
\item \textred{℣.}~Justus ut palma florébit in domo Dómini. \textred{℟.}~Sicut cedrus Líbani multiplicábitur.
\item Orémus. \hfill \textred{Oratio.}
\end{tightlist}
\lettrine{D}eus, qui précibus et méritis beáti Simónis Confessóris tui, Carméli Montis órdinem, per manus Genitrícis Fílii tui Dómini nostri Jesu Christi, singulári privilégio decorásti~: concéde; ut, ipso interveniénte, ad glóriam, quam diligéntibus te præparásti, perveníre valeámus. Per eúmdem Dóminum.

\subsubsection*{\center \textsc{AD VESPERAS.}}
\begin{tightlist}{}
\item \textred{Añæ Laudum, Psalmi de I. Vesperis Apostolorum. [2].}
\item \textred{Capit. Hymn. et ℣. ut in j. Vesp.}
\item \textred{Ad Magnificat, Aña.} Quæsívit verba utília, et conscrípsit sermónes rectíssimos, et veritáte plenos, allelúja.
\item \textred{Et fit comm. Octavæ S. Simonis, Aña.} Magnificávit Dóminus Sanctum suum, et replévit eum spíritu intelligéntiæ ad dandam sciéntiam salútis Ecclésiæ suæ, allelúja.
\item \textred{℣.}~Ora pro nobis, sancte Simon. \textred{℟.}~Ut digni efficiámur promissiónibus Christi.
\item \textred{Oratio} Deus, qui précibus, \textred{ut supra.}
\end{tightlist}

\begin{center}\textcolor{red}{\rule{\linewidth}{0.5pt}} \end{center}

\subsubsection*{\center \textred{Die XXIII. Maii.}}
\subsection*{\center \textsc{In Octava sancti Simonis.}}
\subsubsection*{\center \textred{Duplex minus.}}
\fancyhead[C]{\textsc{23 Maii - Dies Octava S. Simonis.}}
\begin{tightlist}{}
\item \textred{Omnia ut in Festo 943. præter hic signata.}
\item \textred{Ad Vesp. nisi præcedat Festum duplex, Aña.} Decórem et fortitúdinem índuit eum Dóminus, et accépit stolam glóriæ in consummatióne virtútis, allelúja. \textred{quæ sola dicatur.}
\item \textred{Psalmi de Officio præc.}
\item \textred{In singulis Nocturnis dicitur una tantum Aña.}
\item \textred{In j. Noct. Lect. de Script. occurrente, alias ut in die.}
\end{tightlist}

\subsubsection*{\center \textsc{IN SECUNDO NOCTURNO.}}
%\fancyhead[C]{\textsc{Die XXIII Maii - Dies Octava S. Simonis.}}
\subsubsection*{\center De Sermóne sancti Joánnis Chrysóstomi.}
\subsubsection*{\center \textred{De S. Philogonio. \hfill Lectio iv.}}
\lettrineThree{Q}uantum autem sit pietátis arguméntum, quod Deus illum tanto magistrátu dignum judicárit, audi~: Christus excitátus a morte, quid dicit Petro? postquam enim illum interrogásset~: Petre amas me? isque respondísset~: Tu Dómine scis, quod amo te; dicit illi~: Si díligis me, pasce oves meas. Hoc autem dicébat, non solum ut nobis certíssime declaráret, quantópere ipsum dilígeret Petrus; verum étiam, ut palam fáceret, quantam habéret benevoléntiam erga suas oves. Itaque si quis volúerit illi commendátus esse, curam hábeat óvium illíus públicam quærat utilitátem, fratrum suórum salúti prospíciat.

\begin{tightlist}{}
\item \textred{℟.}~Dum sacrum mystérium ágere paráret, urcéolos ambos invénit aqua plenos, \textred{*} Et benedictióne appósita aqua vinum facta est, allelúja. \textred{℣.}~Memóriam fecit mirabílium suórum omnípotens Deus. \textred{E}t.
\end{tightlist}

\subsubsection*{\center \textred{Lectio v.}}
\lettrine{N}ullum enim offícium hoc Deo cárius est, et ob id alícubi dicit~: Simon, Simon, expetívit te sátanas, ut cribráret te tamquam tríticum~: sed ego rogávi pro te, ne defíciat fides tua. Quid ígitur hujus sollicitúdinis, curæque præmium mihi repéndis? Eámdem curam, quam nobis impéndit, repóscit. Nam et tu, inquit, aliquándo convérsus confírma fratres tuos. Neque prorsus ália res est, quæ perínde decláret, doceátque quis sit fidélis et amans Christi, quam si fratrum curam agat, proque illórum salúte gerat sollicitúdinem.

\begin{tightlist}{}
\item \textred{℟.}~Germánus sancti Patris volens frugálem experiri vitam, coctum offert in cibum piscem, quem sanctus Pater in flumen projécit, \textred{*} Qui vivus factus mox natáre cœpit, allelúja. \textred{℣.}~Mirábilis Deus in Sanctis suis, mirabília operátur, dum piscis coctus vitæ restitúitur. \textred{Q}ui.
\end{tightlist}

\subsubsection*{\center \textred{Lectio vj.}}
\lettrine{H}æc omnes áudiant et mónachi, qui móntium occupárunt cacúmina, quique modis ómnibus seípsos mundo crucifixérunt, quo pro víribus Ecclesiárum præféctos ádjuvent, horúmque curas léniant précibus, concórdia, caritáte, sciéntes quod nisi modis ómnibus opituléntur iis, qui tot perículis objiciúntur, per grátiam Dei, totque negotiórum curas sústinent, quibuscúmque possunt ratiónibus auxiliéntur, quamquam procul habitántes, et victus sors ipsis periísset, et tota in scópulum impegísset sapiéntia.

\begin{tightlist}{}
\item \textred{℟.}~Vestítus innocéntiæ toga, Christo famulári stúduit in innocéntia~: \textred{*} Et ad régimen órdinis vocátus, ipsum ampliáre curávit, allelúja. \textred{℣.}~Ad ómnia Christi mandáta dirigébat aspéctum, et omnem viam iniquitátis ódio hábuit. \textred{E}t. Glória Patri. \textred{E}t.
\end{tightlist}

\subsubsection*{\center \textsc{IN TERTIO NOCTURNO.}}
\subsubsection*{\center Léctio sancti Evangélii secúndum Matthæum.}
\subsubsection*{\center \textred{Lectio vij. \hfill Cap. 19.}}
\lettrine{I}n illo témpore~: Dixit Petrus ad Jesum~: Ecce nos relíquimus ómnia, et secúti sumus te~: quid ergo erit nobis? Et réliqua.

\begin{center}Homilía venerábilis Bedæ Presbyteri.\\
\textred{Hom. in Nat. S. Bened.}\end{center}
\lettrineThree{P}erféctus ille est, qui ábiens vendit ómnia, quæ habet, et dat paupéribus, ac véniens séquitur Christum~: habébit enim thesáurum non deficiéntem in cœlis. Unde bene interrogánte Petro, dixit tálibus Jesus~: Amen dico vobis, quod vos, qui secúti estis me, in regeneratióne, cum séderit Fílius hóminis in sede majestátis suæ, sedébitis et vos super sedes duódecim, judicántes duódecim tribus Israel. In hac quippe vita pro ejus nómine laborántes, in ália præmium speráre dócuit, id est, in regeneratióne; cum vidélicet in vitam immortálem fuérimus resurgéndo regeneráti, qui in vitam cadúcam mortáliter erámus géniti.

\begin{tightlist}{}
\item \textred{℟.}~In pastoráli regímine taléntum sibi créditum magnífice dispensávit~: \textred{*} Cum súbditos verbo et exémplo vitæ alléxit, allelúja. \textred{℣.}~Jucunditátem et exsultatiónem thesaurizávit super eum. \textred{C}um.
\end{tightlist}

\subsubsection*{\center \textred{Lectio viij.}}
\lettrine{E}t justa prorsus retribútio, ut qui hic pro Christo humánæ glóriam celsitúdinis neglexérunt, illic a Christo júdices glorificáti singuláriter cum eo assídeant, qui a sequéndis ejus vestígiis nulla ratióne póterant avélli. Nemo autem putet, duódecim tantum Apóstolos, quia pro Juda prævaricánte, Matthías electus est, tunc esse judicatúros; sicut nec duódecim solæ sunt tribus Israel judicándæ, alióquin tribus Levi, quæ tértia décima est, injudicáta recédet.

\begin{tightlist}{}
\item \textred{℟.}~Ampliávit gentem suam vir sanctus, \textred{*} Et in opéribus suis honorávit Deum, allelúja. \textred{℣.}~Opera sacri Patris velut sol in conspéctu Dei. \textred{E}t.
\end{tightlist}

\subsubsection*{\center \textred{Lectio ix.}}
\lettrine{E}t Paulus, qui tértius décimus est Apóstolus, judicándi sorte privábitur, cum ipse dicat~: Nescítis quóniam Angelos judicábimus, quanto magis sæculária? Sciéndum namque est~: omnes, qui ad exémplum Apostolórum sua reliquérunt ómnia, et secúti sunt Christum, júdices cum eo ventúros, sicut étiam omne mortálium genus esse judicándum~: quia enim duodenário sæpe número solet in Scriptúris univérsitas designári, per duódecim sedes Apostolórum, ómnium numerósitas judicántium, et per duódecim tribus Israel, univérsitas eórum, qui judicándi sunt, osténditur.

\begin{tightlist}{}
\item \textred{℟.}~Dǽmonis insídias non ignórans pius Pater, sic Maríam deprecátus est~: O flos Carméli, vitis florígera, splendor cæli, Virgo puérpera, Mater mitis~: \textred{*} Carmelítis da privilégia, allelúja. \textred{℣.}~Ego servus tuus sum, et fílius ancíllæ tuæ. \textred{C}armelítis. Glória Patri. \textred{C}armelítis.
\item \textred{Vesp. a cap. seq. (Die XXIV Majii, B.M.V. Auxilium Christianorum),} 966.
\item \textred{Comm. præc., Aña.} Magnificávit Dóminus Sanctum suum, et replévit eum spíritu intelligéntiæ ad dandam sciéntiam salútis Ecclésiæ suæ, allelúja.
\item \textred{℣.}~Ora pro nobis, sancte Simon. \textred{℟.}~Ut digni efficiámur promissiónibus Christi, allelúja.
\item Orémus. \hfill \textred{Oratio.}
\end{tightlist}
\lettrine{D}eus, qui précibus et méritis beáti Simónis Confessóris tui, Carméli Montis órdinem, per manus Genitrícis Fílii tui Dómini nostri Jesu Christi, singulári privilégio decorásti~: concéde; ut, ipso interveniénte, ad glóriam, quam diligéntibus te præparásti, perveníre valeámus. Per eúmdem Dóminum.

\begin{center}\textcolor{red}{\rule{\linewidth}{0.5pt}} \end{center}

\subsubsection*{\center \textred{Die XXV. Maii.}}
\subsection*{\center \textsc{S. Mariæ Magdalenæ de Pazzis}}
\subsubsection*{\center \textsc{Virginis Ordinis Nostri.}}
\subsubsection*{\center \textred{Duplex Majus 2. Classis.}}

\begin{center}\textcolor{red}{\rule{0.5\linewidth}{0.5pt}} \end{center}

\subsubsection*{\center \textsc{INFRA OCTAVAM}\\ \textred{S. Mariæ Magdalenæ de Pazzis.}}
\fancyhead[C]{\textsc{Infra Octavam S. Mariæ Magdalenæ de Pazzis.}}
\begin{tightlist}{}
\item \textred{Ad Matut. Invit. Hymn. prima Aña primi Noct. Psalmi de Communi} [118], \textred{℣. et ℟. secundum ordinem dierum, ut in die festo, 974.}
\item \textred{Duæ primæ Lect. de Script. occurr. et tertia, ut infra de Proprio.}
\item Te Deum laudámus. 6.
\item \textred{℣.}~Ora pro nobis sancta María Magdaléna. \textred{℟.}~Ut digni efficiámur promissiónibus Christi.
\item \textred{Ad Laudes, prima Aña Laudum:} Mirábilis in Magdaléna Dóminus, quam e vulnéribus suis decórem índuit, et advérsus dǽmonum insídias fortitúdine præcínxit, allelúja.
\item \textred{Reliq. ut in die. 983.}
\item \textred{Ad Horas ut in die.}
\item \textred{Ad Vesperas, prima Aña Laudum; reliqua ut in ij. Vesp. Festi. 986.}
\item \textred{Dom. infra Oct. fit de Dom. cum commem. Oct. per Añas de utrisque Vesp. et Laud. Festi.}
\item \textred{Adveniente Vigilia Pentecostes Oct. cessat.}
\end{tightlist}

\begin{center}\textcolor{red}{\rule{\linewidth}{0.5pt}} \end{center}

\subsubsection*{\center \textred{Die XXX. Maii.\\ Quinta die infra Octav.}}
\subsection*{\center \textsc{S. Mariæ Magdalenæ de Pazzis.}}
\subsubsection*{\center \textred{Semiduplex.}}
\fancyhead[C]{\textsc{30 Maii - Die V in Octava S. Mariæ Magdalenæ.}}
\subsubsection*{\center De Epístola sancti Hierónymi Presbyteri ad Eustóchium.}
\subsubsection*{\center \textred{De Custod. Virg.\\ Lectio iij.}}
\lettrineThree{A}lia fuit in véteri lege felícitas. Ibi dícitur~: Beátus qui habet semen in Sion, et domésticos in Jerúsalem; et~: Maledícta stérilis, quæ non pariébat; et~: Fílii tui sicut novéllæ olivárum in circúitu mensæ tuæ. Vácuus erat orbis; et, ut de typicis táceam, sola erat benedíctio liberórum. Proptérea et Abraham jam senex Cetúræ copulátur, et Jacob mandrágoris redímitur, et conclúsam vulvam in Ecclésiæ figúram Rachel conquéritur. Paulátim vero increscénte ségete messor immíssus est. Virgo Elías, Eliséus virgo, vírgines multi fílii Prophetárum. Jeremíæ dícitur~: Et tu ne accípias uxórem. Sanctificátus in útero, captivitáte propínqua, uxórem prohibétur accípere. Aliis verbis idípsum Apóstolus lóquitur~: Exístimo hoc bonum esse propter instántem necessitátem. Quæ est ista necéssitas, quæ aufert gáudia nuptiárum? Tempus breviátum est, in próximo erit Nabuchodónosor.

\begin{tightlist}{}
\item \textred{Vesp. de seq. (Dies XXXI Maji - B.M.V. Omnium Gratiarum Mediatricis).}
\item \textred{Psal. et com. præc.}
\end{tightlist}

\begin{center}\textcolor{red}{\rule{\linewidth}{0.5pt}} \end{center}

\subsubsection*{\center \textred{Die I. Junii.}}
\subsection*{\center \textsc{In Octava S. Mariæ Magdalenæ de Pazzis.}}
\subsubsection*{\center \textred{Virginis.\\ Duplex minus.}}
\fancyhead[C]{\textsc{1 Junii - Dies Octava S. Mariæ Magdalenæ.}}

\begin{tightlist}{}
\item \textred{Omnia ut in Festo} 974. \textred{exceptis hic signatis.}
\item \textred{Ad Vesp. si non præced. Officium duplex, Aña.} Mirábilis in Magdaléna Dóminus, quam e vulnéribus suis decórem índuit, et advérsus dǽmonum insídias fortitúdine præcínxit, allelúja. 
\item \textred{Psalmi de Off. præc.}
\item \textred{Ad Matutinum in singulis Nocturnis dicitur una tantum Aña. Lectiones primi Nocturni de Scriptura occurrente.}
\item \textred{Alias ut in die~: in secundo et tertio Nocturno.}
\item \textred{Lectiones ut in ejus Officio menstruo, infra.}
\item \textred{In Vesp. com. seq.} 1007.
\end{tightlist}

\begin{center}\textcolor{red}{\rule{\linewidth}{0.5pt}} \end{center}
\clearpage

%\begin{center}\textcolor{red}{\rule{\linewidth}{0.5pt}} \end{center}
\pagenumbering{arabic}
\renewcommand{\thepage}{[\arabic{page}]*}

\section*{\center \textsc{\textred{Officia Frequentius Persolvenda}}}
%\begin{center}\textcolor{red}{\rule{\linewidth}{0.5pt}} \end{center}
\subsection*{\center \textsc{\textred{Officia Hebdomadales}}}
\subsubsection*{\center \textred{Dicenda singulis hebomadibus}}
\begin{center}\textcolor{red}{\rule{\linewidth}{0.5pt}} \end{center}

\subsection*{\center \textsc{Officium Sanctissimi Sacramenti.}}
\subsubsection*{\center \textred{Semiduplex.}}
\begin{center}\textsc{recitandum semel qualibet hebdomada.}\end{center}
\subsubsection*{\center \textsc{AD VESPERAS.}}
\fancyhead[C]{\textsc{Officia Hebdomadales - Feria V - SS. Sacramenti}}
\subsubsection*{\textred{Capitulum. \hfill 1~Cor.~11}}
\lettrineThree{D}ominus Jesus, in qua nocte tradebátur accépit panem, et grátias agens fregit, et dixit~: Accípite, et manducáte~: hoc est Corpus meum, quod pro vobis tradétur.

\subsubsection*{\center\textred{Hymnus.}}
\lettrine{P}ange lingua gloriósi\\ 
Córporis mystérium,\\ 
Sanguinísque pretiósi,\\ 
Quem in mundi prétium\\ 
Fructus ventris generósi\\ 
Rex effúdit Géntium.\\
\\
\textred{N}obis datus, nobis natus\\ 
Ex intácta Vírgine,\\ 
Et in mundo conversátus,\\ 
Sparso verbi sémine,\\ 
Sui moras incolátus\\ 
Miro clausit órdine.\\ 
\\
\textred{I}n suprémæ nocte cœnæ\\ 
Recúmbens cum frátribus,\\ 
Observáta lege plene\\ 
Cibis in legálibus,\\ 
Cibum turbæ duodénæ\\ 
Se dat suis mánibus.\\
\\
\textred{V}erbum caro, panem verum\\ 
Verbo carnem éfficit~:\\ 
Fitque sanguis Christi merum;\\ 
Et si sensus déficit,\\ 
Ad firmándum cor sincérum\\ 
Sola fides súfficit.\\
\\
\textred{T}antum ergo Sacraméntum\\ 
Venerémur cérnui~:\\ 
Et antíquum documéntum\\ 
Novo cedat rítui~:\\ 
Prǽstet fides suppleméntum\\ 
Sénsuum deféctui.\\
\\
\textred{G}enitóri, Genitóque\\ 
Laus, et jubilátio,\\ 
Salus, honor, virtus quoque\\ 
Sit et benedíctio,\\ 
Procedénti ab utróque\\ 
Compar sit laudátio.\\
Amen.

\begin{tightlist}{}
\item \textred{℣.}~Panem de cœlo præstitísti eis. \textred{℟.}~Omne delectaméntum in se habéntem.
\item \textred{Ad Magnificat, Aña.} O quam suávis est Dómine Spíritus tuus, qui ut dulcédinem tuam in fílios demonstráres, pane suavíssimo de cœlo prǽstito, esuriéntes reples bonis, fastidiósos dívites dimíttens inánes!
\end{tightlist}
\subsubsection*{\center \textred{Oratio.}}
\lettrine{D}eus, qui nobis sub Sacraménto mirábili passiónis tuæ memóriam reliquísti~: tríbue, quǽsumus, ita nos córporis, et sánguinis tui sacra mystéria venerári, ut redemptiónis tuæ fructum in nobis júgiter sentiámus. Qui vivis.

\begin{tightlist}{}
\item \textred{Ad Completorium in fine Hymni dicitur ℣.}~Glória tibi Dómine, qui natus es de Vírgine, \textred{et similiter ad Horas.}
\item \textred{Ad Nunc dimittis, Aña.} Allelúja, Panis~: \textred{ut in Psalt. Post Septuagesima dicitur eadem Aña. omisso} Allelúja.
\end{tightlist}

\subsubsection*{\center \textsc{AD MATUTINUM.}}
\begin{tightlist}{}
\item \textred{Invitatorium.} Christum Regem adorémus dominántem Géntibus~: \textred{*} Qui se manducántibus dat spíritus pinguédinem.
\item \textred{Ps.} Veníte, exsultémus. 2.
\end{tightlist}

\subsubsection*{\center\textred{Hymnus.}}
\lettrine{S}acris solémniis juncta sint gáudia,\\ Et ex præcórdiis sonent præcónia;\\ Recédant vétera, nova sint ómnia,\\ Corda, voces, et ópera.\\
\\
\textred{N}octis recólitur cœna novíssima,\\ Qua Christus créditur agnum, et ázyma\\ Dedísse frátribus, juxta legítima\\ Priscis indúlta pátribus.\\
\\
\textred{P}ost agnum týpicum, explétis épulis,\\ Corpus Domínicum datum discípulis,\\ Sic totum ómnibus, quod totum síngulis,\\ Ejus fatémur mánibus.\\
\\
\textred{D}edit fragílibus córporis férculum,\\ Dedit et trístibus sánguinis póculum,\\ Dicens~: Accípite quod trado vásculum,\\ Omnes ex eo bíbite.\\
\\
\textred{S}ic sacrifícium istud instítuit,\\ Cujus offícium commítti vóluit\\ Solis Presbýteris, quibus sic cóngruit,\\ Ut sumant, et dent céteris.\\
\\
\textred{P}anis Angélicus fit panis hóminum~:\\ Dat panis cǽlicus figúris términum~:\\ O res mirábilis! mandúcat Dóminum\\ Pauper servus, et húmilis.\\
\\
\textred{T}e trina Déitas, únaque, póscimus,\\ Sic nos tu vísita, sicut te cólimus,\\ Per tuas sémitas duc nos quo téndimus,\\ Ad lucem, quam inhábitas.\\
Amen.
\subsubsection*{\center \textsc{IN PRIMO NOCTURNO.}}
\textred{Aña.} Fructum salutíferum.
\subsubsection*{\center Psalmus j.}
\lettrineThree{B}eátus vir, qui non ábiit in consílio impiórum, et in via peccatórum non stetit, * et in cáthedra pestiléntiæ non sedit~:

Sed in lege Dómini volúntas ejus, * et in lege ejus meditábitur die ac nocte.

Et erit tamquam lignum, quod plantátum est secus decúrsus aquárum, * quod fructum suum dabit in témpore suo:

Et fólium ejus non défluet~: * et ómnia quæcúmque fáciet, prosperabúntur.

Non sic ímpii, non sic~: * sed tamquam pulvis, quem prójicit ventus a fácie terræ.

Ideo non resúrgent ímpii in judício~: * neque peccatóres in concílio justórum.

Quóniam novit Dóminus viam justórum~: * et iter impiórum períbit.

Glória Patri.

\begin{tightlist}{}
\item \textred{Aña.} Fructum salutíferum gustándum dedit Dóminus mortis suæ témpore.
\item \textred{Aña.} A fructu fruménti.
\end{tightlist}

\subsubsection*{\center Psalmus iv.}
\lettrine{C}um invocárem exaudívit me Deus justítiæ meæ~: * in tribulatióne dilatásti mihi.

Miserére mei, * et exáudi oratiónem meam.

Fílii hóminum, úsquequo gravi corde? * ut quid dilígitis vanitátem, et quǽritis mendácium?

Et scitóte quóniam mirificávit Dóminus sanctum suum~: * Dóminus exáudiet me, cum clamávero ad eum.

Irascímini, et nolíte peccáre~: * quæ dícitis in córdibus vestris, in cubílibus vestris compungímini.

Sacrificáte sacrifícium justítiæ, et speráte in Dómino. * Multi dicunt~: Quis osténdit nobis bona?

Signátum est super nos lumen vultus tui Dómine~: * dedísti lætítiam in corde meo.

A fructu fruménti, vini, et ólei sui * multiplicáti sunt.

In pace in idípsum * dórmiam, et requiéscam.

Quóniam tu Dómine singuláriter in spe * constituísti me.

Glória Patri.

\begin{tightlist}{}
\item \textred{Aña.} A fructu fruménti, et vini multiplicáti fidéles, in pace Christi requiéscunt.
\item \textred{Aña.} Communióne cálicis.
\end{tightlist}
\subsubsection*{\center Psalmus xv.}
\lettrine{C}onsérva me Dómine, quóniam sperávi in te. * Dixi Dómino~: Deus meus es tu, quóniam bonórum meórum non eges.

Sanctis, qui sunt in terra ejus~: * mirificávit omnes voluntátes meas in eis.

Multiplicátæ sunt infirmitátes eórum~: * póstea acceleravérunt.

Non congregábo conventícula eórum de sanguínibus~: * nec memor ero nóminum eórum per lábia mea.

Dóminus pars hereditátis meæ, et cálicis mei~: * tu es, qui restítues hereditátem meam mihi.

Funes cecidérunt mihi in præcláris~: * étenim heréditas mea præclára est mihi.

Benedícam Dóminum, qui tríbuit mihi intelléctum~: * ínsuper et usque ad noctem increpuérunt me renes mei.

Providébam Dóminum in conspéctu meo semper~: * quóniam a dextris est mihi, ne commóvear.

Propter hoc lætátum est cor meum, et exsultávit lingua mea~: * ínsuper et caro mea requiéscet in spe.

Quóniam non derelínques ánimam meam in inférno~: * nec dabis sanctum tuum vidére corruptiónem.

Notas mihi fecísti vias vitæ, adimplébis me lætítia cum vultu tuo~: * delectatiónes in déxtera tua usque in finem.

Gloria Patri.

\begin{tightlist}{}
\item \textred{Aña.} Communióne cálicis, quo Deus ipse súmitur, non vitulórum sánguine, congregávit nos Dóminus.
\item \textred{℣.}~Panem cœli dedit eis. \textred{℟.}~Panem Angelórum manducávit homo.
\item \textred{Lectiones de Scriptura occurrente; si desint, quæ sequuntur.}
\end{tightlist}

\subsubsection*{\center De Epístola prima beáti Pauli Apóstoli ad Corínthios.}
\subsubsection*{\center \textred{Lectio j. \hfill Cap.~11,~℣.~20.}}
\lettrineThree{C}onveniéntibus ergo vobis in unum, jam non est Domínicam cœnam manducáre. Unusquísque enim suam cœnam præsúmit ad manducándum. Et álius quidem ésurit~: álius autem ébrius est. Numquid domos non habétis ad manducándum, et bibéndum? aut Ecclésiam Dei contémnitis, et confúnditis eos, qui non habent? Quid dicam vobis? Laudo vos? in hoc non laudo.

\begin{tightlist}{}
\item \textred{℟.}~Immolábit hœdum multitúdo filiórum Israel ad vésperam Paschæ~: \textred{*} Et edent carnes, et ázymos panes. \textred{℣.}~Pascha nostrum immolátus est Christus~: ítaque epulémur in ázymis sinceritátis, et veritátis. \textred{E}t edent.
\end{tightlist}

\subsubsection*{\center\textred{Lectio ij.}}
\lettrine{E}go enim accépi a Dómino quod et trádidi vobis, quóniam Dóminus Jesus in qua nocte tradebátur, accépit panem, et grátias agens fregit, et dixit~: Accípite, et manducáte~: hoc est corpus meum, quod pro vobis tradétur~: hoc fácite in meam commemoratiónem. Simíliter et cálicem, postquam cœnávit, dicens~: Hic calix novum testaméntum est in meo sánguine~: hoc fácite quotiescúmque bibétis, in meam commemoratiónem. Quotiescúmque enim manducábitis panem hunc, et cálicem bibétis~: mortem Dómini annuntiábitis donec véniat.

\begin{tightlist}{}
\item \textred{℟.}~Comedétis carnes, et saturabímini pánibus. \textred{*} Iste est panis, quem Dóminus dedit vobis ad vescéndum. \textred{℣.}~Non Móyses dedit vobis panem de cœlo, sed Pater meus dat vobis panem de cœlo verum. \textred{I}ste.
\end{tightlist}

\subsubsection*{\center\textred{Lectio iij.}}
\lettrine{I}taque quicúmque manducáverit panem hunc, vel bíberit cálicem Dómini indígne~: reus erit córporis, et sánguinis Dómini. Probet autem seípsum homo~: et sic de pane illo edat, et de cálice bibat. Qui enim mandúcat, et bibit indígne, judícium sibi mandúcat, et bibit; non dijúdicans corpus Dómini. Ideo inter vos multi infírmi et imbecílles, et dórmiunt multi. Quod si nosmetípsos dijudicarémus, non útique judicarémur. Dum judicámur autem, a Dómino corrípimur, ut non cum hoc mundo damnémur.

\begin{tightlist}{}
\item \textred{℟.}~Respéxit Elías ad caput suum subcinerícium panem, qui surgens comédit, et bibit, \textred{*} Et ambulávit in fortitúdine cibi illíus usque ad montem Dei. \textred{℣.}~Si quis manducáverit ex hoc pane, vivet in ætérnum. \textred{E}t ambulávit. Glória Patri. \textred{E}t ambulávit.
\end{tightlist}

\subsubsection*{\center \textsc{IN SECUNDO NOCTURNO}}
\textred{Aña.} Memor sit Dóminus.
\subsubsection*{\center Psalmus xix.}
\lettrineThree{E}xáudiat te Dóminus in die tribulatiónis~: * prótegat te nomen Dei Jacob.

Mittat tibi auxílium de sancto~: * et de Sion tueátur te.

Memor sit omnis sacrifícii tui~: * et holocáustum tuum pingue fiat.

Tríbuat tibi secúndum cor tuum~: * et omne consílium tuum confírmet.

Lætábimur in salutári tuo~: * et in nómine Dei nostri magnificábimur.

Impleat Dóminus omnes petitiónes tuas~: * nunc cognóvi quóniam salvum fecit Dóminus Christum suum.

Exáudiet illum de cælo sancto suo~: * in potentátibus salus déxteræ ejus.

Hi in cúrribus, et hi in equis~: * nos autem in nómine Dómini Dei nostri invocábimus.

Ipsi obligáti sunt, et cecidérunt~: * nos autem surréximus et erécti sumus.

Dómine salvum fac regem~: * et exáudi nos in die, qua invocavérimus te.

Glória Patri.

\begin{tightlist}{}
\item \textred{Aña} Memor sit Dóminus omnis sacrifícii nostri, et holocáustum nostrum pingue fiat.
\item \textred{Aña.} Parátur nobis.
\end{tightlist}

\subsubsection*{\center Psalmus xxij.}
\lettrine{D}óminus regit me, et nihil mihi déerit~: * in loco páscuæ ibi me collocávit.

Super aquam refectiónis educávit me~: * ánimam meam convértit.

Dedúxit me super sémitas justítiæ~: * propter nomen suum.

Nam, et si ambulávero in médio umbræ mortis, non timébo mala~: * quóniam tu mecum es.

Virga tua, et báculus tuus~: * ipsa me consoláta sunt.

Parásti in conspéctu meo mensam, * advérsus eos, qui tríbulant me.

Impinguásti in óleo caput meum~: * et calix meus inébrians quam præclárus est!

Et misericórdia tua subsequétur me * ómnibus diébus vitæ meæ.

Et ut inhábitem in domo Dómini, * in longitúdinem diérum.

Glória Patri.

\begin{tightlist}{}
\item \textred{Aña.} Parátur nobis mensa Dómini, advérsus omnes qui tríbulant nos.
\item \textred{Aña.} In voce exsultatiónis.
\end{tightlist}

\subsubsection*{\center Psalmus xlj.}
\lettrine{Q}uemádmodum desíderat cervus ad fontes aquárum~: * ita desíderat ánima mea ad te, Deus.

Sitívit ánima mea ad Deum fortem vivum~: * quando véniam, et apparébo ante fáciem Dei?

Fuérunt mihi lácrimæ meæ panes die ac nocte~: * dum dícitur mihi quotídie~: Ubi est Deus tuus?

Hæc recordátus sum, et effúdi in me ánimam meam~: * quóniam transíbo in locum tabernáculi admirábilis, usque ad domum Dei.

In voce exsultatiónis, et confessiónis~: * sonus epulántis.

Quare tristis es, ánima mea? * et quare contúrbas me?

Spera in Deo, quóniam adhuc confitébor illi~: * salutáre vultus mei, et Deus meus.

Ad meípsum ánima mea conturbáta est~: * proptérea memor ero tui de terra Jordánis, et Hermóniim a monte módico.

Abýssus abýssum ínvocat, * in voce cataractárum tuárum.

Omnia excélsa tua, et fluctus tui * super me transiérunt.

In die mandávit Dóminus misericórdiam suam~: * et nocte cánticum ejus.

Apud me orátio Deo vitæ meæ, * dicam Deo~: Suscéptor meus es.

Quare oblítus es mei? * et quare contristátus incédo, dum afflígit me inimícus?

Dum confringúntur ossa mea, * exprobravérunt mihi qui tríbulant me inimíci mei.

Dum dicunt mihi per síngulos dies~: Ubi est Deus tuus? * quare tristis es, ánima mea? et quare contúrbas me?

Spera in Deo, quóniam adhuc confitébor illi~: * salutáre vultus mei, et Deus meus.

Glória Patri.

\begin{tightlist}{}
\item \textred{Aña} In voce exsultatiónis résonent epulántes in mensa Dómini.
\item \textred{℣.}~Cibávit eos ex ádipe fruménti. \textred{℟.}~Et de petra melle saturávit eos.
\end{tightlist}

\subsubsection*{\center Sermo sancti Cyrílli Epíscopi Jerosolymitáni.}
\subsubsection*{\center \textred{Lectio iv. \hfill Catechesi myst. 4.}}
\lettrineThree{I}psa beáti Pauli doctrína abúnde suffícere vidétur, ut certam vobis de divínis mystériis fidem fáciat~: quibus digni rédditi, concorpórei, ut ita dicam, et consanguínei Christi facti estis. Ipse enim modo clamábat, quod in nocte, qua tradebátur Dóminus noster Jesus Christus, accípiens panem et grátias agens, fregit, et dedit discípulis suis, dicens~: Accípite et manducáte, hoc est corpus meum. Et accípiens cálicem, et grátias agens, dixit~: Accípite, et bíbite~: hic est sanguis meus. Cum ígitur ipse de pane pronuntiáverit, ac díxerit~: Hoc est corpus meum~: quis audébit deínceps ambígere; et cum idem ipse tam asseveránter díxerit~: Hic est sanguis meus~: quis unquam dubitáverit, ut dicat non esse ejus sánguinem?

\begin{tightlist}{}
\textred{℟.}~Panis, quem ego dabo, caro mea est pro mundi vita. Litigábant ergo Judǽi ad ínvicem, dicéntes~: \textred{*} Quómodo potest hic nobis carnem suam dare ad manducándum? \textred{℣.}~Locútus est pópulus contra Dóminum~: Anima nostra náuseat super cibo isto levíssimo. \textred{Q}uómodo.
\end{tightlist}

\subsubsection*{\center \textred{Lectio v.}}
\lettrine{A}quam olim in vinum convértit in Cana Galilǽæ~: quod habet quamdam cum sánguine propinquitátem~: et eum parum dignum existimábimus, cui credámus quod vinum in sánguinem transmutárit? Ad eas núptias, quibus córpora copulántur, vocátus, præter opiniónem ómnium, hoc fecit miráculum~: et non multo magis sic eum corpus et sánguinem suum fruénda nobis donásse, persuásum fírmiter habébimus; ut ea cum omni certitúdine tamquam corpus ipsíus, et sánguinem sumámus? nam in spécie panis dat nobis corpus, et in spécie vini dat nobis sánguinem~: ut cum súmpseris, gustes corpus, et sánguinem Christi, factus ejúsdem córporis et sánguinis párticeps. Sic enim effícimur Christíferi, hoc est, Christum in corpóribus nostris feréntes, cum corpus ejus et sánguinem in membra nostra recípimus; sic secúndum beátum Petrum, divínæ natúræ consórtes réddimur. 

\begin{tightlist}{}
\item \textred{℟.}~Cœnántibus eis, accépit Jesus panem, et benedíxit, ac fregit, dedítque discípulis suis, et ait~: \textred{*} Accípite et comédite~: hoc est corpus meum. \textred{℣.}~Dixérunt viri tabernáculi mei~: Quis det de cárnibus ejus, ut saturémur? \textred{A}ccípite.
\end{tightlist}

\subsubsection*{\center \textred{Lectio vj.}}
\lettrine{O}lim cum Judǽis Christus dísserens, nisi manducavéritis, inquit, carnem meam, et bibéritis meum sánguinem, non habébitis vitam in vobis. Cum autem illi, quæ dicta fúerant, non spiritáliter accepíssent, offénsi abiérunt retro~: putábant enim, quod eos ad manducándas carnes hortarétur. Erant et in véteri testaménto panes propositiónis; verum illi cum fúerint véteris testaménti, finem jam accepérunt. In novo vero testaménto, panis est cœléstis et calix salutáris~: qui et ánimam et corpus sanctíficant. Quámobrem non sic hæc atténdas velim, tamquam sint nudus et simplex panis, nudum et simplex vinum~: corpus enim sunt et sanguis Christi. Nam etiámsi sensus illud tibi renúntiat, fides tamen te confírmet. Ne júdices rem ex gustu, sed te citra ullam dubitatiónem fides certum reddat, quod sis dignus factus qui córporis et sánguinis Christi párticeps fíeres.

\begin{tightlist}{}
\item \textred{℟.}~Accépit Jesus cálicem, postquam cœnávit, dicens~: Hic calix novum testaméntum est in meo sánguine~: \textred{*} Hoc fácite in meam commemoratiónem. \textred{℣.}~Memória memor ero, et tabéscet in me ánima mea. \textred{H}oc fácite. Glória Patri. \textred{H}oc fácite.
\end{tightlist}

\subsubsection*{\center \textsc{IN TERTIO NOCTURNO.}}
\begin{tightlist}{}
\item \textred{Aña.} Introíbo.
\end{tightlist}
\subsubsection*{\center Psalmus xlij.}
\lettrineThree{J}údica me, Deus, et discérne causam meam de gente non sancta, * ab hómine iníquo, et dolóso érue me.

Quia tu es Deus fortitúdo mea~: * quare me repulísti? et quare tristis incédo, dum afflígit me inimícus?

Emítte lucem tuam, et veritátem tuam~: * ipsa me deduxérunt, et adduxérunt in montem sanctum tuum, et in tabernácula tua.

Et introíbo ad altáre Dei~: * ad Deum, qui lætíficat juventútem meam.

Confitébor tibi in cíthara Deus Deus meus~: * Quare tristis es, ánima mea? et quare contúrbas me?

Spera in Deo, quóniam adhuc confitébor illi~: * salutáre vultus mei, et Deus meus.

Glória Patri.

\begin{tightlist}{}
\item \textred{Aña.} Introíbo ad altáre Dei~: sumam Christum, qui rénovat juventútem meam.
\item \textred{Aña.} Cibávit nos.
\end{tightlist}

\subsubsection*{\center Psalmus lxxx.}
\lettrine{E}xsultáte Deo, adjutóri nostro~: * jubiláte Deo Jacob.

Súmite psalmum, et date týmpanum~: * psaltérium jucúndum cum cíthara.

Buccináte in Neoménia tuba, * in insígni die solemnitátis vestræ.

Quia præcéptum in Israël est~: * et judícium Deo Jacob.

Testimónium in Joseph pósuit illud, cum exíret de terra Ægýpti~: * linguam, quam non nóverat, audívit.

Divértit ab onéribus dorsum ejus~: * manus ejus in cóphino serviérunt.

In tribulatióne invocásti me, et liberávi te~: exaudívi te in abscóndito tempestátis~: * probávi te apud aquam contradictiónis.

Audi, pópulus meus, et contestábor te~: * Israël, si audíeris me, non erit in te deus recens, neque adorábis deum aliénum.

Ego enim sum Dóminus Deus tuus, qui edúxi te de terra Ægýpti~: * diláta os tuum, et implébo illud.

Et non audívit pópulus meus vocem meam~: * et Israël non inténdit mihi.

Et dimísi eos secúndum desidéria cordis eórum~: * ibunt in adinventiónibus suis.

Si pópulus meus audísset me~: * Israël si in viis meis ambulásset:

Pro níhilo fórsitan inimícos eórum humiliássem~: * et super tribulántes eos misíssem manum meam.

Inimíci Dómini mentíti sunt ei~: * et erit tempus eórum in sǽcula.

Et cibávit eos ex ádipe fruménti~: * et de petra, melle saturávit eos.

Glória Patri.

\begin{tightlist}{}
\item \textred{Aña.} Cibávit nos Dóminus ex ádipe fruménti~: et de petra, melle saturávit nos.
\item \textred{Aña.} Ex altári tuo.
\end{tightlist}

\subsubsection*{\center Psalmus lxxxiij.}
\lettrine{Q}uam dilécta tabernácula tua, Dómine virtútum~: * concupíscit, et déficit ánima mea in átria Dómini.

Cor meum, et caro mea~: * exsultavérunt in Deum vivum.

Etenim passer invénit sibi domum~: * et turtur nidum sibi, ubi ponat pullos suos.

Altária tua, Dómine virtútum~: * Rex meus, et Deus meus.

Beáti, qui hábitant in domo tua, Dómine~: * in sǽcula sæculórum laudábunt te.

Beátus vir, cujus est auxílium abs te~: * ascensiónes in corde suo dispósuit, in valle lacrimárum in loco, quem pósuit.

Etenim benedictiónem dabit legislátor, ibunt de virtúte in virtútem~: * vidébitur Deus deórum in Sion.

Dómine, Deus virtútum, exáudi oratiónem meam~: * áuribus pércipe, Deus Jacob.

Protéctor noster, áspice, Deus~: * et réspice in fáciem Christi tui.

Quia mélior est dies una in átriis tuis~: * super míllia.

Elégi abjéctus esse in domo Dei mei~: * magis quam habitáre in tabernáculis peccatórum.

Quia misericórdiam et veritátem díligit Deus~: * grátiam et glóriam dabit Dóminus.

Non privábit bonis eos qui ámbulant in innocéntia~: * Dómine virtútum, beátus homo qui sperat in te.

Glória Patri.

\begin{tightlist}{}
\item \textred{Aña.} Ex altári tuo Dómine, Christum súmimus~: in quem cor, et caro nostra exsúltant.
\item \textred{℣.}~Edúcas panem de terra. \textred{℟.}~Et vinum lætíficet cor hóminis.
\end{tightlist}

\subsubsection*{\center Léctio sancti Evangélii secúndum Joánnem.}
\subsubsection*{\center \textred{Lectio vij. \hfill Cap. 6.}}
\lettrine{I}n illo témpore~: Dixit Jesus turbis Judæórum~: Caro mea vere est cibus, et sanguis meus vere est potus. Et réliqua.

\begin{center}Homilía sancti Cyrílli Epíscopi Alexandríni.\\
\textred{Lib. 4. in Joan. cap. 17.} \end{center}
\lettrineThree{Q}ui mandúcat, inquit, carnem meam et bibit sánguinem meum, in me manet, et ego in illo. Sícuti enim si quis liquefáctæ ceræ áliam ceram infúderit, álteram cum áltera per totum commísceat necésse est; sic qui carnem et sánguinem Dómini récipit, cum ipso ita conjúngitur ut Christus in ipso, et ipse in Christo inveniátur. Símile quodámmodo apud Matthǽum compéries~: símile est, inquit, regnum cœlórum ferménto, quod accéptum abscóndit múlier in farínæ satis tribus. Sicut parum, ut Paulus ait, ferménti totam massam ferméntat; sic párvula benedíctio totum hóminem in seípsam áttrahit, et sua grátia replet~: et hoc modo in nobis Christus manet, et nos in Christo.

\begin{tightlist}{}
\item \textred{℟.}~Melchísedech vero rex Salem, próferens panem et vinum, erat enim sacérdos Dei altíssimi, \textred{*} Benedíxit Abram, et ait. \textred{℣.}~Benedíctus Abram Deo excélso, qui creávit cœlum, et terram. \textred{B}enedíxit.
\end{tightlist}

\subsubsection*{\center \textred{Lectio viij.}}
\lettrine{N}os vero, si vitam ætérnam cónsequi vólumus, si largitórem immortalitátis habére in nobis desiderámus, ad recipiéndam benedictiónem libénter concurrámus~: caveamúsque ne loco láquei damnósam religiónem diábolus nobis præténdat. Recte (inquit) dicis~: scriptum tamen esse non ignorámus, judícium sibi comédere atque bíbere illum, qui de pane cómedit, et de cálice bibit indígne. Ego ígitur probo meípsum, et indígnum invénio. Quando ígitur, quicúmque tu es, qui ista dicis, dignus eris? quando Christo teípsum ófferes? Nam si peccándo indígnus es, et peccáre non désinis (quis enim delícta intélligit? secúndum Psalmístam) expers omníno eris vivíficæ hujus sanctificatiónis.

\begin{tightlist}{}
\item \textred{℟.}~Calix benedictiónis, cui benedícimus, nonne communicátio sánguinis Christi est? \textred{*} Et panis, quem frángimus, nonne participátio córporis Christi est? \textred{℣.}~Quóniam unus panis, et unum corpus multi sumus omnes, qui de uno pane, et de uno cálice participámus. \textred{E}t panis.
\end{tightlist}

\subsubsection*{\center \textred{Lectio ix.}}
\lettrine{Q}uare pias quæso cogitatiónes sucípias, studióse sanctéque vivas, et benedictiónem partícipes~: quæ (mihi crede) non mortem solum, verum étiam morbos omnes depéllit. Sedat enim, cum in nobis máneat Christus, sæviéntem membrórum nostrórum legem, pietátem corróborat, perturbatiónes ánimi exstínguit, ægrótos curat, collísos redíntegrat~: et sicut Pastor bonus, qui ánimam suam pro óvibus pósuit, ab omni nos érigit casu.

\begin{tightlist}{}
\item \textred{℟.}~Unus panis, et unum corpus multi sumus \textred{*} Omnes, qui de uno pane, et de uno cálice participámus. 
\item \textred{℣.}~Parásti in dulcédine tua páuperi Deus, qui habitáre facis unánimes in domo. \textred{O}mnes. Glória Patri. \textred{O}mnes.
\item Te Deum laudámus. 6.
\item \textred{Sacerdotalis ℣.}~Pluit illis manna ad manducándum. \textred{℟.}~Et panem cœli dedit eis.
\end{tightlist}

\subsubsection*{\center \textsc{AD LAUDES.}}
\begin{center}\textred{et per Horas, Aña.}\end{center}
\lettrine{S}apiéntia ædificávit sibi domum, míscuit vinum, et pósuit mensam, allelúja.

\begin{tightlist}{}
\item \textred{Ps. de Dominica.} 27.
\item \textred{Aña.} Angelórum esca nutrivísti pópulum tuum, et panem de cœlo præstitísti eis, allelúja.
\item \textred{Aña.} Pinguis est panis Christi, et præbébit delícias régibus, allelúja.
\item \textred{Aña.} Sacerdótes sancti incénsum, et panes ófferunt Deo, allelúja.
\item \textred{Aña.} Vincénti dabo manna abscónditum, et nomen novum, allelúja.
\end{tightlist}

\subsubsection*{\textred{Capitulum. \hfill 1 Cor. 11.}}
\lettrineThree{D}óminus Jesus in qua nocte tradebátur, accépit panem, et grátias agens fregit, et dixit~: Accípite, et manducáte, hoc est corpus meum quod pro vobis tradétur.

\subsubsection*{\center \textred{Hymnus.}}
\begin{center}\textred{Temp. Pasch. nihil additur, neque mutatur.}\end{center}
\lettrine{V}erbum supérnum pródiens,\\
Nec Patris linquens déxteram.\\
Ad opus suum éxiens\\
Venit ad vitæ vésperam.\\
\\
\textred{I}n mortem a discípulo\\
Suis tradéndus ǽmulis,\\
Prius in vitæ férculo\\
Se trádidit discípulis.\\
\\
\textred{Q}uibus sub bina spécie\\
Carnem dedit et sánguinem~:\\
Ut dúplicis substántiæ\\
Totum cibáret hóminem.\\
\\
\textred{S}e nascens dedit sócium,\\
Convéscens in edúlium,\\
Se móriens in prétium,\\
Se regnans dat in prǽmium.\\
\\
\textred{O} salutáris Hóstia~:\\
Quæ cœli pandis óstium,\\
Bella premunt hostília,\\
Da robur, fer auxílium.\\
\\
\textred{U}ni, trinóque Dómino\\
Sit sempitérna glória~:\\
Qui vitam sine término\\
Nobis donet in pátria.\\
Amen.

\begin{tightlist}{}
\item \textred{℣.}~Pósuit fines tuos pacem. \textred{℟.}~Et ádipe fruménti sátiat te.
\item \textred{Ad Benedictus, Aña.} Ego sum panis vivus, qui de cœlo descéndi. Si quis manducáverit ex hoc pane, vivet in ætérnum, allelúja. 
\end{tightlist}

\subsubsection*{\center \textred{Oratio.}}
\lettrine{D}eus, qui nobis sub Sacraménto mirábili, passiónis tuæ memóriam reliquísti~: tríbue, quǽsumus, ita nos Córporis et Sánguinis tui sacra mystéria venerári, ut redemptiónis tuæ fructum in nobis júgiter sentiámus. Qui vivis.

\begin{tightlist}{}
\item \textred{Ad Horas Psalmi de Dominica, ad Primam tamen ut in Festis.}
\end{tightlist}

\subsubsection*{\center \textsc{AD PRIMAM.}}
\noindent \textred{℣.}~Qui corpus tuum nobis dedísti in cibum, et sánguinem in potum.

\subsubsection*{\center \textsc{AD TERTIAM.}}
\subsubsection*{\textred{Capitulum. \hfill 1 Cor. 11.}}
\lettrine{D}óminus Jesus in qua nocte tradebátur, accépit panem, et grátias agens fregit, et dixit~: Accípite, et manducáte, hoc est corpus meum quod pro vobis tradétur.

\begin{tightlist}{}
%\item \textred{Capitulum.} Dóminus Jesus, \textred{ut supra.}
\item \textred{℟.}~Panem cœli \textred{*} Dedit eis. \textred{P}anem. \textred{℣.}~Panem Angelórum manducávit homo. \textred{D}edit. Glória Patri. \textred{P}anem.
\item \textred{℣.}~Cibávit eos ex ádipe fruménti. \textred{℟.}~Et de petra, melle saturávit eos.
\end{tightlist}

\subsubsection*{\center \textsc{AD SEXTAM.}}
\subsubsection*{\textred{Capitulum. \hfill 1~Cor.~11.}}
\lettrine{Q}uotiescúmque manducábitis panem hunc, et cálicem bibétis, mortem Dómini annuntiábitis, donec véniat. 

\begin{tightlist}{}
\item \textred{℟.}~Cibávit eos \textred{*} Ex ádipe fruménti. \textred{C}ibávit. \textred{℣.}~Et de petra, melle saturávit eos. \textred{E}x. Glória Patri. \textred{C}ibávit.
\item \textred{℣.}~Edúcas panem de terra. \textred{℟.}~Et vinum lætíficet cor hóminis.
\end{tightlist}

\subsubsection*{\center \textsc{AD NONAM.}}
\subsubsection*{\textred{Capitulum. \hfill 1~Cor.~11.}}
\lettrine{Q}uicúmque manducáverit panem hunc, vel bíberit cálicem Dómini indígne~: reus erit córporis et sánguinis Dómini.

\begin{tightlist}{}
\item \textred{℟.}~Edúcas panem \textred{*} De terra. \textred{E}dúcas. \textred{℣.}~Et vinum lætíficet cor hóminis. \textred{D}e. Glória Patri. \textred{E}dúcas.
\item \textred{℣.}~Pósuit fines tuos pacem. \textred{℟.}~Et ádipe fruménti sátiat te.
\end{tightlist}

\subsubsection*{\center \textsc{AD VESPERAS.}}
\begin{tightlist}{}
\item \textred{Prima Aña Laud.} Sapiéntia ædificávit sibi domum, míscuit vinum, et pósuit mensam, allelúja.
\item \textred{Ps. 109.} Dixit Dóminus. [54]. 
\item \textred{Ps. 110.} Confitébor. [55]. 
\item \textred{Ps. 115.} Crédidi. [55]. 
\item \textred{Ps. 127.} Beati omnes, qui. 167. 
\item \textred{Ps. 147.} Lauda, Jerúsalem. [56].
%\item \textred{Capit. Hymn. et ℣. de j. Vesperis.}
\end{tightlist}

\subsubsection*{\textred{Capitulum. \hfill 1~Cor.~11}}
\lettrineThree{D}ominus Jesus, in qua nocte tradebátur accépit panem, et grátias agens fregit, et dixit~: Accípite, et manducáte~: hoc est Corpus meum, quod pro vobis tradétur.

\subsubsection*{\center\textred{Hymnus.}}
\lettrine{P}ange lingua gloriósi\\ 
Córporis mystérium,\\ 
Sanguinísque pretiósi,\\ 
Quem in mundi prétium\\ 
Fructus ventris generósi\\ 
Rex effúdit Géntium.\\ 
\\
\textred{N}obis datus, nobis natus\\ 
Ex intácta Vírgine,\\ 
Et in mundo conversátus,\\ 
Sparso verbi sémine,\\ 
Sui moras incolátus\\ 
Miro clausit órdine.\\ 
\\
\textred{I}n suprémæ nocte cœnæ\\ 
Recúmbens cum frátribus,\\ 
Observáta lege plene\\ 
Cibis in legálibus,\\ 
Cibum turbæ duodénæ\\ 
Se dat suis mánibus.\\ 
\\
\textred{V}erbum caro, panem verum\\ 
Verbo carnem éfficit~:\\ 
Fitque sanguis Christi merum;\\ 
Et si sensus déficit,\\ 
Ad firmándum cor sincérum\\ 
Sola fides súfficit.\\ 
\\
\textred{T}antum ergo Sacraméntum\\ 
Venerémur cérnui~:\\ 
Et antíquum documéntum\\ 
Novo cedat rítui~:\\ 
Prǽstet fides suppleméntum\\ 
Sénsuum deféctui.\\ 
\\
\textred{G}enitóri, Genitóque\\ 
Laus, et jubilátio,\\ 
Salus, honor, virtus quoque\\ 
Sit et benedíctio,\\ 
Procedénti ab utróque\\ 
Compar sit laudátio.\\
Amen.

\begin{tightlist}{}
\item \textred{℣.}~Panem de cœlo præstitísti eis. \textred{℟.}~Omne delectaméntum in se habéntem.
\item \textred{Ad Magnificat, Aña.} O sacrum convívium, in quo Christus súmitur~: recólitur memória passiónis ejus~: mens implétur grátia~: et futúræ glóriæ nobis pignus datur, allelúja.
\item \textred{Ad Nunc dimittis, Aña.} Allelúja~: Panis quem ego dédero, allelúja~: caro mea est pro mundi vita, allelúja, allelúja.
\item \textred{A Septuagesima autem usque ad Pascha:} Panis, quem ego dédero, caro mea est pro mundi vita.
\end{tightlist}

\subsubsection*{\center \textsc{A SEPTUAGESIMA.}}
\noindent \textred{Omnia ut supra; sed omittitur ubique} Allelúja.

\subsubsection*{\center \textsc{TEMPORE PASCHALI.}}
\begin{tightlist}{}
\item \textred{Omnia ut supra; sed tres Psalmi cujuslibet Nocturni dicuntur sub prima sola Antiphona, et in fine Añarum, et Responsoriorum, additur} Allelúja.
\item \textred{Item in fine omnium Hymnorum (exceptis Hymnis Vesp. Matut. et Laud.) dicitur ℣.}~Quǽsumus Auctor. \textred{et ℣.}~Glória tibi Dómine, qui natus es de Vírgine.
\end{tightlist}
\begin{center}\textcolor{red}{\rule{\linewidth}{0.5pt}} \end{center}

\subsection*{\center \textsc{Officium Commem.\\ Beatissimæ V. Mariæ\\ in Sabbato}}
\subsubsection*{\center \textred{Semiduplex.}}
\begin{center}\textsc{recitandum semel qualibet hebdomada.}\end{center}
\subsubsection*{\center \textsc{AD VESPERAS.}}
%\begin{center}\textred{in quibusdam Festis B. Mariæ Virginis.}\end{center}
\begin{tightlist}{}
\item \textred{Omnia de commune B.M.V.} [168] \textred{præter sequentia.}
%\item \textred{In quibusdam Festis B.M.V., in j. Vesperis, Añæ., et Ps. ex j. Vesperis communis B.M.V.}
%\item \textred{Si duplex fieri debet in quibusdam Festis B.M.V. dicitur:}
%\item \textred{℟.}~Christi Virgo dilectíssima, virtútum operátrix, opem fer míseris~: \textred{*} Súbveni Dómine clamántibus ad te júgiter. \textred{℣.}~Quóniam peccatórum mole prémimur, et non est qui ádjuvet. \textred{S}úbveni. Glória Patri. \textred{S}úbveni.
\end{tightlist}

\fancyhead[C]{\textsc{Officia Hebdomadales - Sabbato - B.M.V.}}

\subsubsection*{\center \textred{Oratio.}}
\lettrine{M}ajestátem tuam omnípotens Deus, súpplices deprecámur; ut qui exsultántibus ánimis beatíssimæ, sempérque Vírginis Maríæ commemoratiónem celebráre gaudémus, ejus gloriósis précibus et méritis, ab instántibus malis, et a morte perpétua liberémur. Per Dóminum.

\begin{tightlist}{}
\item \textred{In fine Hymni Completorii et Horarum dicitur ℣.}~Glória tibi Dómine, qui natus es.
\item \textred{Ad Nunc dimittis, Aña.} Sub tuum præsídium.
\end{tightlist}

\begin{center}\textcolor{red}{\rule{\linewidth}{0.5pt}} \end{center}

\subsection*{\center \textred{\textsc{IN ADVENTU.}}}
\begin{tightlist}{}
\item \textred{Omnia ut in die Annuntiationis B.V.M., 25. Martii, ut infra.}
\item \textred{Lect. j. Noct. erunt de Scriptura occurrenti si adsint, alioquin etiam ipsæ accipiantur de eodem Festo Annuntiationis.}
\item \textred{Quum autem non sit dies vacans, quo recitari posset, ideo hic omittitur hoc Officium.}
\end{tightlist}
\fancyhead[C]{\textsc{Officia Hebd. - Sabbato B.V.M. - In Adventu}}
%\begin{center}\textred{\textit{De festo Annuntiationis B. Mariæ. V.}}\end{center}
\subsubsection*{\center \textsc{AD VESPERAS.}}
%\begin{tightlist}{}
%\item \textred{Aña.} Hæc est Regína Vírginum, \textred{cum reliquis Añis. et Psalmis ut in Comm. B.V.M.} [161]
%\end{tightlist}

\subsubsection*{\textred{Capitulum. \hfill Isaiæ 7.}}
\lettrineThree{E}cce Virgo concípiet, et páriet fílium, et vocábitur nomen ejus Emmánuel. Butýrum, et mel cómedet, ut sciat reprobáre malum, et elígere bonum.

\begin{tightlist}{}
%\item \textred{Si fuerit Duplex,}
%\item \textred{℟.}~Gaude, María Virgo, cunctas hǽreses sola interemísti, quæ Gabriélis Archángeli dictis credidísti~: \textred{*} Dum Virgo Deum et hóminem genuísti, et post partum Virgo invioláta permansísti. \textred{℣.}~Gabriélem Archángelum scimus divínitus te esse affátum~: úterum tuum de Spíritu sancto crédimus fecundátum. Erubéscat Judǽus infélix, qui dicit Christum ex Joseph sémine esse natum. \textred{D}um. Glória Patri. \textred{D}um.
\item \textred{Hymn.} Ave, maris stella. [164]
\item \textred{℣.}~Ave María, grátia plena. \textred{℟.}~Dóminus tecum.
\item \textred{Ad Magnif. Aña.} Spíritus sanctus supervéniet in te, María, et virtus Altíssimi obumbrábit tibi.
\end{tightlist}

\subsubsection*{\center \textred{Oratio.}}
\lettrine{D}eus, qui de beátæ Maríæ Vírginis útero, Verbum tuum, Angelo nuntiánte, carnem suscípere voluísti~: præsta supplícibus tuis; ut qui vere eam Genitrícem Dei crédimus, ejus apud te intercessiónibus adjuvémur. Per eúmdem Dóminum.

\begin{tightlist}{}
\item \textred{℣.}~Glória tibi Dómine, Qui natus es de Vírgine, \textred{dicitur in fine Hymnorum Complet. et Horarum.}
\item \textred{Ad Nunc dimittis, Aña.} Ecce ancílla Dómini~: fiat mihi secúndum verbum tuum.
\end{tightlist}

\subsubsection*{\center \textsc{AD MATUTINUM.}}
\begin{tightlist}{}
\item \textred{Invitat.} Ave María, grátia plena~: \textred{*} Dóminus tecum.
\item \textred{Ps.} Veníte, exsultémus. 2.
\item \textred{Hymn.} Quem terra. [172]
\end{tightlist}

\subsubsection*{\center \textsc{IN PRIMO NOCTURNO.}}
\begin{tightlist}{}
\item \textred{Aña.} Ingréssus Angelus ad Maríam dixit~: Ave María, grátia plena, Dóminus tecum.
\item \textred{Ps.} Dómine, Dóminus noster, \textred{cum reliquis de Comm. B.M.V.} [172]
\item \textred{Aña.} Benedícta tu in muliéribus, et benedíctus fructus ventris tui.
\item \textred{Aña.} Angelus Dómini nuntiávit Maríæ, et concépit de Spíritu sancto.
\item \textred{℣.}~Sancta Dei Génitrix, Virgo semper María. \textred{℟.}~Intercéde pro nobis ad Dóminum, Deum nostrum.
\item \textred{Lectiones j. Noct. sumendi sunt de Scriptura occurente, vel si non habeantur, de Festo Annuntiationis, ut infra:}
\item {\small \textred{Bened.} Alma Virgo vírginum intercédat pro nobis ad Dóminum.}
\end{tightlist}

\subsubsection*{\center De Isaía Prophéta}
\subsubsection*{\center \textred{Lectio j. \hfill Cap. 7.}}
\lettrineThree{E}t adjécit Dóminus loqui ad Achaz, dicens:
Pete tibi signum a Dómino Deo tuo in profúndum inférni, sive in excélsum supra.
Et dixit Achaz~: Non petam, et non tentábo Dóminum.
Et dixit~: Audíte ergo domus David~: Numquid parum vobis est moléstos esse homínibus, quia molésti estis et Deo meo?
Propter hoc dabit Dóminus ipse vobis signum. Ecce Virgo concípiet, et páriet fílium, et vocábitur nomen ejus Emmánuel.
Butýrum et mel cómedet, ut sciat reprobáre malum, et elígere bonum. Hæc dicit Dóminus.

\begin{tightlist}{}
\item \textred{℟.}~Missus est Angelus Gábriel ad Maríam Vírginem desponsátam Joseph, núntians ei verbum~: et expavéscit Virgo de lúmine. Ne tímeas María; invenísti enim grátiam apud Deum~: \textred{*} Ecce concípies, et páries, et Fílius Altíssimi vocábitur.  \textred{℣.}~Ave María, grátia plena~: Dóminus tecum. \textred{E}cce.
\item {\small \textred{Bened.} Nos cum prole pia benedícat Virgo María.}
\end{tightlist}

\subsubsection*{\textred{Lectio ij. \hfill Cap. 11.}}
\lettrine{E}t egrediétur virga de radíce Jesse, et flos de radíce ejus ascéndet.
Et requiéscet super eum Spíritus Dómini, spíritus sapiéntiæ, et intelléctus, spíritus consílii et fortitúdinis, spíritus sciéntiæ, et pietátis,
et replébit eum spíritus timóris Dómini. Non secúndum visiónem oculórum judicábit, neque secúndum audítum áurium árguet,
sed judicábit in justítia páuperes, et árguet in æquitáte pro mansuétis terræ, et percútiet terram virga oris sui, et spíritu labiórum suórum interfíciet ímpium.
Et erit justítia cíngulum lumbórum ejus, et fides cinctórium renum ejus.
Hæc dicit Dóminus.

\begin{tightlist}{}
\item \textred{℟.}~Ingréssus Angelus ad Maríam dixit~: \textred{*} Ave María, grátia plena~: Dóminus tecum. \textred{℣.}~Benedícta tu in muliéribus, et benedíctus fructus ventris tui. \textred{A}ve.
\item {\small \textred{Bened.} Sancta Dei Génitrix sit nobis auxiliátrix.}
\end{tightlist}

\subsubsection*{\textred{Lectio iij. \hfill Cap. 35.}}
\lettrine{L}ætábitur desérta et ínvia, et exsultábit solitúdo, et florébit quasi lílium.
Gérminans germinábit, et exsultábit lætabúnda et laudans~: glória Líbani data est ei~: decor Carméli, et Saron, ipsi vidébunt glóriam Dómini, et decórem Dei nostri.
Confortáte manus dissolútas, et génua debília roboráte.
Dícite pusillánimis~: Confortámini, et nolíte timére~: ecce Deus vester ultiónem addúcet retributiónis~: Deus ipse véniet, et salvábit vos.
Tunc aperiéntur óculi cæcórum, et aures surdórum patébunt.
Tunc sáliet sicut cervus claudus, et apérta erit lingua mutórum~: quia scissæ sunt in desérto aquæ, et torréntes in solitúdine.
Et quæ erat árida, erit in stagnum, et sítiens in fontes aquárum. Hæc dicit Dóminus.

\begin{tightlist}{}
\item \textred{℟.}~María cum audísset, turbáta est in sermóne Angeli~: \textred{*} Et cogitábat qualis esset ista salutátio. \textred{℣.}~Quómodo fiet istud, quóniam virum non cognósco? \textred{E}t. Glória Patri. \textred{E}t.
\end{tightlist}

\subsubsection*{\center \textsc{IN SECUNDO NOCTURNO.}}
\begin{tightlist}{}
\item \textred{Aña.} Ne tímeas María, invenísti enim grátiam apud Deum~: ecce concípies in útero, et páries fílium.
%\item \textred{Ps.} Eructávit cor. [169]
\item \textred{Aña.} Dabit illi Dóminus Deus sedem David patris ejus~: et regnábit in ætérnum.
\item \textred{Aña.} Super sólium David, et super regnum ejus sedébit in ætérnum.
\item \textred{℣.}~Post partum, Virgo, invioláta permansísti. \textred{℟.}~Dei Génitrix, intercéde pro nobis.
\item {\small \textred{Bened.} Précibus suæ Matris benedícat nos Fílius Dei Patris.}
\end{tightlist}

\subsubsection*{\center Sermo sancti Leónis Papæ.}
\subsubsection*{\center \textred{Serm. 2. de Nat. Dom.\\
Lectio iv.}}
\lettrineThree{D}eus omnípotens et clemens, cujus natúra bónitas, cujus volúntas poténtia, cujus opus misericórdia est; statim ut nos diabólica malígnitas venéno suæ mortificávit invídiæ, prædestináta renovándis mortálibus suæ pietátis remédia, inter ipsa mundi primórdia præsignávit, denúntians serpénti futúrum semen mulíeris, quod nóxii cápitis elatiónem sua virtúte contéreret, Christum scílicet in carne ventúrum~: Deum hominémque signans, qui natus ex Vírgine, violatórem humánæ propáginis, incorrúpta nativitáte damnáret.

\begin{tightlist}{}
\item \textred{℟.}~Ait Angelus Maríæ~: Ne tímeas, \textred{*} Invenísti enim grátiam apud Deum. \textred{℣.}~Ecce concípies in útero, et páries fílium, et vocábis nomen ejus Jesum. \textred{I}nvenísti.
\item {\small \textred{Bened.} Jesus Maríæ Fílius sit nobis clemens et propítius.}
\end{tightlist}

\subsubsection*{\center \textred{Lectio v.}}
\lettrine{N}am quia gloriabátur diábolus, hóminem sua fraude decéptum divínis caruísse munéribus, et immortalitátis dote nudátum, dura mortis subiísse senténtiam, seque in malis suis quoddam de prævaricatóris consórtio invenísse solátium~: Deum quoque, justæ severitátis exigénte ratióne, erga hóminem, quem in tanto honóre condíderat, antíquam mutásse senténtiam~: opus fuit, dilectíssimi, secréti dispensatióne consílii, ut incommutábilis Deus (cujus volúntas non potest sua benignitáte privári) primam pietátis suæ dispositiónem sacraménto occultióre compléret~: et homo diabólicæ iniquitátis versútia actus in culpam, contra Dei propósitum non períret.

\begin{tightlist}{}
\item \textred{℟.}~Ecce concípies in útero, et páries fílium, \textred{*} Et vocábis nomen ejus Jesum. \textred{℣.}~Hic erit magnus, et Fílius Altíssimi vocábitur, et dabit illi Dóminus Deus sedem David patris ejus. \textred{E}t.
\item {\small \textred{Bened.} Fílius Vírginis Maríæ det nobis gáudia vitæ.}
\end{tightlist}

\subsubsection*{\center \textred{Lectio vj.}}
\lettrine{A}dveniéntibus ergo tempóribus, dilectíssimi, quæ redemptióni hóminum fúerant præstitúta, ingréditur hæc ínfima Jesus Christus Dóminus noster, de cœlésti sede descéndens, et a patérna glória non recédens, novo órdine, nova nativitáte generátus~: novo órdine, quia invisíbilis in suis, visíbilis factus est in nostris, incomprehensíbilis vóluit comprehéndi~: ante témpora manens, esse cœpit ex témpore~: universitátis Dóminus, servílem formam obumbráta majestátis suæ dignitáte suscépit; impassíbilis Deus, non dedignátus est homo esse passíbilis; et immortális, mortis légibus subjacére.

\begin{tightlist}{}
\item \textred{℟.}~Dixit María ad Angelum. Quómodo fiet istud, quóniam virum non cognósco? Et respóndens Angelus dixit ei~: \textred{*} Spíritus sanctus supervéniet in te, et virtus Altíssimi obumbrábit tibi. \textred{℣.}~Ideóque et quod nascétur ex te Sanctum, vocábitur Fílius Dei. \textred{S}píritus. Glória Patri. \textred{S}píritus.
\end{tightlist}

\subsubsection*{\center \textsc{IN TERTIO NOCTURNO.}}
\begin{tightlist}{}
\item \textred{Aña.} Spíritus Sanctus in te descéndet María, ne tímeas habébis in útero Fílium Dei.
%\item \textred{Ps.} Cantáte. [172]
\item \textred{Aña.} Beáta es María, quæ credidísti~: quóniam perficiéntur ea, quæ dicta sunt tibi a Dómino.
\item \textred{Aña.} Beátam me dicent omnes generatiónes~: quia respéxit Dóminus humilitátem ancíllæ suæ.
\item \textred{℣.}~Speciósa facta es et suávis. \textred{℟.}~In delíciis tuis, sancta Dei Génitrix.
\item {\small \textred{Bened.} Fons Evangélii répleat nos dógmate cæli.}
\end{tightlist}

\subsubsection*{\center Léctio sancti Evangélii secúndum Lucam.}
\subsubsection*{\center \textred{Lectio vij. \hfill Cap. 1.}}
\lettrine{I}n illo témpore~: Missus est Angelus Gábriel a Deo in civitátem Galilǽæ, cui nomen Názareth, ad Vírginem desponsátam viro, cui nomen erat Joseph, de domo David, et nomen Vírginis María. Et réliqua.

\begin{center}Homilía sancti Ambrósii Epíscopi.\\
\textred{Lib. 2. in Lucam.}\end{center}
\lettrineThree{L}atent quidem divína mystéria~: nec fácile, juxta prophéticum dictum, quisquam hóminum potest scire consílium Dei~: sed tamen ex céteris factis, atque præcéptis Dómini Salvatóris póssumus intellígere, et hoc propensióris fuísse consílii, quod ea potíssimum elécta est, ut Dóminum páreret, quæ erat desponsáta viro. Cur autem non ántequam desponsarétur, impléta est? Fortásse ne dicerétur, quod concéperat ex adultério.

\begin{tightlist}{}
\item \textred{℟.}~Dabit illi Dóminus Deus sedem David patris ejus~: \textred{*} Et regnábit in domo Jacob in ætérnum. \textred{℣.}~Et regni ejus non erit finis. \textred{E}t.
\item {\small \textred{Bened.} Ab hoste malígno deféndat nos Dei Génitrix Virgo.}
\item \textred{Quando de hac Homilia debent legi tres Lectiones, sequens Lectio dividatur in octavam, et nonam.}
\end{tightlist}

\subsubsection*{\center \textred{Lectio viij.}}
\lettrine{E}t ingréssus ad eam Angelus. Disce vírginem móribus, disce vírginem verecúndia, disce oráculo, disce mystério. Trepidáre vírginum est, et ad omnes viri ingréssus pavére, omnes viri affátus veréri. Discant mulíeres propósitum pudóris imitári. Sola in penetrálibus, quam nemo virórum víderit, solus Angelus repérerit; sola sine cómite, sola sine teste, ne quo degénere depravarétur affátu, ab Angelo salutátur.

¶ Tanti namque mandáti mystérium non hóminis fuit, sed Angeli ore proméndum. Hódie primum audítur~: Spíritus sanctus supervéniet in te. Et audítur, et créditur. Dénique~: Ecce, inquit, ancílla Dómini, contíngat mihi secúndum verbum tuum. Vide humilitátem, vide devotiónem! Ancíllam se dicit Dómini, quæ mater elígitur~: nec repentíno exaltáta promísso est.

\begin{tightlist}{}
\item \textred{℟.}~Dixit autem María ad Angelum~: Ecce ancílla Dómini~: \textred{*} Fiat mihi secúndum verbum tuum. \textred{℣.}~Spíritus Sanctus supervéniet in te, et virtus Altíssimi obumbrábit tibi. \textred{F}iat.
\item \textred{Si ix Lectio legenda sit de alia Homilia, Bened.} Verba Redemptóris pia sint suffrágia nobis.
\item \textred{Alias, Bened.} In omni tribulatióne et angústia succúrrat nobis pia Virgo María.
\item \textred{℟.~ix.} Gaude María Virgo, cunctas hǽreses sola interemísti, quæ Gabriélis Archángeli dictis credidísti~: \textred{*} Dum Virgo Deum et hóminem genuísti, et post partum Virgo invioláta permansísti. \textred{℣.}~Gabriélem Archángelum scimus divínitus te esse affátum~: úterum tuum de Spíritu sancto crédimus fecundátum. Erubéscat Judǽus infélix, qui dicit Christum ex Joseph sémine esse natum. \textred{D}um. Glória Patri. \textred{D}um.
\item Te Deum laudámus. 6.
\item \textred{Sacerd. ℣.}~Ora pro nobis, sancta Dei Génitrix. \textred{℟.}~Ut digni efficiámur promissiónibus Christi.
\end{tightlist}

\subsubsection*{\center \textsc{AD LAUDES.}}
\begin{center}\textred{et per Horas, Aña.}\end{center}
\lettrine{M}issus est Angelus Gábriel a Deo ad Maríam Vírginem desponsátam Joseph.

\begin{tightlist}{}
\item \textred{Ps. de Dominica.} 27.
\item \textred{Aña.} Ave María, grátia plena~: Dóminus tecum~: Benedícta tu in muliéribus.
\item \textred{Aña.} Ne tímeas María, invenísti enim grátiam apud Deum~: ecce concípies in útero, et páries fílium.
\item \textred{Aña.} Dabit illi Dóminus Deus sedem David patris ejus, et regnábit in ætérnum.
\item \textred{Aña.} Ecce ancílla Dómini, fiat mihi secúndum verbum tuum.
\end{tightlist}

\subsubsection*{\textred{Capitulum. \hfill Isa. 7.}}
\lettrineThree{E}cce Virgo concípiet, et páriet fílium, et vocábitur nomen ejus Emmánuel. Butýrum et mel cómedet, ut sciat reprobáre malum, et elígere bonum.

\begin{tightlist}{}
\item \textred{Hymn.} O gloriósa Dómina, \textred{de communi.} [183]
\item \textred{℣.}~Benedícta tu in muliéribus. \textred{℟.}~Et benedíctus fructus ventris tui.
\item \textred{Ad Bened. Aña.} Quómodo fiet istud, Angele Dei, quóniam virum non cognósco? Audi María Virgo~: Spíritus Sanctus supervéniet in te, et virtus Altíssimi obumbrábit tibi.
\end{tightlist}

\subsubsection*{\center \textred{Oratio.}}
\lettrine{D}eus, qui de beátæ Maríæ Vírginis útero, Verbum tuum, Angelo nuntiánte, carnem suscípere voluísti~: præsta supplícibus tuis; ut qui vere eam Genitrícem Dei crédimus, ejus apud te intercessiónibus adjuvémur. Per eúmdem Dóminum.

\begin{tightlist}{}
\item \textred{Ad Horas Añæ de Laudibus, Psalmi de Dominica, ad Primam tamen ut in Festis.}
\end{tightlist}

\subsubsection*{\center \textsc{AD PRIMAM.}}
\noindent \textred{℣.}~Qui de Vírgine dignátus es nasci.

\subsubsection*{\center \textsc{AD TERTIAM.}}
\subsubsection*{\textred{Capitulum. \hfill Isa. 7.}}
\lettrine{E}cce Virgo concípiet, et páriet fílium, et vocábitur nomen ejus Emmánuel. Butýrum et mel cómedet, ut sciat reprobáre malum, et elígere bonum.

\begin{tightlist}{}
\textred{℟.}~Sancta Dei Génitrix, \textred{*} Virgo semper María. \textred{S}ancta. \textred{℣.}~Intercéde pro nobis ad Dóminum Deum nostrum. \textred{V}irgo. Glória Patri. \textred{S}ancta.
\item \textred{℣.}~Post partum Virgo invioláta permansísti. \textred{℟.}~Dei Génitrix intercéde pro nobis.
\end{tightlist}

\subsubsection*{\center \textsc{AD SEXTAM.}}
\subsubsection*{\textred{Capitulum. \hfill Isa. 11.}}
\lettrine{E}grediétur virga de radíce Jesse, et flos de radíce ejus ascéndet. Et requiéscet super eum Spíritus Dómini.

\begin{tightlist}{}
\item \textred{℟.}~Post partum, Virgo \textred{*} Invioláta permansísti. \textred{P}ost. \textred{℣.}~Dei Génitrix intercéde pro nobis. \textred{I}nvioláta. Glória Patri. \textred{P}ost.
\item \textred{℣.}~Speciósa facta es, et suávis. \textred{℟.}~In delíciis tuis, sancta Dei Génitrix.
\end{tightlist}

\subsubsection*{\center \textsc{AD NONAM.}}
\subsubsection*{\textred{Capitulum. \hfill Isa. 45.}}
\lettrine{R}oráte cœli désuper, et nubes pluant justum~: aperiátur terra, et gérminet Salvatórem~: et justítia oriátur simul~: ego Dóminus creávi eum.

\begin{tightlist}{}
\item \textred{℟.}~Speciósa facta es, \textred{*} Et suávis. \textred{S}peciósa. \textred{℣.}~In delíciis tuis, sancta Dei Génitrix. \textred{E}t. Glória Patri. \textred{S}peciósa.
\item \textred{℣.}~Elégit eam Deus, et præelégit eam. \textred{℟.}~Et habitáre eam facit in tabernáculo suo.
\end{tightlist}

\subsubsection*{\center \textsc{IN II. VESPERAS.}}
\begin{tightlist}{}
\item \textred{Añæ. Laudum.}
\item \textred{Ps.} Dixit Dóminus, \textred{cum reliquis de Comm. B.M.V.} [184]
%\item \textred{Capit. Hymn. et ℣. ut in j. Vesp.}
\end{tightlist}

\subsubsection*{\textred{Capitulum. \hfill Isaiæ 7.}}
\lettrineThree{E}cce Virgo concípiet et páriet fílium, et vocábitur nomen ejus Emmánuel. Butyrum, et mel cómedet, ut sciat reprobáre malum, et elígere bonum.

\begin{tightlist}{}
%\item \textred{Si fuerit Duplex,}
%\item \textred{℟.}~Gaude, María Virgo, cunctas hǽreses sola interemísti, quæ Gabriélis Archángeli dictis credidísti~: \textred{*} Dum Virgo Deum et hóminem genuísti, et post partum Virgo invioláta permansísti. \textred{℣.}~Gabriélem Archángelum scimus divínitus te esse affátum~: úterum tuum de Spíritu sancto crédimus fecundátum. Erubéscat Judǽus infélix, qui dicit Christum ex Joseph sémine esse natum. \textred{D}um. Glória Patri. \textred{D}um.
\item \textred{Hymn.} Ave maris stella. [187]
\item \textred{℣.}~Ave María grátia plena. \textred{℟.}~Dóminus tecum.

\item \textred{Ad Magnif. Aña.} O Virgo vírginum, quómodo fiet istud, quia nec primam símilem visa es, nec habére sequéntem? Fíliæ Jerúsalem quid me admirámini? divínum est mystérium hoc, quod cérnitis.
\end{tightlist}

\subsubsection*{\center \textred{Oratio.}}
\lettrine{D}eus, qui de beátæ Maríæ Vírginis útero, Verbum tuum, Angelo nuntiánte, carnem suscípere voluísti~: præsta supplícibus tuis; ut qui vere eam Genitrícem Dei crédimus, ejus apud te intercessiónibus adjuvémur. Per eúmdem Dóminum.

\begin{center}\textcolor{red}{\rule{\linewidth}{0.5pt}} \end{center}

\subsection*{\center \textred{\textsc{A SEPTUAGESIMA}}}
\begin{center}\textred{usque ad Pascha.}\end{center}

\subsubsection*{\center \textsc{AD VESPERAS.}}
\begin{center}\textred{Si fuerit Duplex,}\end{center}
\begin{tightlist}{}
\item \textred{Aña.} O vos omnes, qui transítis per viam, atténdite et vidéte, si est dolor sicut dolor meus.
\item \textred{Psalmi de Officio præc.}
\end{tightlist}
%\fancyhead[C]{\textsc{Sabbato - B.V.M. in Septuagesima}}
\fancyhead[C]{\textsc{Officia Hebd. - Sabbato B.V.M. - Septuagesima}}
\subsubsection*{\textred{Capitulum. \hfill Thren. 1.}}
\lettrineThree{V}ide, Dómine, quóniam tríbulor, conturbátus est venter meus~: subvérsum est cor meum in memetípsa, quóniam amaritúdine plena sum.

\begin{tightlist}{}
\item \textred{Si fuerit Duplex,}
\item \textred{℟.}~Vúlnera Christi moriéntis erant vúlnera matris doléntis; dolóres sævi fuérunt tortóres in ánima matris, \textred{*} Mater erat laniáta morte cari pígnoris.
\item \textred{℣.}~Mente Mater erat percússa cúspide teli, quo membra Christi servi fodérunt iníqui. \textred{M}ater. Glória Patri. \textred{M}ater.
\end{tightlist}

\subsubsection*{\center \textred{Hymnus.}}
\lettrine{L}acrymæ Matris prófluunt,\\
Cordis erúmpunt gémitus,\\
Dum crucis orbis Cónditor\\
Suspénsus est patíbulo.\\
\\
\textred{D}um tot tuétur vúlnera,\\
Rivos fluéntis sánguinis,\\
Fossúmque pectus láncea,\\
Suo crucem fert péctore.\\
\\
\textred{D}olóris ensem séntiunt\\
Piæ paréntis víscera,\\
Cum Jesus dire móritur,\\
Pro reis factus hóstia.\\
\\
\textred{I}lli tantum compátitur,\\
Quod secum fere móritur,\\
Nullus dolor amárior,\\
Nam nulla proles cárior.\\
\\
\textred{O} María, mater pia,\\
Spes única fidélium,\\
Tua prece mellíflua,\\
Natum redde propítium.\\
\\
\textred{Q}uo mortis suæ mérito\\
In hujus vitæ stádio\\
Nos servet, nosque líberet\\
A gehénnæ incéndio.\\
\\
\textred{G}lória tibi, Dómine,\\
Qui natus es de Vírgine,\\
Cum Patre et sancto Spíritu,\\
In sempitérna sǽcula.\\
Amen.

\begin{tightlist}{}
\item \textred{Sic terminantur omnes Hymni.}
\item \textred{℣.}~Sancta Dei Génitrix, Virgo dulcis atque pia. \textred{℟.}~Christum morti tráditum pro nobis semper exóra.
\item \textred{Ad Magnificat, Aña.} Jesu nate de Vírgine sine viríli sémine, qui mortem crucis páteris, suffóssus plagis ásperis~: amóre matris ánxiæ nostræ parcas misériæ; ac ejus piis précibus cœli fruámur gáudiis.
\end{tightlist}

\subsubsection*{\center \textred{Oratio.}}
\lettrine{I}ntervéniat pro nobis, quǽsumus, Dómine Jesu Christe, apud tuam cleméntiam, nunc et in hora mortis nostræ, piíssima Virgo María mater tua~: cujus sacratíssimam ánimam in hora passiónis tuæ dolóris gládius pertransívit, et in gloriósa resurrectióne tua ingens gáudium lætificávit. Qui vivis et regnas.

\begin{tightlist}{}
\item \textred{Ad Nunc dimittis, Aña.} Sub tuum præsídium.
\end{tightlist}

\subsubsection*{\center \textsc{AD MATUTINUM.}}
\begin{tightlist}{}
\item \textred{Invitatorium.} Jesum pro nobis in cruce mórtuum, a Vírgine Matre lamentátum, \textred{*} Veníte adorémus.
\item \textred{Ps.} Venite, exsultémus. 2.
\end{tightlist}

\subsubsection*{\center \textred{Hymnus.}}
\lettrine{M}entes pungant fidélium\\
Dolóres et suspíria\\
Matris Christi, dum víscera\\
Pœna torquéntur áspera.\\
\\
\textred{D}um stat crucis sub stípite,\\
Videns pendéntem Fílium,\\
Præ spíritus angústia\\
Est liquefácta pénitus.\\
\\
\textred{E}xors viri quæ cóncipit,\\
Expers dolóris péperit~:\\
Mœrens dolóres pértulit,\\
Quos páriens non sénserat.\\
\\
\textred{H}oc ille justus Símeon\\
In spíritu prædíxerat,\\
Ampléxibus dum párvulum\\
Fovens piis adstríngeret.\\
\\
\textred{A}uxit dolóres Fílii\\
Matris dolor, ut ánxiæ\\
Matris dolor dolóribus\\
Sic ingravéscit Fílii.\\
\\
\textred{Q}ui Maríam dilígitis,\\
Vos ejus et suspíria\\
Et gémitus recólite,\\
Et dira Christi vúlnera.\\
\\
\textred{G}lória tibi, Dómine,\\
Qui natus es de Vírgine,\\
Cum Patre et Sancto Spíritu,\\
In sempitérna sǽcula.\\
Amen.

\subsubsection*{\center \textsc{IN PRIMO NOCTURNO.}}
\begin{tightlist}{}
\item \textred{Aña.} Surgam et circuíbo civitátem, per vicos et platéas quæram quem díligit ánima mea.
\item \textred{Ps.} Dómine, Dóminus, \textred{cum reliquis de Comm. BMV.} [172]
\item \textred{Aña.} Quæsívi illum, et non invéni, vocávi, et non respóndit mihi.
\item \textred{Aña.} Puer non compáret, et ego quo ibo?
\item \textred{℣.}~Sancta Dei Génitrix, Virgo dulcis atque pia. \textred{℟.}~Christum morti tráditum, pro nobis semper exóra.
\item \textred{Ante Quadragesimam Lectiones primi Nocturni de Scriptura occurrente, sed in Quadragesima sequentes.}
\item {\small \textred{Bened.} Alma Virgo vírginum intercédat pro nobis ad Dóminum.}
\end{tightlist}

\subsubsection*{\center De Jeremía Prophéta.}
\subsubsection*{\center \textred{Lectio j.\\ Thren.~1.~℣.~20.}}
\lettrineThree{V}ide Dómine quóniam tríbulor, conturbátus est venter meus~: subvérsum est cor meum in memetípsa, quóniam amaritúdine plena sum. Foris intérficit gládius, et domi mors símilis est.
Audiérunt quia ingemísco ego, et non est qui consolétur me. Omnes inimíci mei audiérunt malum meum, lætáti sunt, quóniam tu fecísti~: Adduxísti diem consolatiónis, et fient símiles mei. Hæc dicit Dóminus Deus~: Convertímini ad me, et salvi éritis.

\begin{tightlist}{}
\item \textred{℟.}~Adspiciébat dulcem natum Virgo María in cruce pendéntem, et morte turpíssima moriéntem \textred{*} Tantóque dolóre et tristítia vexabátur in mente, ut non possit explicári sermóne. \textred{℣.}~Erat enim Jesus adspéctu dulcis, collóquio suávis, et omni conversatióne benigníssimus. \textred{T}antóque.
\item {\small \textred{Bened.} Nos cum prole pia benedícat Virgo María.}
\end{tightlist}

\subsubsection*{\center \textred{Lectio ij. \hfill Cap.~3.~℣.~48.}}
\lettrine{D}ivisiónes aquárum dedúxit óculus meus, in contritióne fíliæ pópuli mei.
Oculus meus afflíctus est, nec tácuit, eo quod non esset réquies,
donec respíceret et vidéret Dóminus de cælis.
Oculus meus deprædátus est ánimam meam in cunctis filiábus urbis meæ.
Venatióne cepérunt me quasi avem inimíci mei gratis.
Lapsa est in lacum vita mea, et posuérunt lápidem super me.
Inundavérunt aquæ super caput meum~: dixi~: Périi.
Invocávi nomen tuum Dómine de lacu novíssimo.
Vocem meam audísti~: ne avértas aurem tuam a singúltu meo, et clamóribus.
Appropinquásti in die, quando invocávi te~: dixísti~: ne tímeas. Hæc dicit Dóminus.

\begin{tightlist}{}
\item \textred{℟.}~Manábat sanguis ejus ex quátuor pártibus rigántibus undis, ligno mánibus pedibúsque confíxis~: de vultu illíus pulchritúdo efflúxerat omnis, \textred{*} Et qui erat speciósus forma præ fíliis hóminum, videbátur ómnium indecórus. \textred{℣.}~Non erat spécies ei, neque decor; vídimus eum, et non erat adspéctus. \textred{E}t.
\item {\small \textred{Bened.} Sancta Dei Génitrix sit nobis auxiliátrix.}
\end{tightlist}

\subsubsection*{\center \textred{Lectio iij.}}
\lettrine{J}udicásti Dómine causam ánimæ meæ, redémptor vitæ meæ.
Vidísti Dómine iniquitátem illórum advérsum me~: júdica judícium meum.
Vidísti omnem furórem, univérsas cogitatiónes eórum advérsum me.
Audísti oppróbrium eórum Dómine, omnes cogitatiónes eórum advérsum me,
lábia insurgéntium mihi; et meditatiónes eórum advérsum me tota die.
Sessiónem eórum, et resurrectiónem eórum vide~: ego sum psalmus eórum.
Reddes eis vicem, Dómine, juxta ópera mánuum suárum.
Dabis eis scutum cordis labórem tuum.
Persequéris in furóre, et cónteres eos sub cælis, Dómine. Hæc dicit Dóminus.

\begin{tightlist}{}
\item \textred{℟.}~Iste erat máximus dolor Maríæ, quia vidébat se relínqui ab eo quem genúerat, nec supérerat álius, quia illi erat únicus~: \textred{*} Vox triste sonábat foris vulnus denúntians mentis. \textred{℣.}~Amor únice, fili dulcíssime, quis dabit mihi, ut pro te móriar, fili mi? \textred{V}ox. Glória Patri. \textred{V}ox.
\end{tightlist}

\subsubsection*{\center \textsc{IN SECUNDO NOCTURNO.}}
\begin{tightlist}{}
\item \textred{Aña.} Adjúro vos, fíliæ Jerúsalem, si invenéritis diléctum meum, ut nuntiétis ei quia amóre lángueo.
%\item \textred{Ps.} Eructávit. [169]
\item \textred{Aña.} Qualis est diléctus tuus ex dilécto, o pulchérrima mulíerum, quia sic adjurásti nos?
\item \textred{Aña.} Diléctus meus cándidus et rubicúndus, eléctus ex míllibus.
\item \textred{℣.}~Virgo mater Christi juxta Crucem adstitísti. \textred{℟.}~Præparáta mori propter amórem Fílii.
\item {\small \textred{Bened.} Précibus suæ Matris benedícat nos Fílius Dei Patris.}
\end{tightlist}

\subsubsection*{\center Sermo sancti Bernárdi Abbátis.}
\subsubsection*{\center \textred{De Lament. B. Mariæ.\\
Lectio iv.}}
\lettrineThree{J}uxta crucem Christi stabat emórtua María mater, quæ ipsum ex Spíritu sancto concépit. Vox illi non erat, quia dolóre attríta jacens pallébat. Quasi mórtua vivens, vivébat móriens; moriebátur vivens, nec mori póterat, quæ vivens mórtua erat. In illíus ánima dolor sæve sæviébat~: optábat mori magis quam vívere post mortem Christi, quæ male vivens mórtua erat. Ibi stabat dolens, sævo dolóre confécta. O verum elóquium justi Simeónis, quem prædíxit, gládium sentiébat dolóris.

\begin{tightlist}{}
\item \textred{℟.}~Nunc solvis Virgo cum usúra, quod in partu non habuísti a natúra~: \textred{*} Dolórem pariéndo Fílium non sensísti, quem míllies replicátum Fílio moriénte passa fuísti. \textred{℣.}~Juxta Crucem Christi stabat velut emórtua mater, quæ ipsum ex Spíritu sancto concépit. \textred{D}olórem.
\item {\small \textred{Bened.} Jesus Maríæ Fílius sit nobis clemens et propítius.}
\end{tightlist}

\subsubsection*{\center \textred{Lectio v.}}
\lettrine{S}tabat juxta crucem María, íntuens vultu benígno Christum pendéntem in patíbulo, pedúmque summitátibus innítens, manus levábat in altum, ampléctens rubricátam crucem, ac in óscula ejus ruens ea parte, qua unda pretiosíssimi sánguinis defluébat~: sursum manus, nisu quo póterat, extendébat, únicum suum amplécti desíderans, nec valébat. Sperat enim amor multa, quæ nunquam, vel raro fíeri possunt~: impátiens síquidem amor credit quod sibi débeant cédere univérsa. Volébat amplécti Christum in alto pendéntem, sed manus frustra proténsæ, in se complósæ, compléxæ redíbant.

\begin{tightlist}{}
\item \textred{℟.}~Vox illi non erat, quia dolóre attríta intus conflictabátur, quasi mórtua vivens, vivébat móriens, \textred{*} Nec mori póterat, quæ vivens mórtua erat. \textred{℣.}~Optábat mori magis quam vívere post mortem Christi, in cujus ánima dolor sæve sæviébat. \textred{N}ec.
\item {\small \textred{Bened.} Fílius Vírginis Maríæ det nobis gáudia vitæ.}
\end{tightlist}

\subsubsection*{\center \textred{Lectio vj.}}
\lettrine{S}ed cum de cruce corpus ejus fuísset totáliter depósitum, præ dolóris veheméntia, et amóris immensitáte, quasi exánimis facta fuit. Stabat ad caput exstíncti fílii mœstíssima mater, et ejus regálem fáciem mortis obfuscátam pallóribus, magna rigábat affluéntia lacrymárum. Adspiciébat illud reverendíssimum caput corónæ spíneæ diris acúleis perforátum, manus illas et pedes sanctos clavis férreis crudelíssime perforátos, latúsque perfóssum láncea, cum céteris membris lacerátis. Verúmtamen rectum erat amóris et mœróris cóntinens modum. Non desperábat, sed pie et juste dolébat, sperans tamen fírmiter, fortitérque tenens ipsum, secúndum promíssum suum, tértia die resúrgere.

\begin{tightlist}{}
\item \textred{℟.}~O verum elóquium justi Simeónis, quem prædíxit gládium sentiébat dolóris; cum juxta crucem staret, \textred{*} Intuens vultu benígno Christum pendéntem in ligno. \textred{℣.}~Ibi stabat dolens sævo dolóre confécta. \textred{I}ntuens. Glória Patri. \textred{I}ntuens.
\end{tightlist}

\subsubsection*{\center \textsc{IN TERTIO NOCTURNO.}}
\begin{tightlist}{}
\item \textred{Aña.} Jesu Fili mi, Fili mi Jesu, quis mihi tríbuat, ut pro te móriar?
%\item \textred{Ps.} Cantáte. [172]
\item \textred{Aña.} Amor únice, Fili dulcíssime, noli me derelínquere post te, trahe me ad te ipsum, ut ipsa móriar tecum.
\item \textred{Aña.} In tuo me súscipe patíbulo, ut qui una carne vívimus, una morte pereámus.
\item \textred{℣.}~Virgo tuum Fílium in cruce clavátum. \textred{℟.}~Redde tuis précibus míseris placátum. 
\item {\small \textred{Bened.} Fons Evangélii répleat nos dógmate cæli.}
\end{tightlist}

\subsubsection*{\center Lectio sancti Evangélii secúndum Lucam.}
\subsubsection*{\center \textred{Lectio vij. \hfill Cap. 2. ℣.~33.}}
\lettrine{I}n illo témpore~: Erat Pater Jesu et Mater mirántes super his, quæ dicebántur de illo. Et réliqua.

\begin{center}Homilía sancti Ambrósii Epíscopi.\\
\textred{Libro 2. in Luc.}\end{center}
\lettrineThree{V}ides úberem in omnes grátiam Dómini generatióne diffúsam, et prophetíam incrédulis negátam esse, non justis. Ecce et Símeon prophétat in ruínam et resurrectiónem plurimórum venísse Dóminum Jesum Christum, ut justórum iniquorúmque mérita discérnat; et pro nostrórum qualitáte factórum, judex verus et justus, aut supplícia decérnat, aut prǽmia.

\begin{tightlist}{}
\item \textred{℟.}~Quis æstimáre potest quanta tristítia, mœróre et dolóre vulnerátum fúerit cor matérnum, super únici Fílii passióne? \textred{*} Quæ tam caste concépit, tam lætánter péperit, tam dúlciter áluit. \textred{℣.}~Vidit eum ligátum, et sólvere non pótuit~: vidit corpus vulnerátum, et vúlnera ligáre non pótuit. \textred{Q}uæ.
\item {\small \textred{Bened.} Ab hoste malígno deféndat nos Dei Génitrix Virgo.}
\item \textred{Extra Quadrag. et Vigiliam Lectio viij. dividatur in duas,}
\item {\small \textred{et pro ix. Bened.} In omni tribulatióne et angústia succúrrat nobis pia Virgo María.}
\end{tightlist}

\subsubsection*{\center \textred{Lectio viij.}}
\lettrine{E}t tuam, inquit, ipsíus ánimam pertransíbit gládius. Nec líttera, nec história docet, ex hac vita Maríam corporális necis passióne migrásse. Non enim ánima, sed corpus materiáli gládio transverberátur.

¶ Et ídeo prudéntiam Maríæ haud ignáram mystérii cœléstis osténdit~: vivum enim verbum Dei et válidum, et acútius omni gládio acutíssimo, pénetrans usque ad divisiónem ánimæ et spíritus, ártuum quoque et medullárum, cogitatiónes cordis, et secréta scrutátur animórum, quia nuda et apérta sunt ómnia Dei Fílio, quem consciéntiæ secréta non fallunt.

\begin{tightlist}{}
\item \textred{℟.}~Fóderis, o dulcis Jesu, diris clavis extra; sed aspérius intra, ex compassióne Vírginis matris, tecum ómnia participántis. \textred{*} O mirábilis censúræ condítio! peccávit iníquus, et vápulat justus. \textred{℣.}~Delínquit homo reus, et occíditur pro eo homo Deus. \textred{O}.
\item \textred{In Quadrag. vel Vigilia, ix. Lect. et comm. de ea.}
\item {\small \textred{Bened.} Verba Redemptóris pia sint suffrágia nobis.}
\item \textred{℟.~ix.} O dulcíssime Jesu, o benigníssima María Virgo, quid commisístis, quod tantis cruciátibus, tot et tantis afflictiónibus addícti estis? \textred{*} Nos sumus qui peccávimus, qui injúste égimus; et tu innocentíssima Virgo, cum único nato pro nobis dolétis et affligímini. \textred{℣.}~O peccátor, vide quómodo patiúntur, et flent pro te Jesus et María. \textred{N}os. Glória Patri. \textred{N}os.
\item Te Deum laudámus. 6.
\item \textred{Sacerdotalis ℣.}~Tuam ipsíus ánimam dolóris gládius pertransíbit. \textred{℟.}~Ut reveléntur ex multis córdibus cogitatiónes.
\end{tightlist}

\subsubsection*{\center \textsc{AD LAUDES.}}
\begin{center}\textred{et per Horas, Aña.}\end{center}
\lettrine{S}tabant juxta crucem Jesu mater ejus, et soror matris ejus, María Cléophæ, et María Magdaléne.

\begin{tightlist}{}
\item \textred{Ps. de Dominica.} 27.
\item \textred{Aña.} Cum vidísset ergo Jesus matrem et discípulum stantem quem diligébat, dixit matri suæ~: Múlier ecce fílius tuus.
\item \textred{Aña.} Benedíxit illis Símeon, et dixit ad Maríam matrem ejus~: Ecce pósitus est hic in ruínam, et in resurrectiónem multórum in Israel, et in signum, cui contradicétur.
\item \textred{Aña.} Et tuam ipsíus ánimam pertransíbit gládius, ut reveléntur ex multis córdibus cogitatiónes.
\item \textred{Aña.} O vos omnes qui transítis per viam, atténdite et vidéte, si est dolor sicut dolor meus.
\end{tightlist}

\subsubsection*{\textred{Capitulum. \hfill Thren. 1.}}
\lettrineThree{V}ide, Dómine, quóniam tríbulor, conturbátus est venter meus, subvérsum est cor meum in memetípsa, quóniam amaritúdine plena sum.

\subsubsection*{\center \textred{Hymnus.}}
\lettrine{C}ondoléntes mœróribus\\
Maríæ matris Vírginis,\\
Jungámus laudes flétibus\\
Ad laudem sui nóminis.\\
\\
\textred{P}ensémus matris sérvuli\\
Spinas cordis acérrimas,\\
Dum vidit corpus Fílii\\
Pœnas pati tetérrimas.\\
\\
\textred{P}ensémus quid perpénderet;\\
Adstans cruci puérpera,\\
Dum testaméntum cónderet\\
Christus in vitæ véspera.\\
\\
\textred{M}atrem piam discípulo,\\
Patri comméndat spíritum~:\\
Acéti sumpto póculo,\\
Carnis exsólvit débitum.\\
\\
\textred{C}uncti Maríæ sérvuli,\\
Hujus osóres sǽculi~:\\
Jesum occísum cólite,\\
Et cum María plángite.\\
\\
\textred{G}lória tibi Dómine,\\
Qui natus es de Vírgine,\\
Cum Patre et sancto Spíritu,\\
In sempitérna sǽcula.\\
Amen.

\begin{tightlist}{}
\item \textred{℣.}~Sit nobis vita, mors Fílii Dei Patris. \textred{℟.}~Prece devóta suæ piíssimæ matris.
\item \textred{Ad Benedictus, Aña.} De cruce depósitum videns corpus Christi mœsta mater lácrymans, atque vultu tristi, dixit~: O dulcíssime Fili, quid fecísti, quod has pœnas ásperas et mortem subiísti?
\end{tightlist}

\subsubsection*{\center \textred{Oratio.}}
\lettrine{I}ntervéniat pro nobis, quǽsumus, Dómine Jesu Christe, apud tuam cleméntiam, nunc et in hora mortis nostræ, piíssima Virgo María mater tua~: cujus sacratíssimam ánimam in hora passiónis tuæ dolóris gládius pertransívit, et in gloriósa resurrectióne tua ingens gáudium lætificávit. Qui vivis.

\subsubsection*{\center \textsc{AD PRIMAM.}}
\begin{tightlist}{}
\item \textred{℣.}~Qui de Vírgine dignátus es nasci.
\end{tightlist}

\subsubsection*{\center \textsc{AD TERTIAM.}}
\subsubsection*{\textred{Capitulum. \hfill Thren. 1.}}
\lettrine{V}ide, Dómine, quóniam tríbulor, conturbátus est venter meus, subvérsum est cor meum in memetípsa, quóniam amaritúdine plena sum.

\begin{tightlist}{}
%\item \textred{Capitulum.} Vide Dómine, \textred{ut supra.}
\item \textred{℟.}~Sancta Dei Génitrix, \textred{*} Virgo dulcis atque pia. \textred{S}ancta. \textred{℣.}~Christum morti tráditum pro nobis semper exóra. \textred{V}irgo. Glória Patri. \textred{S}ancta.
\item \textred{℣.}~Virgo mater Christi, juxta Crucem adstitísti. \textred{℟.}~Præparáta mori propter amórem Fílii.
\end{tightlist}

\subsubsection*{\center \textsc{AD SEXTAM.}}
\subsubsection*{\textred{Capitulum. \hfill Ruth. 1.}}
\lettrine{N}e vocétis me Noémi, id est pulchram~: sed vocáte me Mara, id est, amáram; quia amaritúdine valde me replévit Omnípotens.

\begin{tightlist}{}
\item \textred{℟.}~Virgo mater Christi, \textred{*} Juxta Crucem adstitísti. \textred{V}irgo. \textred{℣.}~Præparáta mori propter amórem Fílii. \textred{J}uxta. Glória Patri. \textred{V}irgo.
\item \textred{℣.}~Virgo tuum Fílium in cruce clavátum. \textred{℟.}~Redde tuis précibus, míseris placátum.
\end{tightlist}

\subsubsection*{\center \textsc{AD NONAM.}}
\subsubsection*{\textred{Capitulum. \hfill Thren. 1.}}
\lettrine{D}educ, quasi torréntem lácrymas per diem et noctem, neque táceat pupílla óculi tui.

\begin{tightlist}{}
\item \textred{℟.}~Virgo tuum Fílium \textred{*} In cruce clavátum. Virgo. \textred{℣.}~Redde tuis précibus, míseris placátum. \textred{I}n. Glória Patri. \textred{V}irgo.
\item \textred{℣.}~Sit nobis vita mors Fílii Dei Patris. \textred{℟.}~Prece devóta suæ piíssimæ matris.
\item \textred{Ad secundas Vesp. sola prima Aña Laudum, Ps. de Comm. B.V.M. [184]}
\item \textred{Capitul. Hymn. et ℣.~de primis Vesperis.}
\item \textred{Ad Magnificat, Aña.} Ante thronum Trinitátis miserórum miseráta, pia mater pietátis, sis pro nobis advocáta~: et tuo apud Fílium intervéntu véniam peccatórum servis tuis óbtine.
\item \textred{Quando dicitur pro commemoratione, ℣.}~Sancta Dei Génitrix, Virgo dulcis atque pia. \textred{℟.}~Christum morti tráditum pro nobis semper exóra.
\end{tightlist}

\begin{center}\textcolor{red}{\rule{\linewidth}{0.5pt}} \end{center}

\subsubsection*{\center \textsc{TEMP. PASCHALI.}}
\begin{tightlist}{}
\item \textred{Ut in Offic. per annum, sed ante ult. ℣.~Hymnorum (Hymno Vesp. excepto) additur ℣.}~Quǽsumus auctor ómnium, \textred{vel ℣.}~Tu esto nostrum gáudium, \textred{pro Temporis exigentia, et ad Matut. tres Psalmi cujuslibet Noct. dicuntur sub j. sola Aña, et in fine ubique additur} Allelúja, \textred{ut in aliis Officiis.} 
\item \textred{Ad Nunc dimittis, Aña.} Allelúja~: Sancta Dei Génitrix.
\item \textred{Lectiones j. Nocturni de Scriptura occurr. sed cum non habentur, legantur ut per annum,} Ego Sapiéntia hábito. [174]
\end{tightlist}
\fancyhead[C]{\textsc{Officia Hebd. - Sabbato B.V.M. - T.P.}}

\subsubsection*{\center \textsc{IN SECUNDO NOCTURNO.}}
\begin{tightlist}{}
\item {\small \textred{Bened.} Précibus suæ Matris benedícat nos Fílius Dei Patris.}
\end{tightlist}
\subsubsection*{\center Sermo sancti Augustíni Epíscopi.}
\subsubsection*{\center \textred{Serm. 83. In. Append.\\
Lectio iv.}}
\lettrineThree{L}oquámur áliquid, fratres in láudibus sacratíssimæ Vírginis Maríæ. Sed quid nos tantílli, quid actióne pusílli, quid in ejus láudibus referémus~: cum etsi ómnium nostrum membra verteréntur in linguas, eam laudáre suffíceret nullus? Altior cœlo est, de qua lóquimur~: Abýsso profúndior, cui laudes dícere conámur. Hæc est enim quæ sola méruit mater et sponsa vocári. Hæc primæ matris damna resólvit. Hæc hómini pérdito redemptiónem addúxit.

\begin{tightlist}{}
\item \textred{℟.}~Sicut cedrus exaltáta sum in Líbano, et sicut cypréssus in monte Sion~: quasi myrrha elécta, \textred{*} Dedi suavitátem odóris.
\textred{℣.} Et sicut cinnamómum, et bálsamum aromatízans. \textred{D}edi.
\item {\small \textred{Bened.} Jesus Maríæ Fílius sit nobis clemens et propítius.}
\end{tightlist}

\subsubsection*{\center \textred{Lectio v.}}
\lettrine{M}ater enim géneris nostri pœnam íntulit mundo~: Génitrix Dómini nostri salútem édidit mundo~: auctrix peccáti Heva, auctrix mériti María~: Heva occidéndo óbfuit~: María vivificándo prófuit~: illa percússit, ista sanávit. Hæc enim mirábili atque inæstimábili modo ómnium rerum et suum péperit Salvatórem. Quæ hæc Virgo tam sancta, ad quam veníre dignarétur Spíritus sanctus? Quæ tam speciósa, ut eam Deus sibi elígeret Sponsam? Quæ tam casta, ut esse possit Virgo post partum? Hæc est Dei templum, fons ille signátus, et porta in domo Dei clausa. Ad hanc namque, ut dixi, Spíritus sanctus descéndit~: hanc virtus Altíssimi obumbrávit.

\begin{tightlist}{}
\item \textred{℟.}~Quæ est ista, quæ procéssit sicut sol, et formósa tamquam Jerúsalem? \textred{*} Vidérunt eam fíliæ Sion, et beátam dixérunt, et regínæ laudavérunt eam.
\textred{℣.} Et sicut dies verni circúmdabant eam flores rosárum, et lília convállium. \textred{V}idérunt.
\item {\small \textred{Bened.} Fílius Vírginis Maríæ det nobis gáudia vitæ.}
\end{tightlist}

\subsubsection*{\center \textred{Lectio vj.}}
\lettrine{H}æc est immaculáta cóitu, fœcúnda partu, Virgo lactans, Angelórum et hóminum cibum nútriens. Mérito namque beáta singulári a nobis præcónio extóllitur, quæ singuláre commércium mundo prǽbuit. Dénique tantum se ad cœli fastígia sublevávit, ut Verbum in princípio apud Deum de supérna cœli arce suscíperet. O felix María et omni laude digníssima! o Virgo Dei Génitrix gloriósa! o sublímis puérpera! cujus viscéribus auctor cœli terrǽque commíttitur.

\begin{tightlist}{}
\item \textred{℟.}~Ornátam monílibus fíliam Jerúsalem Dóminus concupívit. \textred{*} Et vidéntes eam fíliæ Sion beatíssimam prædicavérunt, dicéntes~: unguéntum effúsum nomen tuum.
\textred{℣.} Adstitit regína a dextris tuis in vestítu deauráto, circúmdata varietáte. \textred{E}t. Glória Patri. \textred{E}t.
\end{tightlist}

\subsubsection*{\center \textsc{IN TERTIO NOCTURNO.}}
\begin{tightlist}{}
\item {\small \textred{Bened.} Fons Evangélii répleat nos dógmate cæli.}
\end{tightlist}
\subsubsection*{\center Léctio sancti Evangélii secúndum Joánnem.}
\subsubsection*{\center \textred{Lectio vij. \hfill Cap. 19.}}
\lettrine{I}n illo témpore~: Stabant juxta crucem Jesu mater ejus, et soror matris ejus María Cléophæ, et María Magdaléne. Et réliqua.

\begin{center}Homilía sancti Ambrósii Epíscopi.\\
\textred{Exhortat. ad Virginit.}\end{center}
\lettrineThree{H}ic Christus in vestris pectóribus alta máneat defíxus radíce, ut et fructus vester exúberet, et liquóre grátiæ spiritális, corpórei vapóris incéndia temperéntur. Hic est, qui venit in nube levi, sicut dixit Prophéta~: Ecce Dóminus sedet super nubem levem, et véniet in Ægýptum. Signíficans quod in Ægýptum, id est, in afflictiónem istam mundi hujus, veníret per Vírginem. Nubem ítaque Maríam dixit, quia carnem gerébat~: levem, quia Virgo erat, nullis onéribus graváta conjúgii. Ipsa est virga gérminans florem, quia pura, et ad Dóminum líbero corde dirécta virgínitas, quæ nullis in hoc sǽculo curárum anfráctibus reflectátur. Ideo illam Dóminus de cruce discípulo suo dilectíssimo sibi trádidit sancto Joánni, qui dixit patri et matri suæ, non novi vos. Dénique vocátus a Christo, relíquit patrem, Verbum secútus.

\begin{tightlist}{}
\item \textred{℟.}~Beátam me dicent omnes generatiónes~: \textred{*} Quia fecit mihi Dóminus magna, qui potens est~: et sanctum nomen ejus.
\textred{℣.} Et misericórdia ejus a progénie in progénies timéntibus eum. \textred{Q}uia.
\item {\small \textred{Bened.} Ab hoste malígno deféndat nos Dei Génitrix Virgo.}
\end{tightlist}

\subsubsection*{\center \textred{Lectio viij.}}
\lettrine{H}uic tráditur Virgo, qui suos nésciat. Huic tráditur Virgo, qui de péctore Christi háuriat sapiéntiam. Huic tráditur Virgo, qui fratres suos non cognóvit, et fílios suos non scivit. Ideóque benedícit eum lex~: date Levi veros ejus, date Levi sortes ejus. Unde et matrem Dómini ipse suscépit; scriptum enim habémus; quia ex illa hora suscépit eam discípulus in sua. Quid est, in sua, cum relíquerit patrem et Christum secútus sit? aut quómodo in sua, cum ipsi Apóstoli díxerint~: ecce nos relíquimus ómnia et secúti sumus te? Quæ sua habébat Joánnes, qui mundána et sæculária non habébat, quia de mundo non erat? quæ ergo habébat sua, nisi ea quæ a Christo accéperat? bonus verbi sapientiǽque posséssor, bonus recéptor grátiæ.

\begin{tightlist}{}
\item \textred{℟.}~Felix namque es, sacra Virgo María, et omni laude digníssima~: \textred{*} Quia ex te ortus est sol justítiæ Christus Deus noster.
\textred{℣.} Ora pro pópulo, intérveni pro clero, intercéde pro devóto femíneo sexu, séntiant omnes tuum juvámen, quicúmque célebrant tuam sanctam commemoratiónem.
\item {\small \textred{Bened.} In omni tribulatióne et angústia succúrrat nobis pia Virgo María.}
\end{tightlist}

\subsubsection*{\center \textred{Lectio ix.}}
\lettrine{A}udíte quæ Apóstoli a Christo accéperint. Accípite, inquit Spíritum sanctum; quorum remiséritis peccáta, remíssa erunt; et quorum detinuéritis, deténta. Neque enim mater Dómini Jesu, nisi ad possessórem grátiæ demigráret, ubi Christus habébat habitáculum. Date ergo et vos fílii mei, vero illi Levi veros ejus. Estóte nubes, sed leves. Eritis profécto, si ábluat virgínitas ónera conditiónis, et ténebras, hujus limósæ carnis illúminet. Ideo dicit illa, nigra sum et decóra, fíliæ Jerúsalem. Nigra per carnem, decóra per Virginitátem.

\begin{tightlist}{}
\item \textred{℟.}~Christi Virgo dilectíssima, virtútum operátrix, opem fer míseris~: \textred{*} Súbveni Dómina clamántibus ad te júgiter.
\textred{℣.} Quóniam peccatórum mole prémimur, et non est qui ádjuvet. \textred{S}úbveni. Glória Patri. \textred{S}úbveni.
\end{tightlist}

\begin{center}\textcolor{red}{\rule{\linewidth}{0.5pt}} \end{center}

\subsection*{\center \textsc{Officium Gloriosissimi Nominis beatæ Virginis Mariæ}}
\subsubsection*{\center \textred{Duplex Majus S.C.}}
\begin{center}\textred{In Festo Septembr. duplex majus S.C. quater in anno duplex minus primæ classis juxta Rubr. 16.}\end{center}
\fancyhead[C]{\textsc{Quater in Anno - Officium de Nomine B.V.M.}}

%\begin{center}\textcolor{red}{\rule{\linewidth}{0.5pt}} \end{center}

\subsection*{\center \textsc{AD VESPERAS.}}
\begin{center}\textred{Antiphona.}\end{center}
\lettrine{A} solis ortu usque ad occásum, laudábile nomen Dómini, et Maríæ Matris ejus.

\begin{tightlist}{}
\item \textred{Psalmi de j. Vesperis de commune B.M.V.} [168]
\item \textred{Aña.} Missus est Angelus Gábriel ad Vírginem desponsátam Joseph, et nomen Vírginis María.
\item \textred{Aña.} Ave María, grátia plena~: Dóminus tecum~: Benedícta tu in muliéribus.
\item \textred{Aña.} Permáneat, et magnificétur nomen tuum usque in sempitérnum, o María.
\item \textred{Aña.} Lauda Jerúsalem Dóminum, lauda ínclytum nomen Maríæ.
\item \textred{Quando hoc Officium recitatur sub ritu duplici minori primæ classis, dicitur in primis Vesperis Aña.} Mirábiles elatiónes maris, sed mirabílior in Nómine Maríæ Dóminus. \textred{prima Aña de Laudibus, Psalmi de Officio præcedente.}
\end{tightlist}

\subsubsection*{\textred{Capitulum. \hfill Judith 13.}}
\lettrineThree{B}enedícta es tu fília a Dómino Deo excélso, præ ómnibus muliéribus super terram~: quia hódie nomen tuum ita magnificávit, ut non recédat laus tua de ore hóminum, qui mémores fúerint virtútis Dómini in ætérnum.

\begin{tightlist}{}
\item \textred{℟.}~Felix namque es sacra Virgo María, et omni laude digníssima~: \textred{*} Quia ex te ortus est Sol justítiæ, Christus Deus noster.
\textred{℣.}~Ora pro pópulo, intérveni pro clero, intercéde pro devóto femíneo sexu~: séntiant omnes tuum juvámen, quicúmque célebrant tui Nóminis commemoratiónem. \textred{Q}uia. Glória Patri. \textred{Q}uia.
\item \textred{Hymnus.} Ave, maris stella. [171]
\item \textred{℣.}~Sit nomen Dóminæ benedíctum. \textred{℟.}~Ex hoc nunc, et usque in sǽculum, allelúja.
\item \textred{Ad Magnificat, Aña.} Maríæ Nomen cunctas illústrat Ecclésias, cui fecit magna qui potens est, et sanctum Nomen ejus.
\end{tightlist}

\subsubsection*{\center \textred{Oratio.}}
\lettrine{D}eus, qui gloriósam Matrem tuam nominári Maríam voluísti~: concéde quǽsumus; ut, qui dulce ejúsdem Maríæ Nomen implórant, perpétuum séntiant tuæ benedictiónis efféctum. Qui vivis et regnas.

\begin{tightlist}{}
\item \textred{Ad Complet. et alias Horas, in fine Hymn. dicitur:} Glória tibi Dómine, Qui natus es.
\item \textred{Ad Nunc dimittis, Aña.} Allelúja~: Sancta Dei Génitrix.
\item \textred{Sed post Septuagesimam, Aña.} Sub tuum præsídium.
\end{tightlist}

\subsubsection*{\center \textsc{AD MATUTINUM.}}
\begin{tightlist}{}
\item \textred{Invitatorium.} Ave María, grátia plena~: \textred{*} Dóminus tecum.
\item \textred{Ps.} Veníte, exsultémus. 2.
\item \textred{Hymn.} Quem terra. [172]
\end{tightlist}

\subsubsection*{\center \textsc{IN PRIMO NOCTURNO.}}
\begin{tightlist}{}
\item \textred{Aña.} In univérsa terra admirábile est Nomen tuum, María.
\item \textred{Ps.} Dómine, Dóminus noster, \textred{cum reliquis de Comm. B.M.V.} [172]
\item \textred{Aña.} In omnem terram diffúsa est laus Nóminis Maríæ.
\item \textred{Aña.} Super mária et flúmina fundávit Maríam, et dedit illi Nomen virtútum Rex glóriæ.
\item \textred{℣.}~Magnificáte Maríam vírginem mecum. \textred{℟.}~Et exaltémus Nomen ejus in ætérnum, allelúja.
\item {\small \textred{Bened.} Alma Virgo vírginum intercédat pro nobis ad Dóminum.}
\end{tightlist}

\subsubsection*{\center De libro Exodi.}
\subsubsection*{\center \textred{Lectio j. \hfill Cap. 15.}}
\lettrineThree{S}umpsit María Prophetíssa, soror Aáron, týmpanum in manu sua, egressǽque sunt omnes mulíeres post eam cum týmpanis et choris, quibus præcinébat, dicens~: Cantémus Dómino; glorióse enim magnificátus est, equum et ascensórem ejus dejécit in mare.

\begin{tightlist}{}
\item \textred{℟.}~Laudémus Nomen sacrum atque insígne Maríæ~: \textred{*} Mundo enim péperit Virgo salútis opem. \textred{℣.}~Læténtur cœli de tali Nómine et orbis. \textred{M}undo.
\item {\small \textred{Bened.} Nos cum prole pia benedícat Virgo María.}
\end{tightlist}

\subsubsection*{\center \textred{Lectio ij.}}
\lettrine{T}ulit autem Móyses Israel de mari Rubro, et egréssi sunt in desértum Sur~: ambulaverúntque tribus diébus per solitúdinem, et non inveniébant aquam. Et venérunt in Mara, nec póterant bíbere aquas de Mara, eo quod essent amáræ~: unde et cóngruum loco nomen impósuit, vocans illum Mara, id est, amaritúdinem.

\begin{tightlist}{}
\item \textred{℟.}~Benedíctum sit Nomen Maríæ, \textred{*} Cujus manus tornátiles plenæ hyacínthis. \textred{℣.}~Implorántibus Maríæ Nomen, María præsto adest. \textred{C}ujus.
\item {\small \textred{Bened.} Sancta Dei Génitrix sit nobis auxiliátrix.}
\end{tightlist}

\subsubsection*{\center \textred{Lectio iij.}}
\lettrine{E}t murmurávit pópulus contra Móysen, dicens~: Quid bibémus? At ille clamávit ad Dóminum, qui osténdit ei lignum~: quod cum misísset in aquas, in dulcédinem versæ sunt. Ibi constítuit ei præcépta atque judícia, et ibi tentávit eum, dicens~: Si audíeris vocem Dómini Dei tui, et quod rectum est coram eo féceris, et obedíeris mandátis ejus, custodierísque ómnia præcépta illíus, cunctum languórem, quem pósui in Ægýpto, non indúcam super te~: ego enim Dóminus sanátor tuus.

\begin{tightlist}{}
\item \textred{℟.}~Salutémus Maríam, pro nobis Fílium exorántem, \textred{*} Et illi dicámus~: Ave María, grátia plena, Dóminus tecum. \textred{℣.}~Hæc est Gedeónis vellus, quæ levíssima oratiónis compressióne, larga éxhibet gratiárum stillicídia. \textred{E}t. Glória Patri. \textred{E}t.
\end{tightlist}

\subsubsection*{\center \textsc{IN SECUNDO NOCTURNO.}}
\begin{tightlist}{}
\item \textred{Aña.} Speciósum Maríæ Nomen deprecabúntur omnes dívites plebis, cujus decórem concupívit Rex.
%\item \textred{Ps.} Eructávit. [169]
\item \textred{Aña.} Refúgium est in tribulatiónibus Maríæ Nomen; ómnibus illud invocántibus.
\item \textred{Aña.} María cívitas Dei sancta, gloriósa dicta sunt de Nómine sancto tuo.
\item \textred{℣.}~Repleátur os meum laude Maríæ. \textred{℟.}~Ut cantem glóriam Nómini sancto ejus, allelúja.
\item {\small \textred{Bened.} Précibus suæ Matris benedícat nos Fílius Dei Patris.}
\end{tightlist}

\subsubsection*{\center De Sermóne sancti Epiphánii Epíscopi.}
\subsubsection*{\center \textred{Orat. de Laud. B.V.M.\\
Lectio iv.}}
\lettrineThree{B}eáta Anna grávida effécta, cœlum, et thronum Cherúbicum péperit sanctam puéllam Maríam~: illa enim reperítur esse cœlum, templum, et thronus; quia Maríam interpretári solémus Dóminam, atque étiam spem. Péperit enim Dóminum, qui est spes totíus mundi, nempe Christum. Rursum, nomen illud, María, interpretátur, myrrha maris; Myrrham vero dicet áliquis (quod ego étiam ajo) de immortalitáte, eo quod paritúra esset gemmam immortálem in mari, hoc est in mundo. Mare autem dicet univérsum mundum, cui Virgo serenitátem, et tranquillitátem cóntulit, dum portum péperit Christum.

\begin{tightlist}{}
\item \textred{℟.}~Benedícta tu a Dómino Deo tuo in omni tabernáculo Jacob~: \textred{*} Quóniam in omni Gente, quæ audíerit nomen tuum, magnificábitur super te Deus. \textred{℣.}~Tu glória Jerúsalem, tu lætítia Israel, tu María honorificéntia pópuli nostri. \textred{Q}uóniam.
\item {\small \textred{Bened.} Jesus Maríæ Fílius sit nobis clemens et propítius.}
\end{tightlist}

\subsubsection*{\center \textred{Lectio v.}}
\lettrine{R}ursum ítaque præcláræ puéllæ Maríæ beátum nomen interpretátur Illustráta, útpote quæ illustráta sit a Fílio Dei, et cunctos usque ad términos terræ credéntes Trinitáti illumináverit. Sancta enim puélla Virgo María est Sponsa Trinitátis, et thesáurus dispensatiónis plane arcánus, cui Gábriel inquit~: Ave grátia plena, Dóminus tecum. Subjúnxit Gábriel~: Et Pater misit arrhabónem de cœlo Spíritum sanctum, præparávit Vírginem unigénito Fílio, cœlésti sponso, quam Pater diléxit, Fílius inhabitávit, Spíritus sanctus percupívit.

\begin{tightlist}{}
\item \textred{℟.}~Mélius est nomen bonum, quam divítiæ multæ~: \textred{*} Turris fortíssima Maríæ nomen. \textred{℣.}~Propugnáculum est ómnibus, qui ad illud confúgiunt. \textred{T}urris.
\item {\small \textred{Bened.} Fílius Vírginis Maríæ det nobis gáudia vitæ.}
\end{tightlist}

\subsubsection*{\center \textred{Lectio vj.}}
\lettrine{I}psa est enim sponsa, et thálamus~: et ex ea Sponsus procédit Christus virgíneum induméntum juxta suávem Davídis Prophétæ cantiónem~: In sole pósuit tabernáculum suum, et ipse tamquam sponsus procédens de thálamo suo. O sponsi desidérium summæ puritátis! própriam ancíllam propensióre afféctu constítuit sponsam et matrem.

\begin{tightlist}{}
\item \textred{℟.}~Nomen sanctum meum notum fáciam in médio pópuli mei~: \textred{*} Et scient, quia ego María. \textred{℣.}~Delíciæ meæ, esse cum fíliis hóminum. \textred{E}t. Glória Patri. \textred{E}t.
\end{tightlist}

\subsubsection*{\center \textsc{IN TERTIO NOCTURNO.}}
\begin{tightlist}{}
\item \textred{Aña.} Nuntiáte inter Gentes glóriam Maríæ, in ómnibus pópulis mirabília nóminis ejus.
%\item \textred{Ps.} Cantáte. [172]
\item \textred{Aña.} Annuntiavérunt cœli nomen Maríæ~: et vidérunt omnes pópuli glóriam ejus.
\item \textred{Aña.} Notum fecit Dóminus sacrum Maríæ nomen, et in conspéctu Géntium revelávit virtútem nóminis ejus.
\item \textred{℣.}~Ave María grátia plena. \textred{℟.}~Dóminus tecum, allelúja.
\item {\small \textred{Bened.} Fons Evangélii répleat nos dógmate cæli.}
\end{tightlist}

\subsubsection*{\center Léctio sancti Evangélii secúndum Lucam.}
\subsubsection*{\center \textred{Lectio vij. \hfill Cap. 1.}}
\lettrine{I}n illo témpore~: Missus est Angelus Gábriel a Deo in civitátem Galilǽæ, cui nomen Názareth, ad Vírginem desponsátam viro, cui nomen erat Joseph, de domo David, et nomen Vírginis María. Et réliqua.

\begin{center}Homilía sancti Bernárdi Abbátis.\\
\textred{Hom. super Missus est.}\end{center}
\lettrineThree{E}t nomen, inquit, Vírginis María. Loquámur pauca et super hoc nómine, quod interpretátum, maris stella dícitur, et matri Vírgini valde conveniénter aptátur. Ipsa namque aptíssime síderi comparátur~: quia sicut sine sui corruptióne sidus suum emíttit rádium, sic absque sui læsióne Virgo partúriit Fílium. Nec síderi rádius suam mínuit claritátem, nec Vírgini Fílius suam integritátem. Ipsa est ígitur nóbilis illa stella ex Jacob orta, cujus rádius univérsum orbem illúminat, cujus splendor et præfúlget in supérnis, et ínferos pénetrat, terras étiam perlústrans et calefáciens magis mentes, quam córpora~: fovet virtútes, éxcoquit vítia. Ipsa, inquam, est præclára et exímia stella, super hoc mare magnum et spatiósum necessário subleváta, micans méritis, illústrans exémplis.

\begin{tightlist}{}
\item \textred{℟.}~Claríssimum Maríæ Nomen ipsa stella est præfúlgens~: \textred{*} Peccatóres illúminat, et dírigit in portum salútis. \textred{℣.}~Consolátur tristes, míseros levat, et jacéntes érigit. \textred{P}eccatóres.
\item {\small \textred{Bened.} Ab hoste malígno deféndat nos Dei Génitrix Virgo.}
\end{tightlist}

\subsubsection*{\center \textred{Lectio viij.}}
\lettrine{O} quisquis te intélligis in hujus sǽculi proflúvio magis inter procéllas et temptestátes fluctuáre, quam per terram ambuláre, ne avértas óculos a fulgóre hujus síderis, si non vis óbrui procéllis. Si insúrgant venti tentatiónum, si incúrras scópulos tribulatiónum, réspice stellam, voca Maríam. Si jactáris supérbiæ undis, si ambitiónis, si detractiónis, si æmulatiónis, réspice stellam, voca Maríam. Si iracúndia, aut avarítia, aut carnis illécebra navículam concússerit mentis, réspice ad Maríam. Si críminum immanitáte turbátus, consciéntiæ fœditáte confúsus, júdicis horróre pertérritus, bárathro incípias absorbéri tristítiæ, desperatiónis abýsso, cógita Maríam.

\begin{tightlist}{}
\item \textred{℟.}~Omnia flúmina intrant in mare, et ipsum non redúndat. \textred{*} Omnia gratiárum charísmata in Maríam, et ipsa non intuméscit. \textred{℣.}~Dixit illi Angelus~: Ave grátia plena. At ipsa respóndit~: Ecce ancílla Dómini. \textred{O}mnia.
\item {\small \textred{Bened.} In omni tribulatióne et angústia succúrrat nobis pia Virgo María.}
\item {\small \textred{Si ix. Lectio fuerit de Hom. Evang., erit}}
\item {\small \textred{Bened.} Verba Redemptóris pia sint suffrágia nobis.}
%\item \textred{Si ix. Lectio fuerit de Hom. Evang. Benedictio.} Verba Redemptóris pia sint suffrágia nobis.
\end{tightlist}

\subsubsection*{\center \textred{Lectio ix.}}
\lettrine{I}n perículis, in angústiis, in rebus dúbiis Maríam cógita, Maríam ínvoca. Non recédat ab ore, non recédat a corde~: et ut ímpetres ejus oratiónis suffrágium, non déseras conversatiónis exémplum. Ipsam sequens, non dévias~: ipsam rogans, non despéras~: ipsam cógitans, non erras~: ipsa tenénte, non córruis~: ipsa protegénte, non métuis~: ipsa duce, non fatigáris~: ipsa propítia, pérvenis~: et sic in temetípso experíris, quam mérito dictum sit~: Et nomen Vírginis María.

\begin{tightlist}{}
\item \textred{℟.}~Felix namque es sacra Virgo María, et omni laude digníssima~: \textred{*} Quia ex te ortus est Sol justítiæ, Christus Deus noster. \textred{℣.}~Ora pro pópulo, intérveni pro clero, intercéde pro devóto femíneo sexu~: séntiant omnes tuum juvámen, quicúmque célebrant tui Nóminis commemoratiónem. \textred{Q}uia. Glória Patri. \textred{Q}uia.
\item Te Deum laudámus. 6.
\item \textred{Sacerdotalis ℣.}~Sit nomen Dóminæ benedíctum. \textred{℟.}~Ex hoc nunc et usque in sǽculum, allelúja.
\end{tightlist}

\subsubsection*{\center \textsc{AD LAUDES.}}
\begin{center}\textred{et per Horas, Aña.}\end{center}
\lettrine{M}irábiles elatiónes maris, sed mirabílior in Nómine Maríæ Dóminus.

\begin{tightlist}{}
\item \textred{Ps. de Dominica.} 27.
\item \textred{Aña.} Laudáte in exsultatióne Maríæ Nomen, quóniam suáve est.
\item \textred{Aña.} Lábia mea laudábunt te, o María, quia dulce est Nomen tuum super mel ori meo.
\item \textred{Aña.} Spíritus et ánimæ justórum, benedícant Maríam in Nómine sancto ejus.
\item \textred{Aña.} Beneplácitum est Dómino in Nómine Maríæ, in quo exaltávit mansuétos in salútem.
\end{tightlist}

\subsubsection*{\textred{Capitulum. \hfill Judith 13.}}
\lettrineThree{B}enedícta es tu fília a Dómino Deo excélso, præ ómnibus muliéribus super terram~: quia hódie nomen tuum ita magnificávit, ut non recédat laus tua de ore hóminum, qui mémores fúerint virtútis Dómini in ætérnum.

\begin{tightlist}{}
\item \textred{Hymn.} O gloriósa Dómina. [183]
\item \textred{℣.}~Lætábor et exsultábo in te, María. \textred{℟.}~Psalmum dicam Nómini tuo, allelúja.
\item \textred{Ad Benedictus, Aña.} Benedíctus Deus, qui Maríæ Vírgini lucis Nomen donávit, ut illumináret eos, qui in ténebris sedent~: et ut stella maris errántes dirígeret in portum pacis.
\end{tightlist}

\subsubsection*{\center \textred{Oratio.}}
\lettrine{D}eus, qui gloriósam Matrem tuam nominári Maríam voluísti~: concéde quǽsumus; ut, qui dulce ejúsdem Maríæ Nomen implórant, perpétuum séntiant tuæ benedictiónis efféctum. Qui vivis et regnas.

\subsubsection*{\center \textsc{AD PRIMAM.}}
\begin{tightlist}{}
\item \textred{Ad Horas Antiphonæ de Laudibus, Psalmi de Dominica, ad Primam tamen ut in Festis.}
\item \textred{℣.}~Qui de Vírgine dignátus es nasci.
\end{tightlist}

\subsubsection*{\center \textsc{AD TERTIAM.}}
\subsubsection*{\textred{Capitulum. \hfill Judith 13.}}
\lettrine{B}enedícta es tu fília a Dómino Deo excélso, præ ómnibus muliéribus super terram~: quia hódie nomen tuum ita magnificávit, ut non recédat laus tua de ore hóminum, qui mémores fúerint virtútis Dómini in ætérnum.

\begin{tightlist}{}
%\item \textred{Capit.} Benedícta es tu, \textred{ut supra ad Laudes.}
\item \textred{℟.}~Magnificáte Maríam Vírginem mecum, \textred{*} Allelúja, allelúja. \textred{M}agnificáte. \textred{℣.}~Et exaltémus Nomen ejus in ætérnum. \textred{A}llelúja, allelúja. Glória Patri. \textred{M}agnificáte.
\item \textred{℣.}~Repleátur os meum laude Maríæ. \textred{℟.}~Ut cantem glóriam nómini sancto ejus, allelúja.
\end{tightlist}

\subsubsection*{\center \textsc{AD SEXTAM.}}
\subsubsection*{\textred{Capitulum. \hfill Judith 13.}}
\lettrine{B}enedícta tu a Deo tuo in omni tabernáculo Jacob; quóniam in omni Gente, quæ audíerit nomen tuum, magnificábitur super te Deus Israel.

\begin{tightlist}{}
\item \textred{℟.}~Repleátur os meum laude Maríæ, \textred{*} Allelúja, alleluja. \textred{R}epleátur. \textred{℣.}~Ut cantem glóriam nómini sancto ejus. \textred{A}llelúja, allelúja. Glória Patri. \textred{R}epleátur.
\item \textred{℣.}~Ave María, grátia plena. \textred{℟.}~Dóminus tecum, allelúja.
\end{tightlist}

\subsubsection*{\center \textsc{AD NONAM.}}
\subsubsection*{\textred{Capitulum. \hfill Ezech. 16.}}
\lettrine{E}gréssum est nomen tuum in Gentes propter spéciem tuam~: quia perfécta eras in decóre meo, quem posúeram super te~: dicit Dóminus Deus.

\begin{tightlist}{}
\item \textred{℟.}~Ave María, grátia plena, \textred{*} Alleluja, alleluja. \textred{A}ve María. \textred{℣.}~Dóminus tecum. \textred{A}llelúja, allelúja. Glória Patri. \textred{A}ve María.
\item \textred{℣.}~Lætábor et exsultábo in te, María. \textred{℟.}~Psalmum dicam Nómini tuo, alleluja.
\end{tightlist}

\subsubsection*{\center \textsc{AD II VESPERAS.}}
\begin{tightlist}{}
\item \textred{Añæ Laudum}
\item \textred{Psalmi de Comm. B.M.V.} [184]
%\item \textred{Capitulum Hymnus, et ℣. ut in j Vesperis.}
\end{tightlist}

\subsubsection*{\textred{Capitulum. \hfill Judith 13.}}
\lettrineThree{B}enedícta es tu fília a Dómino Deo excélso, præ ómnibus muliéribus super terram~: quia hódie nomen tuum ita magnificávit, ut non recédat laus tua de ore hóminum, qui mémores fúerint virtútis Dómini in ætérnum.

\begin{tightlist}{}
%\item \textred{℟.}~Felix namque es sacra Virgo María, et omni laude digníssima~: \textred{*} Quia ex te ortus est Sol justítiæ, Christus Deus noster.
%\textred{℣.}~Ora pro pópulo, intérveni pro clero, intercéde pro devóto femíneo sexu~: séntiant omnes tuum juvámen, quicúmque célebrant tui Nóminis commemoratiónem. \textred{Q}uia. Glória Patri. \textred{Q}uia.
\item \textred{Hymnus.} Ave maris stella. [187]
\item \textred{℣.}~Sit nomen Dóminæ benedíctum. \textred{℟.}~Ex hoc nunc, et usque in sǽculum, allelúja.
\item \textred{Ad Magnificat, Aña.} María Virgo ómnium óculis gratiósa, nominátum est nomen glóriæ ejus usque in extrémum terræ.
\end{tightlist}

%\subsubsection*{\center \textred{Oratio.}}
%\lettrine{D}eus, qui gloriósam Matrem tuam nominári Maríam voluísti~: concéde quǽsumus; ut, qui dulce ejúsdem Maríæ Nomen implórant, perpétuum séntiant tuæ benedictiónis efféctum. Qui vivis et regnas.

\begin{center}\textcolor{red}{\rule{\linewidth}{0.5pt}} \end{center}
\clearpage
\subsection*{\center \textred{\textsc{Officia Mensualia}}
\subsubsection*{\center \textred{Dicenda singulis mensibus}}

\subsection*{\center \textsc{Officium S.P.N. Eliæ.}}
\subsubsection*{\center \textred{Semiduplex.}}
\begin{center}\textred{Recitandum semel in mense, juxta Rubr. 17.}\end{center}
\fancyhead[C]{\textsc{Officia Mensualia - Officium S.P.N. Eliæ}}

%\begin{center}\textcolor{red}{\rule{\linewidth}{0.5pt}} \end{center}
\begin{tightlist}{}
\item \textred{Temp. Pasch. tres Psalmi cujuslibet Nocturni dicuntur sub una sola Aña; et in fine omnium Añarum, et Respons. Versuum additur} allelúja, \textred{et ultimus versus Hymni ad Mat. mutatur.}
\item \textred{Reliqua omnia ut infra.}
\end{tightlist}

\subsubsection*{\center \textsc{AD I. VESPERAS.}}
\subsubsection*{\textred{Capitulum. \hfill Eccli. 48.}}
\lettrineThree{S}urréxit Elías Prophéta quasi ignis, et verbum ipsíus quasi fácula ardébat. Qui indúxit in pópulos famem, et irritántes illum invídia sua pauci facti sunt.

\subsubsection*{\center \textred{Hymnus.}}
\lettrine{N}unc juvat celsi properáre ad alta\\
Mente Carméli juga~: et ore pleno\\
Vatis Elíæ méritos reférre\\
Laudis honóres.\\
\\
\textred{H}ic sacri est ductor, cólumen, decúsque\\
Ordinis nostri, superísque gratus~:\\
O Pater salve, plácidus tuósque\\
Excipe alúmnos.\\
\\
\textred{D}um fames terras prémeret, minístro\\
Uteris corvo; víduæ farínam\\
Súfficis solers; oleóque manat\\
Hýdria pingui.\\
\\
\textred{T}u, Deo fusis précibus, redúctum\\
Fínibus mortis púerum parénti\\
Illius reddis víduæ, senéctæ\\
Dulce levámen.\\
\\
\textred{C}laudis et nubes, aperísque verbo,\\
Imperas post tres properánter annos,\\
Décidant cœlo, sitiénsque tellus\\
Háuriat imbres.\\
\\
\textred{S}it Patri summo, genitǽque Proli,\\
Et tibi compar utriúsque sancte\\
Spíritus, laus, impérium, potéstas\\
Tempus in omne.\\
Amen.

\begin{tightlist}{}
\item \textred{℣.}~Verbo Dómini contínuit cœlum. \textred{℟.}~Et dejécit de cœlo ignem ter.
\item \textred{Ad Magnificat, Aña.} Ecce ego mittam vobis Elíam Prophétam, ántequam véniat dies Dómini magnus, et horríbilis. Et convértet cor patrum ad fílios, et cor filiórum ad patres eórum.
\end{tightlist}

\subsubsection*{\center \textred{Oratio.}}
\lettrine{P}ræsta, quǽsumus omnípotens Deus, ut sicut beátum Elíam Prophétam tuum, et Patrem nostrum ante commúnem mortem curru ígneo aéreum elevásti ad cœlum~: ita nos fácias, eo interveniénte, dum vívimus, a terrénis semper ad cœléstia spíritu sublevári; et cum eo in resurrectióne justórum páriter gaudére. Per Dóminum.

\subsubsection*{\center \textsc{AD MATUTINUM.}}
\begin{tightlist}{}
\item \textred{Invitatorium.} Regem Prophetárum Dóminum, \textred{*} Veníte adorémus.
\item \textred{Ps.} Venite, exsultémus. 2.
\end{tightlist}

\subsubsection*{\center \textred{Hymnus.}}
\lettrine{T}e magne rerum Cónditor,\\
Mens nostra gliscit láudibus\\
In hoc Thesbíte máximo,\\
Quem díligis extóllere.\\
\\
\textred{H}ic namque sacri nóminis\\
Tui zelátor próvocat\\
Vates Baal nequíssimos,\\
Victósque jure intérficit.\\
\\
\textred{I}llo precánte, víctimas\\
Absúmit ignis cœ́litus~:\\
Hi perstrepéntes ácriter,\\
Sunt ómnibus ludíbrio.\\
\\
\textred{T}um Jezabélis ímpiæ\\
Vitans furórem nóxium,\\
Sub junípero dórmiens \\
Adésse cernit Angelum.\\
\\
\textred{P}anem sibi qui próferens,\\
Lymphámque, jussit pérgere\\
Cibo reféctum strénue\\
Ad montis Horeb vérticem.\\
\\
\textred{H}oc in cibo jejúnium\\
Tulit quadragenárium~:\\
Hac vi dapes sunt prǽditæ,\\
Quas dextra Dei pórrigit.\\
\\
\textred{O}mnis tibi sit glória\\
Verbum, Pater, Paráclite,\\
Inseparáta Trínitas,\\
Quæ condidísti ómnia.\\
Amen.

\subsubsection*{\center \textsc{IN PRIMO NOCTURNO.}}
\begin{tightlist}{}
\item \textred{Aña.} Factum est autem, ut ægrotáret fílius mulíeris matris famílias, et erat languor fortíssimus.
\item \textred{Ps.} Beátus vir, \textred{cum reliquis, de Communi unius Conf. Pont.} [76]
\item \textred{Aña.} Dixit múlier ad Elíam~: Quid mihi et tibi, vir Dei? Ingréssus es ad me, ut rememoraréntur iniquitátes meæ, et interfíceres fílium meum?
\item \textred{Aña.} Ait Elías ad víduam~: Da mihi fílium tuum. Et clamávit ad Dóminum; et revíxit puer.
\item \textred{℣.}~Verbo Dómini contínuit cœlum. \textred{℟.}~Et dejécit de cœlo ignem ter.
\item \textred{Lectiones de Scriptura occurrente.}
\item \textred{℟.~j.} Recéde hinc, et vade contra Oriéntem, dixit Dóminus ad Elíam. \textred{*} Qui ábiit, et fecit juxta verbum Dómini. \textred{℣.}~Abscóndere in Torrénte Carith, et ibi de torrénte bibes~: corvísque præcépi, ut pascant te ibi. \textred{Q}ui.
\item \textred{℟.~ij.} Abiit múlier, et fecit juxta verbum Prophétæ~: \textred{*} Et comédit ipse, et illa, et domus ejus. \textred{℣.}~Ex illa die hýdria farínæ non defécit, et lécythus ólei non est imminútus, juxta verbum Elíæ. \textred{E}t.
\item \textred{℟.~iij.} Revérsa est ánima púeri intra eum, et revíxit, et trádidit matri suæ, et ait illi~: \textred{*} En vivit fílius tuus. \textred{℣.}~Tulit Elías púerum, et depósuit eum in inferiórem domum, dicens. \textred{E}n. Glória Patri. \textred{E}n.
\end{tightlist}

\subsubsection*{\center \textsc{IN SECUNDO NOCTURNO.}}
\begin{tightlist}{}
\item \textred{Aña.} Factum est verbum Dómini ad Elíam, dicens~: Vade, et osténde te Achab, ut dem plúviam super fáciem terræ.
%\item \textred{Ps.} Cum invocárem. [67]
\item \textred{Aña.} Ait Elías ad Achab~: Ascénde, cómede et bibe~: quia sonus multæ plúviæ est.
\item \textred{Aña.} Manus Dómini facta est super Elíam, accinctísque lumbis currébat ante Achab, donec veníret in Jézræl.
\item \textred{℣.}~Tu ungis reges ad pœniténtiam. \textred{℟.}~Et Prophétas facis successóres post te.
\end{tightlist}

\subsubsection*{\center Sermo sancti Joánnis Chrysóstomi.}
\subsubsection*{\center \textred{Hom. 2. ad pop. Antioch.\\
Lectio iv.}}
\lettrineThree{I}n ómnibus ígitur philosophémur, et nihil nos inópia lǽdere póterit~: quin et máxime adjuvábit, et clarióres reddet, et ómnibus divítibus copiosióres. Quid enim, dic mihi, Elía paupérius? Sed proptérea omnes dívites vincébat, quóniam pauper hic erat; ipsam vero paupertátem ex mentis opuléntia elégit. Quóniam enim omnem pecuniárum cópiam ánimi sui magnitúdine esse putávit inferiórem, nec philosophía sua dignam, ídeo tantam paupertátem amplexátus est.

\begin{tightlist}{}
\item \textred{℟.}~Ait Abdías ad Elíam~: Cum recéssero a te, Spíritus Dómini asportábit te in locum, quem ego ignóro, et tu dicis~: \textred{*} Dic dómino tuo~: Adest Elías. \textred{℣.}~Dixit Elías~: Vivit Dóminus exercítuum, ante cujus vultum sto, quia hódie apparébo ei. \textred{D}ic.
\end{tightlist}

\subsubsection*{\center \textred{Lectio v.}}
\lettrine{Q}uod si magna præséntia putásset, non solam melótem possedísset~: sed ita vitæ vanitátem condemnávit, et tamquam lutum projéctum, aurum omne despéxit, ut præter amíctum illum, nihil ámplius possidéret. Idcírco rex páuperis indigébat, et verbis nihil plus quam melótem habéntis inhiábat, qui tantum habébat auri, ádeo splendídior trábea erat melóte, et regálibus aulis justi spelúnca. Proptérea et in cœlum ascéndens, nihil áliud quam melótem discípulo relíquit. Cum hac (inquit) cum diábolo pugnávi, et tu contra illum armátus hanc cápies.

\begin{tightlist}{}
\item \textred{℟.}~Cum venísset Elías ad portam civitátis, appáruit ei múlier vídua cólligens ligna, et vocávit eam, dixítque ei~: \textred{*} Da mihi páululum aquæ in vase, ut bibam. \textred{℣.}~Cumque illa pérgeret ut afférret, clamávit post tergum ejus, dicens~: Affer mihi, óbsecro, et buccéllam panis in manu tua. \textred{D}a.
\end{tightlist}

\subsubsection*{\center \textred{Lectio vj.}}
\lettrine{V}alidum enim telum inópia, et inexpugnábile habitáculum, et turris inconcússa. Tamquam máximam hereditátem Eliséus melótem suscépit~: étenim vere máxima fuit heréditas, et omni auro pretiósior. Et erat posthac duplex Elías ille~: et erat sursum Elías, et deórsum Elías. Novi quod justum illum beátum putátis, et vellétis quisque esse ut ille. Quod ígitur si vobis demonstrávero quid áliud, quod illo multo majus omnes sacris mystériis imbúti recípimus? Elías nempe melótem discípulo relíquit~: Fílius autem Dei ascéndens, suam nobis carnem dimísit~: sed Elías quidem exútus, Christus autem et nobis relíquit, et ipsam habens ascéndit.

\begin{tightlist}{}
\item \textred{℟.}~Expándit se, atque mensus est Elías super púerum tribus vícibus, et clamávit ad Dóminum~: \textred{*}~Et exaudívit Dóminus vocem Elíæ. \textred{℣.}~Dómine Deus meus, revertátur óbsecro ánima púeri hujus in víscera ejus. \textred{E}t. Glória Patri. \textred{E}t.
\end{tightlist}

\subsubsection*{\center \textsc{IN TERTIO NOCTURNO.}}
\begin{tightlist}{}
\item \textred{Aña.} Dixit Elías prophétis Baal~: Elígite vobis bovem unum, et invocáte nómina deórum vestrórum. Qui clamábant, et non erat qui respondéret eis.
%\item \textred{Ps.} Dómine, quis. [70]
\item \textred{Aña.} Cum jam tempus esset, ut offerrétur holocáustum, accédens Elías orávit ad Dóminum.
\item \textred{Aña.} Oránte Elía cécidit ignis Dómini, et vorávit holocáustum, et ligna, et lápides, púlverem quoque, et aquam quæ erat in aquædúctu lambens.
\item \textred{℣.}~Beáti sunt qui te vidérunt. \textred{℟.}~Et in amicítia tua decoráti sunt.
\end{tightlist}

\subsubsection*{\center Léctio sancti Evangélii secúndum Lucam.}
\subsubsection*{\center \textred{Lectio vij. \hfill Cap. 9.}}
\lettrine{I}n illo témpore~: Assúmpsit Jesus Petrum, et Jacóbum, et Joánnem, et ascéndit in montem ut oráret, et facta est dum oráret, spécies vultus ejus áltera. Et réliqua. 

\begin{center}Homilía venerábilis Bedæ Presbýteri.\end{center}
\lettrineThree{E}rant autem Móyses et Elías visi in majestáte, et dicébant excéssum ejus, quem completúrus erat in Jerúsalem. Móyses ergo et Elías, qui locúti sunt cum Dómino in monte, et passiónem ejus ac resurrectiónem dicébant, orácula legis et prophetíæ desígnant, quæ in Dómino compléta, et nunc doctis quibúsque patent, et eléctis ómnibus in futúro maniféstius patébunt.

\begin{tightlist}{}
\item \textred{℟.}~Ait Elías núntiis regis Samaríæ~: Numquid non est Deus in Israel, ut eátis ad Beélzebub? Quam ob rem hæc dicit Dóminus~: \textred{*} De léctulo, super quem ascendísti, non descéndes, sed morte moriéris. \textred{℣.}~Revérsi núntii ad Ochozíam dixérunt~: Vir occúrrit nobis, dicens~: Ite, et revertímini ad regem, et dicétis ei. \textred{D}e.
\end{tightlist}

\subsubsection*{\center \textred{Lectio viij.}}
\lettrine{Q}ui bene visi in majestáte dicúntur, quia tunc apértius videbúntur, quando veritátis omnis divinórum eloquiórum non solum sensus, sed et sermo fúerit prolátus~: quamvis cum Móyse et Elía omnes, qui cum Dómino regnatúri sunt, recte intélligi possint. In Móyse quidem, qui mórtuus et sepúltus est, hi qui in judício resuscitándi a morte; in Elía vero, qui necdum mortis débitum solvit, hi qui in advéntu Júdicis, vivi inveniéndi sunt in carne. Qui utríque páriter in uno eodémque moménto rapti in núbibus óbviam Dómino in áera, mox perácto judício ad vitam inducéntur ætérnam.

\begin{tightlist}{}
\item \textred{℟.}~Ascéndit quinquagenárius, et quinquagínta qui erant sub eo, ad Elíam sedéntem in vértice montis. \textred{*} Et descéndit ignis de cœlo, et devorávit eos. \textred{℣.}~Dixit Elías~: Si homo Dei sum, descéndat ignis de cœlo, et dévoret te et quinquagínta tuos. \textred{E}t.
\end{tightlist}

\subsubsection*{\center \textred{Lectio ix.}}
\lettrine{Q}uibus apte cónvenit, quod Móyses et Elías in majestáte fuísse visi perhibéntur~: per excelléntiam namque majestátis, insígne múneris, quo sunt coronándi, monstrátur. Bene cónvenit et hoc, quod excéssum ejus, quem completúrus erat in Jerúsalem, dixísse memorántur~: quia nimírum única laudis matéria fidélibus suis sit pássio Redemptóris, tanto suæ reparatiónis memóriam devóto semper in péctore versant, devóta confessióne testántur.

\begin{tightlist}{}
\item \textred{℟.}~Factum est, dum tólleret Dóminus Elíam per túrbinem in cœlum, \textred{*} Eliséus clamábat, dicens~: Pater mi, pater mi, currus Israel, et auríga ejus. \textred{℣.}~Cumque pérgerent, et incedéntes sermocinaréntur~: ecce currus ígneus et equi ígnei divisérunt utrúmque, et ascéndit Elías per túrbinem in cœlum. \textred{E}liséus. Glória Patri. \textred{E}liséus.
\item Te Deum laudámus. 6.
\item \textred{Sacerd. ℣.}~Surréxit Elías prophéta quasi ignis. \textred{℟.}~Et verbum ipsíus quasi fácula ardébat.
\end{tightlist}

\subsubsection*{\center \textsc{AD LAUDES.}}
\begin{center}\textred{et per Horas, Aña.}\end{center}
\lettrine{E}lías, dum zelat zelum legis, recéptus est in cœlum.

\begin{tightlist}
\item \textred{Ps. de Dominica.} 27.
\item \textred{Aña.} Factum est autem, cum leváre vellet Dóminus Elíam per túrbinem in cœlum, ibant Elías et Eliséus de Gálgalis.
\item \textred{Aña.} Tulit Elías pállium suum, et percússit aquas, et divísæ sunt, et transiérunt ambo per siccum.
\item \textred{Aña.} Cum Elías et Eliséus pérgerent, et incedéntes sermocinaréntur, ecce currus ígneus et equi ígnei divisérunt utrúmque.
\item \textred{Aña.} Ascéndit Elías per túrbinem in cœlum, Eliséus autem vidébat, et clamábat~: Pater mi, pater mi, currus Israel, et auríga ejus.
\end{tightlist}

\subsubsection*{\textred{Capitulum. \hfill Eccli. 48.}}
\lettrineThree{S}urréxit Elías Prophéta quasi ignis, et verbum ipsíus quasi fácula ardébat. Qui indúxit in pópulos famem~: et irritántes illum invídia sua pauci facti sunt.

\subsubsection*{\center \textred{Hymnus.}}
\lettrine{P}ergámus sócii, tóllere cánticis\\
Thesbítæ, exímii Vatis, in ǽthera\\
Res gestas, hílares~: nam decet hoc die\\
Nostri nos mémores ducis.\\
\\
\textred{A}udívit Dómini verba recóndita\\
Auræ sub ténuis múrmure síbili,\\
Jussus sævítiam Jézabel ímprobæ,\\
Atque Achab scélera ínsequi.\\
\\
\textred{S}uccínctus túnica pellícea latus,\\
Carmelíque hábitans antra viréntia,\\
Prædícit mínitans trístia fúnera\\
Regi Ochozíæ ímpio.\\
\\
\textred{H}ujus bis précibus ignis ab ǽthere\\
Absórpsit véniens, flámmaque mílites~:\\
Findit dein próprio flúmina pállio,\\
Et sicco pede prǽterit.\\
\\
\textred{N}os nunquam auxílium destítuat, Pater,\\
Optatúsque tuus perpétuo favor~:\\
Nobis quin pótius múnera Spíritus\\
Augéri prece póstula.\\
\\
\textred{L}aus Patri ingénito, summáque glória,\\
Sit semper génito laus quoque Fílio,\\
Necnon par tibi sit, ínclyte Spíritus,\\
Simplex únaque Trínitas.\\
Amen.

\begin{tightlist}
\item \textred{℣.}~Elías in túrbine tectus est. \textred{℟.}~Et in Eliséo complétus est spíritus ejus.
\item \textred{Ad Benedictus, Aña.} Elías homo erat símilis nobis passíbilis, et oratióne orávit, ut non plúeret super terram; et non pluit annos tres et menses sex. Et rursum orávit~: et cœlum dedit plúviam, et terra dedit fructum suum.
\end{tightlist}

\subsubsection*{\center \textred{Oratio.}}
\lettrine{P}ræsta, quǽsumus omnípotens Deus, ut sicut beátum Elíam Prophétam tuum, et Patrem nostrum ante commúnem mortem curru ígneo aéreum elevásti ad cœlum~: ita nos fácias, eo interveniénte, dum vívimus, a terrénis semper ad cœléstia spíritu sublevári; et cum eo in resurrectióne justórum páriter gaudére. Per Dóminum.

\subsubsection*{\center \textsc{AD TERTIAM.}}
\subsubsection*{\textred{Capitulum. \hfill Eccli. 48.}}
\lettrine{S}urréxit Elías Prophéta quasi ignis, et verbum ipsíus quasi fácula ardébat. Qui indúxit in pópulos famem~: et irritántes illum invídia sua pauci facti sunt.

\begin{tightlist}{}
%\item \textred{Capit.} Surréxit Elías, \textred{ut supra ad Laudes.}
\item \textred{℟.}~Verbo Dómini \textred{*} Contínuit cœlum. \textred{V}erbo. \textred{℣.}~Et dejécit de cœlo ignem ter. \textred{C}ontínuit. Glória Patri. \textred{V}erbo.
\item \textred{℣.}~Tu ungis reges ad pœniténtiam. \textred{℟.}~Et Prophétas facis successóres post te.
\end{tightlist}

\subsubsection*{\center \textsc{AD SEXTAM.}}
\subsubsection*{\textred{Capitulum. \hfill Eccli. 48.}}
\lettrine{V}erbo Dómini contínuit cœlum, et dejécit de cœlo ignem ter~: sic amplificátus est Elías in mirabílibus suis.

\begin{tightlist}
\item \textred{℟.}~Tu ungis reges \textred{*} Ad pœniténtiam. \textred{T}u ungis. \textred{℣.}~Et Prophétas facis successóres post te. \textred{A}d. Glória Patri. \textred{T}u.
\item \textred{℣.}~Beáti sunt, qui te vidérunt. \textred{℟.}~Et in amicítia tua decoráti sunt.
\end{tightlist}

\subsubsection*{\center \textsc{AD NONAM.}}
\subsubsection*{\textred{Capitulum. \hfill Eccli. 48.}}
\lettrine{T}u sustulísti mórtuum ab ínferis de sorte mortis in verbo Dómini Dei. Qui scriptus es in judíciis témporum leníre iracúndiam Dómini.

\begin{tightlist}
\item \textred{℟.}~Beáti sunt, \textred{*} Qui te vidérunt. \textred{B}eáti. \textred{℣.}~Et in amicítia tua decoráti sunt. \textred{Q}ui. Glória Patri. \textred{B}eáti.
\item \textred{℣.}~Elías in túrbine tectus est. \textred{℟.}~Et in Eliséo complétus est spíritus ejus.
\end{tightlist}

\subsubsection*{\center \textsc{AD II. VESPERAS.}}
\begin{tightlist}
\item \textred{Nisi sequatur Festum duplex, dicitur prima Aña. Laudum, sub qua sola dicuntur Psalmi de Dominica; et loco ult.} Lauda, Jerúsalem. [56]
\end{tightlist}

\subsubsection*{\textred{Capitulum. \hfill Eccli. 48.}}
\lettrineThree{S}urréxit Elías Prophéta quasi ignis, et verbum ipsíus quasi fácula ardébat. Qui indúxit in pópulos famem~: et irritántes illum invídia sua pauci facti sunt.

\subsubsection*{\center \textred{Hymnus.}}
\lettrine{A}udiat miras óriens, cadénsque\\
Sol tuas laudes, mare, terra, et aer,\\
Ordinis nostri cólumen, tuǽque\\
Glória gentis.\\
\\
\textred{O} Deo multum (Pater alme) care,\\
O potens signis, meritísque felix,\\
Quem Deus gestis adhibére suévit\\
Grándibus olim.\\
\\
\textred{T}e triumpháli super alta curru\\
Igneis vectum céleres quadrígis\\
Angeli attóllunt, nímium corúsca\\
Luce micántem\\
\\
\textred{U}nico ætérni Patris adfuísti\\
Fílio testis, sócius fidélis\\
Mosis, in summo pósitus Thabóris\\
Vértice montis.\\
\\
\textred{N}os tui præsens clýpeo favóris,\\
Ac tua magna bonitáte fretos,\\
Súpplices dextra fámulos benígna\\
Prótege semper.\\
\\
\textred{S}it Patri summo, genitǽque Proli,\\
Et tibi compar utriúsque sancte\\
Spíritus, laus, impérium, potéstas\\
Tempus in omne.\\
Amen.

\begin{tightlist}
\item \textred{℣.}~Verbo Dómini contínuit cœlum. \textred{℟.}~Et dejécit de cœlo ignem ter.
\item \textred{Ad Magnificat, vel pro commemor. Aña.} Tulit Elías pállium suum, et invólvit illud, et percússit aquas Jordánis, quæ divísæ sunt in utrámque partem, et transiérunt ipse, et Eliséus per siccum~: qui incedéntes, ecce currus ígneus, et equi ígnei divisérunt utrúmque~: et ascéndit Elías per túrbinem in cœlum~: Eliséus autem non vidit eum ámplius.
%\item \textred{Temp. Pasch. tres Psalmi cujuslibet Nocturni dicuntur sub una sola Aña; et in fine omnium Añarum, et Respons. Versuum additur} allelúja, \textred{et ultimus versus Hymni ad Mat. mutatur.}
%\item \textred{Cetera omnia ut supra.}
\end{tightlist}

\begin{center}\textcolor{red}{\rule{\linewidth}{0.5pt}} \end{center}

\subsection*{\center \textsc{Officium S. Andreæ Corsini.}}
\subsubsection*{\center \textred{Ep. et Conf. Ordinis Nostri.\\ Semiduplex.}}
\begin{center}\textred{Recitandum in toto Ordine nostro bis in Mense per Annum, et semel pariter in Mense per Quadragesimam sub ritu Semiduplici, juxta Rubr. 17.}\end{center}
\fancyhead[C]{\textsc{Officia Mensualia - Officium S. Andreæ Corsini}}

\begin{tightlist}{}
\item \textred{Tempore Paschali tres Psalmi cujuslibet Nocturni dicuntur sub una sola Aña, et in fine omnium Añarum etc. additur} Allelúja. \textred{Cetera omnia ut infra.}
\end{tightlist}

\subsubsection*{\center \textsc{AD I. VESPERAS}}
\subsubsection*{\textred{Capitulum. \hfill Eccli. 50.}}
\lettrineThree{E}cce Sacérdos magnus, qui in vita sua suffúlsit domum~: et in diébus suis corroborávit templum. Qui præváluit amplificáre civitátem, qui adéptus est glóriam in conversatióne Gentis.

\subsubsection*{\center \textred{Hymnus.}}
\lettrine{A}ndréam grádibus sídera lúcidis\\
Scandéntem, sócii, et Prǽsule cérnitis\\
Stipátum gémino~: ferte citi lyras,\\
Plectrúmque, et nova cántica.\\
\\
\textred{M}agnis ex átavis éditus áureæ\\
Sortis dum séquitur blanda perícula;\\
Ibat forte puer, qua via flórido\\
Fallit nóxia trámite.\\
\\
\textred{I}nclámat génitrix~: ille pedem loco\\
Festínat rétrahens; áltaque per juga\\
Carméli evólitat, certus inánia\\
Mundi línquere gáudia.\\
\\
\textred{A}st illum cúpide nóbilis ínfula\\
Quærit~: sed látitans défugit ínfulam\\
Frustra~: nam púero vox data trímulo\\
Noléntem ad sólium rapit.\\
\\
\textred{E}rgo divítias fúndere páuperi\\
Et rectum sérere, et prava revéllere;\\
Cœlóque, ac pópulo, non sibi vívere;\\
Hæc laus Pontíficis fuit.\\
\\
\textred{P}astor summe, óvium qui regis ágmina,\\
Pastórum páriter Rector, et Arbiter\\
Christe, ad perpétuæ duc sata Pátriæ\\
Una cum dúcibus gregem.\\
Amen.

\begin{tightlist}{}
\item \textred{℣.}~Ora pro nobis sancte Andréa. \textred{℟.}~Ut digni efficiámur promissiónibus Christi.
\item \textred{Ad Magnificat, Aña.} Natus est firmaméntum gentis, rector fratrum, stabiliméntum pópuli.
\end{tightlist}

\subsubsection*{\center \textred{Oratio.}}
\lettrine{D}eus, qui in Ecclésia tua nova semper instáuras exémpla virtútum~: da pópulo tuo; beáti Andréæ Confessóris tui, atque Pontíficis ita sequi vestígia, ut assequátur et prǽmia. Per Dóminum.

\subsubsection*{\center \textsc{AD MATUTINUM.}}
\begin{tightlist}{}
\item \textred{Invitatorium.} Dóminum misericórdiæ, qui prædestinávit Andréam confórmem fíeri agno immaculáto~: \textred{*} Veníte adorémus.
\item \textred{Ps.} Veníte, exsultémus. 2.
\end{tightlist}

\subsubsection*{\center \textred{Hymnus.}}
\lettrine{V}irgo, Carméli decor et volúptas,\\
Monte subjéctos speculáta campos,\\
Ut lupum vidit; míserans corúsco\\
Térruit arcu.\\
\\
\textred{T}unc feræ (mirum!) cecidére vultus,\\
Ora tunc agni subiére mitis,\\
Tunc datum cursu per amœna ferri\\
Cúlmina montis.\\
\\
\textred{V}írginis cordi pius hæret ardor,\\
Vírginis formam sedet æmulári\\
Móribus puris~: alit illa sacros\\
Altius ignes.\\
\\
\textred{M}ox jubet certam pópulis salútem\\
Ferre commíssis~: mora nulla, late\\
Italas Præsul rádians per oras\\
Lucet, et urit.\\
\\
\textred{N}úminis plenus rétegit futúra,\\
Urbibus rixas fugat, et mináces\\
Artubus morbos; epulásque turbis\\
Dévocat axe.\\
\\
\textred{L}aus Patri, et Nato; tibi laus et alme\\
Spíritus semper; duce quo triúmphat\\
Victor Andréas, Mariána vivens\\
Signa secútus.\\
Amen.

\subsubsection*{\center \textsc{IN PRIMO NOCTURNO.}}
\begin{tightlist}{}
\item \textred{Aña.} Stérilis virtútem in conceptiónem accépit, et unigénitum offerébat.
\item \textred{Ps.} Beátus vir, \textred{cum reliquis de Communi unius Confess. Pontif.} [76]
\item \textred{Aña.} Oculus Dómini respéxit eum in bono, et eréxit eum ab humilitáte sua.
\item \textred{Aña.} Miráti sunt in illo multi, et laudavérunt Deum.
\item \textred{℣.}~Amávit eum Dóminus, et ornávit eum. \textred{℟.}~Stolam glóriæ índuit eum.
\item \textred{Lect. j. Noct. de Script. occurr. quæ si desint, legantur sequentes.}
\end{tightlist}

\subsubsection*{\center De libro Job.}
\subsubsection*{\center \textred{Lectio j. \hfill Cap. 29. ℣.~7.}}
\lettrineThree{Q}uando procedébam ad portam civitátis, et in platéa parábant cáthedram mihi~: auris aúdiens beatificábat me, et óculus videns testimónium reddébat mihi. Eo quod liberássem páuperem vociferántem, et pupíllum, cui non esset adjútor. Benedíctio peritúri super me veniébat, et cor víduæ consolátus sum. Justítia indútus sum~: et vestívi me, sicut vestiménto, et diadémate, judício meo. Oculus fui cæco, et pes claudo. Pater eram páuperum, et causam, quam nesciébam, diligentíssime investigábam. Dicebámque~: in nídulo meo móriar et sicut palma multiplicábo dies. Radix mea apérta est secus aquas, et ros morábitur in messióne mea. Glória mea semper innovábitur.

\begin{tightlist}{}
\item \textred{℟.}~Concépit stérilis et péperit fílium quia magnificábit Dóminus misericórdiam suam cum illa. \textred{*} Et congratulabántur ei. \textred{℣.}~Cœpítque Spíritus Dómini esse cum eo. \textred{E}t.
\end{tightlist}

\subsubsection*{\center\textred{Lectio ij. \hfill Cap. 31. ℣.~1.}}
\lettrine{P}epigi fœdus cum óculis meis, ut ne cogitárem quidem de vírgine. Quam enim partem habéret in me Deus désuper, et hereditátem Omnípotens de excélsis? Numquid non perdítio est iníquo, et alienátio operántibus injustítiam? Si ambulávi in vanitáte, et festinávit in dolo pes meus~: Si declinávit gressus meus de via, et si secútum est óculos meos cor meum, et si mánibus meis adhǽsit mácula~: seram, et álius cómedat. Si decéptum est cor meum super mulíere, et si ad óstium amíci mei insidiátus sum. Hoc enim nefas est, et iníquitas máxima. Ignis est usque ad perditiónem dévorans, et ómnia eradícans genímina. Semper quasi tuméntes super me fluctus tímui Deum, et pondus ejus ferre non pótui. 

\begin{tightlist}{}
\item \textred{℟.}~Humiliátus est póstea, eo quod exaltátum fuísset cor ejus. \textred{*} Et opéruit cilício carnem suam et jejunávit. \textred{℣.}~Ascéndit ut immoláret Dómino hóstiam solémnem et votum suum. \textred{E}t.
\end{tightlist}

\subsubsection*{\center \textred{Lectio iij.}}
\lettrine{S}i putávi aurum robur meum, et obrízo dixi fidúcia mea. Si gavísus sum ad ruínam ejus, qui me óderat~: et exsultávi, quod invenísset eum malum. Non enim dedi ad peccándum guttur meum, ut expéterem maledícens ánimam ejus. Si abscóndi quasi homo peccátum meum, et celávi in sinu meo iniquitátem meam. Si expávi ad multitúdinem nímiam, et despéctio propinquórum térruit me~: et non magis tácui, nec egréssus sum óstium. Si advérsum me terra mea clamat, et cum ipsa sulci ejus deflent. Si fructus ejus comédi absque pecúnia, et ánimam agricolárum ejus afflíxi~: pro fruménto oriátur mihi tríbulus, et pro hórdeo spina. 

\begin{tightlist}{}
\item \textred{℟.}~Invénit hóminem, qui erat cæcus, et ait~: sanat te Dóminus Jesus Christus~: \textred{*} Et conféstim vidit. \textred{℣.}~Est aqua lustratiónis aspérsus. \textred{E}t. Glória Patri. \textred{E}t.
\end{tightlist}

\subsubsection*{\center \textsc{IN SECUNDO NOCTURNO.}}
\begin{tightlist}{}
\item \textred{Aña.} Fecit quod plácuit Deo, et fórtiter ivit in via, quam mandávit Prophéta magnus in conspéctu Dei.
%\item \textred{Ps.} Cum invocárem. [67]
\item \textred{Aña.} Elégit Dóminus servum suum páscere hereditátem suam, et in intelléctibus mánuum suárum dedúxit eos.
\item \textred{Aña.} Dedit eis petitiónem ipsórum, misit saturitátem in ánimas eórum.
\item \textred{℣.}~Justum dedúxit Dóminus per vias rectas. \textred{℟.}~Et osténdit illi regnum Dei.
\end{tightlist}

\subsubsection*{\center Sermo sancti Gregórii Nysséni.}
\subsubsection*{\center \textred{In fun. M. Milet.\\
Lectio iv.}}
\lettrineThree{C}um primum bene moráta, ac modésta Ecclésia beátum Andréam vidit, vidit fáciem ad imáginem Dei vere formátam~: vidit dilectiónem fontis modo scaturiéntem~: vidit grátiam lábiis circumfúsam, ánimi demissiónis summum gradum~: post quem ámplius quidquam cogitári non potest~: vidit mansuetúdinem atque cleméntiam, qualis in Davíde fuit~: qualis in Salomóne, intelligéntiam atque prudéntiam~: qualis in Móyse, bonitátem~: qualis in Samuéle, perfectiónem~: qualis in Josépho, continéntiam pudicitiámque~: qualis in Daniéle, sapiéntiam~: quemádmodum magnus Elías, zelo fídei prǽditum~: sícuti sublímis Joánnes, integritáte córporis ornátum~: sícuti Paulus, inexsuperábili prǽditum dilectióne. Vidit tot bonórum circa unam ánimam concúrsum~: amóre beáto vulneráta est, casto bonóque Sponsum suum amóre atque benevoléntia prosecúta diléxit.

\begin{tightlist}{}
\item \textred{℟.}~In médio fratrum rector illórum loquebátur de regno Dei, et eos qui cura indigébant, sanábat~: \textred{*} Omnibus mirántibus. \textred{℣.}~Ecce quidam hydrópicus erat ante illum; ipse vero sanávit eum. \textred{O}mnibus.
\end{tightlist}

\subsubsection*{\center \textred{Lectio v.}}
\lettrine{S}ed priúsquam cupiditátem expléret, ántequam desidérium recreáret atque sedáret, amóris vi fervens, tentatiónibus athlétam ad certámina vocántibus, sola relícta est. Atque ille quidem in certamínibus pro pietáte suscéptis desudábat~: hæc vero durábat in castitáte matrimónium consérvans. Non ablátus est a nobis Sponsus, in médio nostrum stat, étiam si nos non videámus. In ádytis, ac penetrálibus Sacérdos est. In interióribus veli, quo præcúrsor pro nobis ingréssus est Christus, relíquit carnis teguméntum. Non ámplius signo, et umbræ cœléstium servit; sed in ipsam rerum imáginem intuétur~: non ámplius per spéculum, atque per transénnam, et ænígma; sed ipse fácie cum fácie colláta intercédit apud Deum. 

\begin{tightlist}{}
\item \textred{℟.}~Non seípsum clarificávit, ut Póntifex fíeret \textred{*} sed vocátur a Deo tamquam Aáron. \textred{℣.}~Quasi sol refúlgens, sic ille effúlsit in templo Dei. \textred{S}ed.
\end{tightlist}

\subsubsection*{\center \textred{Lectio vj.}}
\lettrine{I}ntercédit autem pro nobis, et pópuli errátis~: depósuit túnicas pellíceas, neque enim tálibus túnicis opus habent, qui in Paradíso degunt~: sed habet induménta, quæ puritáte vitæ suæ contéxuit, iísque sese ornávit. Honoráta, ac pretiósa coram Dómino talis viri mors est; immo vero non mors, sed ruptúra membrórum est. Dirupísti enim, inquit, víncula mea. Dimíssus est Símeon, liberátus est a vínculis córporis. Láqueus contrítus est, et avícula avolávit. Ingréssus est terram promissiónis~: in monte cum Deo philosophátur. Solvit ánimæ calceaméntum, ut pura planta mentis terram sanctam, ubi conspícitur Deus, conscénderet.

\begin{tightlist}{}
\item \textred{℟.}~Cum turba multa esset, nec habérent quod manducárent, accípiens septem panes fregit, et dabat~: \textred{*} Et manducavérunt, et saturáti sunt. \textred{℣.}~Non sunt fraudáti a desidério suo. \textred{E}t. Glória Patri. \textred{E}t.
\end{tightlist}

\subsubsection*{\center \textsc{IN TERTIO NOCTURNO.}}
\begin{tightlist}{}
\item \textred{Aña.} Minúitur crescens mirabíliter in consummatióne, respléndens glorióse.
%\item \textred{Ps.} Dómine quis. [70]
\item \textred{Aña.} Dóminus dedit illi sedem glóriæ, et propter illum dejécit omnem poténtiam inimicórum.
\item \textred{Aña.} Certe in géntibus multis non erat símilis illi diléctus Deo suo.
\item \textred{℣.}~Justus ut palma florébit in domo Dómini. \textred{℟.}~Sicut cedrus Líbani multiplicábitur.
\end{tightlist}

\subsubsection*{\center Léctio sancti Evangélii secúndum Matthǽum.}
\subsubsection*{\center \textred{Lectio vij.} \hfill \textred{Cap. 25.}}
\lettrine{I}n illo témpore~: Dixit Jesus discípulis suis~: Homo quidam péregre proficíscens, vocávit servos suos, et trádidit illis bona sua. Et réliqua.

\begin{center}Homilía sancti Joánnis Chrysóstomi.\\
\textred{Hom. 79. ad med.}\end{center}
\lettrineThree{T}alénta hic pro eo, quod unusquísque fácere potest accípimus, sive auctoritáte protégere, sive pecúniis juváre, sive doctrína admonére, sive ália quápiam re próximis prodésse queas. Nemo secum dicat~: Cum unum taléntum hábeam, nihil possum effícere~: potest profécto ex ea una re sola comprobári. Non es paupérior illa vídua; non es Petro, atque Joánne rustícior~: qui, quamvis rudes simul atque illiteráti fuérunt, quóniam magno stúdio commúnem utilitátem compléxi fúerant, cœlórum Príncipes facti sunt.

\begin{tightlist}{}
\item \textred{℟.}~In illa hora multos curávit a languóribus et plagis~: \textred{*} Quia virtus de illo exíbat et sanábat. \textred{℣.}~In vita sua fecit monstra, et in morte mirabília operátus est. \textred{Q}uia.
\end{tightlist}

\subsubsection*{\center \textred{Lectio viij.}}
\lettrine{Q}uippe nulla res Deo grátior est, quam ut univérsam vitam ad commúne cómmodum cónferas. Idcírco ratióne, atque oratióne nos Deus decorávit, mentem, atque ingénium concéssit, manus, pedes, vires córporis dedit, ut his ómnibus, et nos ipsos, et próximos tutémur. Non enim ad agéndas grátias solúmmodo confert orátio, verum étiam ad docéndum, et ad monéndum perútilis est~: qua in re, si ea diligénter utémur, Dóminum; sin vero in contráriis, diábolum imitábimur.

\begin{tightlist}{}
\item \textred{℟.}~Repúlsi sunt inimíci præ timóre ejus~: \textred{*} Et dirécta est salus in manu ejus. \textred{℣.}~Cógnitus est in verbis suis fidélis. \textred{E}t.
\end{tightlist}

\subsubsection*{\center \textred{Lectio ix.}}
\lettrine{N}am et Petrus quidem, quando Christum confitebátur, quasi ea dicens, quæ Patris sunt, beatificabátur~: quando autem crucem horrébat, atque repellébat, quasi ea, quæ diáboli sunt sápiens, ácriter increpabátur. Si ergo ignoratióne crucem recusándo reprehénsus est, quam nos sponte peccántes véniam impetrábimus? Ea ergo loquámur, quæ Christi verba esse, conféstim cognoscúntur.

\begin{tightlist}{}
\item \textred{℟.}~Spíritu magno vidit última, consolátus est lugéntes, osténdit futúra, et abscóndita. \textred{*} Quæsívit bona genti suæ, et fecit pacem. \textred{℣.}~Adjúvit páuperem de inópia, et pósuit sicut oves famílias. \textred{Q}uæsívit. Glória Patri. \textred{Q}uæsívit.
\item Te Deum laudámus. 6.
\item \textred{Sacerd. ℣.}~Elégit dum Dóminus sacerdótem sibi. \textred{℟.}~Ad sacrificándum ei hóstiam laudis.
\end{tightlist}

\subsubsection*{\center \textsc{AD LAUDES.}}
\begin{center} \textred{et per Horas, Aña.}\end{center}
\lettrine{D}iléxit Andréam Dóminus in odórem suavitátis.

\begin{tightlist}{}
\item \textred{Ps. de Dominica.} 27.
\item \textred{Aña.} Jucunditátem et exsultatiónem thesaurizávit super illum, et nómine ætérno hereditávit illum.
\item \textred{Aña.} In médio pópuli sui exaltátus est, et in multitúdine sanctórum habet laudem.
\item \textred{Aña.} Perféctus invéntus est, et erit illi glória ætérna.
\item \textred{Aña.} Implétus est benedictiónibus, et vidéntes illum laudábunt.
\end{tightlist}

\subsubsection*{\textred{Capitulum. \hfill Eccli. 50.}}
\lettrineThree{E}cce Sacérdos magnus, qui in vita sua suffúlsit domum~: et in diébus suis corroborávit templum. Qui præváluit amplificáre civitátem, qui adéptus est glóriam in conversatióne Gentis.

\subsubsection*{\center \textred{Hymnus.}}
\begin{center}\textred{In quo Temp. Pasch. nihil additur, neque mutatur.}\end{center}
\lettrine{Q}ua nocte mundi Cónditor\\
Tenéllus infans náscitur,\\
María amánti Prǽsuli\\
Mortem propínquam núntiat.\\
\\
\textred{Q}ua stella nocte ad Párvuli\\
Tres ducit antra Príncipes,\\
María, quæ stella est maris,\\
Ad astra ducit Prǽsulem.\\
\\
\textred{A}rtus relícti léctulum,\\
Seu pulchra lília, aut rosæ,\\
Odóre perflant; ac salus\\
Odor fit unus lánguidis.\\
\\
\textred{Q}uin luce mira prǽnitent,\\
Qua pulsa luce nox fugit~:\\
Hæc nempe sunt vestígia\\
Nitentióris spíritus.\\
\\
\textred{O} qui per alta, spíritus,\\
Spátia vagáris ǽtheris,\\
Torus Deo jam fúlgurans,\\
Ne sperne res mortálium.\\
\\
\textred{T}e fons, Redémptor, lúminis,\\
Te campi odórum lílium,\\
Trahe nos odóre, póscimus,\\
Ad lumen almum glóriæ.\\
Amen.

\begin{tightlist}{}
\item \textred{℣.}~Justus germinábit sicut lílium. \textred{℟.}~Et florébit in ætérnum ante Dóminum.
\item \textred{Ad Benedictus, Aña.} Fecit quod rectum erat in óculis Dómini, et diréxit eum in ómnibus, et adjúvit eum Deus, et divulgátum est nomen ejus.
\end{tightlist}

\subsubsection*{\center \textred{Oratio.}}
\lettrine{D}eus, qui in Ecclésia tua nova semper instáuras exémpla virtútum, da pópulo tuo beáti Andréæ Confessóris tui atque Pontíficis ita sequi vestígia, ut assequátur et prǽmia. Per Dóminum.

\begin{tightlist}{}
\item \textred{Ad Horas Antiphonæ de Laudibus, Psalmi de Dominica, ad Primam tamen ut in Festis. Capit, et ℟℟. de Communi.} [87]
\item \textred{Ad ij. Vesp. Aña. prima Laudum, Psalmi de Dominica, et loco ultimi, Ps. 131.} Meménto Dómine David. [88]
%\item \textred{Capit. Hymn. et ℣. ut in j. Vesperis.}
\end{tightlist}

\subsubsection*{\textred{Capitulum. \hfill Eccli. 50.}}
\lettrineThree{E}cce Sacérdos magnus, qui in vita sua suffúlsit domum~: et in diébus suis corroborávit templum. Qui præváluit amplificáre civitátem, qui adéptus est glóriam in conversatióne Gentis.

\subsubsection*{\center \textred{Hymnus.}}
\lettrine{A}ndréam grádibus sídera lúcidis\\
Scandéntem, sócii, et Prǽsule cérnitis\\
Stipátum gémino~: ferte citi lyras,\\
Plectrúmque, et nova cántica.\\
\\
\textred{M}agnis ex átavis éditus áureæ\\
Sortis dum séquitur blanda perícula;\\
Ibat forte puer, qua via flórido\\
Fallit nóxia trámite.\\
\\
\textred{I}nclámat génitrix~: ille pedem loco\\
Festínat rétrahens; áltaque per juga\\
Carméli evólitat, certus inánia\\
Mundi línquere gáudia.\\
\\
\textred{A}st illum cúpide nóbilis ínfula\\
Quærit~: sed látitans défugit ínfulam\\
Frustra~: nam púero vox data trímulo\\
Noléntem ad sólium rapit.\\
\\
\textred{E}rgo divítias fúndere páuperi\\
Et rectum sérere, et prava revéllere;\\
Cœlóque, ac pópulo, non sibi vívere;\\
Hæc laus Pontíficis fuit.\\
\\
\textred{P}astor summe, óvium qui regis ágmina,\\
Pastórum páriter Rector, et Arbiter\\
Christe, ad perpétuæ duc sata Pátriæ\\
Una cum dúcibus gregem.\\
Amen.

\begin{tightlist}{}
\item \textred{℣.}~Ora pro nobis sancte Andréa. \textred{℟.}~Ut digni efficiámur promissiónibus Christi.
\item \textred{Ad Magnificat, Aña.} Hic est Fratrum amátor, et pópuli~: hic est qui multum orat pro pópulo, et univérsa civitáte.
%\item \textred{Tempore Paschali tres Psalmi cujuslibet Nocturni dicuntur sub una sola Aña, et in fine omnium Añarum etc. additur} Allelúja. \textred{Cetera omnia ut supra.}
\end{tightlist}

\begin{center}\textcolor{red}{\rule{\linewidth}{0.5pt}} \end{center}

\section*{\center \textsc{Officium Sanctæ Teresiæ.}}
\subsubsection*{\center \textred{Virginis Ordinis Nostri.\\ Semiduplex.}}
\begin{center}\textred{Semel in mense sub Ritu semid. juxta Rubr. 17.}\end{center}
\fancyhead[C]{\textsc{Officia Mensualia - Officium S. Teresiæ}}

%\begin{center}\textcolor{red}{\rule{\linewidth}{0.5pt}} \end{center}

\subsubsection*{\center \textsc{AD I. VESPERAS.}}
\subsubsection*{\textred{Capitulum. \hfill Sap. 7.}}
\lettrineThree{O}ptávi, et datus est mihi sensus; et invocávi, et venit in me spíritus sapiéntiæ; et præpósui illam regnis, et sédibus, et divítias nihil esse duxi in comparatióne illíus.

\subsubsection*{\center \textred{Hymnus.}}
\begin{center}\textred{In quo Temp. Pasch. nihil additur neque mutatur.}\end{center}
\lettrine{R}egis supérni núntia\\
Domum patérnam déseris,\\
Terris Terésa bárbaris\\
Christum datúra, aut sánguinem.\\
\\
\textred{S}ed te manet suávior\\
Mors, pœna poscit dúlcior~:\\
Divíni amóris cúspide\\
In vulnus icta cóncides.\\
\\
\textred{O} caritátis víctima!\\
Tu corda nostra cóncrema,\\
Tibíque Gentes créditas\\
Avérni ab igne líbera.\\
\\
\textred{S}it laus Patri cum Fílio,\\
Et Spíritu Paráclito,\\
Tibíque sancta Trínitas\\
Nunc, et per omne sǽculum.\\
Amen.

\begin{tightlist}{}
\item \textred{℣.}~Oblivíscere pópulum tuum et domum patris tui. \textred{℟.}~Et concupíscet Rex decórem tuum.
\item \textred{Ad Magnificat, Aña.} Quæsívi in sponsam mihi eam assúmere. Doctrix enim est disciplínæ Dei, et eléctrix óperum illíus.
\end{tightlist}

\subsubsection*{\center \textred{Oratio.}}
\lettrine{E}xáudi nos Deus salutáris noster, ut sicut de beátæ Terésiæ Vírgnis tuæ commemoratióne gaudémus~: ita cœléstis ejus doctrínæ pábulo nutriámur, et piæ devotiónis erudiámur afféctu. Per Dóminum.

\subsubsection*{\center \textsc{AD MATUTINUM.}}
\begin{tightlist}{}
\item \textred{Invitatorium.} Regem Vírginum Dóminum, \textred{*} Veníte adorémus.
\item \textred{Ps.} Veníte, exsultémus. 2.
\item \textred{Hymnus ut in j. Vesperis.}
\end{tightlist}

\subsubsection*{\center \textsc{IN PRIMO NOCTURNO.}}
\begin{tightlist}{}
\item \textred{Aña.} Quæ est ista, quæ ascéndit per desértum sicut vírgula fumi ex aromátibus myrrhæ et thuris?
\item \textred{Ps.} Dómine, Dóminus noster, \textred{cum reliquis de Communi B.M.V.} [172]
\item \textred{Aña.} En diléctus meus lóquitur mihi~: Surge, própera amíca mea, colúmba mea, et veni.
\item \textred{Aña.} Revértere, revértere Sunamítis~: revértere, revértere, ut intueámur te.
\item \textred{℣.}~Fiat cor meum immaculátum in justificatiónibus tuis. \textred{℟.}~Ut non confúndar.
\item \textred{Lect. de Script. occurr. sed si desunt, legantur sequentes.}
\end{tightlist}

\subsubsection*{\center De libro Sapiéntiæ.}
\subsubsection*{\center \textred{Lectio j.} \hfill \textred{Cap. 7.}}
\lettrineThree{O}ptávi, et datus est mihi sensus~: et invocávi, et venit in me spíritus sapiéntiæ~: et præpósui illam regnis, et sédibus, et divítias nihil esse duxi in comparatióne illíus. Nec comparávi illi lápidem pretiósum~: quóniam omne aurum in comparatióne illíus aréna est exígua, et tamquam lutum æstimábitur argéntum in conspéctu illíus. Super salútem et spéciem diléxi illam, et propósui pro luce habére illam~: quóniam inexstinguíbile est lumen illíus. Venérunt autem mihi ómnia bona páriter cum illa, et innumerábilis honéstas per manus illíus, et lætáta sum in ómnibus~: quóniam antecedébat me ista sapiéntia, et ignorábam, quóniam horum ómnium mater est. Quam sine fictióne dídici, et sine invídia commúnico, et honestátem illíus non abscóndo. Infinítus enim thesáurus est homínibus~: quo qui usi sunt, partícipes facti sunt amicítiæ Dei, propter disciplínæ dona commendáti.

\begin{tightlist}{}
\item \textred{℟.}~Habitábit in solitúdine judícium, et justítia in Carmélo sedébit, et erit opus justítiæ pax, \textred{*} Et cultus justítiæ siléntium, et secúritas usque in sempitérnum. \textred{℣.}~Et sedébit pópulus meus in pulchritúdine pacis, et in tabernáculis fidúciæ. \textred{E}t.
\end{tightlist}

\subsubsection*{\center \textred{Lectio ij.}}
\lettrine{O}mnium enim ártifex dócuit me sapiéntia. Est enim in illa spíritus intelligéntiæ, sanctus, únicus, múltiplex, subtílis, disértus, móbilis, incoinquinátus, certus, suávis, amans bonum, acútus, quem nihil vetat, benefáciens, humánus, benígnus, stábilis, certus, secúrus, omnem habens virtútem, ómnia prospíciens, et qui cápiat omnes spíritus~: intelligíbilis, mundus, subtílis. Omnibus enim mobílibus mobílior est sapiéntia~: attíngit autem ubíque propter suam mundítiam. Vapor est enim virtútis Dei, et emanátio quædam est claritátis omnipoténtis Dei sincéra~: et ídeo nihil inquinátum in eam incúrrit. Candor est enim lucis ætérnæ, et spéculum sine mácula Dei majestátis, et imágo bonitátis illíus. Et cum sit una, ómnia potest, et in se pérmanens ómnia ínnovat, et per natiónes in ánimas sanctas se transfert, amícos Dei, et Prophétas constítuit.

\begin{tightlist}{}
\item \textred{℟.}~Diffúsa est grátia in lábiis tuis~: \textred{*} Proptérea benedíxit te Deus in ætérnum. \textred{℣.}~Spécia tua, et pulchritúdine tua inténde, próspere procéde, et regna. \textred{P}roptérea.
\end{tightlist}

\subsubsection*{\center \textred{Lectio iij.}}
\lettrine{N}eminem enim díligit Deus, nisi eum, qui cum sapiéntia inhábitat. Est enim hæc speciósior sole, et super omnem dispositiónem stellárum, luci comparáta invenítur prior. Illi enim succédit nox, sapiéntiam autem non vincit malítia. Attíngit ergo a fine usque ad finem fórtiter, et dispónit ómnia suáviter. Hanc amávi et exquisívi a juventúte mea, et quæsívi sponsam mihi eam assúmere, et amátrix facta sum formæ illíus. Generositátem illíus gloríficat, contubérnium habens Dei~: sed et ómnium Dóminus diléxit illam. Doctrix enim est disciplínæ Dei, et eléctrix óperum illíus. Et si divítiæ appetúntur in vita, quid sapiéntia locuplétius, quæ operátur ómnia? Si autem sensus operátur~: quis horum, quæ sunt, magis, quam illa est ártifex? Et si justítiam quis díligit~: labóres hujus magnas habent virtútes~: sobrietátem enim, et prudéntiam docet, et justítiam, et virtútem, quibus utílius nihil est in vita homínibus.

\begin{tightlist}{}
\item \textred{℟.}~Audi fília, et vide, et inclína aurem tuam; \textred{*} Quia deínceps ut vera sponsa meum zelábis honórem. \textred{℣.}~Sponsábo te mihi in justítia, et judício, et in misericórdia, et miseratiónibus. \textred{Q}uia. Glória Patri. \textred{Q}uia.
\end{tightlist}

\subsubsection*{\center \textsc{IN SECUNDO NOCTURNO.}}
\begin{tightlist}{}
\item \textred{Aña.} Vidi Dóminum, qui dedit corónam decóris in cápite meo.
%\item \textred{Ps.} Eructávit. [169]
\item \textred{Aña.} Induit me Génitrix Dómini vestiméntis salútis, et datum est mihi, ut coopériar býssino splendénti, et cándido.
\item \textred{Aña.} Ego autem in Dómino gaudébo, et exsultábo in Deo Jesu meo.
\item \textred{℣.}~Sitívit in te ánima mea. \textred{℟.}~Quam multiplíciter tibi caro mea.
\end{tightlist}

\subsubsection*{\center Sermo sancti Máximi Epíscopi.}
\subsubsection*{\center \textred{Hom. in nat. S. Agnetis.\\
Lectio iv.}}
\lettrineThree{C}um in toto mundo virgíneus flos Maríæ immarcescíbiles corónas innéctat, et sceptrígeram pudóris aulam immaculáto consérvet afféctu, eoúsque intégritas perseverávit ad palmam, ut in puéllis trophǽum sanctitátis arríperet, et per vestígia matris Vírginis ad cœléstem thálamum perveníret. O quam imitábile virgínibus, Virgo, tui amóris tradidísti exémplum! O quam sanctum pectóribus dandum virgíneis tribuísti respónsum, contemnéndo mundi divítias, exsecrándo carnis concupiscéntias, et solam Christi pulchritúdinem adamándo! Accédite ad puéllam puéllæ, et díscite, qualis in péctore suo circa amórem Christi flammas accéndit. Díscite, vírgines, Christi amórem in ea fervéntem, et omnes mundi divítias velut stércora recusántem.

\begin{tightlist}{}
\item \textred{℟.}~Dixit mihi Dóminus meus~: Ne tímeas fília~: \textred{*} Non rápiet te quisquam de manu mea. \textred{℣.}~Neque mors, neque vita, neque creatúra ália póterit me separáre a caritáte Dei. \textred{N}on.
\end{tightlist}

\subsubsection*{\center \textred{Lectio v.}}
\lettrine{I}deo enim sanctórum per lítteras tradúntur exémpla, ut unusquísque pro qualitáte sexus modum sui propósiti téneat, et sicut per vírginem vírgines, ita per cónjugem cónjuges víncere sǽculum pro Christi amóre stúdeant, ac repugnatióne ánimi, mundi dúlcia recusántes, et ejus non recusántes austéra, ad delectatiónes perpétuas, et ad gáudia ætérna pertíngant. Hæc dicéntes, non vituperámus cónjugum casta commércia, sed vírginum perseverántiam adornámus. Non enim virgínitas malis rebus superponénda est, sed est rebus óptimis præferénda~: quæ enim mundis múndior, et sanctis est sánctior, non índiget malárum rerum vituperatióne laudári.

\begin{tightlist}{}
\item \textred{℟.}~Gaudens gaudébo in Dómino, et exsultábit ánima mea in Deo meo, quia índuit me vestiméntis salútis; \textred{*} Et induménto justítiæ circúmdedit me. \textred{℣.}~Quasi sponsam decorátam coróna, et quasi sponsam ornátam monílibus suis. \textred{E}t.
\end{tightlist}

\subsubsection*{\center \textred{Lectio vj.}}
\lettrine{V}eni jam, Virgo, ad thálamum, quem quæsísti, ut sponsam te virgínitas immaculáta cómprobet Christi. Felix catérva vírginum, quibus stúdium est tuæ sequi confessiónis vestígia. Habébunt enim tecum corónam in cœlo, quæ tecum superavérunt inimícum in mundo. Nam sicut tibi vestígia Maríæ prosequénti participátio ætérnæ remuneratiónis appáruit~: ita eis te imitántibus ætérna crédimus gáudia non negánda. Itaque o spléndida Christo Dei Fílio, et ómnibus Angelis, atque Archángelis grata! ut nostri meminísse dignéris, quibus póssumus précibus exorámus, ut nobis ille tríbuat indulgéntiam, qui tibi tuórum ómnium labórum trádidit palmam. 

\begin{tightlist}{}
\item \textred{℟.}~Pulchra fácie, sed púlchrior fide beáta es Terésia virgo, réspuens mundum, lætáberis cum Angelis~: \textred{*} Intercéde pro ómnibus nobis. \textred{℣.}~Gaudeámus, et exsultémus, et demus glóriam ei, quia venérunt núptiæ Agni. \textred{I}ntercéde. Glória Patri. \textred{I}ntercéde.
\end{tightlist}

\subsubsection*{\center \textsc{IN TERTIO NOCTURNO.}}
\begin{tightlist}{}
\item \textred{Aña.} Diléctus meus mihi, et ego illi, qui páscitur inter lília.
%\item \textred{Ps.} Cantáte. [172]
\item \textred{Aña.} Dómine prævenísti eam in benedictiónibus dulcédinis~: posuísti in cápite ejus corónam de lápide pretióso.
\item \textred{Aña.} Os suum apéruit sapiéntiæ, et lex cleméntiæ in lingua ejus. 
\item \textred{℣.}~Deus præcínxit me virtúte. \textred{℟.}~Et pósuit immaculátam viam meam.
\end{tightlist}

\subsubsection*{\center Lectio sancti Evangélii secúndum Matthǽum.}
\subsubsection*{\center \textred{Lectio vij. \hfill Cap. 11.}}
\lettrine{I}n illo témpore~: Respóndens Jesus dixit~: Confíteor tibi Pater, Dómine cœli et terræ, quia abscondísti hæc a sapiéntibus, et prudéntibus, et revelásti ea párvulis. Et réliqua.

\begin{center}Homilía sancti Augustíni Epíscopi.\\
\textred{Serm. 9. de Verbis Dom.}\end{center}
\lettrineThree{N}on immérito ille vas electiónis cum ingénti lætítia dixit~: Non sunt condígnæ passiónes hujus témporis ad futúram glóriam, quæ revelábitur in nobis. Ecce unde illud jugum suáve est, et sárcina levis, et si angústa est paucis eligéntibus, fácilis tamen ómnibus diligéntibus. Dixit Psalmísta~: Propter verba labiórum tuórum ego custodívi vias duras. Sed quæ dura sunt laborántibus, éadem ipsa mitéscunt amántibus.

\begin{tightlist}{}
\item \textred{℟.}~Habitáre fecit eam Deus stérilem in domo~: \textred{*} Matrem filiórum lætántem. \textred{℣.}~Cibávit eam pane vitæ, et intelléctus, et aqua sapiéntiæ salutáris potávit illam. \textred{M}atrem.
\end{tightlist}

\subsubsection*{\center \textred{Lectio viij.}}
\lettrine{P}ropter quod ita divínæ pietátis dispensatióne actum est, ut intérior homo, qui renovátur de die in diem, non adhuc sub lege pósitus, sed jam sub grátia exonerátus sárcinis innumerabílium observatiónum, quod erat revéra grave jugum, sed duræ cervíci conveniénter impósitum, facilitáte símplicis fídei, et bonæ spei, et sanctæ caritátis, quidquid molestiárum exterióri hómini forínsecus intulísset ille princeps, qui missus est foras, interióri gáudio fíeret leve. Nihil enim tam fácile bonæ voluntáti est, quam ipsa sibi~: et hæc súfficit Deo.

\begin{tightlist}{}
\item \textred{℟.}~Tu glória Jerúsalem, tu lætítia Israel, tu honorificéntia pópuli nostri, quia fecísti viríliter, et confortátum est cor tuum~: ídeo et manus Dómini confortábit te~: \textred{*} Et eris benedícta in ætérnum. \textred{℣.}~Surrexérunt fílii tui, et. beatíssimam prædicavérunt. \textred{E}t.
\end{tightlist}

\subsubsection*{\center \textred{Lectio ix.}}
\lettrine{Q}uantúmlibet ergo sǽviat iste mundus, veríssime Angeli, nato in carne Dómino clamavérunt~: Glória in excélsis Deo, et in terra pax homínibus bonæ voluntátis; quia, qui natus erat, suáve jugum est, et sárcina levis. Et sicut dicit Apóstolus~: Fidélis Deus, qui nos non sinit tentári supra id, quod póssumus ferre, sed facit cum tentatióne étiam éxitum, ut possímus sustinére.

\begin{tightlist}{}
\item \textred{℟.}~Super salútem et spéciem diléxi sapiéntiam, et propósui pro luce habére illam~: \textred{*} Quóniam inexstinguíbile est lumen illíus. \textred{℣.}~Lætáta sum in ómnibus, quóniam antecedébat me ista sapiéntia. \textred{Q}uóniam. Glória Patri. \textred{Q}uóniam.
\item Te Deum laudámus. 6.
\item \textred{Sacerdotalis ℣.}~Veni, Sponsa Christi, áccipe corónam. \textred{℟.}~Quam tibi Dóminus præparávit in ætérnum.
\end{tightlist}

\subsubsection*{\center \textsc{AD LAUDES.}}
\begin{center}\textred{et per Horas, Aña.}\end{center}
\lettrine{H}æc est virgo sápiens, et una de número prudéntum. 

\begin{tightlist}{}
\item \textred{Ps. de Dominica.} 27.
\item \textred{Aña.} Date ei de fructu mánuum suárum, et laudent eam in portis ópera ejus.
\item \textred{Aña.} Sitívit in te Dómine ánima sponsæ tuæ, et lætáta est in his, quæ dicta sunt illi.
\item \textred{Aña.} Trahe me post te, in odórem currémus unguentórum tuórum~: óleum effúsum nomen tuum.
\item \textred{Aña.} Lætáre Terésia in eo, qui te fecit~: quóniam benedíxit tibi Dóminus, et ádipe fruménti sátiat te.
\end{tightlist}

\subsubsection*{\center \textred{Capitulum. \hfill Sap. 7.}}
\lettrineThree{O}ptávi, et datus est mihi sensus; et invocávi, et venit in me spíritus sapiéntiæ, et præpósui illam regnis et sédibus, et divítias nihil esse duxi in comparatióne illíus.

\subsubsection*{\center \textred{Hymnus.}}
\begin{center}\textred{In quo Temp. Pasch. nihil additur, neque mutatur.}\end{center}
\lettrine{F}elix dies, qua cándidæ\\
Instar colúmbæ, cœ́litum\\
Ad sacra templa spíritus\\
Se tránstulit Terésiæ;\\
\\
\textred{S}ponsíque voces áudiit~:\\
Veni soror de vértice\\
Carméli, ad Agni núptias;\\
Veni ad corónam glóriæ.\\
\\
\textred{T}e sponse Jesu Vírginum\\
Beáti adórent órdines,\\
Et nuptiáli cántico\\
Laudent per omne sǽculum.\\
Amen.

\begin{tightlist}{}
\item \textred{℣.}~Adjuvábit eam Deus vultu suo. \textred{℟.}~Deus in médio ejus non commovébitur.
\item \textred{Ad Benedictus, Aña.} Dedit ei Dóminus sapiéntiam et prudéntiam multam nimis, et latitúdinem cordis, quasi arénam, quæ est in líttore maris.
\end{tightlist}

\subsubsection*{\center \textred{Oratio.}}
\lettrine{E}xáudi nos Deus salutáris noster, ut sicut de beátæ Terésiæ Vírginis tuæ commemoratióne gaudémus~: ita cœléstis ejus doctrínæ pábulo nutriámur, et piæ devotiónis erudiámur afféctu. Per Dóminum nostrum.

\begin{tightlist}{}
\item \textred{Ad Horas Antiphonæ de Laudibus, Psalmi de Dominica, ad Primam tamen ut in Festis.}
\end{tightlist}

\subsubsection*{\center \textsc{AD TERTIAM.}}
\subsubsection*{\center \textred{Capitulum. \hfill Sap. 7.}}
\lettrine{O}ptávi, et datus est mihi sensus; et invocávi, et venit in me spíritus sapiéntiæ, et præpósui illam regnis et sédibus, et divítias nihil esse duxi in comparatióne illíus.

\begin{tightlist}{}
%\item \textred{Capit. Optávi, ut supra.}
\item \textred{℟.}~Tibi dixit cor meum, \textred{*} Exquisívit te fácies mea. \textred{T}ibi dixit. \textred{℣.}~Fáciem tuam Dómine requíram. \textred{E}xquisívit. Glória Patri. \textred{T}ibi.
\item \textred{℣.}~Anima mea exsultábit in Dómino. \textred{℟.}~Et delectábitur super salutári suo.
\end{tightlist}

\subsubsection*{\center \textsc{AD SEXTAM.}}
\subsubsection*{\textred{Capitulum. \hfill Sap. 7.}}
\lettrine{S}uper salútem et spéciem diléxi illam, et propósui pro luce habére illam, quóniam inexstinguíbile est lumen illíus.

\begin{tightlist}{}
\item \textred{℟.}~Anima mea exsultábit \textred{*} In Dómino. \textred{A}nima mea. \textred{℣.}~Et delectábitur super salutári suo. \textred{I}n. Glória Patri. \textred{A}nima.
\item \textred{℣.}~Apud me orátio Deo vitæ meæ. \textred{℟.}~Dicam Deo~: suscéptor meus es.
\end{tightlist}

\subsubsection*{\center \textsc{AD NONAM.}}
\subsubsection*{\textred{Capitulum. \hfill Sap. 7.}}
\lettrine{V}enérunt autem mihi ómnia bona páriter cum illa, et innumerábilis honéstas per manus illíus, et lætáta sum in ómnibus~: quóniam antecedébat me ista sapiéntia.

\begin{tightlist}{}
\item \textred{℟.}~Apud me orátio \textred{*} Deo vitæ meæ. \textred{A}pud me. \textred{℣.}~Dicam Deo~: suscéptor meus es. \textred{D}eo. Glória Patri. \textred{A}pud.
\item \textred{℣.}~Passer invénit sibi domum. \textred{℟.}~Et turtur nidum sibi ubi ponat pullos suos.
\end{tightlist}

\subsubsection*{\center \textsc{AD VESPERAS.}}
\begin{tightlist}{}
\item \textred{Aña j. Laud.} Hæc est virgo sápiens, et una de número prudéntum.
\item \textred{Psalmi de Comm. B.M.V.} [184]
%\item \textred{Capit. Hymn. et ℣. de j. Vesperis.}
\end{tightlist}

\subsubsection*{\textred{Capitulum. \hfill Sap. 7.}}
\lettrineThree{O}ptávi, et datus est mihi sensus; et invocávi, et venit in me spíritus sapiéntiæ; et præpósui illam regnis, et sédibus, et divítias nihil esse duxi in comparatióne illíus.

\subsubsection*{\center \textred{Hymnus.}}
\begin{center}\textred{In quo Temp. Pasch. nihil additur, neque mutatur.}\end{center}
\lettrine{R}egis supérni núntia\\
Domum patérnam déseris,\\
Terris Terésa bárbaris\\
Christum datúra, aut sánguinem.\\
\\
\textred{S}ed te manet suávior\\
Mors, pœna poscit dúlcior~:\\
Divíni amóris cúspide\\
In vulnus icta cóncides.\\
\\
\textred{O} caritátis víctima!\\
Tu corda nostra cóncrema,\\
Tibíque Gentes créditas\\
Avérni ab igne líbera.\\
\\
\textred{S}it laus Patri cum Fílio,\\
Et Spíritu Paráclito,\\
Tibíque sancta Trínitas\\
Nunc, et per omne sǽculum.\\
Amen.

\begin{tightlist}{}
\item \textred{℣.}~Oblivíscere pópulum tuum et domum patris tui. \textred{℟.}~Et concupíscet Rex decórem tuum.
\item \textred{Ad Magnificat, Aña.} Sapiéntiam ejus enarrábunt gentes, et laudem ejus enuntiábit Ecclésia.
\item \textred{Tempore Paschali tres Psalmi cujuslibet Nocturni dicuntur sub una sola Aña; et in fine Antiphonar. Responsor. et ℣.~additur} allelúja, \textred{cetera omnia ut supra.}
\end{tightlist}

\begin{center}\textcolor{red}{\rule{\linewidth}{0.5pt}} \end{center}

\subsection*{\center \textsc{Officium Sanctæ Mariæ Magdalenæ de Pazzis.}}
\subsubsection*{\center \textred{Virginis Ordinis Nostri.\\ Semiduplex}}
\begin{center}\textred{Semel in mense sub Ritu semid. juxta Rubr. 17.}\end{center}
\fancyhead[C]{\textsc{Officia Mensualia - Officium S. Mariæ Magdalenæ}}

%\begin{center}\textcolor{red}{\rule{\linewidth}{0.5pt}} \end{center}

\subsubsection*{\textred{Capitulum. \hfill Cant. 3.}}
\lettrineThree{I}n léctulo meo per noctes quæsívi quem díligit ánima mea~: quæsívi illum, et non invéni. Surgam, et circuíbo civitátem~: per vicos et platéas quæram quem díligit ánima mea.

\subsubsection*{\center \textred{Hymnus.}}
\lettrine{P}ange Carméli specióse vertex\\
Vírginis sanctæ memoránda gesta,\\
Inclytum cujus decus univérsus\\
Cóncinit orbis.\\
\\
\textred{C}ándido velans scápulas amíctu,\\
Síngulos Christi pátitur dolóres~:\\
Pérvigil cunctos súperat furéntis\\
Dǽmonis astus.\\
\\
\textred{I}nter urtícas, rigidásque spinas\\
Innocens puro rútilat cruóre~:\\
Cándida sacris mánibus Maríæ\\
Veste beátur.\\
\\
\textred{C}ordis humáni pénetrat secréta,\\
Orbis evéntus mémorat futúros,\\
Imperat morbis, fébribus medétur,\\
Ulcera sanat.\\
\\
\textred{P}róprium nunquam plácitum secúta,\\
In nono vitæ moritúra lustro,\\
Spíritum jussu cóhibens patérno,\\
Vincit agónem.\\
\\
\textred{I}n lacu quinos Erebi per annos\\
Sústinet larvas, móritur; nec unquam\\
Gnara quid fœdet sacra virginális\\
Lília cordis.\\
\\
\textred{H}orret in tumba júvenem protérvum~:\\
Cláruit vivens, moriénsque signis~:\\
Integrum miros hódie profúndit\\
Corpus odóres.\\
\\
\textred{L}aus Deo summo, Genitóque Patris,\\
Flámini sancto sit honor perénnis~:\\
Et Deum trinum venerémur uno\\
Semper honóre.\\
Amen.

\begin{tightlist}{}
\item \textred{℣.}~Ora pro nobis sancta María Magdaléna. \textred{℟.}~Ut digni efficiámur promissiónibus Christi.
\item \textred{Ad Magnificat, Aña.} Aquæ multæ non potuérunt exstínguere caritátem ejus, nec flúmina óbruent illam; quia fortis est ut mors diléctio.
\end{tightlist}

\subsubsection*{\center \textred{Oratio.}}
\lettrine{D}eus virginitátis amátor, qui beátam Maríam Magdalénam vírginem, tuo amóre succénsam, cœléstibus donis decorásti; da, ut quam devóta commemoratióne venerámur, puritáte, et caritáte imitémur. Per Dóminum nostrum.

\subsubsection*{\center \textsc{AD MATUTINUM.}}
\begin{tightlist}{}
\item \textred{Invitat.} Jesum sponsum Vírginum, quem Magdaléna diléxit multum, \textred{*} Veníte adorémus.
\item \textred{Ps.} Veníte, exsultémus. 2.
\end{tightlist}

\subsubsection*{\center \textred{Hymnus.}}
\begin{center}\textred{in quo Temp. Pach. nihil additur neque mutatur.}\end{center}
\lettrine{H}æc est Etrúsca Mágdala,\\
Magístra veri spíritus~:\\
Partem prehéndit óptimam,\\
Non auferéndam sǽculis.\\
\\
\textred{C}um Rege sacras cœ́litum\\
Vovens decénnis núptias,\\
Exspéctat indelébili\\
Sponsum coléndum lámpade.\\
\\
\textred{M}anúque dotem própria\\
Crucem, flagélla, vúlnera,\\
Cor, ánnulum, spinas, dapes\\
Dat Sponsus Eucharísticas.\\
\\
\textred{S}criptóris augústi manu\\
In corde descríptum suo,\\
Verbum caro factum est, legit,\\
Verbúmque grande intélligit.\\
\\
\textred{M}irábili sub schémate\\
Ignis, colúmnæ, flúminis,\\
Linguæ, colúmbæ, fláminis,\\
Sanctum recépit Spíritum.\\
\\
\textred{A}móris æstro pércita,\\
Compúlsat æra túrrium,\\
Ut totus orbis cónvolet,\\
Deúmque amórem díligat.\\
\\
\textred{T}e Sponse Jesu Vírginum\\
Beáti adórent Ordines,\\
Et nuptiáli cántico\\
Laudent per omne sǽculum.\\
Amen.

\subsubsection*{\center \textsc{IN PRIMO NOCTURNO.}}
\begin{tightlist}{}
\item \textred{Aña.} Ex ore puéllæ perfecísti laudem Dómine, quia novit quam admirábile est nomen tuum.
\item \textred{Ps.} Dómine, Dóminus noster, \textred{cum reliquis de Comm. B.M.V.} [172]
\item \textred{Aña.} Dóminus prǽstitit sapiéntiam párvulæ desiderábilem super aurum et lápidem pretiósum multum.
\item \textred{Aña.} Accépit benedictiónem a Dómino, quia fuit ínnocens, et mundo corde.
\item \textred{℣.}~Diffúsa est grátia in lábiis tuis. \textred{℟.}~Proptérea benedíxit te Deus in ætérnum.
\item \textred{Lectiones de Scriptura occurr. si desint, legantur sequentes.}
\end{tightlist}

\subsubsection*{\center De Cánticis Canticórum.}
\subsubsection*{\center \textred{Lectio j. \hfill Cant. 3.}}
\lettrineThree{I}n léctulo meo per noctes quæsívi quem díligit ánima mea~: quæsívi illum, et non invéni. Surgam, et circuíbo civitátem~: per vicos et platéas quæram quem díligit ánima mea~: quæsívi illum, et non invéni. Invenérunt me vígiles, qui custódiunt civitátem~: Num quem díligit ánima mea vidístis? Páululum cum pertransíssem eos, invéni quem díligit ánima mea~: ténui eum, nec dimíttam, donec introdúcam illum in domum matris meæ, et in cubículum genitrícis meæ.

\begin{tightlist}{}
\item \textred{℟.}~Flos campi diléctus meus Jesus, cujus adspéctu lætántur Angeli~: \textred{*} Cujus pulchritúdine pascúntur beáti, cujus odóre revivíscunt mórtui. \textred{℣.}~Sicut malus inter ligna silvárum, sic diléctus meus inter fílias. \textred{C}ujus.
\end{tightlist}

\subsubsection*{\center \textred{Lectio ij. \hfill Cant. 8.}}
\lettrine{Q}uis mihi det te fratrem meum sugéntem úbera matris meæ, ut invéniam te foris, et deósculer te, et jam me nemo despíciat? Apprehéndam te, et ducam in domum matris meæ~: ibi me docébis, et dabo tibi póculum ex vino condíto, et mustum malórum granatórum meórum. Læva ejus sub cápite meo, et déxtera illíus amplexábitur me. Adjúro vos, fíliæ Jerúsalem, ne suscitétis, neque evigiláre faciátis diléctam, donec ipsa velit.

\begin{tightlist}{}
\item \textred{℟.}~Sponsa Jesu ardet caritáte, fervet desidério, \textred{*} Amóris æstu rapta quo se vertat nescit. \textred{℣.}~Fulcíte me flóribus, stipáte me malis, quia amóre lángueo. \textred{A}móris.
\end{tightlist}

\subsubsection*{\center \textred{Lectio iij.}}
\lettrine{Q}uæ est ista, quæ ascéndit de desérto delíciis áffluens inníxa super diléctum suum? Sub árbore malo suscitávi te~: ibi corrúpta est mater tua, ibi violáta est génitrix tua. Pone me ut signáculum super cor tuum, ut signáculum super bráchium tuum~: quia fortis est ut mors diléctio, dura sicut inférnus æmulátio~: lámpades ejus lámpades ignis, atque flammárum. Aquæ multæ non potuérunt exstínguere caritátem, nec flúmina óbruent illam. Si déderit homo omnem substántiam domus suæ pro dilectióne, quasi nihil despíciet eam.

\begin{tightlist}{}
\item \textred{℟.}~O felícem Magdalénæ ánimam caritátis ádeo succénsam ardóribus, \textred{*} Ut præ amóre languens defíceret. \textred{℣.}~Déficit cor meum, et caro mea, Deus cordis mei, et pars mea Deus in ætérnum. \textred{U}t. Glória Patri. \textred{U}t.
\end{tightlist}

\subsubsection*{\center \textsc{IN SECUNDO NOCTURNO.}}
\begin{tightlist}{}
\item \textred{Aña.} Magdaléna oblíta est pópulum suum, et domum patris sui, et concupívit Rex decórem illíus.
%\item \textred{Ps.} Eructávit. [169]
\item \textred{Aña.} Deus, adjútor in tribulatiónibus, fuit cum Magdaléna~: ídeo non tímuit, cum transferrétur in lacum leónum.
\item \textred{Aña.} Magdalénæ fundaménta in Monte Carmélo, et diléxit eam Dóminus super ómnia tabernácula Jacob.
\item \textred{℣.}~Spécie tua et pulchritúdine tua. \textred{℟.}~Inténde, próspere procéde, et regna.
\end{tightlist}

\subsubsection*{\center Sermo sancti Ambrósii Epíscopi.}
\subsubsection*{\center \textred{Lectio iv. \hfill Lib. 1. de Virg.}}
\lettrineThree{E}go licet laudatiónem non suscéperim Virginitátis, sed expressiónem, ad rem tamen pertinére árbitror, ut quæ sit ei pátria, et quis auctor appáreat. Ac prius ubi sit pátria definiámus. Si enim ibi est pátria, ubi genitále domicílium~: in cœlo profécto est pátria castitátis~: ítaque hic ádvena, ibi íncola est. Quid autem est cástitas virginális, nisi expers contagiónis intégritas? atque ejus auctórem, quem póssumus existimáre, nisi immaculátum Dei Fílium, cujus caro non vidit corruptiónem, divínitas non est expérta contagiónem?

\begin{tightlist}{}
\item \textred{℟.}~Cum adhuc júnior essem quæsívi sapiéntiam in oratióne mea~: \textred{*} Animam meam diréxi ad illam, et excépi eam. \textred{℣.}~Spíritus tuus bonus dedúxit me in terram rectam, et docuísti me facere voluntátem tuam. \textred{A}nimam.
\end{tightlist}

\subsubsection*{\center \textred{Lectio v.}}
\lettrine{V}idéte ígitur quanta virginitátis mérita sint. Christus ante Vírginem, Christus ex Vírgine~: a Patre quidem natus ante sǽcula, sed ex Vírgine natus ob sǽcula. Illud natúræ suæ, hoc nostræ utilitátis est. Illud erat semper, hoc esse vóluit. Spectáte áliud virginitátis méritum, Christus virgo vírginis Sponsus, et si dici potest, Christus virgíneæ castitátis~: virgínitas enim Christi, non virginitátis est Christus. Virgo est ergo quæ nupsit, Virgo, quæ nos suo útero portávit, Virgo, quæ génuit, Virgo, quæ próprio lacte nutrívit. Sed jam a matre descendámus ad fílias.

\begin{tightlist}{}
\item \textred{℟.}~Certámen forte dedit illi ut vínceret, et sciret quóniam ómnium poténtior est sapiéntia, descendítque cum illa in fóveam, \textred{*} Et non derelíquit eam. \textred{℣.}~Secúndum multitúdinem dolórum meórum consolatiónes tuæ lætificavérunt ánimam meam. \textred{E}t.
\end{tightlist}

\subsubsection*{\center \textred{Lectio vj.}}
\lettrine{D}e virgínibus, ínquit Apóstolus, præcéptum Dómini non hábeo. Si Doctor géntium non hábuit, habére quis pótuit? Et præcéptum quidem non hábuit, sed hábuit exémplum. Non enim imperári potest virgínitas; sed optári. Nam quæ super nos sunt, in voto magis, quam in magistério sunt. Sed volo vos, inquit, sine sollicitúdine esse~: nam qui sine uxóre est, sollícitus est quæ Dómini sunt, quómodo pláceat Deo~: et virgo cógitat, quæ Dómini sunt, ut sit sancta córpore et spíritu. Non ego quidem dissuádeo matrimónium, sed virginitátis attéxo benefícium.

\begin{tightlist}{}
\item \textred{℟.}~Absit mihi gloriári nisi in Cruce Dómini nostri Jesu Christi, \textred{*} Per quem mihi mundus crucifíxus est, et ego mundo. \textred{℣.}~Vivo ego, jam non ego, vivit vero in me Christus. \textred{P}er. Glória Patri. \textred{P}er.
\end{tightlist}

\subsubsection*{\center \textsc{IN TERTIO NOCTURNO.}}
\begin{tightlist}{}
\item \textred{Aña.} Introívit in átria sancta Religiónis, et Dómino hóstiam própriæ voluntátis óbtulit.
%\item \textred{Ps.} Cantáte. [172]
\item \textred{Aña.} De manu dǽmonum liberávit Magdalénam Dóminus, et vidérunt omnes pópuli glóriam sanctificatiónis ejus.
\item \textred{Aña.} Mons Carméli exsúltet in conspéctu Dómini, qui mirabília fecit in Magdaléna.
\item \textred{℣.}~Adjuvábit eam Deus vultu suo. \textred{℟.}~Deus in médio ejus, non commovébitur.
\end{tightlist}

\subsubsection*{\center Lectio sancti Evangélii secúndum Matthǽum.}
\subsubsection*{\center \textred{Lectio vij. \hfill Cap. 11.}}
\lettrine{I}n illo témpore~: Respóndens Jesus dixit~: Confíteor tibi Pater, Dómine cœli et terræ, quia abscondísti hæc a sapiéntibus et prudéntibus, et revelásti ea párvulis. Et réliqua.

\begin{center}Homilía sancti Augustíni Epíscopi.\\
\textred{Serm. 9. de Verb. Dom.}\end{center}
\lettrineThree{O}mnia sæva, et immánia prorsus facília et prope nulla éfficit amor. Quanto ergo cértius ac facílius ad veram beatitúdinem cáritas facit, quod ad misériam, quantum pótuit, cupíditas fecit! quam fácile tolerátur quǽlibet advérsitas temporális, ut ætérna pœna vitétur, ætérna réquies comparétur. Non immérito ille vas electiónis cum ingénti lætítia dixit~: Non sunt condígnæ passiónes hujus témporis ad superventúram glóriam, quæ revelábitur in nobis. Ecce unde illud jugum suáve est, et sárcina levis, et si angústa est paucis eligéntibus, fácilis tamen ómnibus diligéntibus. Dixit Psalmísta~: Propter verba labiórum tuórum ego custodívi vias duras.

\begin{tightlist}{}
\item \textred{℟.}~Diléxi decórem domus tuæ, et locum habitatiónis glóriæ tuæ, \textred{*} Elégi abjécta esse in domo tua, magis quam habitáre in tabernáculis peccatórum. \textred{℣.}~Quam dilécta tabernácula tua Dómine virtútum, concupíscit, et déficit ánima mea in átria sanctæ Religiónis. \textred{E}légi.
\end{tightlist}

\subsubsection*{\center \textred{Lectio viij.}}
Sed quæ dura sunt laborántibus, éadem ipsa mitéscunt amántibus. Propter quod ita divínæ pietátis dispensatióne actum est, ut intérior homo, qui renovátur de die in diem, non adhuc sub lege pósitus, sed jam sub grátia exonerátus sárcinis innumerabílium observatiónum, quod erat revéra grave jugum, sed duræ cervíci conveniénter impósitum, facilitáte símplicis fídei, et bonæ spei et sanctæ caritátis, quidquid molestiárum exterióri hómini forínsecus intulísset ille princeps, qui missus est foras, interióri gáudio fíeret leve.

\begin{tightlist}{}
\item \textred{℟.}~Fascículus myrrhæ diléctus meus mihi, inter úbera mea commorábitur~: \textred{*} Pósui eum ut signáculum, super cor meum, et super bráchium meum. \textred{℣.}~Adjúro vos fíliæ Jerúsalem, si invenéritis diléctum meum, ut nuntiétis ei, quia amóre lángueo. \textred{P}ósui.
\end{tightlist}

\subsubsection*{\center \textred{Lectio ix.}}
Nihil enim tam fácile bonæ voluntáti est quam ipsa sibi, et hæc súfficit Deo. Quantúmlibet ergo sǽviat iste mundus, veríssime Angeli, nato in carne Dómino, clamavérunt~: Glória in excélsis Deo, et in terra pax homínibus bonæ voluntátis. Quia, qui natus erat, suáve jugum est, et sárcina levis. Et sicut dicit Apóstolus~: Fidélis Deus, qui nos non sinit tentári supra id, quod póssumus ferre, sed facit cum tentatióne étiam éxitum, ut possímus sustinére.

\begin{tightlist}{}
\item \textred{℟.}~Egredímini fíliæ Sion, et vidéte Sponsam in diadémate, quo coronávit eam Sponsus suus~: \textred{*} In die desponsatiónis illíus, et lætítiæ cordis ejus. \textred{℣.}~Collocávit super caput sponsæ suæ corónam spíneam, quæ pósita fuit super caput suum. \textred{I}n. Glória Patri. \textred{I}n.
\item Te Deum laudámus. 6.
\item \textred{Sacerdotalis ℣.}~Ora pro nobis sancta María Magdaléna. \textred{℟.}~Ut digni efficiámur promissiónibus Christi, allelúja.
\end{tightlist}

\subsubsection*{\center \textsc{AD LAUDES.}}
\begin{center}\textred{et per Horas, Aña.}\end{center}
\lettrine{M}irábilis in Magdaléna Dóminus, quam e vulnéribus suis decórem índuit, et advérsus dǽmonum insídias fortitúdine præcínxit.

\begin{tightlist}{}
\item \textred{Ps. de Dominica.} 27.
\item \textred{Aña.} Relíctis sǽculi pompis, introívit in conspéctu Dei Magdaléna, quam ánnulo suo subarrhávit Dóminus.
\item \textred{Aña.} In terra desérta tentatiónum appáruit Magdaléna, ut vidéret virtútem et glóriam tuam.
\item \textred{Aña.} Patrem, et Fílium cum sancto Spíritu puérula benedicébat, dum sanctæ Trinitátis fidem párvulos docébat.
\item \textred{Aña.} In týmpano et choro, in chordis, et órgano omnis spíritus laudet Dóminum, cujus in Magdaléna nomen exaltátum est.
\end{tightlist}

\subsubsection*{\textred{Capitulum. \hfill Cant. 3.}}
\lettrineThree{I}n léctulo meo per noctes quæsívi quem díligit ánima mea~: quæsívi illum, et non invéni. Surgam, et circuíbo civitátem~: per vicos et platéas quæram quem díligit ánima mea.

\subsubsection*{\center \textred{Hymnus.}}
\lettrine{C}orporis nexus óbitu rescíndens,\\
Plena virtútum, méritis onústa,\\
Scandit ætérnas, comitánte Sponso,\\
Mágdala sedes.\\
\\
\textred{A}bdito ventris látitans recéssu,\\
Non gravat matris grémium, nec unquam\\
Críminis pressit própriam molésto\\
Póndere mentem.\\
\\
\textred{O}btegit nudos, docet imperítos,\\
Cárceres visit, juvat univérsos~:\\
Plúrimas fusis ánimas ab orco\\
Flétibus arcet.\\
\\
\textred{D}éserens mundum, crucifíxa Christo,\\
Corpus afflíxit précibus, caténis\\
Férreis, flagris, fame, felle, spinis,\\
Frígore, flamma.\\
\\
\textred{C}um suis Horas cécinit Patrónis,\\
Exstases, raptus pátitur frequénter,\\
Osculo Sponsi recreáta, Jesum\\
Portat in ulnis.\\
\\
\textred{N}unc opem nobis miseráta confer,\\
Qui tuos læti cánimus triúmphos;\\
Erudi mentem, refluátque amóris\\
Flamma perénnis.\\
\\
\textred{L}aus Deo summo, Genitóque Patris,\\
Flámini sancto sit honor perénnis~:\\
Et Deum Trinum venerémur uno\\
Semper honóre.\\
Amen.

\begin{tightlist}{}
\item \textred{℣.}~Elégit eam Deus, et præelégit eam. \textred{℟.}~Et habitáre facit in tabernáculo suo.
\item \textred{Ad Benedictus, Aña.} Magdaléna Virgo óptimam partem elégit, quæ non auferétur ab ea~: puríssima in vita, adhuc post mortem ab omni corruptióne manet immúnis.
\end{tightlist}

\subsubsection*{\center \textred{Oratio.}}
\lettrine{D}eus Virginitátis amátor, qui beátam Maríam Magdalénam Vírginem, tuo amóre succénsam, cœléstibus donis decorásti~: da, ut quam devóta commemoratióne venerámur, puritáte et caritáte imitémur. Per Dóminum nostrum.

\begin{tightlist}{}
\item \textred{Ad Horas Antiphonæ de Laudibus, Psalmi de Dominica, ad Primam tamen ut in Festis.}
\end{tightlist}

\subsubsection*{\center \textsc{AD TERTIAM.}}
\subsubsection*{\textred{Capitulum. \hfill Cant. 3.}}
\lettrine{I}n léctulo meo per noctes quæsívi quem díligit ánima mea~: quæsívi illum, et non invéni. Surgam, et circuíbo civitátem~: per vicos et platéas quæram quem díligit ánima mea.

\begin{tightlist}{}
%\item \textred{Capit. In léctulo meo.}
\item \textred{℟.}~Diffúsa est grátia \textred{*} In lábiis tuis. \textred{D}iffúsa. \textred{℣.}~Proptérea benedíxit te Deus in ætérnum. \textred{I}n. Glória Patri. \textred{D}iffúsa.
\item \textred{℣.}~Spécie tua et pulchritúdine tua. \textred{℟.}~Inténde, próspere procéde, et regna.
\end{tightlist}

\subsubsection*{\center \textsc{AD SEXTAM.}}
\subsubsection*{\textred{Capitulum. \hfill Cant. 3.}}
\lettrine{I}nvenérunt me vígiles, qui custódiunt civitátem~: num quem díligit ánima mea vidístis? Páululum cum pertransíssem eos, invéni quem díligit ánima mea~: ténui eum nec dimíttam.

\begin{tightlist}{}
\item \textred{℟.}~Spécie tua \textred{*} Et pulchritúdine tua. \textred{S}pécie. \textred{℣.}~Inténde, próspere procéde, et regna. \textred{E}t. Glória Patri. \textred{S}pécie.
\item \textred{℣.}~Adjuvábit eam Deus vultu suo. \textred{℟.}~Deus in médio ejus, non commovébitur.
\end{tightlist}

\subsubsection*{\center \textsc{AD NONAM.}}
\subsubsection*{\textred{Capitulum. \hfill Cant. 3.}}
\lettrine{A}djúro vos, fíliæ Jerúsalem, si invenéritis diléctum meum, ut nuntiétis ei, quia amóre lángueo.

\begin{tightlist}{}
\item \textred{℟.}~Adjuvábit eam \textred{*} Deus vultu suo. \textred{A}djuvábit. \textred{℣.}~Deus in médio ejus, non commovébitur. \textred{D}eus. Glória Patri. \textred{A}djuvábit.
\item \textred{℣.}~Elégit eam Deus, et præelégit eam. \textred{℟.}~Et habitáre eam facit in tabernáculo suo.
\end{tightlist}

\subsubsection*{\center \textsc{AD II. VESPERAS.}}
\begin{tightlist}{}
\item \textred{Aña j. Laud.} Mirábilis in Magdaléna Dóminus, quam e vulnéribus suis decórem índuit, et advérsus dǽmonum insídias fortitúdine præcínxit. 
\item \textred{Psalmi de Comm. B.M.V.} [184]
%\item \textred{Capit. Hymn. et ℣. ut in j. Vesperis.}
\end{tightlist}

\subsubsection*{\textred{Capitulum. \hfill Philipp. 2.}}
\lettrineThree{I}n léctulo meo per noctes quæsívi quem díligit ánima mea~: quæsívi illum, et non invéni. Surgam, et circuíbo civitátem~: per vicos et platéas quæram quem díligit ánima mea.

\subsubsection*{\center \textred{Hymnus.}}
\lettrine{P}ange Carméli specióse vertex\\
Vírginis sanctæ memoránda gesta,\\
Inclytum cujus decus univérsus\\
Cóncinit orbis.\\
\\
\textred{C}ándido velans scápulas amíctu,\\
Síngulos Christi pátitur dolóres~:\\
Pérvigil cunctos súperat furéntis\\
Dǽmonis astus.\\
\\
\textred{I}nter urtícas, rigidásque spinas\\
Innocens puro rútilat cruóre~:\\
Cándida sacris mánibus Maríæ\\
Veste beátur.\\
\\
\textred{C}ordis humáni pénetrat secréta,\\
Orbis evéntus mémorat futúros,\\
Imperat morbis, fébribus medétur,\\
Ulcera sanat.\\
\\
\textred{P}róprium nunquam plácitum secúta,\\
In nono vitæ moritúra lustro,\\
Spíritum jussu cóhibens patérno,\\
Vincit agónem.\\
\\
\textred{I}n lacu quinos Erebi per annos\\
Sústinet larvas, móritur; nec unquam\\
Gnara quid fœdet sacra virginális\\
Lília cordis.\\
\\
\textred{H}orret in tumba júvenem protérvum~:\\
Cláruit vivens, moriénsque signis~:\\
Integrum miros hódie profúndit\\
Corpus odóres.\\
\\
\textred{L}aus Deo summo, Genitóque Patris,\\
Flámini sancto sit honor perénnis~:\\
Et Deum trinum venerémur uno\\
Semper honóre.\\
Amen.

\begin{tightlist}{}
\item \textred{℣.}~Ora pro nobis sancta María Magdaléna. \textred{℟.}~Ut digni efficiámur promissiónibus Christi.
\item \textred{Ad Magnificat, Aña.} Veni Sponsa de Líbano, veni, coronáberis de cápite Amana, de vértice Sanir, et Hermon, de cubílibus leónum, et móntibus pardórum.
\item \textred{Completorium de Dominica.}
\item \textred{Tempore Paschali tres Psalmi cujuslibet Nocturni dicuntur sub una sola Antiphona, et in fine Añarum etc. additur} Allelúja. \textred{Cetera omnia ut supra.}
\end{tightlist}
\begin{center}\textcolor{red}{\rule{\linewidth}{0.5pt}} \end{center}
\clearpage
%\begin{center}\textcolor{red}{\rule{\linewidth}{0.5pt}} \end{center}
\section*{\center \textred{\textsc{OFFICIA VOTIVA PER ANNUM.}}}
\subsubsection*{\center \textred{\textsc{Dicenda ad libitum}}}

\lettrineThree{O}\textred{fficia hæc Votiva per annum ritu Semiduplici pro singulis Hebdomadæ Feriis a Summo Pontifice Leone XIII. per Decretum \textsc{Urbis et Orbis} Sacrorum Rituum Congregationis die 5 Julii 1883 generali Indulto concessa, et ad Nostrum Ordinem extensa Decr. 20 Martii 1884, persolvi possunt, juxta idem Indultum, loco Officiorum Ferialium occurrentium, quocumque anni tempore; exceptis Feria IV. Cinerum, Feriis totius temporis Passionis, ac Sacri Adventus a die 17 ad 24 Decembris inclusive.}

\begin{center}\textcolor{red}{\rule{\linewidth}{0.5pt}} \end{center}

\subsection*{\center \textred{\textsc{PRO FERIA II.}}}
\subsection*{\center \textsc{\textred{OFFICIUM VOTIVUM}\\
De Sanctis Angelis.}}
\subsubsection*{\center \textred{Semiduplex.}}
\fancyhead[C]{\textsc{Officia Votiva - Feria II - De Angelis.}}


\subsubsection*{\center \textsc{AD VESPERAS.}}
\subsubsection*{\textred{Capitulum. \hfill Exod. 23.}}
\lettrineThree{E}cce ego mittam Angelum meum, qui præcédat te, et custódiat in via, et introdúcat in locum quem parávi. Obsérva eum, et audi vocem ejus.

\subsubsection*{\center \textred{Hymnus.}}
\lettrine{T}ibi, Christe, splendor Patris,\\
Vita, virtus córdium,\\
In conspéctu Angelórum,\\
Votis, voce psállimus~:\\
Alternántes concrepándo\\
Melos damus vócibus.\\
\\
\textred{C}ollaudámus venerántes\\
Omnes cæli mílites;\\
Sed præcípue primátem\\
Cæléstis exércitus\\
Michaélem, in virtúte,\\
Conteréntem zábulum.\\
\\
\textred{Q}uo custóde, procul pelle\\
Rex Christe piíssime,\\
Omne nefas inimíci,\\
Mundo corde, et córpore,\\
Paradíso redde tuo\\
Nos sola cleméntia.\\
\\
\textred{G}lóriam Patri melódis\\
Personémus vócibus,\\
Glóriam Christo canámus,\\
Glóriam Paráclito;\\
Qui Trinus, et unus Deus\\
Exstat ante sǽcula.\\
Amen.

\begin{tightlist}{}
\item \textred{℣.}~In conspéctu Angelórum psallam tibi, Deus meus. \textred{℟.}~Adorábo ad templum sanctum tuum, et confitébor nómini tuo.
\item \textred{Ad Magnif. Aña.} Omnes sunt administratórii spíritus, in ministérium missi propter eos, qui hereditátem cápiunt salútis.
\end{tightlist}

\subsubsection*{\center \textred{Oratio.}}
\lettrine{D}eus, qui miro órdine Angelórum ministéria, hominúmque dispénsas~: concéde propítius; ut a quibus tibi ministrántibus in cœlo semper assístitur, ab his in terra vita nostra muniátur. Per Dóminum.

\subsubsection*{\center \textsc{AD MATUTINUM.}}
\fancyhead[C]{\textsc{Officia Votiva - Feria II - De Angelis.}}
\begin{tightlist}{}
\item \textred{Invit.} Regem Angelórum Dóminum, \textred{*} Veníte, adorémus.
\item \textred{Ps.} Veníte, exsultémus. 2.
\item \textred{Hymnus.} Tibi Christe, \textred{ut supra.}
\end{tightlist}

\subsubsection*{\center \textsc{IN PRIMO NOCTURNO.}}
\fancyhead[C]{\textsc{Officia Votiva - Feria II - De Angelis.}}
\begin{tightlist}{}
\item \textred{Aña.} Multa magnália de Michaéle Archángelo, qui fortis in prælio fecit victóriam.
\item \textred{Ps. 8.} Dómine, Dóminus. [25] 
\item \textred{Aña.} Gábriel Archángelus locútus est Maríæ dicens~: Ecce concípies in útero et páries Fílium et vocábis nomen ejus Jesum.
\item \textred{Ps. 10.} In Dómino confído. [27]
\item \textred{Aña.} Benedícite Deum cœli, dixit Angelus Ráphael, et coram ómnibus vivéntibus confitémini ei, quia fecit vobíscum misericórdiam suam.
\item \textred{Ps. 14.} Dómine, quis. [28]
\item \textred{℣.}~In conspéctu Angelórum psallam tibi, Deus meus. \textred{℟.}~Adorábo ad templum sanctum tuum, et confitébor nómini tuo.
\item \textred{Lectiones de Scriptura occurr. sed si desint, legantur sequentes.}
\end{tightlist}

\subsubsection*{\center De libro Exodi.}
\subsubsection*{\center \textred{Lectio j. \hfill Cap. 23.}}
\lettrineThree{E}cce ego mittam Angelum meum, qui præcédat te, et custódiat in via, et introdúcat in locum quem parávi.
Obsérva eum, et audi vocem ejus, nec contemnéndum putes~: quia non dimíttet cum peccáveris, et est nomen meum in illo.
Quod si audíeris vocem ejus, et féceris ómnia quæ loquor, inimícus ero inimícis tuis, et afflígam affligéntes te. Præcedétque te Angelus meus.

\begin{tightlist}{}
\item \textred{℟.}~Venit Míchael Archángelus cum multitúdine Angelórum, cui trádidit Deus ánimas sanctórum, \textred{*} Ut perdúcat eas in paradísum exsultatiónis. \textred{℣.}~Emítte, Dómine, Spíritum sanctum tuum de cœlis, Spíritum sapiéntiæ et intelléctus. \textred{U}t.
\end{tightlist}

\subsubsection*{\center Ex Zacharía Prophéta.}
\subsubsection*{\center \textred{Lectio ij. \hfill Cap. 1.}}
\lettrine{F}actum est verbum Dómini ad Zacharíam, fílium Barachíæ fílii Addo, prophétam, dicens~: Vidi per noctem, et ecce vir ascéndens super equum rufum, et ipse stabat inter myrtéta, quæ erant in profúndo~: et post eum equi rufi, várii, et albi.
Et dixi~: Quid sunt isti, Dómine mi? Et dixit ad me Angelus, qui loquebátur in me~: Ego osténdam tibi quid sint hæc.
Et respóndit vir, qui stabat inter myrtéta, et dixit~: Isti sunt, quos misit Dóminus ut perámbulent terram.
Et respondérunt Angelo Dómini, qui stabat inter myrtéta, et dixérunt~: Perambulávimus terram, et ecce omnis terra habitátur, et quiéscit. Hæc dicit Dóminus.

\begin{tightlist}{}
\item \textred{℟.}~Respóndit Angelus Dómini, et dixit~: Dómine exercítuum, \textred{*} Usquequo tu non miseréberis Jerúsalem et úrbium Juda, quibus irátus es? \textred{℣.}~Iste enim septuagésimus annus est. \textred{U}squequo.
\end{tightlist}

\subsubsection*{\center \textred{Lectio iij. \hfill Cap. 2.}}
\lettrine{E}t levávi óculos meos, et vidi~: et ecce vir, et in manu ejus funículus mensórum.
Et dixi~: Quo tu vadis? Et dixit ad me~: Ut métiar Jerúsalem, et vídeam quanta sit latitúdo ejus, et quanta longitúdo ejus.
Et ecce Angelus, qui loquebátur in me, egrediebátur, et Angelus álius egrediebátur in occúrsum ejus.
Et dixit ad eum~: Curre, lóquere ad púerum istum, dicens~: Absque muro habitábitur Jerúsalem præ multitúdine hóminum, et jumentórum in médio ejus.
Et ego ero ei, ait Dóminus, murus ignis in circúitu~: et in glória ero in médio ejus. Hæc dicit Dóminus.

\begin{tightlist}{}
\item \textred{℟.}~In conspéctu Géntium nolíte timére~: vos enim in córdibus vestris adoráte et timéte Dóminum~: \textred{*} Angelus enim ejus vobíscum est. \textred{℣.}~Stetit Angelus juxta aram templi, habens thuríbulum áureum in manu sua. \textred{A}ngelus. Glória Patri. \textred{A}ngelus.
\end{tightlist}

\subsubsection*{\center \textsc{IN SECUNDO NOCTURNO.}}
%\fancyhead[C]{\textsc{Officia Votiva - Feria II - De Angelis.}}
\begin{tightlist}{}
\item \textred{Aña.} Ascéndit fumus arómatum in conspéctu Dómini de manu Angeli.
\item \textred{Ps. 18.} Cœli enárrant. [173]
\item \textred{Aña.} Immíttet Angelus Dómini in circúitu timéntium eum, et erípiet eos.
\item \textred{Ps. 23.} Dómini est terra. [174]
\item \textred{Aña.} Laudémus Dóminum, quem laudant Angeli, quem Chérubim et Séraphim, Sanctus, Sanctus, Sanctus, proclámant.
\item \textred{Ps. 33.} Benedícam Dóminum. [6]
\item \textred{℣.}~Stétit Angelus juxta aram templi. \textred{℟.}~Habens thuríbulum áureum in manu sua.
\end{tightlist}

\subsubsection*{\center De Sermóne sancti Gregórii Papæ.}
\subsubsection*{\center \textred{Ex. Homilia 34. in Evang.\\
Lectio iv.}}
\lettrineThree{Q}uia Angelórum nómina interpretándo perstrínximus, súperest, ut ipsa officiórum vocábula bréviter exsequámur. Virtútes étenim vocántur illi nimírum spíritus, per quos signa et mirácula frequéntius fiunt. Potestátes étiam vocántur hi, qui hoc poténtius céteris in suo órdine percepérunt, ut eórum ditióni virtútes advérsæ subjéctæ sint, quorum potestáte refrænántur, nec corda hóminum tantum tentáre præváleant, quantum volunt.

\begin{tightlist}{}
\item \textred{℟.}~Omnes Angeli stabant in circúitu throni, et seniórum, et quátuor animálium~: \textred{*} Et cecidérunt in conspéctu throni in fácies suas, et adoravérunt Deum. \textred{℣.}~Adoráte Dóminum omnes Angeli ejus. \textred{E}t.
\end{tightlist}

\subsubsection*{\center \textred{Lectio v.}}
\lettrine{P}rincipátus étiam vocántur, qui ipsis quoque bonis Angelórum spirítibus præsunt, qui subjéctis áliis dum quæquæ sunt agénda dispónunt, eis ad explénda divína ministéria principántur. Dominatiónes autem vocántur, qui étiam potestátes principátuum dissimilitúdine alta transcéndunt. Nam principári est inter réliquos priórem exsístere~: dominári vero est étiam subjéctos quosque possidére. Throni quoque illa ágmina sunt vocáta, quibus ad exercénda judícia semper Deus omnípotens prǽsidet. Unde et per Psalmístam dícitur~: Sedes super thronum, qui júdicas æquitátem.

\begin{tightlist}{}
\item \textred{℟.}~Stetit Angelus juxta aram templi, habens thuríbulum áureum in manu sua, et data sunt ei incénsa multa~: \textred{*} Et ascéndit fumus arómatum de manu Angeli in conspéctu Dómini. \textred{℣.}~In conspéctu Angelórum psallam tibi~: adorábo ad templum sanctum tuum, et confitébor nómini tuo, Dómine. \textred{E}t.
\end{tightlist}

\subsubsection*{\center \textred{Lectio vj.}}
\lettrine{C}hérubim quoque plenitúdo sciéntiæ dícitur; et sublimióra illa ágmina idcírco Chérubim vocáta sunt, quia tanto perfectióri sciéntia plena sunt, quanto claritátem Dei vicínius contemplántur. Séraphim étiam vocántur illa spirítuum sanctórum ágmina, quæ ex singulári propinquitáte Conditóris sui incomparábili ardent amóre. Séraphim namque ardéntes vel incendéntes vocántur; quorum profécto flamma amor est~: quia quo subtílius claritátem divinitátis ejus adspíciunt, eo valídius in ejus amóre flamméscunt.

\begin{tightlist}{}
\item \textred{℟.}~In conspéctu Angelórum psallam tibi, et adorábo ad templum sanctum tuum~: \textred{*} Et confitébor nómini tuo, Dómine. \textred{℣.}~Super misericórdia tua et veritáte tua~: quóniam magnificásti super nos nomen sanctum tuum. \textred{E}t. Glória Patri. \textred{E}t confitébor.
\end{tightlist}

\subsubsection*{\center \textsc{IN TERTIO NOCTURNO.}}
\begin{tightlist}{}
\item \textred{Aña.} Misit Dóminus Angelum suum, qui percússit omnem virum robústum, et bellatórem, et príncipem exércitus regis Assyriórum.
\item \textred{Ps. 95.} Cantáte Dómino. [179]
\item \textred{Aña.} Adoráte Dóminum omnes Angeli ejus~: audívit, et lætáta est Sion.
\item \textred{Ps. 96.} Dóminus regnávit. [180]
\item \textred{Aña.} Benedícite Dómino omnes Angeli ejus, poténtes virtúte, faciéntes verbum illíus, ad audiéndam vocem sermónum ejus.
\item \textred{Ps. 102} Bénedic, ánima. 288.
\item \textred{℣.}~Ascendit fumus arómatum in conspéctu Dómini. \textred{℟.}~De manu Angeli.
\end{tightlist}

\subsubsection*{\center Léctio sancti Evangélii secúndum Joánnem.}
\subsubsection*{\center \textred{Lectio vij. \hfill Cap. 1.}}
\lettrine{I}n illo témpore~: Vidit Jesus Nathánael veniéntem ad se, et dicit de eo~: Ecce vere Israelíta, in quo dolus non est. Et réliqua.

\begin{center}Homilía sancti Bernárdi Abbátis.\\
\textred{In Ps. Qui habitat.\\
Sermo 11.}\end{center}
\lettrineThree{Q}uæ vero sunt sanctórum Angelórum viæ? Profécto illæ, quas Unigénitus enarrávit, dicens~: Vidébitis Angelos ascendéntes et descendéntes super Fílium hóminis. Ascénsio ígitur et descénsio viæ illórum; ascénsio propter se, descénsio, vel pótius condescénsio propter nos. Sic beáti illi spíritus ascéndunt per contemplatiónem Dei, descéndunt per compassiónem tui, ut custódiant te in ómnibus viis tuis. Ascéndunt ad vultum ejus, descéndunt ad nutum ejus~: quóniam Angelis suis mandávit de te. Nec tamen vel descendéndo visióne glóriæ fraudántur, quia semper vident fáciem Patris.

\begin{tightlist}{}
\item \textred{℟.}~Angelus Dómini descéndit cum Azaría et sóciis ejus in fornácem, et excússit flammam ignis de fornáce~: \textred{*} Et non tétigit eos omníno ignis, neque contristávit. \textred{℣.}~Benedíctus Deus eórum, qui misit Angelum suum, et éruit servos suos, qui credidérunt in eum. \textred{E}t.
\end{tightlist}

\subsubsection*{\center \textred{Lectio viij.}}
\lettrine{C}um ascéndunt ad contemplatiónem, inquírunt veritátem, de qua et desiderándo satiántur, et satiándo desíderant. Cum vero descéndunt, fáciunt nobíscum misericórdiam, ut custódiant nos in ómnibus viis nostris. Administratórii enim spíritus sunt, missi in ministérium propter nos. Plane minístri nostri, non dómini nostri. Et in hoc Unigéniti formam imitántur, qui non venit ministrári, sed ministráre, qui stetit inter discípulos tamquam qui minístrat. Fructus angelicárum viárum, quod ad ipsos spectat, sua ipsórum beatitúdo, et obediéntia est caritátis~: quod autem ad nos pértinet, inde quidem obténtus divínæ grátiæ, hinc vero custódia viæ nostræ~: síquidem Angelis suis mandávit de te, ut custódiant te in ómnibus indigéntiis tuis.

\begin{tightlist}{}
\item \textred{℟.}~Tu, Dómine, qui misísti Angelum tuum sub Ezechía, rege Juda, et interfecísti de castris Senácherib centum octogínta quinque míllia~: \textred{*} Et nunc Dominátor cœlórum mitte Angelum tuum bonum ante nos, in timóre et tremóre magnitúdinis bráchii tui. \textred{℣.}~Ut métuant qui cum blasphémia véniunt advérsus sanctum pópulum tuum. \textred{E}t.
\end{tightlist}

\subsubsection*{\center \textred{Lectio ix.}}
\lettrine{I}nterim ergo mandávit Angelis suis Deus, non quidem ut amóveant te a viis tuis, sed in eis ipsis te custódiant, et quasi per vias suas, vias tuas dírigant in vias ipsíus. Quonam modo, inquis? Nempe ut quod Angelus ex sola caritáte púrius agit, ut própria saltem necessitáte compúlsus et admónitus, descéndas, et condescéndas; exhibére scílicet próximo misericórdiam~: et rursum cum eódem Angelo levans desidéria tua, tota ánimi cupiditáte ascéndere stúdeas ad summa et perpétuam Veritátem. Hinc enim monémur leváre corda nostra cum mánibus; hinc audímus quotídie~: Sursum corda. Hinc quoque negligéntes argúimur, et dícitur nobis~: Fílii hóminum, úsquequo gravi corde? ut quid dilígitis vanitátem et quǽritis mendácium? Exonerátum enim cor et leve levátur magis, ut quærat et díligat veritátem.

\begin{tightlist}{}
\item \textred{℟.}~Te sanctum Dóminum in excélsis laudant omnes Angeli, dicéntes~: \textred{*} Te decet laus, et honor Dómine. \textred{℣.}~Chérubim quoque, ac Séraphim, Sanctus proclámant, et omnis cœlicus ordo, dicéntes. \textred{T}e. Glória Patri. \textred{T}e.
\item Te Deum laudámus. 6.
\item \textred{Sacerd. ℣.}~Sancti Angeli, Custódes nostri, deféndite nos in prǽlio. \textred{℟.}~Ut non pereámus in treméndo judício.
\end{tightlist}

\subsubsection*{\center \textsc{AD LAUDES.}}
\begin{center}\textred{et per Horas, Aña.}\end{center}
\lettrine{A}ngeli Dómini Dóminum benedícite in ætérnum.

\begin{tightlist}{}
\item \textred{Ps. de Dominica.} 27.
\item \textred{Aña.} Angelis suis Deus mandávit de te, ut custódiant te in ómnibus viis tuis.
\item \textred{Aña.} Angeli eórum semper vident fáciem Patris mei, qui est in cœlis.
\item \textred{Aña.} Laudáte Deum, omnes Angeli ejus, laudáte eum, omnes virtútes ejus.
\item \textred{Aña.} Angeli, Archángeli, Throni et Dominatiónes, Principátus et Potestátes, Virtútes cœlórum, laudáte Dóminum de cœlis, allelúja.
\end{tightlist}

\subsubsection*{\textred{Capitulum. \hfill Exod. 23.}}
\lettrineThree{E}cce ego mittam Angelum meum, qui præcédat te, et custódiat in via, et introdúcat in locum quem parávi. Obsérva eum, et audi vocem ejus.

\subsubsection*{\center \textred{Hymnus.}}
\lettrine{C}hriste, sanctórum decus Angelórum,\\
Rector humáni géneris, et auctor,\\
Nobis ætérnum tríbue benígnus\\
Scándere cœlum.\\
\\
\textred{A}ngelum pacis Míchael ad istam\\
Cœ́litus mitti rogitámus aulam~:\\
Nobis ut crebro veniénte, crescant\\
Próspera cuncta.\\
\\
\textred{A}ngelus fortis Gábriel, ut hostem\\
Pellat antíquum, vólitet ab alto,\\
Sǽpius templum véniens ad istud\\
Vísere nostrum.\\
\\
\textred{A}ngelum nobis médicum salútis\\
Mite de cœlis Ráphael, ut omnes\\
Sanet ægrótos, paritérque nostros\\
Dírigat actus.\\
\\
\textred{H}inc Dei nostri Génitrix María,\\
Totus et nobis chorus Angelórum\\
Semper assístat, simul et beáta\\
Cóncio tota.\\
\\
\textred{P}ræstet hoc nobis Déitas beáta\\
Patris, ac Nati, paritérque sancti\\
Spíritus cujus réboat in omni\\
Glória mundo.\\
Amen.

\begin{tightlist}{}
\item \textred{℣.}~Adoráte Deum. \textred{℟.}~Omnes Angeli ejus.
\item \textred{Ad Benedictus, Aña.} Angelus, qui loquebátur in me, et suscitávit me, quasi virum qui suscitátur a somno suo.
\end{tightlist}

\subsubsection*{\center \textred{Oratio.}}
\lettrine{D}eus, qui miro órdine Angelórum ministéria, hominúmque dispénsas~: concéde propítius; ut a quibus tibi ministrántibus in cœlo semper assístitur, ab his in terra vita nostra muniátur. Per Dóminum.

\begin{tightlist}{}
\item \textred{Ad Horas Añæ de Laudibus, Psalmi de Dominica, ad Primam tamen ut in Festis.}
\end{tightlist}

\subsubsection*{\center \textsc{AD TERTIAM.}}
\subsubsection*{\textred{Capitulum. \hfill Exod. 23.}}
\lettrine{E}cce ego mittam Angelum meum, qui præcédat te, et custódiat in via, et introdúcat in locum quem parávi. Obsérva eum, et audi vocem ejus.

\begin{tightlist}{}
\item \textred{℟.}~In conspéctu Angelórum \textred{*} Psallam tibi, Deus meus. \textred{I}n. \textred{℣.}~Adorábo ad templum sanctum tuum, et confitébor nómini tuo. \textred{P}sallam. Glória Patri. \textred{I}n.
\item \textred{℣.}~Stetit Angelus juxta aram templi. \textred{℟.}~Habens thuríbulum áureum in manu sua.
\end{tightlist}

\subsubsection*{\center \textsc{AD SEXTAM.}}
\subsubsection*{\textred{Capitulum. \hfill Apoc. 5.}}
\lettrine{A}udívi vocem Angelórum multórum in circúitu throni, et animálium, et seniórum~: et erat númerus eórum míllia míllium, voce magna dicéntium~: Salus Deo nostro.

\begin{tightlist}{}
\item \textred{℟.}~Stetit Angelus \textred{*} Juxta aram templi. \textred{S}tetit. \textred{℣.}~Habens thuríbulum áureum in manu sua. \textred{J}uxta. Glória Patri. \textred{S}tetit.
\item \textred{℣.}~Ascéndit fumus arómatum in conspéctu Dómini. \textred{℟.}~De manu Angeli.
\end{tightlist}

\subsubsection*{\center \textsc{AD NONAM.}}
\subsubsection*{\textred{Capitulum. \hfill Apoc. 12.}}
\lettrine{F}actum est prǽlium magnum in cœlo~: Míchael et Angeli ejus præliabántur cum dracóne, et draco pugnábat, et ángeli ejus~: et non prævaluérunt, neque locus invéntus est eórum ámplius in cœlo.

\begin{tightlist}{}
\item \textred{℟.}~Ascéndit fumus arómatum \textred{*} In conspéctu Dómini. \textred{A}scéndit. \textred{℣.}~De manu Angeli. \textred{I}n. Glória Patri. \textred{A}scéndit.
\item \textred{℣.}~Adoráte Deum. \textred{℟.}~Omnes Angeli ejus.
\end{tightlist}

\subsubsection*{\center \textsc{AD II. VESPERAS.}}
\begin{tightlist}{}
\item \textred{Aña. prima de Laudibus.} Angeli Dómini, Dóminum benedícite in ætérnum. 
\item \textred{Ps. de Dominica et loco ultimi, Ps. 137.} Confitébor tibi, Dómine. 205.
%\item \textred{Capit. Hymnus et ℣. ut in j. Vesperis.}
\end{tightlist}

\subsubsection*{\textred{Capitulum. \hfill Exod. 23.}}
\lettrineThree{E}cce ego mittam Angelum meum, qui præcédat te, et custódiat in via, et introdúcat in locum quem parávi. Obsérva eum, et audi vocem ejus.

\subsubsection*{\center \textred{Hymnus.}}
\lettrine{T}ibi, Christe, splendor Patris,\\
Vita, virtus córdium,\\
In conspéctu Angelórum\\
Votis, voce psállimus~:\\
Alternantes concrepándo\\
Melos damus vócibus.\\
\\
\textred{C}ollaudámus venerántes\\
Omnes Cæli mílites;\\
Sed præcípue Primátem\\
Cæléstis exércitus\\
Michaélem, in virtúte\\
Conteréntem zábulum.\\
\\
\textred{Q}uo custóde, procul pelle\\
Rex, Christe, piíssime,\\
Omne nefas inimíci,\\
Mundo corde et córpore,\\
Paradíso redde tuo\\
Nos sola cleméntia.\\
\\
\textred{G}lóriam Patri melódis\\
Personémus vócibus,\\
Glóriam Christo canámus,\\
Glóriam Paráclito~:\\
Qui Trinus et unus Deus\\
Exstat ante sǽcula.\\
Amen.

\begin{tightlist}{}
\item \textred{℣.}~In conspéctu Angelórum psallam tibi, Deus meus. \textred{℟.}~Adorábo ad templum sanctum tuum, et confitébor nómini tuo.
\item \textred{Ad Magnif. Aña.} Sancti Angeli, qui Deo in cœlis semper assístitis, deféndite nos in prǽlio, ut non pereámus in treméndo judício.
\end{tightlist}

\begin{center}\textcolor{red}{\rule{\linewidth}{0.5pt}} \end{center}

\subsection*{\center \textred{\textsc{PRO FERIA III.}}}
\subsection*{\center \textsc{\textred{OFFICIUM VOTIVUM}\\ De Sanctis Apostolis.}}
\subsubsection*{\center \textred{Semiduplex}}

\begin{tightlist}{}
\item \textred{Omnia de Communi Apostolorum} [2] \textred{juxta diversitatem temporis, præter sequentia.}
\end{tightlist}

\subsubsection*{\center \textred{Oratio.}}
\lettrine{D}eus, qui nos per beátos Apóstolos tuos ad agnitiónem tui nóminis veníre tribuísti~: da nobis eórum glóriam sempitérnam et proficiéndo celebráre, et celebrándo profícere. Per Dóminum.

\subsubsection*{\center \textsc{IN PRIMO NOCTURNO.}}
\fancyhead[C]{\textsc{Officia Votiva - Feria III - De Apostolis.}}
\begin{tightlist}{}
\item \textred{Lectiones de Scriptura occurrente. In Quadragesima vero:} Sic nos exístimet homo, \textred{de Communi.} [8]
\end{tightlist}

\subsubsection*{\center \textsc{IN SECUNDO NOCTURNO.}}

\subsubsection*{\center De Sermóne sancti Augustíni Epíscopi.}
\subsubsection*{\center \textred{Serm. 43. De Sanctis.\\
Lectio iv.}}
\lettrineThree{I}sti sunt lux mundi~: quia per ipsos lumen fídei, et veræ sciéntæ primum Dóminus huic mundo trádidit, et ab errórum ac peccatórum ténebris gentes et pópulos éruit. Isti sunt sal terræ~: quia per eos terrígenæ condiméntum sapóris vitæ ætérnæ percepérunt, ut restríngerent carnis lascíviam, et a putrédine peccatórum et vitiórum vérmibus servaréntur illǽsi. Isti sunt lápides pretiósi, quos in fundaménto cœléstis ædifícii pósitos Joánnes in Apocalýpsi sua descríbit~: quia prædicátio eórum Ecclésiæ fundaménta locávit. Unde et Paulus ait~: Vos estis cives Sanctórum et doméstici Dei, superædificáti super fundaméntum Apostolórum et Prophetárum.

\begin{tightlist}{}
\item \textred{℟.}~Vidi conjúnctos viros habéntes spléndidas vestes~: et Angelus Dómini locútus est ad me dicens~: \textred{*} Isti sunt viri sancti, facti amíci Dei.
\textred{℣.}~Vidi Angelum Dei fortem volántem per médium cœli, voce magna clamántem et dicéntem. \textred{I}sti.
\end{tightlist}

\subsubsection*{\center \textred{Lectio v.}}
\lettrine{H}oc ergo considerántes, fratres caríssimi, quæ tanti duces pópuli nos docuérunt, studeámus factis implére, quæ illi jussérunt. Discámus eórum exémplo mundi divítias contémnere, præséntis sǽculi voluptátem non amáre, regnum cœléste desideráre, Christo nihil præpónere, sed ejus mandátis in ómnibus obedíre, paupertátem præséntium rerum dilígere, virtútum divítias habére, sapiéntiæ thesáurum appétere, spiritáles delícias quǽrere, nulli invidére, sed omnes hómines dilígere, amícos in Deo et inimícos propter Deum; quia hæc est vera et sola diléctio.

\begin{tightlist}{}
\item \textred{℟.}~Isti sunt triumphatóres, et amíci Dei, qui contemnéntes jussa príncipum, meruérunt prǽmia ætérna~: \textred{*} Modo coronántur, et accípiunt palmam.
\textred{℣.}~Isti sunt, qui venérunt ex magna tribulatióne, et lavérunt stolas suas in sánguine Agni. \textred{M}odo
\end{tightlist}

\subsubsection*{\center \textred{Lectio vj.}}
\lettrine{I}sti ergo Príncipes nostri, in amóre Dei perfectíssimi, et in proximórum dilectióne impléti, hinc potuérunt mundi ímpetum víncere, et cruéntum sǽculum domáre, quia nihil amavérunt in omni re, nisi Dei voluntátem. Ita et nos, fratres, diligámus in ómnibus rebus Dei fácere voluntátem, et creatórem nostrum amémus in semetípso, creatúram vero in conditóre suo, et sic habébimus ordinatíssimam caritátem~: quia Deus cáritas est, et qui díligit hanc dilectiónem, díligit Deum. Et si ita dilexérimus, díligit nos ipse Deus, díligunt nos et sancti Apóstoli júdices nostri, et exórant pro nobis, ut in universáli judício Christi cum ipsis in perpétuum coronémur.

\begin{tightlist}{}
\item \textred{℟.}~Fuérunt sine queréla ante Dóminum, et ab ínvicem non sunt separáti~: \textred{*} Cálicem Dómini bibérunt, et amíci Dei facti sunt.
\item \textred{℣.}~Tradidérunt córpora sua propter Deum ad supplícia~: ídeo coronántur, et accípiunt palmam. \textred{C}álicem.
\end{tightlist}

\subsubsection*{\center \textsc{IN TERTIO NOCTURNO.}}
\subsubsection*{\center Léctio sancti Evangélii secúndum Matthǽum.}
\subsubsection*{\center \textred{Lectio vij. \hfill Cap. 19.}}
\lettrine{I}n illo témpore~: Dixit Petrus ad Jesum~: Ecce nos relíquimus ómnia, et secúti sumus te~: quid ergo erit nobis? Et réliqua.

\begin{center}Homilía sancti Hilárii Epíscopi.\\
\textred{Comm. in Matth. cap. 20.}\end{center}
\lettrineThree{A}póstoli dictis Dómini hæc reddunt, reliquísse se ómnia, et cum ipso esse. Quibus Dóminus, cum séderit in majestátis suæ sede, sessúros super sedes duódecim, ac tótidem tribus Israel judicatúros spopóndit; omnibúsque qui univérsa relíquerint propter nomen ejus, fructum céntupli prǽmii reservátum~: multos autem ex novíssimis primos, et ex primis futúros novíssimos. Multa sunt quæ non sinunt nos símplici intelléctu dicta Evangélica suscípere.

\begin{tightlist}{}
\item \textred{℟.}~Isti sunt, qui vivéntes in carne, plantavérunt Ecclésiam sánguine suo. Non sunt de terris córpora eórum separáta~: \textred{*} Quorum mérita sunt in cælo glorióse cumuláta. \textred{℣.}~Sancti per fidem vicérunt regna, operáti sunt justítiam, adépti sunt repromissiónes. \textred{Q}uorum.
\end{tightlist}

\subsubsection*{\center \textred{Lectio viij.}}
\lettrine{I}nterpósitis nonnúllis rebus, quæ ex natúra humáni sensus sibi contráriæ sunt, ratiónem quǽrere cœléstis intelligéntiæ admonémur. Apóstoli dicunt, et sequi se Christum, et se ómnia reliquísse. Quómodo ígitur fiunt tristes, et quómodo métuunt dicéntes, salvum esse néminem posse? Namque et ab áliis fíeri póterat, si quid fecíssent ipsi. Deínde cum fecíssent ipsi; quare metus, vel unde suscéptus est? Additur étiam in responsióne Dómini, hæc apud hómines impossibília, possibília apud Deum. Numquid apud hóminem impossibília erant, quæ et Apóstoli fecísse se gloriántur, et fecísse eos Dóminus agnóscit?

\begin{tightlist}{}
\item \textred{℟.}~Isti sunt viri sancti, quos elégit Dóminus in caritáte non ficta, et dedit illis glóriam sempitérnam~: \textred{*} Quorum doctrína fulget Ecclésia, ut sole luna. \textred{℣.}~In omnem terram exívit sonus eórum, et in fines orbis terræ verba eórum. \textred{Q}uorum.
\end{tightlist}

\subsubsection*{\center \textred{Lectio ix.}}
\lettrine{O}mnis ítaque hic sermo est spiritális. Quid autem Dei virtúti tam possíbile est, quam ut per fidem salvet? ut per aquam regéneret? ut per crucem vincat? ut per Evangélia adóptet? ut per resurrectiónem ex morte vivíficet? Quibus audítis Apóstoli cito credunt, seséque ómnia reliquísse profiténtur. Sed hanc eórum obediéntiam cito Dóminus remunerátur, omnem difficultátem superióris quæstiónis absólvens, cum dicit~: Vos, qui secúti estis me, in regeneratióne, judicábitis duódecim tribus Israel. Hæc enim illa regenerátio est, quam Apóstoli sunt secúti, quam lex indulgére non pótuit, quæ eos super duódecim thronos in judicándis duódecim tríbubus Israel, in duódecim Patriarchárum glória copulávit.

\begin{tightlist}{}
\item \textred{℟.}~Cives Apostolórum, et doméstici Dei advenérunt hódie, \textred{*} Portántes facem, et illuminántes pátriam, dare pacem Géntibus, et liberáre pópulum Dómini. \textred{℣.}~Audíte preces súpplicum, vitæ ætérnæ poscéntes prǽmia, qui fertis in dextris manípulos justítiæ, quique gaudéntes advenístis hódie. \textred{P}ortántes. Glória Patri. \textred{P}ortántes.
\item Te Deum laudámus. 6.
\item \textred{Sacerdotalis ℣.}~Dedísti hereditátem. \textred{℟.}~Timéntibus nomen tuum Dómine.
\end{tightlist}

\begin{center}\textcolor{red}{\rule{\linewidth}{0.5pt}} \end{center}

\subsection*{\center \textred{PRO FERIA IV.}}
\subsection*{\center \textsc{\textred{OFFICIUM VOTIVUM}\\ De S. Joseph,\\ Sponso B.M.V.}}
\subsubsection*{\center \textred{Confess. et Catholicæ Ecclesiæ Patrono ac Protectore Primario Ordinis Nostri.\\ Semiduplex.}}
\fancyhead[C]{\textsc{Officia Votiva - Feria IV - De S. Joseph.}}

\subsubsection*{\center \textsc{AD VESPERAS.}}
\subsubsection*{\textred{Capitulum. \hfill Gen. 49.}}
\lettrineThree{B}enedictiónes patris tui confortátæ sunt benedictiónibus patrum ejus, donec veníret desidérium cóllium æternórum~: fiant in cápite Joseph, et in vértice Nazarǽi inter fratres suos.

\subsubsection*{\center \textred{Hymnus.}}
\lettrine{T}e, Joseph, célebrent ágmina cœ́litum~:\\
Te cuncti résonent Christíadum chori,\\
Qui clarus méritis, junctus es ínclytæ\\
Casto fœ́dere Vírgini.\\
\\
\textred{A}lmo cum túmidam gérmine cónjugem\\
Admírans, dúbio tángeris ánxius,\\
Afflátu súperi Fláminis Angelus\\
Concéptum púerum docet.\\
\\
\textred{T}u natum Dóminum stringis, ad éxteras\\
Ægýpti prófugum tu séqueris plagas;\\
Amíssum Sólymis, quæris, et ínvenis,\\
Miscens gáudia flétibus.\\
\\
\textred{P}ost mortem réliquos mors pia cónsecrat,\\
Palmámque eméritos glória súscipit~:\\
Tu vivens, Súperis par, frúeris Deo,\\
Mira sorte beátior.\\
\\
\textred{N}obis, summa Trias, parce precántibus,\\
Da Joseph méritis sídera scándere~:\\
Ut tandem líceat nos tibi pérpetim\\
Gratum prómere cánticum.\\
Amen.

\begin{tightlist}{}
\item \textred{℣.}~Constítuit eum dóminum domus suæ. \textred{℟.}~Et príncipem omnis possessiónis suæ.
\item \textred{Ad Magnif. Aña.} Cum esset desponsáta Mater Jesu María Joseph, ántequam convenírent, invénta est in útero habens de Spíritu sancto.
\end{tightlist}

\subsubsection*{\center \textred{Oratio.}}
\lettrine{D}eus, qui ineffábili providéntia beátum Joseph sanctíssimæ Genitrícis tuæ sponsum elígere dignátus es~: præsta, quǽsumus; ut quem protectórem venerámur in terris, intercessórem habére mereámur in cœlis~: Qui vivis et regnas.

\subsubsection*{\center \textsc{AD MATUTINUM.}}
\begin{tightlist}{}
\item \textred{Invit.} Laudémus Deum nostrum \textred{*} In veneratióne beáti Joseph, protectóris nostri.
\item \textred{Ps.} Veníte, exsultémus. 2.
\item \textred{Hymnus.} Te, Joseph, \textred{ut supra.}
\end{tightlist}

\subsubsection*{\center \textsc{IN PRIMO NOCTURNO.}}
\begin{tightlist}{}
\item \textred{Aña.} Ascéndit Joseph a Galilǽa de civitáte Názareth in Judǽam, in civitátem David, quæ vocátur Béthlehem, ut profiterétur cum María.
\item \textred{Ps.} Beátus vir, \textred{cum reliquis de Communi unius Conf. Pont.} [76]
\item \textred{Aña.} Venérunt pastóres festinántes~: et invenérunt Maríam, et Joseph, et infántem pósitum in præsépio.
\item \textred{Aña.} Ecce Angelus Dómini appáruit in somnis Joseph, dicens~: Surge, et áccipe púerum, et matrem ejus, et fuge in Ægýptum.
\item \textred{℣.}~Constítuit eum Dóminum domus suæ. \textred{℟.}~Et príncipem omnis possessiónis suæ.
\item \textred{Lectiones de Scriptura occurrente; si vero desint, legantur sequentes.}
\end{tightlist}

\subsubsection*{\center De libro Génesis.}
\subsubsection*{\center \textred{Lectio j. \hfill Cap. 39.}}
\lettrineThree{I}gitur Joseph ductus est in Ægýptum, emítque eum Pútiphar eunúchus Pharaónis, princeps exércitus, vir Ægýptius, de manu Ismaëlitárum, a quibus perdúctus erat.
Fuítque Dóminus cum eo, et erat vir in cunctis próspere agens~: habitavítque in domo dómini sui,
qui óptime nóverat Dóminum esse cum eo, et ómnia quæ géreret, ab eo dírigi in manu illíus.
Invenítque Joseph grátiam coram dómino suo, et ministrábat ei~: a quo præpósitus ómnibus, gubernábat créditam sibi domum, et univérsa quæ ei trádita fúerant~:
benedixítque Dóminus dómui Ægýptii propter Joseph, et multiplicávit tam in ǽdibus quam in agris cunctam ejus substántiam~:
nec quidquam áliud nóverat, nisi panem quo vescebátur. Erat autem Joseph pulchra fácie, et decórus aspéctu.

\begin{tightlist}{}
\item \textred{℟.}~Clamávit pópulus ad regem, aliménta petens~: \textred{*} Quibus ille respóndit~: Ite ad Joseph. \textred{℣.}~Salus nostra in manu tua est~: réspice nos tantum, et læti serviémus regi. \textred{Q}uibus.
\end{tightlist}

\subsubsection*{\center \textred{Lectio ij. \hfill Cap. 41.}}
\lettrine{P}lácuit Pharaóni consílium et cunctis minístris ejus~:
locutúsque est ad eos~: Num inveníre potérimus talem virum, qui spíritu Dei plenus sit?
Dixit ergo ad Joseph~: Quia osténdit tibi Deus ómnia quæ locútus es, numquid sapientiórem et consímilem tui inveníre pótero?
Tu eris super domum meam, et ad tui oris impérium cunctus pópulus obédiet~: uno tantum regni sólio te præcédam.
Dixítque rursus Phárao ad Joseph~: Ecce, constítui te super univérsam terram Ægýpti.
Tulítque ánnulum de manu sua, et dedit eum in manu ejus~: vestivítque eum stola býssina, et collo torquem áuream circumpósuit.
Fecítque eum ascéndere super currum suum secúndum, clamánte præcóne, ut omnes coram eo genu flécterent, et præpósitum esse scirent univérsæ terræ Ægýpti.

\begin{tightlist}{}
\item \textred{℟.}~Fecit me Deus quasi patrem regis, et dóminum univérsæ domus ejus~: \textred{*} Exaltávit me, ut salvos fáceret multos pópulos. \textred{℣.}~Veníte ad me, et ego dabo vobis ómnia bona Ægýpti, ut comedátis medúllam terræ. \textred{E}xaltávit.
\end{tightlist}

\subsubsection*{\center \textred{Lectio iij.}}
\lettrine{D}ixit quoque rex ad Joseph~: Ego sum Phárao~: absque tuo império non movébit quisquam manum aut pedem in omni terra Ægýpti.
Vertítque nomen ejus, et vocávit eum, lingua ægyptíaca, Salvatórem mundi. Dedítque illi uxórem Aseneth, fíliam Pútiphare sacerdótis Heliopóleos. Egréssus est ítaque Joseph ad terram Ægýpti
(trigínta autem annórum erat quando stetit in conspéctu regis Pharaónis), et circuívit omnes regiónes Ægýpti.
Venítque fertílitas septem annórum~: et in manípulos redáctæ ségetes congregátæ sunt in hórrea Ægýpti.
Omnis étiam frugum abundántia in síngulis úrbibus cóndita est.
Tántaque fuit abundántia trítici, ut arénæ maris coæquarétur, et cópia mensúram excéderet.

\begin{tightlist}{}
\item \textred{℟.}~Jam lætus móriar quia vidi fáciem tuam, et supérstitem te relínquo. Non sum fraudátus adspéctu tuo~: \textred{*} Insuper osténdit mihi Dóminus semen tuum. \textred{℣.}~Qui pascit me ab adolescéntia mea, benedícat púeris istis, et invocétur super eos nomen meum. \textred{I}nsuper. Glória Patri. \textred{I}nsuper.
\end{tightlist}

\subsubsection*{\center \textsc{IN SECUNDO NOCTURNO.}}
\begin{tightlist}{}
\item \textred{Aña.} Consúrgens Joseph, accépit púerum et matrem ejus nocte, et secéssit in Ægýptum~: et erat ibi usque ad óbitum Heródis.
%\item \textred{Ps.} Cum invocárem. [67]
\item \textred{Aña.} Defúncto Heróde, Angelus Dómini appáruit in somnis Joseph in Ægýpto, dicens~: Surge, et áccipe púerum et matrem ejus, et vade in terram Israel~: defúncti sunt enim qui quærébant ánimam púeri.
\item \textred{Aña.} Accépit Joseph púerum et matrem ejus, et venit in terram Israel.
\item \textred{℣.}~De ventre matris meæ tu es protéctor meus. \textred{℟.}~In te confirmátus sum ex útero.
\end{tightlist}

\subsubsection*{\center Sermo sancti Bernárdi Abbátis.}
\subsubsection*{\center \textred{Homilia 2. super Missus.\\
Lectio iv.}}
\lettrineThree{D}esponsáta est María Joseph, vel pótius (sicut ponit Evangelísta) viro, cui nomen erat Joseph~: virum nóminat, non quia marítus, sed quod homo virtútis erat~: vel pótius quia juxta álium Evangelístam, non vir simplíciter, sed vir ejus dictus est, mérito appellátur, quod necessário putátur. Débuit ergo vir ejus appellári, quia necésse fuit et putári~: sicut et pater Salvatóris non quidem esse, sed dici méruit, ut putarétur esse, dicénte hoc ipso Evangelísta~: Et ipse Jesus erat incípiens quasi annórum trigínta, ut putabátur, fílius Joseph.

\begin{tightlist}{}
\item \textred{℟.}~Dedísti mihi protectiónem salútis tuæ, et déxtera tua suscépit me~: \textred{*} Protéctor meus, et cornu salútis meæ, et suscéptor meus. \textred{℣.}~Ego protéctor tuus sum, et merces tua magna nimis. \textred{P}rotéctor.
\end{tightlist}

\subsubsection*{\center \textred{Lectio v.}}
\lettrine{N}on est dúbium, quin bonus et fidélis homo fúerit iste Joseph, cui Mater desponsáta est Salvatóris. Fidélis, inquam, servus et prudens, quem constítuit Dóminus suæ Matris solátium, suæ carnis nutrítium, solum dénique in terris magni consílii coadjutórem fidíssimum. Huc accédit, quod dícitur fuísse de domo David. Vere enim de domo David. Vere de régia stirpe descéndit vir iste Joseph, nóbilis génere, mente nobílior. Plane fílius David, non degénerans a patre suo David; prorsus, inquam, fílius David, non tantum carne, sed fide, sed sanctitáte, sed devotióne, quem tamquam álterum David Dóminus invénit secúndum cor suum, cui tuto commítteret sacratíssimum atque secretíssimum sui cordis arcánum~: cui, tamquam álteri David, incérta et occúlta sapiéntiæ suæ manifestávit, et dedit illi non ignárum esse mystérii, quod nemo príncipum hujus sǽculi agnóvit.

\begin{tightlist}{}
\item \textred{℟.}~Státuet fílios suos sub tégmine illíus, et sub ramis ejus morábitur~: protegétur sub tégmine illíus a fervóre~: \textred{*} Et in glória ejus requiéscet. \textred{℣.}~Speráte in eo, omnis congregátio pópuli~: effúndite coram illo corda vestra. \textred{E}t.
\end{tightlist}

\subsubsection*{\center \textred{Lectio vj.}}
\lettrine{C}ui dénique datum est, quod multi reges et prophétæ cum vellent vidére, non vidérunt, audíre et non audiérunt, non solum vidére et audíre, sed étiam portáre, dedúcere, amplécti, deosculári, et nutríre et custodíre. Non tantum autem Joseph, sed et María descendísse credénda est de domo David~: alióquin non esset desponsáta viro de domo David, si non esset et ipsa de domo David. Ambo ígitur erant de domo David~: sed in áltera compléta est véritas, quam jurávit Dóminus David, áltero tamen cónscio et teste adimplétæ promissiónis.

\begin{tightlist}{}
\item \textred{℟.}~Si consístant advérsum me castra, non timébit cor meum~: \textred{*} Si exsúrgat advérsum me prǽlium, in hoc ego sperábo. \textred{℣.}~In te cantátio mea semper, quóniam tu adjútor fortis. \textred{S}i. Glória Patri. \textred{S}i.
\end{tightlist}

\subsubsection*{\center \textsc{IN TERTIO NOCTURNO.}}
\begin{tightlist}{}
\item \textred{Aña.} Audiens Joseph quod Archeláus regnáret in Judǽa pro Heróde, patre suo, tímuit illo ire.
%\item \textred{Ps.} Dómine, quis habitábit. [70]
\item \textred{Aña.} Admónitus in somnis, Joseph secéssit in partes Galilǽæ, et véniens habitávit in civitáte quæ vocátur Názareth; ut adimplerétur quod dictum est per Prophétas~: Quóniam Nazarǽus vocábitur.
\item \textred{Aña.} Erat pater Jesu et mater mirántes super his quæ dicebántur de illo; et benedíxit illis Símeon.
\item \textred{℣.}~Invocávi Dóminum, Patrem Dómini mei. \textred{℟.}~Ut non derelínquat me in die tribulatiónis.
\end{tightlist}

\subsubsection*{\center Lectio sancti Evangélii secúndum Lucam.}
\subsubsection*{\center \textred{Lectio vij. \hfill Cap. 3.}}
\lettrine{I}n illo témpore~: Factum est autem, cum baptizarétur omnis pópulus et Jesu baptizáto et oránte, apértum est cœlum. Et réliqua.

\begin{center}Homilía sancti Augustíni Epíscopi.\\
\textred{Lib. 23. contra Faust. cap. 7. 8.}\end{center}
\lettrineThree{S}ic de cœlo dictum est super aquam Jordánis~: Hic est Fílius meus diléctus, in quo mihi bene complácui~: quemádmodum dictum est et in monte. Neque enim quia et ibi de cœlo vox ipsa sónuit, Fílius Dei ante non fuit~: quandóquidem ex útero Vírginis ille accépit formam servi, qui cum in forma Dei esset, non rapínam arbitrátus est, esse se æquálem Deo. Dénique idem Apóstolus Paulus álio loco apertíssime dicit~: Cum autem venit plenitúdo témporis, misit Deus Fílium suum factum ex mulíere, factum sub lege, ut eos, qui sub lege erant, redímeret, ut filiórum adoptiónem reciperémus. Ipse ergo est Fílius Dei, qui et Dóminus David secúndum divinitátem, et idem ipse Fílius David ex sémine David secúndum carnem.

\begin{tightlist}{}
\item \textred{℟.}~Joseph, fili David, noli timére accípere Maríam, cónjugem tuam; quod enim in ea natum est, de Spíritu sancto est~: páriet autem fílium, \textred{*} Et vocábis nomen ejus Jesum. \textred{℣.}~Ipse enim salvum fáciet pópulum suum a peccátis eórum. \textred{E}t.
\end{tightlist}

\subsubsection*{\center \textred{Lectio viij.}}
\lettrine{Q}uod si nobis crédere non prodésset, non hoc tam atténte idem Apóstolus Timótheo commendáret dicens~: Memor esto Christum Jesum resurrexísse a mórtuis, ex sémine David, secúndum Evangélium meum. Quid ergo jam móveat sancti Evangélii sectatórem, quod sine concúbitu Joseph Christus natus ex Vírgine, fílius tamen David appellátur, cum generatiónum sériem non usque ad Maríam, sed usque ad Joseph Matthǽus Evangelísta perdúcat? Primo quia maríti ejus fúerat propter virílem sexum pótius honoránda persóna~: neque enim quia concúbitu non permíxtus, ídeo non marítus, cum ipse Matthǽus narret ab Angelo Maríam cónjugem ipsíus appellátam, qui narrat, quod de Spíritu sancto concéperat.

\begin{tightlist}{}
\item \textred{℟.}~Surge, et áccipe púerum et matrem ejus, et fuge in Ægýptum~: \textred{*} Et esto ibi, usque dum dicam tibi. \textred{℣.}~Ut adimplerétur quod dictum est a Dómino per Prophétam dicéntem~: Ex Ægýpto vocávi fílium meum. \textred{E}t.
\end{tightlist}

\subsubsection*{\center \textred{Lectio ix.}}
\lettrine{C}um vero unus idémque narrátor utrúmque dicat, utrúmque comméndet, et virum Maríæ Joseph, et Christi Vírginem Matrem, et Christum ex sémine David, et Joseph in série progeneratórum Christi ex David~: quid restat, nisi et Maríam non fuísse extráneam a cognatióne David, et eam Joseph cónjugem non frustra appellátam propter órdinem sexus et animórum confœderatiónem; et Joseph pótius propter dignitátem virílem ab órdine generatiónum illárum non fuísse separándum, ne hoc ipso viderétur ab illa fémina separátus, cui eum conjungébat mentis afféctus?

\begin{tightlist}{}
\item \textred{℟.}~Valde honorándus est, et laudábilis nimis beátus Joseph, \textred{*} Quem constítuit Dóminus magni consílii coadjutórem fidíssimum. \textred{℣.}~Fidélis servus et prudens, Vírginis Matris Sponsus, et Fílii Dei Nutrítius. \textred{Q}uem. Glória Patri. \textred{Q}uem.
\item Te Deum laudámus. 6.
\item \textred{Sacerdotalis ℣.}~Ora pro nobis sancte Joseph. \textred{℟.}~Ut digni efficiámur promissiónibus Christi.
\end{tightlist}

\subsubsection*{\center \textsc{AD LAUDES.}}
\begin{center}\textred{et per Horas, Aña.}\end{center}
\lettrine{J}acob autem génuit Joseph, virum Maríæ, de qua natus est Jesus, qui vocátur Christus.

\begin{tightlist}{}
\item \textred{Ps. de Dominica.} 27.
\item \textred{Aña.} Missus est Angelus Gábriel a Deo in civitátem Galilǽæ, cui nomen Názareth, ad Vírginem desponsátam viro, cui nomen erat Joseph.
\item \textred{Aña.} Ascéndit autem Joseph a Galilǽa de civitáte Názareth, in Judǽam, in civitátem David, quæ vocátur Béthlehem.
\item \textred{Aña.} Et venérunt festinántes, et invenérunt Maríam, et Joseph, et infántem pósitum in præsépio.
\item \textred{Aña.} Et ipse Jesus erat incípiens quasi annórum trigínta, ut putabátur fílius Joseph.
\end{tightlist}

\subsubsection*{\textred{Capitulum. \hfill Gen. 49.}}
\lettrineThree{B}enedictiónes patris tui confortátæ sunt benedictiónibus patrum ejus, donec veníret desidérium cóllium æternórum~: fiant in cápite Joseph, et in vértice Nazarǽi inter fratres suos.

\subsubsection*{\center \textred{Hymnus.}}
\lettrine{C}œ́litum, Joseph, decus, atque nostræ\\
Certa spes vitæ, columénque mundi,\\
Quas tibi læti cánimus, benígnus\\
Súscipe laudes.\\
\\
\textred{T}e Sator rerum státuit pudícæ\\
Vírginis sponsum, voluítque Verbi\\
Te patrem dici, dedit et minístrum\\
Esse salútis.\\
\\
\textred{T}u Redemptórem stábulo jacéntem,\\
Quem chorus Vatum cécinit futúrum,\\
Adspicis gaudens, humilísque natum\\
Numen adóras.\\
\\
\textred{R}ex Deus regum, Dominátor orbis,\\
Cujus ad nutum tremit inferórum\\
Turba, cui pronus famulátur æther,\\
Se tibi subdit.\\
\\
\textred{L}aus sit excélsæ Tríadi perénnis,\\
Quæ tibi præbens súperos honóres,\\
Det tuis nobis méritis beátæ\\
Gáudia vitæ.\\
Amen.

\begin{tightlist}{}
\item \textred{℣.}~Dedísti mihi protectiónem salútis tuæ. \textred{℟.}~Et déxtera tua suscépit me.
\item \textred{Ad Bened. Aña.} Joseph, fili David, noli timére accípere Maríam cónjugem tuam; quod enim in ea natum est, de Spíritu sancto est.
\end{tightlist}

\subsubsection*{\center \textred{Oratio.}}
\lettrine{D}eus, qui ineffábili providéntia beátum Joseph sanctíssimæ Genitrícis tuæ sponsum elígere dignátus es~: præsta, quǽsumus; ut quem protectórem venerámur in terris, intercessórem habére mereámur in cœlis~: Qui vivis.

\begin{tightlist}{}
\item \textred{Ad Horas Añæ de Laudibus, Psalmi de Dominica, ad Primam tamen ut in Festis.}
\end{tightlist}

\subsubsection*{\center \textsc{AD TERTIAM.}}
\subsubsection*{\textred{Capitulum. \hfill Gen. 49.}}
\lettrine{B}enedictiónes patris tui confortátæ sunt benedictiónibus patrum ejus, donec veníret desidérium cóllium æternórum~: fiant in cápite Joseph, et in vértice Nazarǽi inter fratres suos.

\begin{tightlist}{}
\item \textred{℟.}~Constítuit eum \textred{*} Dóminum domus suæ. \textred{C}onstítuit. \textred{℣.}~Et príncipem omnis possessiónis suæ. \textred{D}óminum. Glória Patri. \textred{C}onstítuit.
\item \textred{℣.}~De ventre matris meæ, tu es protéctor meus. \textred{℟.}~In te confirmátus sum ex útero.
\end{tightlist}

\subsubsection*{\center \textsc{AD SEXTAM.}}
\subsubsection*{\textred{Capitulum. \hfill Gen. 49.}}
\lettrine{D}eus patris tui erit adjútor tuus, et Omnípotens benedícet tibi benedictiónibus cœli désuper.

\begin{tightlist}{}
\item \textred{℟.}~De ventre matris meæ, \textred{*} Tu es protéctor meus. \textred{D}e. \textred{℣.}~In te confirmátus sum ex útero. \textred{T}u. Glória Patri. \textred{D}e.
\item \textred{℣.}~Invocávi Dóminum Patrem Dómini mei. \textred{℟.}~Ut non derelínquat me in die tribulatiónis.
\end{tightlist}

\subsubsection*{\center \textsc{AD NONAM.}}
\subsubsection*{\textred{Capitulum. \hfill Gen. 49.}}
\lettrine{F}ilius accréscens Joseph, fílius accréscens, et decórus adspéctu~: fíliæ discurrérunt super murum.

\begin{tightlist}{}
\item \textred{℟.}~Invocávi Dóminum \textred{*} Patrem Dómini mei. \textred{I}nvocávi. \textred{℣.}~Ut non derelínquat me in die tribulatiónis. \textred{P}atrem. Glória Patri. \textred{I}nvocávi.
\item \textred{℣.}~Dedísti mihi protectiónem salútis tuæ. \textred{℟.}~Et déxtera tua suscépit me.
\end{tightlist}

\subsubsection*{\center \textsc{AD II. VESPERAS.}}
\begin{tightlist}{}
\item \textred{Prima Aña Laudum.} Jacob autem génuit Joseph, virum Maríæ, de qua natus est Jesus, qui vocátur Christus. 
\item \textred{Ps. de j. Vesp. Comm. Apost. [2]}
%\item \textred{Capit. Hymnus et ℣. ut in j. Vesperis.}
\end{tightlist}

\subsubsection*{\textred{Capitulum. \hfill Gen. 49.}}
\lettrineThree{B}enedictiónes patris tui confortátæ sunt benedictiónibus patrum ejus, donec veníret desidérium cóllium æternórum~: fiant in cápite Joseph, et in vértice Nazarǽi inter fratres suos.

\subsubsection*{\center \textred{Hymnus.}}
\lettrine{T}e, Joseph, célebrent ágmina cœ́litum~:\\
Te cuncti résonent Christíadum chori,\\
Qui clarus méritis, junctus es ínclytæ\\
Casto fœ́dere Vírgini.\\
\\
\textred{A}lmo cum túmidam gérmine cónjugem\\
Admírans, dúbio tángeris ánxius,\\
Afflátu súperi Fláminis Angelus\\
Concéptum púerum docet.\\
\\
\textred{T}u natum Dóminum stringis, ad éxteras\\
Ægýpti prófugum tu séqueris plagas;\\
Amíssum Sólymis, quæris, et ínvenis,\\
Miscens gáudia flétibus.\\
\\
\textred{P}ost mortem réliquos mors pia cónsecrat,\\
Palmámque eméritos glória súscipit~:\\
Tu vivens, Súperis par, frúeris Deo,\\
Mira sorte beátior.\\
\\
\textred{N}obis, summa Trias, parce precántibus,\\
Da Joseph méritis sídera scándere~:\\
Ut tandem líceat nos tibi pérpetim\\
Gratum prómere cánticum.\\
Amen.

\begin{tightlist}{}
\item \textred{℣.}~Constítuit eum dóminum domus suæ. \textred{℟.}~Et príncipem omnis possessiónis suæ.
\item \textred{Ad Magnif. Aña.} Fili, quid fecísti nobis sic? Ecce pater tuus et ego doléntes quærebámus te.
\end{tightlist}

\begin{center}\textcolor{red}{\rule{\linewidth}{0.5pt}} \end{center}

\subsection*{\center \textred{PRO FERIA VI.}}
\subsection*{\center \textsc{\textred{OFFICIUM VOTIVUM}\\ De Passione D.N.J.C.}}
\subsubsection*{\center \textred{Semiduplex.}}
\fancyhead[C]{\textsc{Officia Votiva - Feria VI - De Passione D.N.J.C.}}

\subsubsection*{\center \textsc{AD I. VESPERAS.}}
\subsubsection*{\textred{Capitulum. \hfill Philipp. 2.}}
\lettrineThree{F}ratres~: Hoc enim sentíte in vobis, quod et in Christo Jesu~: qui, cum in forma Dei esset, non rapínam arbitrátus est esse se æquálem Deo~: sed semetípsum exinanívit, formam servi accípiens, in similitúdinem hóminum factus, et hábitu invéntus ut homo. Humiliávit semetípsum, factus obédiens usque ad mortem, mortem autem crucis.

\subsubsection*{\center \textred{Hymnus.}}
\lettrine{M}œréntes óculi spárgite lácrymas,\\
Et luctu résonent íntima córdium,\\
Illátas réfero Númini ab ímpiis\\
Pœnas, et fera vúlnera.\\
\\
\textred{A}ccíncta, heu! gládiis turba satéllitum\\
Arréptum Dóminum fústibus ímpetit,\\
Nunc cædit cólaphis, nunc quatit hórridis\\
Divínum caput íctibus.\\
\\
\textred{H}aud finis scéleri~: tráditur ímprobo\\
Christus carnífici; nec mora, bárbarus\\
In Regem súperum non timet ímpio\\
Ausu vértere déxteram.\\
\\
\textred{A}udíte, o pópuli, Numen amábile,\\
Manánte ex húmeris úndique sánguine,\\
Lictóris rábidi sústinet ímpetum,\\
Et vocem premit ínnocens.\\
\\
\textred{Q}uis non illácrimet? Jam nova cóndidit\\
Torménta indómitæ gentis iníquitas;\\
Infígit cérebro, proh dolor! éfferis\\
Sertum vépribus ásperum.\\
\\
\textred{H}eu crimen! tráhitur fúnibus ímprobis\\
Funésti Dóminus supplícii ad locum;\\
Illic oppétiit fúnere, spíritum\\
Patri restítuens suo.\\
\\
\textred{P}asso pro míseris tam fera vúlnera\\
In terris résonet débita glória;\\
Sacrúmque assídue Nomen in ǽthera\\
Humánum genus éfferat.\\
Amen.

\begin{tightlist}{}
\item \textred{℣.}~Oblátus est quia ipse vóluit. \textred{℟.}~Et livóre ejus sanáti sumus.
\item \textred{Ad Magnificat, Aña.} O vos omnes, qui transítis per viam, atténdite, et vidéte si est dolor sicut dolor meus.
\end{tightlist}

\subsubsection*{\center \textred{Oratio.}}
\lettrine{D}omine Jesu Christe, qui de cœlis ad terram de sinu Patris descendísti et sánguinem tuum pretiósum in remissiónem peccatórum nostrórum fudísti~: te humíliter deprecámur; ut in die judícii, ad déxteram tuam audíre mereámur~: Veníte benedícti. Qui vivis et regnas.

\begin{tightlist}{}
\item \textred{Ad Complet. et per Horas, in fine Hymnorum dicitur:} ℣.~Glória tibi, Dómine, Qui nos salvásti in Cruce.
\item \textred{Completorium de Dominica.}
\item \textred{Ad Nunc dimittis, Aña.} Hoc signum Crucis * erit in cælo, cum Dóminus ad judicándum vénerit.
\end{tightlist}

\subsubsection*{\center \textsc{AD MATUTINUM.}}
\begin{tightlist}{}
\item \textred{Invit.} Christum Regem crucifíxum \textred{*} Veníte, adorémus.
\item \textred{Ps.} Veníte, exsultémus. 2. 
\end{tightlist}
\subsubsection*{\center\textred{Hymnus.}}
\lettrine{A}dspice, infámi Deus ipse ligno\\
Pendet, effúso mádidus cruóre;\\
Adspice, immíti manus alma clavo\\
Fínditur alte.\\
\\
\textred{H}unc, velut pravi scéleris minístrum,\\
Inter indígnos médium latrónes\\
Cerne~: crudélis fuit ista genis\\
Dira volúntas.\\
\\
\textred{P}allet heu! vultus; caput ecce lassum\\
Fléctitur~: clausis óculis Redémptor\\
Spíritum sacro méritis onústum\\
Fundit ab ore.\\
\\
\textred{O} cor, æs durum súperas, geméndo\\
Ni scelus tergas; tua culpa Christum\\
Stípiti affíxit~: tua culpa morti\\
Súbdidit atræ.\\
\\
\textred{S}it Deo ætérnum decus omne in ævum,\\
Qui pius nostri géneris Redémptor,\\
Críminis labem míseris nocéntem\\
Sánguine tersit.\\
Amen.

\subsubsection*{\center \textsc{IN PRIMO NOCTURNO.}}
\begin{tightlist}{}
\item \textred{Aña.} Adstitérunt reges terræ, et príncipes convenérunt in unum advérsus Dóminum, et advérsus Christum ejus.
\item \textred{Ps. 2.} Quare fremuérunt. 22.
\item \textred{Aña.} Multiplicáti sunt, qui tríbulant me; multi insúrgunt advérsum me.
\item \textred{Ps. 3.} Dómine, quid. 22.
\item \textred{Aña.} Divisérunt sibi vestiménta mea, et super vestem meam misérunt sortem.
\item \textred{Ps. 21.} Deus, Deus meus. 230.
\item \textred{℣.}~Oblátus est, quia ipse vóluit. \textred{℟.}~Et livóre ejus sanáti sumus.
\item \textred{Lectiones de Scriptura occurr. si vero desint, legantur sequentes.}
\end{tightlist}

\subsubsection*{\center De Epístola beáti Pauli Apóstoli ad Romános.}
\subsubsection*{\center \textred{Lectio j. \hfill Cap.~5.~et~6.}}
\lettrineThree{C}omméndat caritátem suam Deus in nobis~: quóniam cum adhuc peccatóres essémus, secúndum tempus, Christus pro nobis mórtuus est~: multo ígitur magis nunc justificáti in sánguine ipsíus, salvi érimus ab ira per ipsum. Si enim cum inimíci essémus, reconciliáti sumus Deo per mortem Fílii ejus, multo magis reconciliáti, salvi érimus in vita ipsíus. Non solum autem, sed et gloriámur in Deo per Dóminum nostrum Jesum Christum, per quem nunc reconciliatiónem accépimus. Si enim, uníus delícto mors regnávit per unum~: multo magis abundántiam grátiæ, et donatiónis, et justítiæ accipiéntes, in vita regnábunt per unum Jesum Christum.

\begin{tightlist}{}
\item \textred{℟.}~Vínea mea elécta, ego te plantávi, dicit Dóminus, et tu facta es mihi nimis amára; \textred{*} Quia parásti crucem Salvatóri tuo. \textred{℣.}~Pópule meus, quid feci tibi, aut in quo contristávi te? respónde mihi. \textred{Q}uia.
\end{tightlist}

\subsubsection*{\center\textred{Lectio ij.}}
\lettrine{I}gitur, sicut per uníus delíctum in omnes hómines in condemnatiónem, sic et per uníus justítiam, in omnes hómines in justificatiónem vitæ. Sicut enim per inobediéntiam uníus hóminis, peccatóres constitúti sunt multi~: ita et per uníus obeditiónem, justi constituéntur multi. Lex autem subintrávit ut abundáret delíctum. Ubi autem abundávit delíctum, superabundávit grátia. Ut sicut regnávit peccátum in mortem, ita et grátia regnet per justítiam in vitam ætérnam, per Jesum Christum Dóminum nostrum.

\begin{tightlist}{}
\item \textred{℟.}~Ego edúxi te de Ægýpto, demérso Pharaóne in mare rubrum, et ante te præívi in colúmna nubis. \textred{*} Et tu me tradidísti princípibus sacerdótum, et me duxísti ad prætórium Piláti. \textred{℣.}~Pópule meus, quid feci tibi, aut in quo contristávi te? respónde mihi. \textred{E}t.
\end{tightlist}

\subsubsection*{\center\textred{Lectio iij.}}
\lettrine{Q}uid ergo dicémus? Permanébimus in peccáto ut grátia abúndet? Absit. Qui enim mórtui sumus peccáto, quómodo adhuc vivémus in illo? An ignorátis, quia quicúmque baptizáti sumus in Christo Jesu, in morte ipsíus baptizáti sumus? Consepúlti enim sumus cum illo per baptísmum in mortem~: ut quómodo Christus surréxit a mórtuis per glóriam Patris, ita et nos in novitáte vitæ ambulémus. Si enim complantáti facti sumus similitúdini mortis ejus; simul et resurrectiónis érimus. Hoc sciéntes, quia vetus homo noster simul crucifíxus est, ut destruátur corpus peccáti, et ultra non serviámus peccáto. Qui enim mórtuus est, justificátus est a peccáto. Si autem mórtui sumus cum Christo, crédimus quia simul étiam vivémus cum Christo.

\begin{tightlist}{}
\item \textred{℟.}~Ego propter te flagellávi Ægýptum cum primogénitis suis; \textred{*} Et tu me flagellátum tradidísti, qui, tamquam agnus, coram tondénte se, obmútui, et non apérui os meum. \textred{℣.}~Pópule meus, quid feci tibi, aut in quo contristávi te? respónde mihi. \textred{E}t.  Glória Patri. \textred{E}t.
\end{tightlist}

\subsubsection*{\center \textsc{IN SECUNDO NOCTURNO.}}
\begin{tightlist}{}
\item \textred{Aña.} Inimíci mei dixérunt mala mihi, quando moriétur, et períbit nomen ejus?
\item \textred{Ps. 40.} Beátus qui intélligit. 124.
\item \textred{Aña.} Aliéni insurrexérunt in me, et fortes quæsiérunt ánimam meam.
\item \textred{Ps. 53.} Deus, in nómine. 160.
\item \textred{Aña.} Conculcavérunt me inimíci mei tota die, advérsum me omnes cogitatiónes eórum in malum.
\item \textred{Ps. 55.} Miserére mei. 162.
\item \textred{℣.}~Próprio Fílio suo non pepércit Deus. \textred{℟.}~Sed pro ómnibus trádidit illum.
\end{tightlist}

\subsubsection*{\center Sermo sancti Bernárdi Abbátis.}
\subsubsection*{\center \textred{Feria 4. Hebdom. S. Serm. de Passione Domini. \hfill Lectio iv.}}
\lettrineThree{I}n passióne domínica, fratres, tria speciáliter cónvenit intuéri~: opus, modum, causam. Nam in ópere quidem patiéntia, in modo humílitas, in causa cáritas commendátur. Patiéntia autem singuláris, quod vidélicet~: cum supra dorsum ejus fabricárent peccatóres; cum sic extenderétur in ligno, ut dinumeraréntur ómnia ossa ejus; cum fortíssimum illud propugnáculum, quod custódit Israel, úndique forarétur; cum foderéntur manus ejus et pedes, sicut agnus ad occisiónem ductus sit et tamquam ovis coram tondénte non apéruit os suum. Non advérsus Patrem múrmurans, a quo missus fúerat; non advérsus humánum genus, pro quo, quæ non rápuit, exsolvébat; non dénique vel contra pópulum ipsum peculiárem sibi, a quo pro tantis benefíciis tanta mala recipiébat.

\begin{tightlist}{}
\item \textred{℟.}~Ego propter te Chananæórum reges percússi, et dedi tibi sceptrum regále; \textred{*} Et tu dedísti cápiti meo spíneam corónam, et percussísti arúndine caput meum. \textred{℣.}~Pópule meus, quid feci tibi, aut in quo contristávi te? respónde mihi. \textred{E}t.
\end{tightlist}

\subsubsection*{\center \textred{Lectio v.}}
\lettrine{M}odum ipsum, si diligénter atténdas, non modo mitem, sed et húmilem corde cognósces. Nempe in humilitáte judícium ejus sublátum est, cum nec ad tantas blasphémias, nec ad falsíssima, quæ sibi objiciebántur crímina respondéret. Vídimus, inquit, eum et non erat ei adspéctus, nec speciósum forma præ fíliis hóminum, sed oppróbrium hóminum et tamquam leprósum, novíssimum virórum, plane virum dolórum a Deo percússum et humiliátum; ita ut nulla esset ei spécies, neque decor. O novíssimum et altíssimum! O húmilem et sublímem! O oppróbrium hóminum et glóriam Angelórum! Nemo illo sublímior, neque humílior. Dénique sputis íllitus est, oppróbriis saturátus est, morte turpíssima condemnátus est, cum scelerátis deputátus est. Nihílque merébitur vel ista humílitas, quæ hunc habet modum; imo quæ tam est ultra modum? Sicut est patiéntia singuláris, sic et humílitas admirábilis; útraque sine exémplo.

\begin{tightlist}{}
\item \textred{℟.}~Ego duxi te per desértum quadragínta annis, et manna cibávi te; \textred{*} Et tu me cecidísti álapis, et flagéllis. \textred{℣.}~Pópule meus, quid feci tibi, aut in quo contristávi te? respónde mihi. \textred{E}t.
\end{tightlist}

\subsubsection*{\center \textred{Lectio vj.}}
\lettrine{U}trámque tamen magnífice causa ipsa comméndat, nimírum cáritas est. Propter nímiam enim caritátem suam, qua diléxit nos Deus, ut servum redímeret, nec Pater Fílio, nec sibi Fílius ipse pepércit. Vere nímiam, quia et hæc mensúram excédit, modum súperat, plane superéminens univérsis. Majórem, inquit, caritátem nemo habet, quam ut ánimam suam ponat quis pro amícis suis. Tu majórem habuísti, Dómine, ponens eam étiam pro inimícis. Cum enim adhuc inimíci essémus, per mortem tuam et tibi reconciliáti sumus et Patri. Quænam ergo ália vidébitur esse, vel fuísse, vel fore huic símilis caritáti? Vix pro justo quis móritur. Tu pro injústis passus es, móriens propter delícta nostra, qui venísti justificáre gratis peccatóres, servos fácere fratres, captívos coherédes, éxsules reges.

\begin{tightlist}{}
\item \textred{℟.}~Ego exaltávi te magna virtúte~: \textred{*} Et tu me suspendísti in patíbulo crucis, in qua expándi manus meas ad pópulum non credéntem et contradicéntem. \textred{℣.}~Pópule meus, quid feci tibi, aut in quo contristávi te? respónde mihi. \textred{E}t. Glória Patri. \textred{E}t.
\end{tightlist}

\subsubsection*{\center \textsc{IN TERTIO NOCTURNO.}}
\begin{tightlist}{}
\item \textred{Aña.} Fílii hóminum dentes eórum arma, et sagíttæ, et lingua eórum gládius acútus.
\item \textred{Ps. 56.} Miserére mei, Deus. 163.
\item \textred{Aña.} Intendérunt arcum rem amáram, ut sagíttent in occúltis immaculátum.
\item \textred{Ps. 63.} Exáudi, Deus. 263.
\item \textred{Aña.} Factus sum sicut homo sine adjutório inter mórtuos liber.
\item \textred{Ps. 87.} Dómine, Deus salútis. 287.
\item \textred{℣.}~Factus est obédiens usque ad mortem. \textred{℟.}~Mortem autem crucis.
\end{tightlist}

\subsubsection*{\center Léctio sancti Evangélii secúndum Joánnem.}
\subsubsection*{\center \textred{Lectio vij. \hfill Cap. 19.}}
\lettrine{I}n illo témpore~: Sciens Jesus, quia ómnia consummáta sunt, ut consummarétur Scriptúra, dixit~: Sítio. Et réliqua.

\begin{center}Homilía sancti Cyrílli Epíscopi Alexandríni.\\
\textred{Comment. in Joannem.\\ Lib. 12. cap. 25.}\end{center}
\lettrineThree{C}um truculénter ómnia impietáte Judæórum essent perácta, nec quidquam ad summam crudelitátem deésset, tandem sanctíssima caro Christi naturále quid pátitur. Multis enim, ac váriis dolóribus exsiccáta, siti torquétur. Multum enim possunt dolóres ad commovéndam sitim; innátum enim calórem exagitántes, vel quod in profunditáte húmidum est, consúmunt et ígneis ardóribus doléntis víscera urunt. Sicut ígitur céteras passiónes, sic hanc étiam sponte Christus perpéssus est. Bíbere ígitur petit~: illi autem ádeo longe ab omni humanitáte áberant, ut pro juvánte potu atque jucúndo nocéntem et acérbum attúlerint, et humanitátem, qua movéri videbántur, in impietátem convérterint. Nam cum peténti dabant, caritátis hábitum suscípere videbántur; sed non est possíbile mentíri Scriptúram, in qua ex persóna Christi dícitur~: Et dedérunt in escam meam fel, et in siti mea potavérunt me acéto.

\begin{tightlist}{}
\item \textred{℟.}~Ego te potávi aqua salútis de petra~: \textred{*} Et tu me potásti felle et acéto. \textred{℣.}~Pópule meus, quid feci tibi, aut in quo contristávi te? respónde mihi. \textred{E}t.
\end{tightlist}

\subsubsection*{\center\textred{Lectio viij. \hfill Idem. Cap. 36.}}
\lettrine{C}um ergo accepísset Jesus acétum, dixit~: Consummátum est~: et inclináto cápite trádidit spíritum. Consummátum jam esse Salvátor dicit, quia sæviéndi Judæórum potéstas in ipsum per mortem suam términum hábuit. Quid enim non excogitárunt, aut quid suprémæ inhumanitátis relíctum sibi est? Quis cædéndi, quis contumeliándi modus non adhíbitus? Mérito ígitur consummátum esse, ait. Cumque tempus jam esset, ut étiam spirítibus apud ínferos prædicáret (venit enim, ut tam vivórum quam mortuórum dominarétur), mortem pro nobis suscépit~: et hanc natúræ nostræ própriam passiónem sponte súbiit secúndum carnem, quamvis, ut Deus vita naturáliter sit, ut spoliátis ínferis, primítiæ dormiéntium, et primogénitus, (ut Scriptúræ dicunt) mortuórum factus, recúrsum ad vitam natúræ nostræ largirétur.

\begin{tightlist}{}
\item \textred{℟.}~Ego ante te apérui mare~: \textred{*} Et tu aperuísti láncea latus Salvatóri tuo, qui redémit mundum in sánguine suo. \textred{℣.}~Pópule meus, quid feci tibi, aut in quo contristávi te? respónde mihi. \textred{E}t.
\end{tightlist}

\subsubsection*{\center\textred{Lectio ix.}}
\lettrine{I}nclinávit ígitur caput. Accidit enim moriéntibus, vigóre carnis jam consúmpto et ánima, quæ corpus rétinet, abeúnte, inclinátio cápitis. Quámobrem verbis hoc modo usus Evangelísta est. Sic et tradidísse spíritum, a consuetúdine nostra sibi dictum est; multi enim sic exstíngui, et mori dicúntur. Quamvis mihi videátur ídeo non simplíciter mórtuum esse, sed tradidísse spíritum Evangelísta dixísse; quia in manus Dei Patris ánimam suam commendávit, ut ipse dixit~: Pater! in manus tuas comméndo spíritum meum. Quod nobis magnæ spei fundaméntum atque oríginem præbet.

\begin{tightlist}{}
\item \textred{℟.}~Crucem tuam adorámus, Dómine, sanctámque recólimus passiónem~: \textred{*} Ecce enim propter lignum venit gáudium in univérso mundo. \textred{℣.}~Deus misereátur nostri, et benedícat nobis, illúminet vultum suum super nos, et misereátur nostri. \textred{E}cce. Glória Patri. \textred{E}cce.
\item Te Deum laudámus. 6.
\item \textred{Sacerd. ℣.}~Adorámus te, Christe, et benedícimus tibi. \textred{℟.}~Quia per mortem tuam redemísti mundum.
\end{tightlist}

\subsubsection*{\center \textsc{AD LAUDES.}}
\begin{center}\textred{et per Horas, Aña.}\end{center}
\lettrine{F}ui flagellátus tota die, et castigátio mea in matutínis.

\begin{tightlist}{}
\item \textred{Ps. de Dominica.} 27. 
\item \textred{Aña.} Veni ad montem myrrhæ et ad collem thuris; tamquam agnus ad occisiónem ductus, obmútui, et non apérui os meum.
\item \textred{Aña.} Fodérunt manus meas, et pedes meos, dinumeravérunt ómnia ossa mea.
\item \textred{Aña.} Consolántem me quæsívi, et non invéni~: dedérunt in escam meam fel, et in siti mea potavérunt me acéto.
\item \textred{Aña.} Cum accepísset Jesus acétum, dixit~: Consummátum est, et inclináto cápite, emísit spíritum.
\end{tightlist}

\subsubsection*{\textred{Capitulum. \hfill Philipp. 2.}}
\lettrineThree{F}ratres~: Hoc enim sentíte in vobis, quod et in Christo Jesu, qui cum in forma Dei esset, non rapínam arbitrátus est esse se æquálem Deo, sed semetípsum exinanívit, formam servi accípiens, in similitúdinem hóminum factus, et hábitu invéntus ut homo. Humiliávit semetípsum factus obédiens usque ad mortem, mortem autem crucis.

\subsubsection*{\center \textred{Hymnus.}}
\lettrine{S}ævo dolórum túrbine\\
Jactátur, atris óbrutus\\
Pœnis, acérba sústinens\\
Redémptor affíxus cruci.\\
\\
\textred{P}edes manúsque horríbili\\
Clavi cruéntant vúlnere,\\
Cor, vultus, artus, péctora\\
Sacro redúndant sánguine.\\
\\
\textred{F}let, orat, et clamans obit~:\\
Cor Matris ictum cóncidit~:\\
Heu, Mater! Heu, Fili! dolor\\
Ingráta frangat péctora.\\
\\
\textred{M}ontes, sepúlcra, sáxaque\\
Scindúntur; arva, flúmina,\\
Rupes et æquor cóntremunt,\\
Templíque velum scínditur.\\
\\
\textred{S}ol, luna, cœlum, sídera\\
Plangunt, et orbis íngemit~:\\
O vos viri, vos párvuli,\\
Nuptæ, puéllæ, plángite.\\
\\
\textred{A}dstáre mœréntes cruci,\\
Pedes beátos úngite,\\
Laváte fletu, térgite\\
Comis, et ore lámbite.\\
\\
\textred{T}u caritátis víctima,\\
Ut nostra tollas crímina,\\
Nobis salúbri pérficis\\
Adoptiónem sánguine.\\
\\
\textred{N}ostra ergo pax, et gáudium\\
Sis vita, Jesu, et prǽmium;\\
Sis ductor, et lux in via,\\
Merces, coróna in Pátria.\\
Amen.

\begin{tightlist}{}
\item \textred{℣.}~Ipse passus est propter iniquitátes nostras. \textred{℟.}~Et attrítus est propter scélera nostra.
\item \textred{Ad Benedictus, Aña.} Ad Jesum autem cum veníssent, ut vidérunt eum jam mórtuum, non fregérunt ejus crura, sed unus mílitum láncea latus ejus apéruit, et contínuo exívit sanguis et aqua.
\end{tightlist}

\subsubsection*{\center\textred{Oratio.}}
\lettrine{D}omine Jesu Christe, qui de cœlis ad terram de sinu Patris descendísti et sánguinem tuum pretiósum in remissiónem peccatórum nostrórum fudísti~: te humíliter deprecámur; ut in die judícii, ad déxteram tuam audíre mereámur~: Veníte benedícti. Qui vivis et regnas.

\begin{tightlist}{}
\item \textred{Ad Horas Añæ de Laudibus, Psalmi de Dominica, ad Primam tamen ut in Festis.}
\end{tightlist}

\subsubsection*{\center \textsc{AD PRIMAM.}}
\begin{tightlist}{}
\item \textred{℣.}~Qui nos salvásti in Cruce.
\end{tightlist}

\subsubsection*{\center \textsc{AD TERTIAM.}}
\subsubsection*{\textred{Capitulum. \hfill Philipp. 2.}}
\lettrine{F}ratres~: Hoc enim sentíte in vobis, quod et in Christo Jesu, qui cum in forma Dei esset, non rapínam arbitrátus est esse se æquálem Deo~: sed semetípsum exinanívit, formam servi accípiens, in similitúdinem hóminum factus, et hábitu invéntus ut homo. Humiliávit semetípsum factus obédiens usque ad mortem, mortem autem crucis.

\begin{tightlist}{}
\item \textred{℟.}~Oblátus est, \textred{*} Quia ipse vóluit. \textred{O}blátus. \textred{℣.}~Et livóre ejus sanáti sumus. \textred{Q}uia. Glória Patri. \textred{O}blátus 
\item \textred{℣.}~Próprio Fílio suo non pepércit Deus. \textred{℟.}~Sed pro nobis ómnibus trádidit illum.
\end{tightlist}

\subsubsection*{\center \textsc{AD SEXTAM.}}
\subsubsection*{\textred{Capitulum. \hfill Zach. 13.}}
\lettrine{Q}uid sunt plagæ istæ in médio mánuum tuárum? Et dicet~: His plagátus sum in domo eórum, qui diligébant me.

\begin{tightlist}{}
\item \textred{℟.}~Próprio Fílio suo \textred{*} Non pepércit Deus. \textred{P}róprio. \textred{℣.}~Sed pro nobis ómnibus trádidit illum. \textred{N}on. Glória Patri. \textred{P}róprio.
\item \textred{℣.}~Factus est obédiens usque ad mortem. \textred{℟.}~Mortem autem crucis.
\end{tightlist}

\subsubsection*{\center \textsc{AD NONAM.}}
\subsubsection*{\textred{Capitulum. \hfill Is. 53.}}
\lettrine{G}eneratiónem ejus quis enarrábit? Quia abscíssus est de terra vivéntium, propter scelus pópuli mei percússi eum.

\begin{tightlist}{}
\item \textred{℟.}~Factus est obédiens \textred{*} Usque ad mortem. \textred{F}actus. \textred{℣.}~Mortem autem crucis. \textred{U}sque. Glória Patri. \textred{F}actus.
\item \textred{℣.}~Ipse passus est propter iniquitátes nostras. \textred{℟.}~Et attrítus est propter scélera nostra.
\end{tightlist}

\subsubsection*{\center \textsc{AD II. VESPERAS.}}
\begin{tightlist}{}
\item \textred{Prima Aña Laudum.} Fui flagellátus tota die, et castigátio mea in matutínis.
\item \textred{Psalm. 115} Crédidi. 91. 
\item \textred{Ps. 119.} Ad Dóminum. 92. 
\item \textred{Ps. 139.} Eripe me. 244. 
\item \textred{Ps. 140.} Dómine, clamávi. 244. 
\item \textred{Ps. 141.} Voce mea. 245.
%\item \textred{Capit, Hymn. et ℣. ut in j. Vesperis.}
\end{tightlist}

\subsubsection*{\textred{Capitulum. \hfill Philipp. 2.}}
\lettrineThree{F}ratres~: Hoc sentíte in vobis, quod et in Christo Jesu~: qui, cum in forma Dei esset, non rapínam arbitrátus est esse se æquálem Deo~: sed semetípsum exinanívit, formam servi accípiens, in similitúdinem hóminum factus, et hábitu invéntus ut homo. Humiliávit semetípsum, factus obédiens usque ad mortem, mortem autem crucis.

\subsubsection*{\center \textred{Hymnus.}}
\lettrine{M}œréntes óculi, spárgite lácrymas,\\
Et luctu résonent íntima córdium~:\\
Illátas réfero Númini ab ímpiis\\
Pœnas, et fera vúlnera.\\
\\
\textred{A}ccíncta, heu! gládiis turba satéllitum\\
Arréptum Dóminum fústibus ímpetit,\\
Nunc cædit cólaphis, nunc quatit hórridis\\
Divínum caput íctibus.\\
\\
\textred{H}aud finis scéleri~: tráditur ímprobo\\
Christus carnífici; nec mora, bárbarus\\
In Regem súperum non timet ímpio\\
Ausu vértere déxteram.\\
\\
\textred{A}udíte, o pópuli, Numen amábile,\\
Manánte ex húmeris úndique sánguine,\\
Lictóris rábidi sústinet ímpetum,\\
Et vocem premit ínnocens.\\
\\
\textred{Q}uis non illácrimet? Jam nova cóndidit\\
Torménta indómitæ gentis iníquitas~:\\
Infígit cérebro, proh dolor! éfferis\\
Sertum vépribus ásperum.\\
\\
\textred{H}eu crimen! tráhitur fúnibus ímprobis\\
Funésti Dóminus supplícii ad locum;\\
Illic oppétiit fúnere, spíritum\\
Patri restítuens suo.\\
\\
\textred{P}asso pro míseris tam fera vúlnera,\\
In terris résonet débita glória;\\
Sacrúmque assídue Nomen in ǽthera\\
Humánum genus éfferat.\\
Amen.

\begin{tightlist}{}
\item \textred{℣.}~Oblátus est quia ipse vóluit. \textred{℟.}~Et livóre ejus sanáti sumus.
\item \textred{Ad Magnif. Aña.} Depónens Joseph corpus Jesu invólvit síndone, et pósuit eum in monuménto excíso. Fuit planctus magnus in illa die; mulíeres autem sedéntes contra sepúlcrum dolébant, sicut doléri solet in morte primogéniti.
\end{tightlist}

\subsubsection*{\center\textred{Oratio.}}
\lettrine{D}omine Jesu Christe, qui de cœlis ad terram de sinu Patris descendísti et sánguinem tuum pretiósum in remissiónem peccatórum nostrórum fudísti~: te humíliter deprecámur; ut in die judícii, ad déxteram tuam audíre mereámur~: Veníte benedícti. Qui vivis et regnas.

\begin{tightlist}{}
\item \textred{Completorium de Dominica.}
\item \textred{Ad Nunc dimittis, Aña.} Hoc signum Crucis * erit in cælo, cum Dóminus ad judicándum vénerit.
\end{tightlist}

\begin{center}\textcolor{red}{\rule{\linewidth}{0.5pt}} \end{center}

\subsection*{\center \textred{ROMÆ\\ PRO FERIA III.}}
\subsection*{\center \textsc{\textred{OFFICIUM VOTIVUM}\\ De SS. Apostolis Petro et Paulo.}}
\subsubsection*{\center \textred{Semiduplex.}}
\fancyhead[C]{\textsc{Offic. Vot. - Feria III Romæ - SS. Petro et Paulo}}

\begin{tightlist}{}
\item \textred{Omnia de Comm. Apost.} [2] \textred{pro diversitate temporis, præter sequentia.}
\end{tightlist}

\subsubsection*{\center \textsc{IN UTRISQUE VESPERIS.}}
\begin{tightlist}{}
\item \textred{℣.}~Constítues eos príncipes super omnem terram. \textred{℟.}~Mémores erunt nóminis tui, Dómine.
\item \textred{Ad Magnificat, Aña.} Petrus Apóstolus, et Paulus Doctor géntium, ipsi nos docuérunt legem tuam, Dómine.
\end{tightlist}

\subsubsection*{\center \textred{Oratio.}}
\lettrine{D}eus, cujus déxtera beátum Petrum ambulántem in flúctibus, ne mergerétur, eréxit, et co-apóstolum ejus Paulum, tértio naufragántem, de profúndo pélagi liberávit~: exáudi nos propítius, et concéde; ut ambórum méritis æternitátis glóriam consequámur~: Qui vivis.

\subsubsection*{\center \textsc{IN PRIMO NOCTURNO.}}
\begin{tightlist}{}
\item \textred{Lect. de Script. occurr.}
\item \textred{In Quadrag. vero dicuntur sequentes.}
\end{tightlist}

\subsubsection*{\center Incipit Epístola prima beáti Petri Apóstoli.}
\subsubsection*{\center \textred{Lectio j. \hfill Cap. 1.}}
\lettrineThree{P}etrus, Apóstolus Jesu Christi, eléctis ádvenis dispersiónis Ponti, Galátiæ, Cappadóciæ, Asiæ, et Bithýniæ,
secúndum præsciéntiam Dei Patris, in sanctificatiónem Spíritus, in obediéntiam, et aspersiónem sánguinis Jesu Christi. Grátia vobis, et pax multiplicétur.
Benedíctus Deus et Pater Dómini nostri Jesu Christi, qui secúndum misericórdiam suam magnam regenerávit nos in spem vivam, per resurrectiónem Jesu Christi ex mórtuis,
in hereditátem incorruptíbilem, et incontaminátam, et immarcescíbilem, conservátam in cælis in vobis,
qui in virtúte Dei custodímini per fidem in salútem, parátam revelári in témpore novíssimo.

\begin{tightlist}{}
\item \textred{℟.}~Ecce, ego mitto vos sicut oves in médio lupórum, dicit Dóminus~: \textred{*} Estóte ergo prudéntes sicut serpéntes, et símplices sicut colúmbæ. \textred{℣.}~Dum lucem habétis, crédite in lucem, ut fílii lucis sitis. \textred{E}stóte.
\end{tightlist}

\subsubsection*{\center \textred{Lectio ij.}}
\lettrine{I}n quo exsultábitis, módicum nunc si opórtet contristári in váriis tentatiónibus:
ut probátio vestræ fídei multo pretiósior auro (quod per ignem probátur) inveniátur in laudem, et glóriam, et honórem in revelatióne Jesu Christi:
quem, cum non vidéritis, dilígitis~: in quem nunc quoque non vidéntes créditis~: credéntes autem exsultábitis lætítia inenarrábili, et glorificáta:
reportántes finem fídei vestræ, salútem animárum.

\begin{tightlist}{}
\item \textred{℟.}~Tóllite jugum meum super vos, dicit Dóminus~: et díscite a me, quia mitis sum, et húmilis corde~: \textred{*} Jugum enim meum suáve est, et onus meum leve. \textred{℣.}~Et inveniétis réquiem animábus vestris. \textred{J}ugum.
\end{tightlist}

\subsubsection*{\center \textred{Lectio iij.}}
\lettrine{D}e qua salúte exquisiérunt, atque scrutáti sunt prophétæ, qui de futúra in vobis grátia prophetavérunt:
scrutántes in quod, vel quale tempus significáret in eis Spíritus Christi~: prænúntians eas quæ in Christo sunt passiónes, et posterióres glórias:
quibus revelátum est quia non sibimetípsis, vobis autem ministrábant ea, quæ nunc nuntiáta sunt vobis per eos, qui evangelizavérunt vobis, Spíritu sancto misso de cælo, in quem desíderant ángeli prospícere.

\begin{tightlist}{}
\item \textred{℟.}~Dum stetéritis ante reges et prǽsides, nolíte cogitáre quómodo, aut quid loquámini~: \textred{*} Dábitur enim vobis in illa hora, quid loquámini. \textred{℣.}~Non enim vos estis, qui loquímini~: sed Spíritus Patris vestri, qui lóquitur in vobis. \textred{D}ábitur. Glória Patri. \textred{D}ábitur.
\end{tightlist}

\subsubsection*{\center \textsc{IN SECUNDO NOCTURNO.}}
\subsubsection*{\center De Sermóne sancti Leónis Papæ.}
\subsubsection*{\center \textred{Ex hom. 82. in Natal. SS. Apost. Petri et Pauli.\\ Lectio iv.}}
\lettrineThree{N}erónis furor per ómnium vitiórum inflammátus excéssum, in hunc eum usque torréntem suæ præcipitávit insániæ, ut primus nómini Christiáno atrocitátem generális persecutiónis inférret, quasi per sanctórum Apostolórum Petri et Pauli neces grátia Dei posset exstíngui; quibus hoc ipsum erat máximum lucrum, ut contémptus vitæ hujus occíduæ, percéptio fíeret felicitátis ætérnæ. Pretiósa est ergo in conspéctu Dómini mors sanctórum ejus~: nec ullo crudelitátis génere déstrui potest sacraménto crucis Christi fundáta relígio. Non minúitur persecutiónibus Ecclésia, sed augétur~: et semper domínicus ager ségete ditióri vestítur, dum grana, quæ síngula cadunt, multiplicáta nascúntur. Unde duo ista præclára divíni séminis gérmina in quantam sóbolem germinárint, beatórum míllia Mártyrum protestántur, qui apostolicórum æmuli triumphórum, Urbem nostram purpurátis, et longe latéque rutilántibus pópulis ambiérunt, et quasi ex multárum honóre gemmárum consérto uno diadémate coronárunt.

\begin{tightlist}{}
\item \textred{℟.}~Vidi conjúnctos viros habéntes spléndidas vestes~: et Angelus Dómini locútus est ad me, dicens~: \textred{*} Isti sunt viri sancti, facti amíci Dei. \textred{℣.}~Vidi Angelum Dei fortem volántem per médium cœli, voce magna clamántem et dicéntem. \textred{I}sti.
\end{tightlist}

\subsubsection*{\center \textred{Lectio v.}}
\lettrine{D}e quorum præsídio, dilectíssimi, divínitus nobis ad exémplum patiéntiæ et confirmatiónem fídei præparáto, universáliter quidem in ómnium sanctórum commemoratióne lætándum est~: sed in horum excelléntia Patrum mérito est excelléntius gloriándum, quos grátia Dei in tantum ápicem inter ómnia Ecclésiæ membra provéxit, ut eos in córpore, cui caput est Christus, quasi géminum constitúeret lumen oculórum. De quorum méritis atque virtútibus, quæ omnem loquéndi súperant facultátem, nihil divérsum, nihil debémus sentíre discrétum~: quia illos et eléctio pares, et labor símiles, et finis fecit æquáles.

\begin{tightlist}{}
\item \textred{℟.}~Isti sunt triumphatóres, et amíci Dei, qui contemnéntes jussa príncipum, meruérunt prǽmia ætérna~: \textred{*} Modo coronántur, et accípiunt palmam. \textred{℣.}~Isti sunt, qui venérunt ex magna tribulatióne, et lavérunt stolas suas in sánguine Agni. \textred{M}odo.
\end{tightlist}

\subsubsection*{\center \textred{Lectio vj.}}
\lettrine{S}icut autem et nos expérti sumus, et nostri probavére majóres, crédimus atque confídimus, inter omnes labóres istíus vitæ, ad obtinéndam misericórdiam Dei, semper nos speciálium patronórum oratiónibus adjuvándos~: ut, quantum própriis peccátis deprímimur, tantum apostólicis méritis erigámur. Per Dóminum nostrum Jesum Christum, cui est cum Patre et sancto Spíritu éadem potéstas, una divínitas, in sǽcula sæculórum. Amen.

\begin{tightlist}{}
\item \textred{℟.}~Fuérunt sine queréla ante Dóminum, et ab ínvicem non sunt separáti~: \textred{*} Cálicem Dómini bibérunt, et amíci Dei facti sunt.
\textred{℣.}~Tradidérunt córpora sua propter Deum ad supplícia~: ídeo coronántur, et accípiunt palmam. \textred{C}álicem. Glória Patri. \textred{C}álicem.
\end{tightlist}

\subsubsection*{\center \textsc{IN TERTIO NOCTURNO.}}
\subsubsection*{\center Léctio sancti Evangélii secúndum Matthǽum.}
\subsubsection*{\center \textred{Lectio vij. \hfill Cap. 19.}}
\lettrine{I}n illo témpore~: Dixit Petrus ad Jesum~: Ecce nos relíquimus ómnia, et secúti sumus te~: quid ergo erit nobis? Et réliqua.

\begin{center}Homilía sancti Ambrósii Epíscopi.\\
\textred{In Ps. 118. Serm. 8.}\end{center}
\lettrineThree{E}cce nos ómnia derelíquimus, et secúti sumus te. Hoc est, non quæsívimus quæ sǽculi sunt, non quæsívimus partem de possessiónibus, sed te elégimus portiónem. Reliquísti ergo ómnia, Petre, quæ habébas; unde habes quod habére te dicis? Surgit claudus, et sono tui sermónis erígitur~: sanitátem áliis donas, qui indigébas ipse tuæ salútis auxílio. Reliquísti ergo quæ habébas, et cepísti, quæ non habébas. Christus tibi pórtio, Christus tibi posséssio est~: illíus nomen tibi muníficum, illíus nomen tibi est fructuósum, illíus nomen tibi tribúta depéndit, et bona tribúta, non pecúniæ, sed grátiæ.

\begin{tightlist}{}
\item \textred{℟.}~Isti sunt, qui vivéntes in carne, plantavérunt Ecclésiam sánguine suo. Non sunt de terris córpora eórum separáta~: \textred{*} Quorum mérita sunt in cælo glorióse cumuláta. \textred{℣.}~Sancti per fidem vicérunt regna, operáti sunt justítiam, adépti sunt repromissiónes. \textred{Q}uorum.
\end{tightlist}

\subsubsection*{\center \textred{Lectio viij.}}
\lettrine{Q}uicúmque relíquerit patrem, aut matrem, aut domum, aut fratres, aut soróres, céntuplum accípiet in hoc sǽculo, et vitam ætérnam possidébit. Multos quidem in hoc sǽculo, qui sua paupéribus donavérunt, uberióre thesáuro locupletátos fuísse cognóvimus. Est pulchérrimus locus, ut étiam ad misericórdiam provocéntur. Sed non debent a Dómino sæculárem hujúsmodi exígere mercédem, nec speráre quæ mundi sunt; sed magis, quia iis, qui ómnia relíquerint, Deus pórtio est. Ipse est útique merces perfécta virtútum, qui non céntupli enumeratióne, sed perféctæ plenitúdinis æstimatióne censétur. Ego, inquit, Deus tuus sum. Non dixit~: Ero, sed~: Jam sum, jam inhábito, jam possídeo.

\begin{tightlist}{}
\item \textred{℟.}~Isti sunt viri sancti, quos elégit Dóminus in caritáte non ficta, et dedit illis glóriam sempitérnam~: \textred{*} Quorum doctrína fulget Ecclésia, ut sole luna. \textred{℣.}~In omnem terram exívit sonus eórum, et in fines orbis terræ verba eórum. \textred{Q}uorum.
\end{tightlist}

\subsubsection*{\center \textred{Lectio ix.}}
\lettrine{H}abes, homo, inter multa propósita, étiam hujúsmodi portiónem. Propósuit tibi Dóminus in portióne possessiónes, in portióne aurum, in portióne argéntum, in portióne honóres, in portióne nobilitátem, propósuit étiam in portióne seípsum. Habes ígitur plúrimas portiónes; élige quam putas; non te númerus pertúrbet, sed éxcitet grátia; non te labor avértat, sed fructus invítet.

\begin{tightlist}{}
\item \textred{℟.}~Cives Apostolórum et doméstici Dei advenérunt hódie, \textred{*} Portántes facem et illuminántes pátriam, dare pacem géntibus et liberáre pópulum Dómini. \textred{℣.}~Audíte preces súpplicum, vitæ ætérnæ poscéntes prǽmia, qui fertis in dextris manípulos justítiæ, quique gaudéntes advenístis hódie. \textred{P}ortántes. Glória Patri. \textred{P}ortántes.
\item Te Deum laudámus. 6.
\item \textred{Sacerdotalis ℣.}~Dedísti hereditátem. \textred{℟.}~Timéntibus nomen tuum, Dómine.
\end{tightlist}

\subsubsection*{\center \textsc{AD LAUDES.}}
\begin{tightlist}{}
\item \textred{℣.}~In omnem terram exívit sonus eórum. \textred{℟.}~Et in fines orbis terræ verba eórum.
\item \textred{Ad Bened. Aña.} Gloriósi príncipes terræ, quómodo in vita sua dilexérunt se, ita et in morte non sunt separáti.
\end{tightlist}

\begin{center}\textcolor{red}{\rule{\linewidth}{0.5pt}} \end{center}
\clearpage
\section*{\center \textsc{SUFFRAGIA CONSUETA}}
\subsection*{\center \textsc{Seu Commemorationes Communes.}}

\fancyhead[C]{\textsc{Commemorationes Communes - In Adventu}}

\begin{tightlist}{}
\item \textred{Dicendæ in fine Vesperarum et Laudum ab Octava Epiphaniæ usque ad Dominicam Passionis exclusive, et ab Octava Corporis Christi usque ad Vigiliam Nativitatis Domini exclusive, in omnibus Dominicis, in quibus non fit Officium Duplex, in Feriis, comprehensis etiam Vigiliis, in Festis Semiduplicibus, Simplicibus, et Simplicissimis extra quascumque Octavas celebratis, juxta Rubricas Generales, Rubr. 41. prout dictæ commemorationes pro suo respectivo tempore inferius distributæ et designatæ sunt.}
\end{tightlist}

\begin{center}\textcolor{red}{\rule{\linewidth}{0.5pt}} \end{center}

\subsection*{\center \textsc{IN ADVENTU.}}
\begin{tightlist}{}
\item \textred{Exceptis utrisque Vesperis et Laudibus primæ Dominicæ, et Vesp. in quibus dicitur Aña.} O. \textred{ac Laudibus Feriæ quartæ Quatuor Temporum.}
\end{tightlist}

\subsubsection*{\center \textsc{DE B. VIRGINE.}}
\begin{tightlist}{}
\item \textred{Quando Officium diei non est de B. V. etiamsi in Choro recitetur ejus Officium parvum.}
\item \textred{Ad Vesperas et Laudes. Aña.} Spíritus sanctus in te descéndet, María, ne tímeas; habébis in útero Fílium Dei, allelúja.
\item \textred{℣.}~Ave María, grátia plena. \textred{℟.}~Dóminus tecum.
\item Orémus. \hfill \textred{Oratio.}
\end{tightlist}
\lettrine{D}eus qui de beátæ Maríæ Vírginis útero, Verbum tuum, Angelo nuntiánte, carnem suscípere voluísti~: præsta supplícibus tuis; ut, qui vere eam Genitrícem Dei crédimus, ejus apud te intercessiónibus adjuvémur. Per eúmdem Christum Dóminum nostrum. \textred{℟.}~Amen.

\subsubsection*{\center \textsc{DE OMNIBUS SANCTIS.}}
\begin{tightlist}{}
\item \textred{Aña.} Ecce Dóminus véniet, et omnes Sancti ejus cum eo~: et erit in die illa lux magna, allelúja.
\item \textred{℣.}~Ecce apparébit Dóminus super nubem cándidam. \textred{℟.}~Et cum eo Sanctórum míllia.
\item Orémus. \hfill \textred{Oratio.}
\end{tightlist}
\lettrine{C}onsciéntias nostras, quǽsumus Dómine, visitándo purífica~: ut véniens Jesus Christus Fílius tuus Dóminus noster, cum ómnibus Sanctis, parátam sibi in nobis invéniat mansiónem. Qui tecum vivit et regnat in unitáte Spíritus Sancti Deus. \textred{℟.}~Amen.

\begin{center}\textcolor{red}{\rule{\linewidth}{0.5pt}} \end{center}

\subsection*{\center \textsc{PER ANNUM.}}
\subsubsection*{\center \textsc{DE SS. SACRAMENTO.}}
\begin{tightlist}{}
\item \textred{Fit comm. ab Oct. Epiph. usque ad Quadrag. et ab Oct. Corp. Christi usque ad Adventum; solis Feriis quartis ad Vesperas, et Feriis quintis ad Laudes et Vesperas, nisi fit Officium de Venerabili Sacramento, hoc modo:}
\item \textred{Feria quarta ad Vesp., nisi præcedat Festum Duplex, Aña.} O quam suávis est, Dómine, Spíritus tuus~: qui, ut dulcédinem tuam in fílios demonstráres, pane suavíssimo de cœlo prǽstito, esuriéntes reples bonis, fastidiósos dívites dimíttens inánes.
\item \textred{℣.}~Panem de cœlo præstitísti eis. \textred{℟.}~Omne delectaméntum in se habéntem.
\item \textred{Feria quinta ad Laudes, Aña.} Ego sum panis vivus, qui de cœlo descéndi~: si quis manducáverit ex hoc pane, vivet in ætérnum, allelúja.
\item \textred{℣.}~Pósuit fines tuos pacem. \textred{℟.}~Et ádipe fruménti sátiat te.
\item \textred{Feria quinta ad Vesperas, nisi sequatur Festum Duplex, Aña.} O sacrum convívium, in quo Christus súmitur~: recólitur memória passiónis ejus~: mens implétur grátia, et futúræ glóriæ nobis pignus datur, allelúja.
\item \textred{℣.}~Panem de cœlo præstitísti eis. \textred{℟.}~Omne delectaméntum in se habéntem.
\item Orémus. \hfill \textred{Oratio.}
\end{tightlist}
\lettrine{D}eus, qui nobis sub Sacraménto mirábili passiónis tuæ memóriam reliquísti~: tríbue quǽsumus, ita nos córporis et sánguinis tui sacra mystéria venerári; ut redemptiónis tuæ fructum in nobis júgiter sentiámus. Qui vivis et regnas in sǽcula sæculórum. \textred{℟.}~Amen.

\subsubsection*{\center \textsc{DE RESURRECTIONE.}}
\fancyhead[C]{\textsc{Commemorationes Communes - Per Annum}}
\begin{tightlist}{}
\item \textred{Fit com. ab Oct. Corp. Christi usque ad Adventum, Dominicis diebus ad Laudes tantum, quando de Dominica agitur, etiam si sit infra Octav. et hæc præcedere debet omnes alias commem. tam occurrentes proprias, quam communes.}
\item \textred{Aña.} Et valde mane una Sabbatórum véniunt ad monuméntum, orto jam sole, allelúja.
\item \textred{℣.}~Surréxit Dóminus de sepúlcro. \textred{℟.}~Qui pro nobis pepéndit in ligno.
\item Orémus. \hfill \textred{Oratio.}
\end{tightlist}
\lettrine{P}ræsta quǽsumus omnípotens Deus, ut qui Resurrectiónis Domínicæ commemoratiónem ágimus, ereptiónis nostræ lætítiam suscípere mereámur. Per eúmdem Christum Dóminum nostrum. \textred{℟.}~Amen.

\subsubsection*{\center \textsc{DE CRUCE.}}
\begin{tightlist}{}
\item \textred{Ferialibus diebus tantum, (comprehensis etiam Vigiliis) ab Oct. Epiph. usque ad Dom Passionis, et ab Oct. Corp. Christi usque ad Adventum.}
\item \textred{Ad Vesp. et Laud. Aña.} Per signum Crucis de inimícis nostris líbera nos, Deus noster.
\item \textred{℣.}~Omnis terra adóret te Deus, et psallat tibi. \textred{℟.}~Psalmum dicat nómini tuo Dómine.
\item Orémus. \hfill \textred{Oratio.}
\end{tightlist}
\lettrine{P}erpétua nos quǽsumus Dómine pace custódi~: quos per lignum sanctæ Crucis redímere dignátus es. Qui vivis et regnas in sǽcula sæculórum. \textred{℟.}~Amen.

\begin{center}\textcolor{red}{\rule{\linewidth}{0.5pt}} \end{center}

\subsection*{\center \textsc{POST EPIPHANIAM USQUE AD ADVENTUM.}}
\begin{tightlist}{}
\item \textred{Ab Oct. Epiph. usque ad Domin. Passionis, et ab Oct. Corp. Christi usque ad Adventum, fiunt sequentes commemorationes.}
\end{tightlist}

\subsubsection*{\center \textsc{DE TITULARI.}}
\textred{Vel Patrono Ecclesiæ fit commemoratio ante vel post comm. de Sanctis Ordinis Nostri, juxta ejus dignitatem.}

\subsubsection*{\center \textsc{DE BEATA VIRGINE.}}
\begin{tightlist}{}
\item \textred{Quando Officium diei non est de B.V. etiamsi in Choro recitetur ejus Officium parvum.}
\item \textred{Ad Vesperas et Laudes, Aña.} Sancta María, succúrre míseris, juva pusillánimes, réfove flébiles~: ora pro pópulo, intérveni pro clero, intercéde pro devóto femíneo sexu; séntiant omnes tuum juvámen, quicúmque célebrant tuam sanctam Commemoratiónem.
\item \textred{Ab Oct. Epiph. usque ad Purificationem dicitur:}
\item \textred{℣.}~Post partum, Virgo, invioláta permansísti. \textred{℟.}~Dei Génitrix intercéde pro nobis.
\item Orémus. \hfill \textred{Oratio.}
\end{tightlist}
\lettrine{D}eus, qui salútis ætérnæ, beátæ Maríæ virginitáte fecúnda, humáno géneri prǽmia præstitísti~: tríbue quǽsumus; ut ipsam pro nobis intercédere sentiámus, per quam merúimus auctórem vitæ suscípere, Dóminum nostrum Jesum Christum Fílium tuum. \textred{℟.}~Amen.

\begin{tightlist}{}
\item \textred{A Purificatione usque ad Adventum:}
\item \textred{℣.}~Dignáre me laudáre te, Virgo sacráta. \textred{℟.}~Da mihi virtútem contra hostes tuos.
\item Orémus. \hfill \textred{Oratio.}
\end{tightlist}
\lettrine{C}oncéde miséricors Deus fragilitáti nostræ præsídium; ut qui sanctæ Dei Genitrícis memóriam ágimus, intercessiónis ejus auxílio, a nostris iniquitátibus resurgámus. Per eúmdem Christum Dóminum nostrum. \textred{℟.}~Amen.

\subsubsection*{\center \textsc{DE S. JOSEPH.}}
\begin{tightlist}{}
\item \textred{Fit post commem. B. V. Mariæ et ante quoscumque alios Sanctos, exceptis Angelis et S. Joanne Baptista.}
\item \textred{Ad Vesperas, Aña.} Ecce fidélis servus et prudens, quem constítuit Dóminus super famíliam suam.
\item \textred{℣.}~Glória et divítiæ in domo ejus. \textred{℟.}~Et justítia ejus manet in sǽculum sǽculi.
\item \textred{Ad Laudes, Aña.} Ipse Jesus erat incípiens quasi annórum trigínta ut putabátur fílius Joseph.
\item \textred{℣.}~Os justi meditábitur sapiéntiam. \textred{℟.}~Et lingua ejus loquétur judícium.
\item Orémus. \hfill \textred{Oratio.}
\end{tightlist}
\lettrine{D}eus, qui ineffábili providéntia beátum Joseph sanctíssimæ Genitrícis tuæ sponsum elígere dignátus es~: præsta quǽsumus, ut quem protectórem venerámur in terris intercessórem habére mereámur in cœlis. Qui vivis et regnas in sǽcula sæculórum. \textred{℟.}~Amen.

\subsubsection*{\center \textsc{DE SANCTIS}}
\begin{center}\textsc{Nostri Ordinis.}\end{center}
\begin{tightlist}{}
\item \textred{Ad Vesperas et Laudes, Aña.} Christi pia grátia Sanctos sublimávit, quos beátæ Vírginis Ordo propagávit. Nos eórum méritis pétimus adjuvári, atque suis précibus Deo commendári.
\item \textred{℣.}~Sancti tui, Dómine, benedícent te. \textred{℟.}~Glóriam regni tui dicent, et poténtiam tuam loquéntur.
\item Orémus. \hfill \textred{Oratio.}
\end{tightlist}
\lettrine{O}mnípotens et clementíssime Deus, qui Montis Carméli Ordinem gloriosíssimæ Vírginis Maríæ, Genitrícis Fílii tui Dómini nostri Jesu Christi, sacráto título insignítum, sanctórum tuórum Patris nostri Elíæ et Eliséi Prophetárum, Angeli et Anastásii Mártyrum, Cyrílli et Albérti Confessórum, Euphrásiæ, Terésiæ, et Maríæ Magdalénæ Vírginum, et aliórum plurimórum méritis decorásti~: tríbue nobis quǽsumus; ut per eórum suffrágia, ab instántibus malis ánimæ et córporis liberáti, ad te verum Carméli vérticem gaudéntes perveníre valeámus. Per eúmdem Christum Dóminum nostrum. \textred{℟.}~Amen.

\subsubsection*{\center \textsc{DE OMNIBUS SANCTIS.}}
\begin{tightlist}{}
\item \textred{Ad Vesperas et Laudes, Aña.} Omnes elécti Dei, nostri memorámini ante Deum, ut vestris précibus adjúti mereámur vobis adjúngi.
\item \textred{℣.}~Oráte pro nobis omnes Sancti Dei. \textred{℟.}~Ut digni efficiámur promissiónibus Christi.
\item Orémus. \hfill \textred{Oratio.}
\end{tightlist}
\lettrine{O}mnium Sanctórum tuórum, quǽsumus Dómine, intercessióne placátus, et véniam nobis delictórum tríbue, et remédia sempitérna concéde. Per Christum Dóminum nostrum. \textred{℟.}~Amen.

\subsubsection*{\center \textsc{DE PACE.}}
\begin{tightlist}{}
\item \textred{Ad Vesperas et Laudes, Aña.} Da Pacem, Dómine, in diébus nostris~: quia non est álius qui pugnet pro nobis nisi tu Deus noster.
\item \textred{℣.}~Fiat pax in virtúte tua. \textred{℟.}~Et abundántia in túrribus tuis.
\item Orémus. \hfill \textred{Oratio.}
\end{tightlist}
\lettrine{D}eus, a quo sancta desidéria, recta consília, et justa sunt ópera~: da servis tuis illam, quam mundus dare non potest, pacem; ut et corda nostra mandátis tuis dédita, et hóstium subláta formídine, témpora sint tua protectióne tranquílla. Per Dóminum nostrum. \textred{℟.}~Amen.

\begin{center}\textcolor{red}{\rule{\linewidth}{0.5pt}} \end{center}

\subsection*{\center \textsc{TEMPORE PASCHALI.}}
\fancyhead[C]{\textsc{Commemorationes Communes - T.P.}}
\begin{tightlist}{}
\item \textred{Solum fit commem. sequens de Cruce, quæ dicatur etiam in Dominicis, et Festis Semiduplicibus, Simplicibus, et Simplicissimis, extra Octavas celebratis.}
\item \textred{Ad Vesperas et Laudes, Aña.} Crucifíxus surréxit a mórtuis, et redémit nos Sánguine suo, allelúja.
\item \textred{℣.}~Dícite in natiónibus. \textred{℟.}~Quia Dóminus regnávit a ligno, allelúja.
\item Orémus. \hfill \textred{Oratio.}
\end{tightlist}
\lettrine{D}eus, qui pro nobis Fílium tuum Crucis patíbulum subíre voluísti, ut inimíci a nobis expélleres potestátem; concéde nobis fámulis tuis, ut Resurrectiónis grátiam consequámur. Per eúmdem Dóminum nostrum. \textred{℟.}~Amen.

\begin{center}\textcolor{red}{\rule{\linewidth}{0.5pt}} \end{center}

\end{multicols}
\vfill
\begin{center} \includegraphics[width=0.8 \textwidth]{Images/Cross_divider} \end{center}

\cleardoublepage

\section*{\center \textsc{PRECES VESPERTINÆ}}
\fancyhead[C]{\textsc{Preces Vespertinæ OCD}}
\subsection*{\center \textsc{Ad usum Carmelitarum Discalceatorum.}}
\subsubsection*{\center \textred{Post examen conscientiæ vel communem disciplinam quotidie recitandæ, Præside incipiente.}}

\begin{multicols}{2}
\subsubsection*{\center \textsc{De Beata Virgine}}
\begin{tightlist}{}
\item \textred{Ant.} Ait Dóminus Deus * ad serpentem~: Inimicítias ponam inter te et mulíerem, et semen tuum et semen illíus~: Ipsa cónteret caput tuum.
\item {\small \textred{Et dicitur ter:} Ave, María.}
\item \textred{℣.}~Tota pulchra es, María. \textred{℟.}~Et mácula originális non est in te.
\item Orémus. \hfill \textred{Oratio}
\end{tightlist}
\lettrine{D}eus, qui per immaculátam Vírginis Conceptiónem dignum Fílio tuo habitáculum præparásti~: quǽsumus; ut, qui ex morte ejúsdem Fílii tui prævísa, eam ab omni labe præservásti, nos quoque mundos ejus intercessióne ad te perveníre concédas.

\lettrine{E}cclésiæ tuæ, quǽsumus, Dómine, preces placátus admítte~: ut, destrúctis adversitátibus et erróribus univérsis, secúra tibi sérviat libertáte.

\lettrine{D}eus, ómnium fidélium pastor et rector, fámulum tuum N., quem pastórem Ecclésiæ tuæ præésse voluísti, propítius réspice~: da ei, quǽsumus, verbo et exémplo quibus præest profícere; ut ad vitam, una cum grege sibi crédito, pervéniat sempitérnam. Per Christum Dóminum nostrum.\\
\textred{℟.}~Amen.

\begin{center}\textcolor{red}{\rule{\linewidth}{0.5pt}} \end{center}

\subsubsection*{\center \textsc{De S. P. N. Joseph}}
\begin{tightlist}{}
\item \textred{Ant.} Joseph, fili David, * noli timére accípere Maríam cónjugem tuam~: quod enim in ea natum est, de Spíritu Sancto est~: páriet autem Fílium et vocábis nomen ejus Jesum.
\item \textred{℣.}~Ora pro nobis, sancte Pater Joseph. \textred{℟.}~Ut digni efficiámur promissiónibus Christi.
\item Orémus. \hfill \textred{Oratio}
\end{tightlist}
\lettrine{S}anctíssimæ Genitrícis tuæ sponsi, quǽsumus, Dómine, méritis adjuvémur~: ut, quod possibílitas nostra non óbtinet, ejus nobis intercessióne donétur~: Qui vivis et regnas in sǽcula sæculórum. \textred{℟.}~Amen.

\begin{center}\textcolor{red}{\rule{\linewidth}{0.5pt}} \end{center}

\subsection*{\center \textsc{\textred{Officia Mensualia}}}
\subsubsection*{\center \textred{Dicenda singulis mensibus}}

\begin{center}\textcolor{red}{\rule{\linewidth}{0.5pt}} \end{center}

\subsubsection*{\center \textsc{De S. P. N. Elia}}
\begin{tightlist}{}
\item \textred{Ant.} Sancte Prophéta Dei Elía, * Carmelitárum Dux et Pater~: intercéde pro nostra omniúmque salúte.
\item \textred{℣.}~Ora pro nobis, sancte Pater Elía. \textred{℟.}~Ut digni efficiámur promissiónibus Christi.
\item Orémus. \hfill \textred{Oratio}
\end{tightlist}
\lettrine{P}ræsta, quǽsumus, omnípotens Deus~: ut, qui beátum Elíam Prophétam tuum et Patrem nostrum, ígneo curru mirabíliter elevátum esse crédimus; eódem interveniénte ad cæléstia sublevémur, et Sanctórum tuórum consórtio gaudeámus. Per Christum Dóminum nostrum.\\ 
\textred{℟.}~Amen.

\begin{center}\textcolor{red}{\rule{\linewidth}{0.5pt}} \end{center}

\subsubsection*{\center \textsc{De S. M. N. Teresia}}
\begin{tightlist}{}
\item \textred{Ant.} Sancta Mater Terésia, * réspice de cælo, et vide, et vísita víneam istam~: et pérfice eam, quam plantávit déxtera tua.
\item \textred{℣.}~Ora pro nobis, sancta Mater Terésia. \textred{℟.}~Ut digni efficiámur promissiónibus Christi.
\item Orémus. \hfill \textred{Oratio}
\end{tightlist}
\lettrine{E}xáudi nos, Deus salutáris noster~: ut sicut de beátæ Terésiæ Vírginis tuæ et Matris nostræ Commemoratióne gaudémus; ita cæléstis ejus doctrínæ pábulo nutriámur, et piæ devotiónis erudiámur afféctu. Per Christum Dóminum nostrum. \textred{℟.}~Amen.

\begin{center}\textcolor{red}{\rule{\linewidth}{0.5pt}} \end{center}

\subsubsection*{\center \textsc{De S. P. N. Joanne}}
\begin{tightlist}{}
\item \textred{Ant.} Fílii, * confortámini, et viríliter ágite in lege~: quia in ipsa gloriósi éritis.
\item \textred{℣.}~Ora pro nobis, sancte Pater Joánnes. \textred{℟.}~Ut digni efficiámur promissiónibus Christi.
\item Orémus. \hfill \textred{Oratio}
\end{tightlist}
\lettrine{D}eus, qui sanctum Joánnem Confessórem tuum atque Doctórem, Patrem nostrum, perféctæ sui abnegatiónis, et crucis amatórem exímium effecísti~: concéde; ut ejus imitatióni júgiter inhæréntes, glóriam assequámur ætérnam. Per Christum Dóminum nostrum. \textred{℟.}~Amen.

\begin{center}\textcolor{red}{\rule{\linewidth}{0.5pt}} \end{center}
\end{multicols}

\end{document}